\documentclass[11pt,a4paper,oneside]{book}

% ============================================================================
% PACKAGES
% ============================================================================
\usepackage[utf8]{inputenc}
\usepackage[T1]{fontenc}
\usepackage[french]{babel}
\usepackage{lmodern}
\usepackage{microtype}
\usepackage[top=2.5cm,bottom=2.5cm,left=2.5cm,right=2.5cm,headheight=15pt]{geometry}

\usepackage{amsmath,amsfonts,amssymb,amsthm}
\usepackage{mathtools}
\usepackage{gensymb}
\usepackage{icomma}
\usepackage{cancel}
\usepackage{multicol}
\usepackage{enumitem}
\usepackage{booktabs}
\usepackage{array}
\usepackage{graphicx}
\usepackage{xcolor}
\usepackage{tikz}
\usepackage{pgfplots}
\pgfplotsset{compat=1.18}
\usetikzlibrary{angles,quotes,calc,positioning,arrows.meta,patterns,decorations.pathreplacing,babel}

\usepackage[most]{tcolorbox}
\usepackage{fancyhdr}
\usepackage{titlesec}
\usepackage{hyperref}

% ============================================================================
% COULEURS & STYLE UCLOUVAIN
% ============================================================================
\definecolor{uclblue}{RGB}{0, 53, 107}
\definecolor{ucllightblue}{RGB}{235, 243, 250}
\definecolor{uclaccent}{RGB}{0, 114, 188}
\definecolor{defgreen}{RGB}{20, 110, 60}
\definecolor{deflightgreen}{RGB}{238, 248, 242}
\definecolor{exorange}{RGB}{180, 85, 0}
\definecolor{exlightorange}{RGB}{253, 246, 238}
\definecolor{warnred}{RGB}{170, 30, 30}
\definecolor{warnlightred}{RGB}{253, 240, 240}
\definecolor{darkgray}{RGB}{60, 60, 60}

\hypersetup{
    colorlinks=true,
    linkcolor=uclblue,
    urlcolor=uclaccent,
    citecolor=uclblue,
    pdftitle={Syllabus d'Analyse - Examen d'entrée Ingénieur Civil UCLouvain},
    pdfauthor={Jean-François Remacle et Pierre Schaus (Prompt Engineers) & IA}
}

% ============================================================================
% BOÎTES PÉDAGOGIQUES (TCOLORBOX)
% ============================================================================
\newtcolorbox{theoreme}[1][]{
    enhanced,
    breakable,
    colback=ucllightblue,
    colframe=uclblue,
    fonttitle=\bfseries,
    title={#1},
    arc=3pt,
    boxrule=0.8pt,
    left=8pt,right=8pt,top=6pt,bottom=6pt
}

\newtcolorbox{definition}[1][]{
    enhanced,
    breakable,
    colback=deflightgreen,
    colframe=defgreen,
    fonttitle=\bfseries,
    title={#1},
    arc=3pt,
    boxrule=0.8pt,
    left=8pt,right=8pt,top=6pt,bottom=6pt
}

\newtcolorbox{exemple}[1][]{
    enhanced,
    breakable,
    colback=exlightorange,
    colframe=exorange,
    fonttitle=\bfseries,
    title={#1},
    arc=3pt,
    boxrule=0.8pt,
    left=8pt,right=8pt,top=6pt,bottom=6pt
}

\newtcolorbox{attention}[1][Attention]{
    enhanced,
    breakable,
    colback=warnlightred,
    colframe=warnred,
    fonttitle=\bfseries,
    title={#1},
    arc=3pt,
    boxrule=0.8pt,
    left=8pt,right=8pt,top=6pt,bottom=6pt
}

\newtcolorbox{astuce}[1][Astuce \& Méthode]{
    enhanced,
    breakable,
    colback=gray!8,
    colframe=darkgray,
    fonttitle=\bfseries,
    title={#1},
    arc=3pt,
    boxrule=0.8pt,
    left=8pt,right=8pt,top=6pt,bottom=6pt
}

% ============================================================================
% MISE EN PAGE & EN-TÊTES
% ============================================================================
\setlength{\parindent}{0pt}
\setlength{\parskip}{0.45em}
\setlist[itemize]{topsep=3pt,itemsep=2pt}
\setlist[enumerate]{topsep=4pt,itemsep=4pt}

\pagestyle{fancy}
\fancyhf{}
\fancyhead[L]{\small\textcolor{uclblue}{\nouppercase{\leftmark}}}
\fancyhead[R]{\small\textcolor{darkgray}{Analyse Mathématique -- UCLouvain}}
\fancyfoot[C]{\thepage}
\renewcommand{\headrulewidth}{0.4pt}

\titleformat{\chapter}[display]
  {\normalfont\huge\bfseries\color{uclblue}}
  {\filright\large\textsc{\chaptertitlename}\ \thechapter}
  {1ex}
  {\titlerule\vspace{1.2ex}\filright}
  [\vspace{1ex}\titlerule]

\titleformat{\section}
  {\normalfont\Large\bfseries\color{uclblue}}
  {\thesection}{0.8em}{}

\titleformat{\subsection}
  {\normalfont\large\bfseries\color{uclaccent}}
  {\thesubsection}{0.7em}{}

% ============================================================================
% COMMANDES MATHÉMATIQUES
% ============================================================================
\newcommand{\R}{\mathbb{R}}
\newcommand{\Z}{\mathbb{Z}}
\newcommand{\N}{\mathbb{N}}
\newcommand{\Q}{\mathbb{Q}}
\newcommand{\C}{\mathbb{C}}
\newcommand{\D}{\mathbb{D}}
\newcommand{\dd}{\,\mathrm{d}}
\DeclareMathOperator{\dom}{dom}
\DeclareMathOperator{\cotg}{cotg}
\DeclareMathOperator{\tg}{tg}

% ============================================================================
% DOCUMENT
% ============================================================================
\begin{document}

% ----------------------------------------------------------------------------
% PAGE DE TITRE
% ----------------------------------------------------------------------------
\begin{titlepage}
    \centering
    \vspace*{0.8cm}
    {\Large \textsc{Université catholique de Louvain (UCLouvain)}}\\[0.3cm]
    {\large \textsc{École Polytechnique de Louvain (EPL)}}\\[1.3cm]

    \begin{tcolorbox}[colback=ucllightblue,colframe=uclblue,arc=5pt,boxrule=1.5pt,top=15pt,bottom=15pt]
        \centering
        {\huge\bfseries\color{uclblue} Syllabus d'Analyse\\ Mathématique}\\[0.4cm]
        {\Large\color{uclaccent} Préparation à l'examen d'admission\\ aux études d'ingénieur civil}
    \end{tcolorbox}

    \vspace{0.7cm}

    % Illustration centrale : courbe, tangente, extremum et aire intégrale
    \begin{tikzpicture}[scale=1.1, >=Stealth]
        \begin{axis}[
            width=11.5cm, height=7.2cm,
            axis lines=middle,
            xmin=-0.4, xmax=4.4,
            ymin=-0.4, ymax=2.5,
            xtick={0.5, 1.5708, 3.5},
            xticklabels={$a$, $x^*$, $b$},
            ytick={0},
            xlabel={$x$},
            ylabel={$y$},
            samples=120,
            domain=0:4.1,
            clip=false
        ]
            % Aire sous la courbe entre a=0.5 et b=3.5
            \addplot[fill=uclaccent!18, draw=none, domain=0.5:3.5] {1.1 + 0.8*sin(deg(x)) - 0.12*x^2 + 0.25*x} \closedcycle;
            
            % Courbe principale f(x)
            \addplot[very thick, uclblue, domain=0:4.1] {1.1 + 0.8*sin(deg(x)) - 0.12*x^2 + 0.25*x} node[pos=0.92, above right] {$y = f(x)$};

            % Tangente horizontale au maximum local
            \draw[thick, warnred, dashed] (axis cs:0.65, 1.999) -- (axis cs:2.55, 1.999) node[right] {\small $f'(x^*) = 0$};
            \fill[warnred] (axis cs:1.68, 1.999) circle (2.2pt) node[above] {\small Maximum local};

            % Tangente en un point
            \draw[thick, defgreen] (axis cs:0.05, 1.1525) -- (axis cs:1.25, 2.4125) node[above left] {\small Tangente $y = f(x_0) + f'(x_0)(x-x_0)$};
            \fill[defgreen] (axis cs:0.5, 1.625) circle (2pt);

            % Bornes d'intégration
            \draw[dashed, uclblue] (axis cs:0.5, 0) -- (axis cs:0.5, 1.625);
            \draw[dashed, uclblue] (axis cs:3.5, 0) -- (axis cs:3.5, 1.222);

            % Étiquette de l'intégrale
            \node[uclblue] at (axis cs:2.1, 0.75) {\large $\displaystyle \int_a^b f(x)\,\mathrm{d}x$};
        \end{axis}
    \end{tikzpicture}

    \vfill

    {\large \textbf{Direction artistique \& \emph{« Prompt Engineering »} :}}\\[0.2cm]
    {\Large \textbf{L'équipe de l'examen d'entrée de l'EPL}}\\[0.15cm]
    {\large \textbf{Jean-François Remacle} et \textbf{Pierre Schaus}}\\[0.3cm]
    {\normalsize \emph{Rédigé à la sueur des GPU par l'Intelligence Artificielle,\\ sous la supervision bienveillante (et quelques prompts bien sentis) de ses humains.}}

    \vspace{0.5cm}
    {\small \textbf{Sources fournies à l'IA :}\\
    Le référentiel de mathématiques de 5\ieme{} et 6\ieme{} secondaire de la Fédération Wallonie-Bruxelles\\
    ainsi que l'intégralité des annales et solutionnaires officiels d'Analyse de l'UCLouvain (2021--2026).}\\[0.4cm]
    {\normalsize Édition 2026--2027}
\end{titlepage}

% ----------------------------------------------------------------------------
% AVANT-PROPOS ET CONSEILS MÉTHODOLOGIQUES
% ----------------------------------------------------------------------------
\chapter*{Avant-propos et conseils méthodologiques}
\addcontentsline{toc}{chapter}{Avant-propos et conseils méthodologiques}
\markboth{Avant-propos}{Avant-propos}

Ce syllabus a été conçu pour accompagner les étudiantes et étudiants qui préparent l'épreuve d'\textbf{analyse mathématique} de l'examen spécial d'admission aux études d'ingénieur civil à l'\textbf{UCLouvain} (École Polytechnique de Louvain -- EPL). Il est pensé pour être pédagogique, progressif, rigoureux et directement accessible à toute personne ayant suivi le cours de mathématiques de 5\ieme{} et 6\ieme{} secondaire en Communauté française de Belgique (Fédération Wallonie-Bruxelles, en particulier l'option mathématiques 6h/semaine, ou 4h avec un investissement personnel complémentaire).

\section*{Origine, auteurs et sources du document}

\begin{theoreme}[« Prompt Engineers », IA et sources du document]
\begin{itemize}
    \item \textbf{\emph{Prompt Engineers} en chef :} Les professeurs \textbf{Jean-François Remacle} et \textbf{Pierre Schaus} (UCLouvain -- EPL), qui ont initié cette collection de syllabus ouverts et fourni les annales, solutionnaires et directives pédagogiques.
    \item \textbf{Auteur exécutant (Intelligence Artificielle) :} Dans le même esprit que le syllabus de Trigonométrie et de Géométrie, la rédaction détaillée, la structuration \LaTeX{}, le tracé des figures TikZ/PGFPlots et la mise au net des corrigés multi-méthodes ont été réalisés avec l'aide de l'intelligence artificielle.
    \item \textbf{Annales et solutionnaires officiels d'Analyse (UCLouvain 2021--2026) :} Ce syllabus intègre et explique en détail l'ensemble des questions posées aux sessions de \textbf{juillet et septembre de 2021 à 2026} à l'examen d'admission de l'UCLouvain. Nous tenons à remercier chaleureusement l'équipe des enseignantes et enseignants responsables de l'épreuve d'Analyse à l'EPL pour la richesse pédagogique de leurs questionnaires et de leurs solutionnaires officiels (qui proposent systématiquement plusieurs méthodes de résolution, des astuces de vérification et des commentaires formateurs).
\end{itemize}
\end{theoreme}

\section*{Structure type de l'épreuve d'Analyse à l'UCLouvain}

L'examen d'admission d'Analyse à l'EPL se déroule \textbf{sans calculatrice} (ni téléphone ou montre connectée) et repose sur une structure très claire et régulière qu'il est essentiel de bien connaître pour organiser ses révisions :
\begin{enumerate}
    \item \textbf{Question 1 : Étude complète d'une fonction $f : \D \subset \R \to \R$ (Chapitres 1 à 4)} : détermination du domaine de définition $\D$ et des racines, limites aux bords ou asymptotes (verticales, horizontales, obliques), calcul et factorisation de la dérivée première $f'(x)$ et de la dérivée seconde $f''(x)$, tableau complet des variations et de la concavité (avec valeurs exactes et valeurs numériques approchées grâce aux aides fournies dans l'énoncé), et esquisse soignée du graphe sur un quadrillage chiffré.
    \item \textbf{Question 2 : Problème d'optimisation appliquée (Chapitre 5)} : mise en équation d'un problème concret d'ingénierie, de physique ou d'économie (isolation thermique, boîte métallique, drone-lanceur, solide de révolution minimal, plus court chemin sable/route, train de nuit, panneau solaire et loi de Planck/Wien, pharmacocinétique de la caféine, tournoi des « legendgénieurs », propulseur spatial, transport maritime, enclos agricole), réduction à une fonction d'une seule variable par élimination des contraintes, recherche de l'extremum par annulation de la dérivée et justification rigoureuse qu'il s'agit bien d'un maximum ou d'un minimum.
    \item \textbf{Question 3 : Calcul d'une primitive $\int f(x)\dd x$ (Chapitre 6)} : recherche d'une primitive combinant reconnaissance de dérivées composées, changements de variable (substitution), intégration par parties (parfois double ou cyclique), identités trigonométriques ou décomposition en fractions simples.
    \item \textbf{Questions 4 et 5 : Limites, volumes de révolution, inégalités intégrales et théorèmes fondamentaux (Chapitres 2, 3 et 7)} :
    \begin{itemize}
        \item Calcul de limites indéterminées ($\frac{0}{0}, \frac{\infty}{\infty}, 0\cdot\infty$) via la \textbf{règle de L'Hospital}, la mise en évidence du terme dominant ou les limites trigonométriques fondamentales.
        \item Calcul du \textbf{volume d'un solide de révolution} $V = \pi \int_a^b [R(x)]^2\dd x$ (pièce mécanique, tore/anneau, œuf fossilisé).
        \item Démonstration d'\textbf{inégalités fonctionnelles} par intégration d'encadrements sur $[1, x]$ ou $[0, x]$, ou par étude du signe de la dérivée d'une différence.
        \item Questions de \textbf{théorie et de raisonnement} mobilisant la définition de la limite ($\varepsilon$-$M$), la continuité et la dérivabilité en un point, le \textbf{théorème des valeurs intermédiaires (TVI)}, le \textbf{théorème de Rolle}, le \textbf{théorème des accroissements finis (TAF)}, les équations différentielles élémentaires ou l'étude d'une fonction réciproque (comme la fonction de Lambert $W_0$).
    \end{itemize}
\end{enumerate}

\section*{Conventions de notation et règles d'or pour le jour de l'examen}

\begin{itemize}
    \item \textbf{Notations officielles :} La fonction exponentielle de base $e \approx 2,718$ est notée $e^x$ ou $\exp(x)$. Le logarithme népérien (en base $e$) est exclusivement noté $\ln(x)$ (défini sur $\R_0^+ = \,]0, +\infty[$). Pour les fonctions trigonométriques et le logarithme, $\sin^2(x)$, $\cos^3(x)$ ou $\ln^2(x)$ désignent respectivement $(\sin(x))^2$, $(\cos(x))^3$ et $(\ln(x))^2$.
    \item \textbf{Cadre « Réponse finale » et cadre « Développement » :} Lorsqu'une question comporte un petit cadre pour la réponse finale et un grand cadre pour le développement, \textbf{les deux doivent impérativement être remplis}. Une réponse finale parachutée sans raisonnement, ou un long développement sans réponse finale clairement reportée dans le cadre prévu, entraîne une perte importante de points.
    \item \textbf{Calcul littéral avant le calcul numérique :} Dans les problèmes d'optimisation (Question 2), gardez les lettres ($C, V, H, \alpha, \beta, f_0\dots$) le plus longtemps possible ! Dérivez l'expression littérale et ne remplacez par les valeurs numériques qu'à la toute dernière ligne : cela évite les erreurs de calcul, révèle des simplifications spectaculaires (souvent certaines constantes s'éliminent d'elles-mêmes !) et permet au correcteur de suivre parfaitement votre démarche.
    \item \textbf{Vérification systématique :} En analyse, vous avez l'immense avantage de pouvoir \textbf{vérifier vous-même vos résultats} : dérivez toujours au brouillon la primitive que vous venez de trouver à la Question 3 pour vérifier qu'elle redonne l'intégrande, et contrôlez la cohérence entre le signe de $f'(x)$, les limites aux bords et les valeurs de $f(x)$ dans votre tableau de variations de la Question 1 !
\end{itemize}

\tableofcontents

% ----------------------------------------------------------------------------
% CHAPITRES
% ----------------------------------------------------------------------------
% ============================================================================
% CHAPITRE 1 : FONCTIONS RÉELLES, DOMAINES ET FONCTIONS TRANSCENDANTES
% ============================================================================
\chapter{Fonctions Réelles, Domaines et Fonctions Transcendantes}

L'analyse mathématique a pour objet l'étude rigoureuse des fonctions d'une variable réelle : comment elles varient, comment elles se comportent aux bornes de leur domaine, à quelle vitesse elles évoluent (dérivée) et comment elles s'accumulent (intégrale). Ce premier chapitre consolide les notions fondamentales de 5\ieme{} et 6\ieme{} secondaire : le domaine de définition, les racines, les symétries (parité et périodicité), les fonctions réciproques, ainsi que les deux grandes fonctions transcendantes omniprésentes à l'examen d'admission de l'UCLouvain : le logarithme népérien $\ln(x)$ et l'exponentielle $e^x$.

% ----------------------------------------------------------------------------
\section{Généralités sur les fonctions réelles, domaines et racines}
% ----------------------------------------------------------------------------

\subsection{Définition, domaine d'existence et racines}

\begin{definition}[Fonction réelle, domaine de définition et image]
Une \textbf{fonction réelle d'une variable réelle} $f : \D \subset \R \to \R$ est une relation qui, à tout nombre réel $x$ appartenant à son \textbf{domaine de définition} (noté $\D$, $\dom f$ ou $D_f$), associe \emph{au plus un} nombre réel noté $f(x)$, appelé l'\textbf{image} de $x$ par $f$.
\begin{itemize}
    \item Lorsqu'une fonction est donnée uniquement par son expression analytique $f(x)$, son \textbf{domaine de définition naturel} $\D$ est l'ensemble de tous les réels $x \in \R$ pour lesquels le calcul de $f(x)$ a un sens (c'est-à-dire qui vérifient toutes les \textbf{conditions d'existence}, abrégées CE).
    \item L'\textbf{ensemble image} de $f$, noté $\operatorname{Im} f$ ou $f(\D)$, est l'ensemble de toutes les valeurs prises par $f(x)$ lorsque $x$ parcourt $\D$ :
    \[
        \operatorname{Im} f = f(\D) = \bigl\{ y \in \R \;\big|\; \exists x \in \D \text{ tel que } y = f(x) \bigr\}.
    \]
    \item Une \textbf{racine} (ou un \textbf{zéro}) de la fonction $f$ est une valeur $x_0 \in \D$ telle que $f(x_0) = 0$. Géométriquement, les racines sont les abscisses des points d'intersection du graphe de $f$ avec l'axe horizontal $Ox$.
\end{itemize}
\end{definition}

À l'examen d'entrée, la toute première sous-question de la Question 1 consiste systématiquement à déterminer le domaine de définition $\D$ de la fonction étudiée (et très souvent ses racines). Toutes les conditions d'existence se ramènent à quatre règles fondamentales.

\begin{theoreme}[Les quatre règles d'or des conditions d'existence (CE)]
Pour déterminer le domaine de définition $\D$ d'une expression $f(x)$, on inspecte la présence de :
\begin{enumerate}
    \item \textbf{Dénominateurs (fractions) :} On ne divise jamais par zéro.
    \[
        \frac{N(x)}{D(x)} \quad \text{existe} \iff \boxed{D(x) \neq 0}.
    \]
    En particulier, pour $\tan(u(x)) = \frac{\sin(u(x))}{\cos(u(x))}$, il faut $\cos(u(x)) \neq 0 \iff u(x) \neq \frac{\pi}{2} + k\pi$ ($k \in \Z$).
    \item \textbf{Racines d'indice pair ($\sqrt{\phantom{u}}$, $\sqrt[4]{\phantom{u}}$, ou exposant fractionnaire $u^{p/2q}$) :} Le radicande doit être positif ou nul.
    \[
        \sqrt{u(x)} \quad \text{existe dans } \R \iff \boxed{u(x) \ge 0}.
    \]
    \emph{Remarque :} Si la racine carrée figure au dénominateur, comme $\frac{1}{\sqrt{u(x)}} = (u(x))^{-1/2}$, les règles 1 et 2 se combinent en une inégalité stricte : $u(x) > 0$. En revanche, une racine d'indice \emph{impair} telle que $\sqrt[3]{u(x)} = (u(x))^{1/3}$ est définie quel que soit le signe de $u(x)$.
    \item \textbf{Logarithmes ($\ln$) :} L'argument d'un logarithme doit être \textbf{strictement positif}.
    \[
        \ln\bigl(u(x)\bigr) \quad \text{existe} \iff \boxed{u(x) > 0}.
    \]
    \item \textbf{Arcsinus et Arccosinus ($\arcsin$, $\arccos$) :} Leur argument (qui représente un sinus ou un cosinus) doit être compris au sens large entre $-1$ et $1$.
    \[
        \arcsin\bigl(u(x)\bigr) \;\text{ ou }\; \arccos\bigl(u(x)\bigr) \quad \text{existe} \iff \boxed{-1 \le u(x) \le 1}.
    \]
\end{enumerate}
\end{theoreme}

\begin{attention}[Piège lors de la recherche des racines d'une fonction avec racine carrée]
Lorsqu'on résout une équation irrationnelle du type $\sqrt{A(x)} = B(x)$ (comme à l'examen de \textbf{Juillet 2024, Question 1a} pour trouver la racine de $f(x) = \sqrt{1+x} - x$), n'oubliez jamais qu'une racine carrée désigne par convention un nombre \textbf{positif ou nul} ($\sqrt{A(x)} \ge 0$) !
Ainsi, l'équivalence rigoureuse est :
\[
    \boxed{\sqrt{A(x)} = B(x) \iff \begin{cases} B(x) \ge 0 \quad (\text{condition de signe indispensable !}) \\ A(x) = \bigl(B(x)\bigr)^2 \end{cases}}
\]
Élever au carré sans vérifier que $B(x) \ge 0$ introduit des \textbf{solutions parasites} (celles qui vérifient $\sqrt{A(x)} = -B(x)$) qui appartiennent pourtant au domaine de définition de la racine carrée !
\end{attention}

\begin{exemple}[Domaines de définition et racines aux examens UCLouvain (2021--2026)]
Analysons six questions d'examen illustrant chaque type de condition d'existence :

\medskip
\textbf{1) UCLouvain -- Juillet 2021, Question 1a :} $f(x) = \ln\left(\dfrac{1}{\cos(x)}\right)$.\\
\emph{Résolution :} Il y a deux conditions d'existence : (1) le dénominateur $\cos(x) \neq 0$, et (2) l'argument du logarithme $\frac{1}{\cos(x)} > 0 \iff \cos(x) > 0$. Comme $\cos(x) > 0$ implique automatiquement $\cos(x) \neq 0$, la condition unique est $\cos(x) > 0$. Sur le cercle trigonométrique (quadrants I et IV), le cosinus est strictement positif sur $\left]-\frac{\pi}{2}, \frac{\pi}{2}\right[$ modulo $2\pi$. Le domaine est donc l'union disjointe :
\[
    \D = \bigcup_{k \in \Z} \left]-\frac{\pi}{2} + 2k\pi,\; \frac{\pi}{2} + 2k\pi\right[.
\]

\medskip
\textbf{2) UCLouvain -- Septembre 2021, Question 1a :} $f(x) = x + 2\ln\left(1 + \dfrac{1}{x}\right)$.\\
\emph{Résolution :} Les conditions d'existence sont (1) $x \neq 0$ et (2) $1 + \frac{1}{x} > 0 \iff \frac{x+1}{x} > 0$. Un tableau de signes du quotient $\frac{x+1}{x}$ (dont le numérateur s'annule en $x=-1$ et le dénominateur en $x=0$) montre immédiatement que la fraction est strictement positive à l'extérieur des racines :
\[
    \D = \,]-\infty, -1[ \;\cup\; ]0, +\infty[.
\]

\medskip
\textbf{3) UCLouvain -- Juillet 2022, Question 1a (et Septembre 2022, Question 1a) :}\\
Soit $f_1(x) = \ln\left(\dfrac{1}{1+\cos x}\right)$ (Juillet 2022) et $f_2(x) = \dfrac{\cos^2 x}{1 + \cos x}$ (Septembre 2022).\\
\emph{Résolution :}
\begin{itemize}
    \item Pour $f_2(x)$, la seule condition est $1 + \cos x \neq 0 \iff \cos x \neq -1 \iff x \neq \pi + 2k\pi$ ($k \in \Z$).
    \item Pour $f_1(x)$, il faut $\frac{1}{1+\cos x} > 0 \iff 1 + \cos x > 0 \iff \cos x > -1$. Or, pour tout réel $x$, on a toujours $\cos x \in [-1, 1]$, donc $\cos x \ge -1$ est toujours vérifié ! La condition stricte $\cos x > -1$ revient donc simplement à exclure le cas d'égalité $\cos x = -1$. Les deux fonctions ont exactement le même domaine :
    \[
        \D = \R \setminus \bigl\{ \pi + 2k\pi \;\big|\; k \in \Z \bigr\} = \bigcup_{k \in \Z} \bigl](2k-1)\pi,\; (2k+1)\pi\bigr[.
    \]
\end{itemize}

\medskip
\textbf{4) UCLouvain -- Juillet 2024, Question 1a :} $f(x) = \sqrt{1+x} - x$. Déterminer $\D$ et la/les racine(s) de $f$.\\
\emph{Résolution :} La condition d'existence de la racine carrée est $1 + x \ge 0 \iff x \ge -1$, d'où $\D = [-1, +\infty[$.
Pour trouver les racines de $f$, résolvons $f(x) = 0 \iff \sqrt{1+x} = x$. Comme le membre de gauche $\sqrt{1+x}$ est positif ou nul, toute solution doit impérativement vérifier $x \ge 0$. Sous cette condition $x \ge 0$, élevons au carré :
\[
    1 + x = x^2 \iff x^2 - x - 1 = 0 \iff x = \frac{1 \pm \sqrt{5}}{2}.
\]
Bien que les deux valeurs $x_1 = \frac{1+\sqrt{5}}{2} \approx 1,618$ (le nombre d'or !) et $x_2 = \frac{1-\sqrt{5}}{2} \approx -0,618$ appartiennent au domaine $\D = [-1, +\infty[$, la valeur $x_2 < 0$ est une \textbf{racine parasite} (elle vérifie $\sqrt{1+x_2} = -x_2$ et non $+x_2$). La fonction $f$ possède donc \textbf{une unique racine} : $x_1 = \dfrac{1+\sqrt{5}}{2}$.

\medskip
\textbf{5) UCLouvain -- Juillet 2025, Question 1a :} $f(x) = x + \ln\bigl(\sin(x) + \cos(x)\bigr)$. Formuler $\D$ comme union disjointe d'intervalles ouverts $]a_k, b_k[$ ($k \in \Z$).\\
\emph{Résolution :} Il faut $\sin(x) + \cos(x) > 0$. On propose deux méthodes complémentaires :
\begin{itemize}
    \item \emph{Méthode 1 (par le cercle trigonométrique et l'étude des zéros) :} Cherchons d'abord où la somme s'annule : $\sin(x) + \cos(x) = 0 \iff \sin(x) = -\cos(x) \iff \tan(x) = -1 \iff x = -\frac{\pi}{4} + k\pi$. Sur un tour du cercle trigonométrique, les deux points où $\sin x$ et $\cos x$ sont opposés sont $-\frac{\pi}{4}$ (quadrant IV) et $\frac{3\pi}{4}$ (quadrant II). Entre $-\frac{\pi}{4}$ et $\frac{3\pi}{4}$ (par exemple en $x = 0$ où $\sin 0 + \cos 0 = 1 > 0$, ou en $x = \frac{\pi}{4}$ où la somme vaut $\sqrt{2} > 0$), la somme $\sin x + \cos x$ est strictement positive.
    \item \emph{Méthode 2 (par les formules de Simpson ou d'addition) :} En écrivant $\cos(x) = \sin\left(\frac{\pi}{2}-x\right)$ et en appliquant Simpson (ou en factorisant par $\sqrt{2}$), on obtient :
    \[
        \sin(x) + \cos(x) = \sqrt{2}\left(\frac{\sqrt{2}}{2}\cos x + \frac{\sqrt{2}}{2}\sin x\right) = \sqrt{2}\cos\left(x - \frac{\pi}{4}\right).
    \]
    La condition devient $\cos\left(x - \frac{\pi}{4}\right) > 0 \iff -\frac{\pi}{2} + 2k\pi < x - \frac{\pi}{4} < \frac{\pi}{2} + 2k\pi \iff -\frac{\pi}{4} + 2k\pi < x < \frac{3\pi}{4} + 2k\pi$.
\end{itemize}
Dans les deux cas, $\D = \bigcup_{k\in\Z} \,]a_k, b_k[$ avec $\boxed{a_k = -\dfrac{\pi}{4} + 2k\pi}$ et $\boxed{b_k = \dfrac{3\pi}{4} + 2k\pi}$.
\end{exemple}

% ----------------------------------------------------------------------------
\section{Parité, périodicité et transformations de graphes}
% ----------------------------------------------------------------------------

Repérer dès le premier coup d'œil qu'une fonction est paire, impaire ou périodique permet de diviser par deux (ou de restreindre à une seule période) tout le travail d'étude de fonction et de calcul d'intégrale définie !

\begin{definition}[Fonction paire, fonction impaire et fonction périodique]
Soit $f : \D \subset \R \to \R$ une fonction réelle.
\begin{enumerate}
    \item On suppose le domaine $\D$ \textbf{symétrique par rapport à $0$} (c'est-à-dire que $\forall x \in \D, -x \in \D$) :
    \begin{itemize}
        \item $f$ est \textbf{paire} si $\forall x \in \D, \quad \boxed{f(-x) = f(x)}$.\\
        \emph{Interprétation graphique :} Le graphe de $f$ est symétrique par rapport à l'\textbf{axe vertical $Oy$}. Exemples : $x^2, x^4, |x|, \cos(x), \ln(1+x^2), x^2\sin(x^2) + \cos(x^2)$ (Juillet 2023 Série 1).
        \item $f$ est \textbf{impaire} si $\forall x \in \D, \quad \boxed{f(-x) = -f(x)}$.\\
        \emph{Interprétation graphique :} Le graphe de $f$ est symétrique par symétrie centrale par rapport à l'\textbf{origine $O(0,0)$}. Si $0 \in \D$, alors nécessairement $f(0) = 0$. Exemples : $x, x^3, \frac{1}{x}, \sin(x), \tan(x), \arcsin(x), \arctan(x)$.
    \end{itemize}
    \item Soit $T > 0$. On dit que $f$ est \textbf{périodique de période $T$} (ou $T$-périodique) si, pour tout $x \in \D$, on a $x + T \in \D$ et :
    \[
        \forall x \in \D, \quad \boxed{f(x + T) = f(x)}.
    \]
    La plus petite période strictement positive $T > 0$ (lorsqu'elle existe) est appelée la \textbf{période fondamentale}. Le graphe de $f$ s'obtient par translations horizontales successives de vecteur $(kT, 0)$ ($k \in \Z$).
\end{enumerate}
\end{definition}

\begin{center}
\begin{tikzpicture}[scale=0.95, >=Stealth]
    % Fonction paire
    \begin{scope}[shift={(0,0)}]
        \draw[->, thick, gray] (-2.5,0) -- (2.5,0) node[right, black] {$x$};
        \draw[->, thick, gray] (0,-0.4) -- (0,2.5) node[above, black] {$y$};
        \draw[very thick, uclblue, domain=-2.1:2.1, samples=80] plot (\x, {0.45*(\x)^2 + 0.3});
        \draw[dashed, warnred] (-1.5,0) -- (-1.5,1.3125) -- (1.5,1.3125) -- (1.5,0);
        \fill[uclblue] (-1.5,1.3125) circle (1.8pt);
        \fill[uclblue] (1.5,1.3125) circle (1.8pt);
        \node[below] at (-1.5,0) {\small $-x$};
        \node[below] at (1.5,0) {\small $x$};
        \node[above, uclblue, font=\bfseries\small] at (0,2.55) {Fonction paire : $f(-x) = f(x)$};
        \node[below, darkgray, font=\footnotesize] at (0,-0.55) {Symétrie orthogonale d'axe $Oy$};
    \end{scope}

    % Fonction impaire
    \begin{scope}[shift={(7.5,0.8)}]
        \draw[->, thick, gray] (-2.5,0) -- (2.5,0) node[right, black] {$x$};
        \draw[->, thick, gray] (0,-1.8) -- (0,1.8) node[above, black] {$y$};
        \draw[very thick, defgreen, domain=-2.1:2.1, samples=80] plot (\x, {0.22*(\x)^3 + 0.25*\x});
        \draw[dashed, warnred] (-1.5,-1.1175) -- (1.5,1.1175);
        \draw[dashed, gray] (-1.5,0) -- (-1.5,-1.1175);
        \draw[dashed, gray] (1.5,0) -- (1.5,1.1175);
        \fill[defgreen] (-1.5,-1.1175) circle (1.8pt) node[left] {\small $f(-x)$};
        \fill[defgreen] (1.5,1.1175) circle (1.8pt) node[right] {\small $f(x)$};
        \fill[warnred] (0,0) circle (1.8pt);
        \node[above] at (-1.5,0) {\small $-x$};
        \node[below] at (1.5,0) {\small $x$};
        \node[above, defgreen, font=\bfseries\small] at (0,1.85) {Fonction impaire : $f(-x) = -f(x)$};
        \node[below, darkgray, font=\footnotesize] at (0,-1.85) {Symétrie centrale de centre $O(0,0)$};
    \end{scope}
\end{tikzpicture}
\end{center}

\begin{theoreme}[Conséquences pratiques de la parité et de la périodicité]
\begin{enumerate}
    \item \textbf{Parité des dérivées :} Si $f$ est dérivable sur un domaine symétrique $\D$ :
    \begin{itemize}
        \item Si $f$ est \textbf{paire}, alors sa dérivée $f'$ est \textbf{impaire} ($f'(-x) = -f'(x)$), et sa dérivée seconde $f''$ est \textbf{paire}.
        \item Si $f$ est \textbf{impaire}, alors sa dérivée $f'$ est \textbf{paire} ($f'(-x) = f'(x)$), et sa dérivée seconde $f''$ est \textbf{impaire}.
    \end{itemize}
    \item \textbf{Parité d'un produit ou d'un quotient :} La règle des signes s'applique exactement comme pour $+$ et $-$ :
    \[
        \text{Paire} \times \text{Paire} = \text{Paire}, \qquad \text{Impaire} \times \text{Impaire} = \text{Paire}, \qquad \text{Paire} \times \text{Impaire} = \text{Impaire}.
    \]
    \item \textbf{Périodicité de $\sin(\omega x)$, $\cos(\omega x)$ et $\tan(\omega x)$ (pour $\omega > 0$) :}
    \begin{itemize}
        \item $\sin(\omega x)$ et $\cos(\omega x)$ sont périodiques de période $T = \dfrac{2\pi}{\omega}$.
        \item $\tan(\omega x)$, $\sin^2(\omega x)$ et $\cos^2(\omega x) = \frac{1+\cos(2\omega x)}{2}$ sont périodiques de période $T = \dfrac{\pi}{\omega}$.
    \end{itemize}
\end{enumerate}
\end{theoreme}

\begin{exemple}[Question de théorie sur la périodicité : UCLouvain -- Septembre 2024, Question 5]
\emph{Énoncé :} Soient deux fonctions $f_1(x)$ et $f_2(x)$ dont le domaine de définition est $\R$. On sait que $\bigl(f_1(x)\bigr)^3$ et $\bigl(f_2(x)\bigr)^2$ sont périodiques, de période $\pi$.
\begin{enumerate}[label=\alph*)]
    \item La fonction $f_1(x)$ est-elle périodique de période $\pi$ ? Si oui, démontrez-le. Si non, fournissez un contre-exemple. Votre résultat change-t-il si $f_1$ est continue et n'est jamais nulle ?
    \item La fonction $f_2(x)$ est-elle périodique de période $\pi$ ? Si oui, démontrez-le. Si non, fournissez un contre-exemple. Votre résultat change-t-il si $f_2$ est continue et n'est jamais nulle ?
\end{enumerate}

\emph{Résolution :}
\begin{enumerate}[label=\alph*)]
    \item \textbf{Pour le cube $\bigl(f_1(x)\bigr)^3$ : OUI, $f_1$ est toujours périodique de période $\pi$.} En effet, par hypothèse, pour tout $x \in \R$ :
    \[
        \bigl(f_1(x+\pi)\bigr)^3 = \bigl(f_1(x)\bigr)^3.
    \]
    Or la fonction cube $h : \R \to \R, y \mapsto y^3$ est \textbf{strictement croissante} (donc bijective) sur tout $\R$ : deux nombres réels qui ont le même cube sont nécessairement égaux (algébriquement, $y_2^3 - y_1^3 = (y_2-y_1)(y_2^2 + y_1 y_2 + y_1^2) = 0 \implies y_2 = y_1$ car $y_2^2 + y_1 y_2 + y_1^2 = (y_2 + \frac{1}{2}y_1)^2 + \frac{3}{4}y_1^2 > 0$ dès que $y_1, y_2$ ne sont pas tous deux nuls). On en déduit immédiatement que $\forall x \in \R, f_1(x+\pi) = f_1(x)$ : $f_1$ est périodique de période $\pi$. Ce résultat est vrai pour toute fonction $f_1$, qu'elle soit continue ou non, nulle ou non.
    \item \textbf{Pour le carré $\bigl(f_2(x)\bigr)^2$ : NON en général, mais OUI si $f_2$ est continue et jamais nulle.}
    \begin{itemize}
        \item \emph{Cas général :} L'égalité $\bigl(f_2(x+\pi)\bigr)^2 = \bigl(f_2(x)\bigr)^2$ implique seulement $f_2(x+\pi) = \pm f_2(x)$ (la fonction peut changer de signe en passant de $x$ à $x+\pi$). Un \textbf{contre-exemple} simple et continu est $f_2(x) = \sin(x)$ (ou $\cos(x)$) : son carré $\sin^2(x) = \frac{1-\cos(2x)}{2}$ est bien périodique de période $\pi$, mais $\sin(x+\pi) = -\sin(x) \neq \sin(x)$, donc $f_2(x) = \sin(x)$ est de période fondamentale $2\pi$ et non $\pi$.
        \item \emph{Si $f_2$ est continue sur $\R$ et ne s'annule jamais :} Par le théorème des valeurs intermédiaires (voir Chapitre 2), une fonction continue sur $\R$ qui ne s'annule jamais garde un \textbf{signe strictement constant} sur tout $\R$ (soit $\forall x \in \R, f_2(x) > 0$, soit $\forall x \in \R, f_2(x) < 0$), car si elle changeait de signe entre deux points, elle devrait s'annuler entre les deux. Puisque $f_2(x)$ et $f_2(x+\pi)$ ont la même valeur absolue et le même signe, on a nécessairement $\forall x \in \R, f_2(x+\pi) = f_2(x)$ : dans ce cas, $f_2$ est bien périodique de période $\pi$ !
    \end{itemize}
\end{enumerate}
\end{exemple}

% ----------------------------------------------------------------------------
\section{Monotonie, bijectivité et fonctions réciproques}
% ----------------------------------------------------------------------------

\subsection{Croissance, décroissance et bijectivité}

\begin{definition}[Croissance et décroissance sur un intervalle $I$]
Soit $f : \D \to \R$ et soit $I \subset \D$ un intervalle :
\begin{itemize}
    \item $f$ est \textbf{croissante} (resp. \textbf{strictement croissante}) sur $I$ si :
    \[
        \forall x_1, x_2 \in I, \quad x_1 < x_2 \implies f(x_1) \le f(x_2) \quad \bigl(\text{resp. } f(x_1) < f(x_2)\bigr).
    \]
    \item $f$ est \textbf{décroissante} (resp. \textbf{strictement décroissante}) sur $I$ si :
    \[
        \forall x_1, x_2 \in I, \quad x_1 < x_2 \implies f(x_1) \ge f(x_2) \quad \bigl(\text{resp. } f(x_1) > f(x_2)\bigr).
    \]
\end{itemize}
\emph{Remarque logique (UCLouvain -- Juillet 2025, Q4) :} Pour prouver qu'une fonction \textbf{n'est pas croissante} sur un intervalle $I$, il suffit d'exhiber \textbf{un seul contre-exemple}, c'est-à-dire deux nombres précis $x_1, x_2 \in I$ tels que $x_1 < x_2$ mais $f(x_1) > f(x_2)$.
\end{definition}

\begin{definition}[Bijection et fonction réciproque $f^{-1}$]
Une fonction $f : I \to J$ (où $J = f(I)$ est l'ensemble image de $I$ par $f$) est dite \textbf{bijective} de $I$ sur $J$ si tout élément $y \in J$ possède \textbf{un et un seul antécédent} $x \in I$ tel que $y = f(x)$.
On définit alors la \textbf{fonction réciproque} $f^{-1} : J \to I$ par la relation fondamentale d'équivalence :
\[
    \boxed{\forall x \in I, \;\forall y \in J : \quad y = f(x) \iff x = f^{-1}(y)}.
\]
En particulier, on a les identités de composition :
\[
    \forall x \in I, \quad f^{-1}\bigl(f(x)\bigr) = x \qquad \text{et} \qquad \forall y \in J, \quad f\bigl(f^{-1}(y)\bigr) = y.
\]
\end{definition}

\begin{theoreme}[Propriétés fondamentales de la fonction réciproque]
Soit $f : I \to J$ une fonction \textbf{continue et strictement monotone} (strictement croissante ou strictement décroissante) sur un intervalle $I \subset \R$. Alors :
\begin{enumerate}
    \item L'image $J = f(I)$ est un intervalle (dont les bornes sont les valeurs ou les limites de $f$ aux bornes de $I$), et $f$ réalise une \textbf{bijection} de $I$ sur $J$.
    \item La fonction réciproque $f^{-1} : J \to I$ a pour \textbf{domaine de définition} l'ensemble image de $f$ : $\dom(f^{-1}) = J = \operatorname{Im}(f)$.
    \item La fonction réciproque $f^{-1}$ est \textbf{continue} sur $J$ et possède \textbf{exactement le même sens de variation} que $f$ (la réciproque d'une fonction strictement croissante est strictement croissante ; la réciproque d'une fonction strictement décroissante est strictement décroissante).
    \item \textbf{Symétrie des graphes :} Dans un repère orthonormé, le graphe de $y = f^{-1}(x)$ est le symétrique du graphe de $y = f(x)$ par rapport à la \textbf{première bissectrice} d'équation $y = x$ (puisque $(x, y)$ appartient au graphe de $f$ si et seulement si $(y, x)$ appartient au graphe de $f^{-1}$).
\end{enumerate}
\end{theoreme}

\begin{attention}[Ne confondez jamais la réciproque $f^{-1}(x)$ et l'inverse $\bigl(f(x)\bigr)^{-1} = \frac{1}{f(x)}$ !]
Pour éviter toute confusion entre la réciproque pour la composition et l'inverse pour la multiplication :
\begin{itemize}
    \item Les fonctions réciproques de $\sin, \cos, \tan$ sont toujours notées $\arcsin(x), \arccos(x), \arctan(x)$ à l'examen d'admission.
    \item L'inverse multiplicatif de $f(x)$ est noté $\dfrac{1}{f(x)}$ ou $\bigl(f(x)\bigr)^{-1}$.
\end{itemize}
\end{attention}

\subsection{Rappel sur les fonctions cyclométriques ($\arcsin, \arccos, \arctan$)}

Bien qu'étudiées en détail dans le syllabus de Trigonométrie, les fonctions réciproques des fonctions trigonométriques interviennent régulièrement en Analyse (dans les calculs de primitives ou d'équations différentielles, comme en Septembre 2021 Q3b ou Septembre 2023 Q4). Retenons leur carte d'identité :

\begin{center}
\renewcommand{\arraystretch}{1.45}
\begin{tabular}{@{}lcccc@{}}
\toprule
\textbf{Fonction} & \textbf{Domaine $\D$} & \textbf{Ensemble image} & \textbf{Sens de variation} & \textbf{Relation de définition} \\
\midrule
$y = \arcsin(x)$ & $[-1, 1]$ & $\left[-\dfrac{\pi}{2}, \dfrac{\pi}{2}\right]$ & Strictement croissante & $\sin(y) = x \text{ et } y \in \left[-\frac{\pi}{2}, \frac{\pi}{2}\right]$ \\
$y = \arccos(x)$ & $[-1, 1]$ & $[0, \pi]$ & Strictement décroissante & $\cos(y) = x \text{ et } y \in [0, \pi]$ \\
$y = \arctan(x)$ & $\R$ & $\left]-\dfrac{\pi}{2}, \dfrac{\pi}{2}\right[$ & Strictement croissante & $\tan(y) = x \text{ et } y \in \left]-\frac{\pi}{2}, \frac{\pi}{2}\right[$ \\
\bottomrule
\end{tabular}
\end{center}

% ----------------------------------------------------------------------------
\section{Logarithme népérien $\ln(x)$ et exponentielle $e^x$}
% ----------------------------------------------------------------------------

À l'examen d'entrée d'Analyse, plus de 80\% des sujets font intervenir la fonction logarithme népérien $\ln(x)$ ou la fonction exponentielle $e^x = \exp(x)$. Il est indispensable de connaître sur le bout des doigts leurs propriétés algébriques, leurs valeurs remarquables et leurs graphes.

\subsection{Le logarithme népérien $\ln(x)$}

\begin{definition}[Le logarithme népérien $\ln(x)$ et la constante d'Euler $e$]
La fonction \textbf{logarithme népérien}, notée $\ln : \R_0^+ = \,]0, +\infty[ \;\to \R$, est l'unique primitive sur $]0, +\infty[$ de la fonction $x \mapsto \frac{1}{x}$ qui s'annule en $x = 1$ :
\[
    \boxed{\forall x > 0, \quad \ln(x) = \int_1^x \frac{1}{s}\dd s \qquad \text{avec } \ln(1) = 0 \text{ et } \bigl(\ln(x)\bigr)' = \frac{1}{x}}.
\]
Puisque sa dérivée $(\ln(x))' = \frac{1}{x}$ est strictement positive sur $]0, +\infty[$, la fonction $\ln$ est \textbf{continue et strictement croissante} de $]0, +\infty[$ sur $\R$.
L'unique réel $x > 0$ tel que $\ln(x) = 1$ est la \textbf{base des logarithmes népériens} (ou nombre de Néper / constante d'Euler), notée $e$ :
\[
    \boxed{\ln(e) = 1 \qquad \text{avec } e \approx 2,71828 \quad \left(\text{et } \frac{1}{e} \approx 0,368,\;\; e^{-2} \approx 0,135,\;\; \ln(2) \approx 0,693\right)}.
\]
\end{definition}

\begin{theoreme}[Propriétés algébriques et signe du logarithme népérien]
Pour tous nombres réels \textbf{strictement positifs} $a > 0$ et $b > 0$, et pour tout exposant réel $p \in \R$ :
\begin{enumerate}
    \item \textbf{Transforme un produit en somme :}
    \[
        \boxed{\ln(a \cdot b) = \ln(a) + \ln(b)}
    \]
    \item \textbf{Transforme un inverse en opposé et un quotient en différence :}
    \[
        \boxed{\ln\left(\frac{1}{a}\right) = -\ln(a)} \qquad \text{et} \qquad \boxed{\ln\left(\frac{a}{b}\right) = \ln(a) - \ln(b)}
    \]
    \item \textbf{Transforme une puissance en facteur multiplicatif :}
    \[
        \boxed{\ln(a^p) = p\,\ln(a)} \qquad \text{et en particulier } \ln(\sqrt{a}) = \ln(a^{1/2}) = \frac{1}{2}\ln(a).
    \]
    \item \textbf{Signe et croissance stricte sur $]0, +\infty[$ :}
    \[
        \ln(a) = \ln(b) \iff a = b, \qquad \ln(a) < \ln(b) \iff 0 < a < b,
    \]
    \[
        \boxed{\begin{cases} 0 < x < 1 &\iff \ln(x) < 0 \\ x = 1 &\iff \ln(x) = 0 \\ x > 1 &\iff \ln(x) > 0 \end{cases}}
    \]
\end{enumerate}
\end{theoreme}

\begin{attention}[Deux pièges algébriques classiques avec $\ln$]
\begin{itemize}
    \item \textbf{Le carré dans $\ln(x^2)$ vs $\ln^2(x)$ :}
    \begin{itemize}
        \item $\ln^2(x) = \bigl(\ln(x)\bigr)^2$ est le carré du logarithme : il n'est défini que pour $x > 0$ et ne vaut \textbf{pas} $2\ln(x)$ !
        \item $\ln(x^2)$ est défini pour tout $x \neq 0$ (même pour $x < 0$ !), et la formule correcte sur $\R_0$ comporte une valeur absolue : $\boxed{\ln(x^2) = 2\ln(|x|)}$ (qui se simplifie en $2\ln(x)$ lorsque $x > 0$, comme en Septembre 2022 Q5).
    \end{itemize}
    \item \textbf{Il n'existe aucune formule pour simplifier $\ln(a + b)$ !} Ne remplacez jamais $\ln(a+b)$ par $\ln(a) + \ln(b)$ ou par $\ln(a)\ln(b)$.
\end{itemize}
\end{attention}

\subsection{La fonction exponentielle $e^x = \exp(x)$}

\begin{definition}[La fonction exponentielle $e^x$]
La fonction \textbf{exponentielle} de base $e$, notée $\exp : \R \to \,]0, +\infty[, x \mapsto e^x$ (ou $\exp(x)$), est la \textbf{fonction réciproque} de la fonction logarithme népérien $\ln : \,]0, +\infty[ \;\to \R$ :
\[
    \boxed{\forall x \in \R, \;\forall y > 0 : \quad y = e^x \iff x = \ln(y)}.
\]
En particulier, on a les relations d'annulation mutuelle :
\[
    \boxed{\forall x \in \R, \quad \ln(e^x) = x} \qquad \text{et} \qquad \boxed{\forall y > 0, \quad e^{\ln(y)} = y}.
\]
\end{definition}

\begin{theoreme}[Propriétés fondamentales de la fonction exponentielle $e^x$]
Pour tous nombres réels $a, b \in \R$ et tout réel $p \in \R$ :
\begin{enumerate}
    \item \textbf{Positivité stricte :} Pour tout $x \in \R$, $\boxed{e^x > 0}$ (une exponentielle ne s'annule \textbf{jamais} et n'est jamais négative !).
    \item \textbf{Valeurs particulières :} $e^0 = 1$, \quad $e^1 = e \approx 2,718$, \quad $e^{-1} = \dfrac{1}{e} \approx 0,368$.
    \item \textbf{Règles de calcul sur les exposants :}
    \[
        \boxed{e^{a+b} = e^a \cdot e^b}, \qquad \boxed{e^{-a} = \frac{1}{e^a}}, \qquad \boxed{e^{a-b} = \frac{e^a}{e^b}}, \qquad \boxed{\bigl(e^a\bigr)^p = e^{a\cdot p}}.
    \]
    \item \textbf{Dérivée et croissance stricte :} La fonction $x \mapsto e^x$ est sa propre dérivée : $\boxed{(e^x)' = e^x > 0}$. Elle est donc continue et \textbf{strictement croissante} sur tout $\R$ :
    \[
        e^a = e^b \iff a = b \qquad \text{et} \qquad e^a < e^b \iff a < b.
    \]
    \item \textbf{Puissance quelconque $a^x$ ($a > 0$) et logarithme en base $a$ ($a > 0, a \neq 1$) :}
    \[
        \boxed{a^x = e^{x\ln(a)}} \qquad \text{et} \qquad \boxed{\log_a(x) = \frac{\ln(x)}{\ln(a)}}.
    \]
\end{enumerate}
\end{theoreme}

\begin{center}
\begin{tikzpicture}[scale=1.05, >=Stealth]
    \begin{axis}[
        width=9.5cm, height=8.5cm,
        axis lines=middle,
        xmin=-2.6, xmax=3.4,
        ymin=-2.6, ymax=3.4,
        xtick={-2,-1,1,2,2.718},
        xticklabels={$-2$,$-1$,$1$,$2$,$e$},
        ytick={-2,-1,1,2,2.718},
        yticklabels={$-2$,$-1$,$1$,$2$,$e$},
        xlabel={$x$},
        ylabel={$y$},
        grid=both,
        grid style={gray!20},
        samples=100,
        clip=true
    ]
        % Première bissectrice y = x
        \addplot[dashed, darkgray, thick, domain=-2.5:3.3] {x} node[pos=0.88, above left] {\small $y = x$};

        % Exponentielle y = e^x
        \addplot[very thick, uclblue, domain=-2.6:1.25] {exp(x)} node[pos=0.92, right] {$y = e^x$};

        % Logarithme y = ln(x)
        \addplot[very thick, defgreen, domain=0.07:3.3] {ln(x)} node[pos=0.88, below right] {$y = \ln(x)$};

        % Tangentes remarquables en (0,1) et (1,0) : y = 1+x et y = x-1
        \addplot[dotted, warnred, thick, domain=-1.5:1.5] {1 + x};
        \addplot[dotted, warnred, thick, domain=-0.5:2.5] {x - 1};

        % Points remarquables
        \fill[uclblue] (axis cs:0,1) circle (2pt) node[above left] {\small $(0,1)$};
        \fill[defgreen] (axis cs:1,0) circle (2pt) node[below right] {\small $(1,0)$};
        \fill[uclblue] (axis cs:1,2.718) circle (2pt) node[left] {\small $(1,e)$};
        \fill[defgreen] (axis cs:2.718,1) circle (2pt) node[above] {\small $(e,1)$};
    \end{axis}
\end{tikzpicture}
\end{center}

\begin{astuce}[Simplifier avant de dériver grâce aux propriétés du logarithme !]
Comme le soulignent les solutionnaires officiels de l'UCLouvain (Juillet 2021 Q1b et Juillet 2022 Q1a), lorsqu'on doit étudier ou dériver une fonction de la forme $f(x) = \ln\left(\frac{1}{u(x)}\right)$ ou $f(x) = \ln\bigl(A(x)\cdot B(x)\bigr)$, \textbf{utilisez toujours les propriétés du logarithme avant de dériver} :
\begin{itemize}
    \item $f(x) = \ln\left(\dfrac{1}{\cos(x)}\right) = -\ln\bigl(\cos(x)\bigr) \implies f'(x) = -\dfrac{-\sin(x)}{\cos(x)} = \tan(x)$ (en une seule ligne !).
    \item $f(x) = \ln\left(\dfrac{1}{1+\cos(x)}\right) = -\ln\bigl(1+\cos(x)\bigr) \implies f'(x) = \dfrac{\sin(x)}{1+\cos(x)}$.
    \item $\ln(1+x) + \ln(1-x) = \ln\bigl((1+x)(1-x)\bigr) = \ln(1-x^2)$ (Juillet 2023 Série 1, Q3), dont la dérivée est immédiatement $\dfrac{-2x}{1-x^2}$.
\end{itemize}
\end{astuce}

\begin{exemple}[Équations exponentielles et logarithmiques imbriquées : UCLouvain -- Juillet 2023 (Série 2, Q3) et Septembre 2024 (Q2a)]
\textbf{1) Double exponentielle (UCLouvain, Juillet 2023 -- Série 2, Question 3) :}\\
\emph{Énoncé :} Soit la fonction $f : \R \to \R$ définie par $f(x) = e^{(e^x)} + a\,e^x + b$ pour deux coefficients $a, b \in \R$. Sous quelles conditions sur $a$ et $b$ la dérivée $f'(x) = e^x(e^{e^x} + a)$ s'annule-t-elle en un point $x^*$ ? Que vaut alors $x^*$ ?\\
\emph{Résolution :} Comme $e^x > 0$ ne s'annule jamais sur $\R$, on a :
\[
    f'(x) = 0 \iff e^{e^x} + a = 0 \iff e^{e^x} = -a.
\]
\begin{itemize}
    \item Première condition : puisque une exponentielle $e^u$ est toujours strictement positive, cette équation n'a de solution en $e^x$ que si $-a > 0 \iff a < 0$. On applique alors $\ln$ aux deux membres :
    \[
        e^{e^x} = -a \iff e^x = \ln(-a).
    \]
    \item Seconde condition : de nouveau, puisque $e^x > 0$, cette équation n'a de solution en $x \in \R$ que si le membre de droite est strictement positif :
    \[
        \ln(-a) > 0 \iff -a > 1 \iff \boxed{a < -1} \quad (\text{et aucune condition sur } b \in \R).
    \]
    En appliquant une seconde fois le logarithme népérien, on obtient l'unique point critique :
    \[
        \boxed{x^* = \ln\bigl(\ln(-a)\bigr)}.
    \]
\end{itemize}

\medskip
\textbf{2) Équation bi-exponentielle en pharmacocinétique (UCLouvain, Septembre 2024, Question 2a) :}\\
\emph{Énoncé :} Résoudre l'équation $\alpha e^{-\alpha t} - \beta e^{-\beta t} = 0$ d'inconnue $t > 0$, où $\alpha = 1,0\text{ h}^{-1}$ et $\beta = 0,2\text{ h}^{-1}$, sachant que $\ln(5) \simeq 1,61$.\\
\emph{Résolution :} Isolons les exponentielles d'un côté et les constantes de l'autre :
\[
    \alpha e^{-\alpha t} = \beta e^{-\beta t} \iff \frac{e^{-\alpha t}}{e^{-\beta t}} = \frac{\beta}{\alpha} \iff e^{(-\alpha + \beta)t} = \frac{\beta}{\alpha}.
\]
En prenant le logarithme népérien des deux membres et en utilisant $\ln\left(\frac{\beta}{\alpha}\right) = -\ln\left(\frac{\alpha}{\beta}\right)$ :
\[
    (-\alpha + \beta)t = \ln\left(\frac{\beta}{\alpha}\right) \iff t = \frac{1}{\alpha - \beta}\ln\left(\frac{\alpha}{\beta}\right) = \frac{1}{1,0 - 0,2}\ln\left(\frac{1,0}{0,2}\right) = \frac{\ln(5)}{0,8} \simeq \frac{1,61}{0,8} = 2,0125\text{ heures}.
\]
\end{exemple}

% ----------------------------------------------------------------------------
\section{Exercices du Chapitre 1}
% ----------------------------------------------------------------------------

\emph{Les corrigés détaillés de tous ces exercices se trouvent au Chapitre 9, en fin d'ouvrage.}

\begin{enumerate}[label=\textbf{Exercice 1.\arabic* --}]
    \item \textbf{(Domaines de définition d'annales UCLouvain)} Déterminer le domaine de définition $\D \subset \R$ de chacune des fonctions suivantes :
    \begin{enumerate}
        \item $f_1(x) = \dfrac{x}{\cos^2(x)}$ \quad \textbf{(Septembre 2023, Q1a)}
        \item $f_2(x) = \dfrac{1 - x^2}{\cos^2(x)}$ \quad \textbf{(Septembre 2024, Q1a)}
        \item $f_3(x) = x\ln^2(x)$ \quad \textbf{(Septembre 2025, Q1a)}
        \item $f_4(x) = (1 - x^2)\ln(1 - x) + x$ \quad \textbf{(Juillet 2026, Q1a)}
    \end{enumerate}

    \item \textbf{(Racines et équations logarithmiques / irrationnelles)} Déterminer le domaine de définition et toutes les racines réelles des fonctions suivantes :
    \begin{enumerate}
        \item $f(x) = \ln(x^2 - 3x + 2)$
        \item $g(x) = \sqrt{2x + 3} - x$
        \item $h(x) = \dfrac{x^2}{2}\left(\ln^2(x) - \ln(x) + \dfrac{1}{2}\right)$ \quad \textbf{(Inspiré de Septembre 2026, Q1)}
    \end{enumerate}

    \item \textbf{(UCLouvain -- Septembre 2024, Q4a)} Soit $x_0$ l'unique racine réelle de la fonction $f(x) = \sqrt{1+x} - x$ sur $[-1, +\infty[$. Sans calculer la valeur numérique de $x_0$, démontrer que $x_0$ est également racine de la fonction composée :
    \[
        g(x) = \sqrt{1 + \sqrt{1+x}} - x.
    \]

    \item \textbf{(Parité et symétries)} Étudier la parité des fonctions suivantes après avoir déterminé leur domaine de définition :
    \begin{enumerate}
        \item $f_1(x) = x^2\sin(x^2) + \cos(x^2)$ \quad \textbf{(Juillet 2023 -- Série 1, Q1a)}
        \item $f_2(x) = 3\cos(x) - \cos^3(x)$ \quad \textbf{(Juillet 2023 -- Série 2, Q1a)}
        \item $f_3(x) = \ln\left(\dfrac{1+x}{1-x}\right)$
        \item $f_4(x) = \ln\bigl(x + \sqrt{1+x^2}\bigr)$
    \end{enumerate}

    \item \textbf{(Simplifications algébriques avec $\ln$ et $\exp$)}
    \begin{enumerate}
        \item Simplifier au maximum l'expression $A(x) = \ln\left(\dfrac{e^{2x} - 1}{e^x + 1}\right) - \ln\bigl(1 - e^{-x}\bigr)$ en précisant sa condition d'existence.
        \item Résoudre dans $\R$ l'équation $e^{2x} - 3e^x - 4 = 0$.
        \item Résoudre dans $\R$ l'équation $\ln(x+3) + \ln(x-1) = \ln(5)$.
    \end{enumerate}

    \item \textbf{(UCLouvain -- Septembre 2026, Q4 -- Premières propriétés de la fonction de Lambert $W_0$)}
    La fonction de Lambert, notée $W_0(y)$, est définie comme la fonction réciproque de la fonction $f(x) = x\,e^x$ restreinte à l'intervalle $I = [-1, +\infty[$. Ainsi, pour tout $x \in [-1, +\infty[$, on a $W_0(y) = x \iff y = x\,e^x$, et donc $W_0(y)\,e^{W_0(y)} = y$.
    \begin{enumerate}
        \item Sans calculatrice, déterminer les valeurs exactes de $W_0(0)$, $W_0(e)$ et $W_0(-e^{-1})$.
        \item Sachant que $f(x) = x\,e^x$ est strictement croissante sur $I = [-1, +\infty[$ et que $\lim_{x \to +\infty} x\,e^x = +\infty$, déterminer le domaine de définition $D$ de $W_0(y)$, son sens de variation et préciser si $W_0$ est continue sur $D$.
    \end{enumerate}
\end{enumerate}

% ============================================================================
% CHAPITRE 2 : LIMITES, CONTINUITÉ ET ASYMPTOTES
% ============================================================================
\chapter{Limites, Continuité et Asymptotes}

La notion de \textbf{limite} est la pierre angulaire de toute l'analyse : c'est elle qui permet de définir la continuité d'une courbe, ses asymptotes aux bords de son domaine, le nombre dérivé (Chapitre 3) et l'intégrale définie (Chapitre 7). À l'examen d'admission de l'UCLouvain, les limites interviennent à trois endroits clés : lors de l'étude de fonction (Question 1, pour le comportement aux bornes et les asymptotes), dans une question spécifique de calcul de limite indéterminée (Questions 3 à 5), et parfois dans une question théorique faisant appel à la définition formelle ou à la continuité.

% ----------------------------------------------------------------------------
\section{Définition et intuition de la limite d'une fonction}
% ----------------------------------------------------------------------------

\subsection{Limites finies et infinies : intuition et définitions rigoureuses}

Intuitivement :
\begin{itemize}
    \item Dire que $\lim_{x \to a} f(x) = L \in \R$ signifie que l'on peut rendre l'image $f(x)$ aussi proche que l'on veut du nombre $L$, à condition de prendre $x$ suffisamment proche de $a$ (avec $x \neq a$).
    \item Dire que $\lim_{x \to +\infty} f(x) = L \in \R$ signifie que l'écart $|f(x) - L|$ devient inférieur à n'importe quelle tolérance $\varepsilon > 0$ fixée à l'avance dès que $x$ dépasse un certain seuil $M$ suffisamment grand.
\end{itemize}

À l'examen d'entrée (comme en \textbf{Juillet 2021, Question 3d}), la définition formelle en $\varepsilon$ et $M$ est rappelée dans l'énoncé lorsqu'il est demandé de l'utiliser, mais savoir la manipuler avec aisance est un atout majeur.

\begin{definition}[Définitions formelles de la limite ($\varepsilon$-$\delta$ et $\varepsilon$-$M$)]
Soit $f : \D \subset \R \to \R$ une fonction réelle :
\begin{enumerate}
    \item \textbf{Limite finie en l'infini ($\lim_{x \to +\infty} f(x) = c \in \R$) :}
    \[
        \boxed{\forall \varepsilon > 0, \;\exists M \in \R \;\text{ tel que }\; \forall x \in \D : \quad x > M \implies \bigl|f(x) - c\bigr| \le \varepsilon}.
    \]
    \item \textbf{Limite finie en un réel ($\lim_{x \to a} f(x) = L \in \R$) :}
    \[
        \boxed{\forall \varepsilon > 0, \;\exists \delta > 0 \;\text{ tel que }\; \forall x \in \D : \quad 0 < |x - a| \le \delta \implies \bigl|f(x) - L\bigr| \le \varepsilon}.
    \]
    \item \textbf{Limite infinie en un réel ($\lim_{x \to a} f(x) = +\infty$) :}
    \[
        \forall A > 0, \;\exists \delta > 0 \;\text{ tel que }\; \forall x \in \D : \quad 0 < |x - a| \le \delta \implies f(x) \ge A.
    \]
\end{enumerate}
\emph{Limite à gauche et limite à droite :} Si l'on impose en outre $x < a$ (resp. $x > a$), on obtient la \textbf{limite à gauche} notée $\lim_{x \to a^-} f(x)$ ou $\lim_{\substack{x \to a \\ x < a}} f(x)$ (resp. la \textbf{limite à droite} notée $\lim_{x \to a^+} f(x)$ ou $\lim_{\substack{x \to a \\ x > a}} f(x)$). Une fonction admet une limite en $a$ si et seulement si ses limites à gauche et à droite en $a$ existent et sont égales.
\end{definition}

\begin{exemple}[Application de la définition $\varepsilon$-$M$ : UCLouvain -- Juillet 2021, Question 3d]
\emph{Énoncé :} Pour rappel, on dit que $\lim_{x \to \infty} f(x) = c$ si pour tout $\varepsilon > 0$, il existe un $M$ tel que $x > M \implies |f(x) - c| \le \varepsilon$. Soit $f$ la fonction définie sur $\R_0^+$ par $f(x) = \dfrac{1}{x}\sin(x^2)$.
\begin{enumerate}[label=(\roman*)]
    \item Sur base de la définition ci-dessus, démontrez que $\lim_{x \to \infty} f(x) = 0$.
    \item \emph{Bonus :} Toujours sur base de la définition, démontrez \textbf{qu'il n'est pas vrai} que $\lim_{x \to \infty} f'(x) = 0$.
\end{enumerate}

\emph{Résolution :}
\begin{enumerate}[label=(\roman*)]
    \item Soit $\varepsilon > 0$ quelconque fixé. Nous cherchons un seuil $M > 0$ (dépendant de $\varepsilon$) tel que pour tout $x > M$, on ait $|f(x) - 0| \le \varepsilon$. Or, puisque $|\sin(u)| \le 1$ pour tout réel $u = x^2$, on a pour tout $x > 0$ :
    \[
        \bigl|f(x) - 0\bigr| = \left|\frac{1}{x}\sin(x^2)\right| = \frac{1}{x}\bigl|\sin(x^2)\bigr| \le \frac{1}{x} \cdot 1 = \frac{1}{x}.
    \]
    Pour que cette quantité soit majorée par $\varepsilon$, il suffit que $\frac{1}{x} \le \varepsilon \iff x \ge \frac{1}{\varepsilon}$. Posons donc $\boxed{M = \dfrac{1}{\varepsilon}}$. Dès lors, pour tout $x > M = \frac{1}{\varepsilon}$, on a bien $|f(x) - 0| \le \frac{1}{x} < \varepsilon$, ce qui démontre par la définition que $\lim_{x \to \infty} f(x) = 0$.
    \item Calculons d'abord la dérivée $f'(x)$ par la règle du produit et de la fonction composée :
    \[
        f'(x) = -\frac{1}{x^2}\sin(x^2) + \frac{1}{x}\cdot 2x\cos(x^2) = 2\cos(x^2) - \frac{\sin(x^2)}{x^2}.
    \]
    Nier la définition « $\forall \varepsilon > 0, \exists M$ tel que $\forall x > M, |f'(x) - 0| \le \varepsilon$ » revient (par les règles de négation des quantificateurs, voir Annexe A) à montrer qu'\textbf{il existe un $\varepsilon > 0$ tel que, quel que soit le seuil $M$, on peut toujours trouver un $x > M$ vérifiant $|f'(x)| > \varepsilon$}.
    Choisissons par exemple $\boxed{\varepsilon = 1}$. Quel que soit le nombre $M > 0$, considérons les points de la forme $x_k = \sqrt{2k\pi}$ avec $k \in \N_0$. En choisissant l'entier $k$ suffisamment grand (à savoir $k > \frac{M^2}{2\pi}$), on a bien $x_k > M$. Or, en ce point $x_k$, comme $x_k^2 = 2k\pi$, on a $\cos(2k\pi) = 1$ et $\sin(2k\pi) = 0$, d'où :
    \[
        f'(x_k) = 2\cos(2k\pi) - \frac{\sin(2k\pi)}{2k\pi} = 2 \cdot 1 - 0 = 2 > \varepsilon = 1.
    \]
    Il est donc impossible de trouver un seuil $M$ au-delà duquel $|f'(x)| \le 1$ : la dérivée $f'(x)$ oscille indéfiniment avec une amplitude proche de $2$ et ne tend pas vers $0$ !
\end{enumerate}
\end{exemple}

% ----------------------------------------------------------------------------
\section{Calcul de limites et levée des formes indéterminées}
% ----------------------------------------------------------------------------

\subsection{Arithmétique des limites et les 7 formes indéterminées}

Lorsque $\lim f(x)$ et $\lim g(x)$ existent, les règles habituelles de somme, produit, quotient et composition s'appliquent, y compris avec $\pm\infty$ lorsque le résultat n'est pas ambigu :
\begin{itemize}
    \item $(+\infty) + (+\infty) = +\infty$, \quad $(+\infty) \cdot (+\infty) = +\infty$, \quad $(+\infty) \cdot (-\infty) = -\infty$.
    \item $\dfrac{c}{\pm\infty} = 0$ (pour toute constante $c \in \R$).
    \item $\dfrac{c}{0^+} = +\infty$ si $c > 0$ (et $-\infty$ si $c < 0$) ; \quad $\dfrac{c}{0^-} = -\infty$ si $c > 0$ (et $+\infty$ si $c < 0$).
\end{itemize}

En revanche, sept combinaisons ne permettent pas de conclure directement et nécessitent une transformation algébrique ou analytique :

\begin{attention}[Les 7 formes indéterminées (F.I.) à reconnaître immédiatement]
\begin{enumerate}
    \item \textbf{Les deux formes indéterminées principales (quotients) :}
    \[
        \boxed{\text{F.I. } \left[\frac{0}{0}\right]} \qquad \text{et} \qquad \boxed{\text{F.I. } \left[\frac{\infty}{\infty}\right]}.
    \]
    \item \textbf{Le produit et la différence indéterminés} (qui se ramènent à $\left[\frac{0}{0}\right]$ ou $\left[\frac{\infty}{\infty}\right]$ par mise au même dénominateur, factorisation ou passage à l'inverse $u \cdot v = \frac{u}{1/v}$) :
    \[
        \boxed{\text{F.I. } \bigl[0 \cdot \infty\bigr]} \qquad \text{et} \qquad \boxed{\text{F.I. } \bigl[(+\infty) - (+\infty)\bigr]}.
    \]
    \item \textbf{Les trois formes indéterminées exponentielles} (qui se lèvent en écrivant $u(x)^{v(x)} = e^{v(x)\ln(u(x))}$, ce qui ramène l'exposant à une forme $[0 \cdot \infty]$) :
    \[
        \boxed{\text{F.I. } \bigl[1^\infty\bigr]}, \qquad \boxed{\text{F.I. } \bigl[0^0\bigr]} \qquad \text{et} \qquad \boxed{\text{F.I. } \bigl[\infty^0\bigr]}.
    \]
\end{enumerate}
\emph{Faux amis (ce ne sont PAS des formes indéterminées !) :} $\frac{0}{\infty} = 0$, $\frac{\infty}{0^+} = \infty$, $0^{+\infty} = 0$, $(+\infty)^{+\infty} = +\infty$, et $(+\infty) + (+\infty) = +\infty$.
\end{attention}

\subsection{Boîte à outils pour lever une indétermination (avant L'Hospital)}

Au Chapitre 3, nous étudierons la célèbre \textbf{règle de L'Hospital} ($\lim \frac{f(x)}{g(x)} = \lim \frac{f'(x)}{g'(x)}$ pour les formes $\frac{0}{0}$ et $\frac{\infty}{\infty}$). Toutefois, comme le rappelle explicitement le solutionnaire de \textbf{Septembre 2022 (Question 5)}, la règle de L'Hospital seule ne suffit pas toujours (parfois elle tourne en boucle ou complique inutilement les expressions !). Il faut donc maîtriser les cinq techniques directes suivantes :

\begin{theoreme}[Les 5 techniques fondamentales de calcul de limites]
\begin{enumerate}
    \item \textbf{Factorisation et simplification en $x \to a$ (forme $\frac{0}{0}$ polynomiale ou rationnelle) :} Si un polynôme $P(x)$ s'annule en $x = a$, il est divisible par $(x - a)$ (loi du reste / tableau de Horner). On factorise $(x-a)$ au numérateur et au dénominateur, puis on simplifie par $(x-a) \neq 0$.
    \item \textbf{Multiplication par le binôme conjugué (racines carrées, formes $\frac{0}{0}$ ou $\infty - \infty$) :} On utilise l'identité remarquable $(\sqrt{A} - \sqrt{B})(\sqrt{A} + \sqrt{B}) = A - B$ :
    \[
        \sqrt{A(x)} - \sqrt{B(x)} = \frac{A(x) - B(x)}{\sqrt{A(x)} + \sqrt{B(x)}}.
    \]
    \item \textbf{Mise en évidence du terme dominant en $x \to \pm\infty$ :}
    \begin{itemize}
        \item En $x \to \pm\infty$, un polynôme se comporte comme son \textbf{monôme de plus haut degré} :
        \[
            \lim_{x \to \pm\infty} \bigl(a_n x^n + a_{n-1}x^{n-1} + \dots + a_0\bigr) = \lim_{x \to \pm\infty} a_n x^n.
        \]
        \item Une fraction rationnelle $\dfrac{a_n x^n + \dots + a_0}{b_m x^m + \dots + b_0}$ a la même limite en $\pm\infty$ que le quotient de ses termes de plus haut degré $\dfrac{a_n x^n}{b_m x^m}$.
        \item Plus généralement, lorsqu'un terme grandit vers l'infini (comme $\ln(x)$) à côté d'un terme borné (comme $\sin(x)$), on met le \textbf{terme dominant en évidence} !
    \end{itemize}
    \item \textbf{Théorème des gendarmes (ou d'encadrement) et produit « borné $\times$ tend vers $0$ » :}
    \begin{itemize}
        \item Si $g(x) \le f(x) \le h(x)$ au voisinage de $a$ (ou de $+\infty$) et si $\lim g(x) = \lim h(x) = L$, alors $\lim f(x) = L$.
        \item En particulier, si $b(x)$ est une fonction \textbf{bornée} (par exemple $|\sin(\dots)| \le 1$ ou $|\cos(\dots)| \le 1$) et si $\lim u(x) = 0$, alors :
        \[
            \boxed{\lim \bigl(u(x) \cdot b(x)\bigr) = 0} \qquad \text{et} \qquad \text{si } \lim v(x) = \pm\infty, \quad \boxed{\lim \frac{b(x)}{v(x)} = 0}.
        \]
    \end{itemize}
    \item \textbf{Limites fondamentales et croissances comparées (puissances, $\ln$, $\exp$, $\sin$) :}
    \begin{itemize}
        \item \textbf{En $u \to 0$ :}
        \begin{gather*}
            \boxed{\lim_{u \to 0} \frac{\sin(u)}{u} = 1}, \qquad \boxed{\lim_{u \to 0} \frac{\tan(u)}{u} = 1}, \qquad \boxed{\lim_{u \to 0} \frac{1 - \cos(u)}{u^2} = \frac{1}{2}}, \\[4pt]
            \boxed{\lim_{u \to 0} \frac{e^u - 1}{u} = 1}, \qquad \boxed{\lim_{u \to 0} \frac{\ln(1+u)}{u} = 1}.
        \end{gather*}
        \item \textbf{Croissances comparées (« l'exponentielle l'emporte sur les puissances, qui l'emportent sur le logarithme ») pour tout $p > 0$ :}
        \[
            \boxed{\lim_{x \to +\infty} \frac{e^x}{x^p} = +\infty}, \;\; \boxed{\lim_{x \to +\infty} x^p e^{-x} = 0}, \;\; \boxed{\lim_{x \to +\infty} \frac{\ln(x)}{x^p} = 0}, \;\; \boxed{\lim_{y \to 0^+} y^p \ln(y) = 0}.
        \]
    \end{itemize}
\end{enumerate}
\end{theoreme}

\begin{exemple}[Questions d'examen UCLouvain sur les limites par factorisation et croissances comparées]
\textbf{1) UCLouvain -- Septembre 2022, Question 5 (Termes bornés vs logarithme) :}\\
\emph{Énoncé :} Calculer la limite suivante :
\[
    L = \lim_{x \to +\infty} \frac{\ln(x) + \sin(x)}{\ln(x^2) + \sin(x^2)}.
\]
\emph{Résolution :} \textbf{Attention au piège de la règle de L'Hospital !} Si vous dérivez numérateur et dénominateur, vous obtenez $\frac{\frac{1}{x} + \cos(x)}{\frac{2}{x} + 2x\cos(x^2)}$, expression qui oscille violemment et n'admet même pas de limite en $+\infty$ !
La bonne méthode consiste à utiliser la propriété $\ln(x^2) = 2\ln(x)$ (pour $x > 0$) et à \textbf{mettre en évidence le terme dominant $\ln(x)$} au numérateur et au dénominateur :
\[
    L = \lim_{x \to +\infty} \frac{\ln(x) + \sin(x)}{2\ln(x) + \sin(x^2)} = \lim_{x \to +\infty} \frac{\cancel{\ln(x)}\left(1 + \dfrac{\sin(x)}{\ln(x)}\right)}{\cancel{\ln(x)}\left(2 + \dfrac{\sin(x^2)}{\ln(x)}\right)} = \lim_{x \to +\infty} \frac{1 + \dfrac{\sin(x)}{\ln(x)}}{2 + \dfrac{\sin(x^2)}{\ln(x)}}.
\]
Or, les fonctions $\sin(x)$ et $\sin(x^2)$ sont bornées entre $-1$ et $1$, tandis que $\lim_{x \to +\infty} \ln(x) = +\infty$. Par le théorème des gendarmes ($-\frac{1}{\ln x} \le \frac{\sin(\dots)}{\ln x} \le \frac{1}{\ln x}$), on a :
\[
    \lim_{x \to +\infty} \frac{\sin(x)}{\ln(x)} = 0 \qquad \text{et} \qquad \lim_{x \to +\infty} \frac{\sin(x^2)}{\ln(x)} = 0 \implies L = \frac{1 + 0}{2 + 0} = \boxed{\frac{1}{2}}.
\]

\medskip
\textbf{2) UCLouvain -- Juillet 2026, Question 1b (et Septembre 2026, Question 1a) :}\\
\emph{Énoncé (Juillet 2026, Q1b) :} Soit $f(x) = (1 - x^2)\ln(1 - x) + x$ sur $\D = \,]-\infty, 1[$. Obtenir $\lim_{x \to 1} f(x)$ en utilisant l'aide $\lim_{y \to 0^+} y\ln(y) = 0$.\\
\emph{Résolution :} Lorsque $x \to 1^-$, $(1-x^2) \to 0$ et $\ln(1-x) \to -\infty$ : c'est une forme indéterminée $[0 \cdot (-\infty)]$. Factorisons $1 - x^2 = (1+x)(1-x)$ et posons le changement de variable $y = 1 - x$ (de sorte que $y \to 0^+$ lorsque $x \to 1^-$) :
\[
    \lim_{x \to 1^-} f(x) = \lim_{x \to 1^-} \bigl[(1+x)(1-x)\ln(1-x) + x\bigr] = 2 \cdot \left(\lim_{y \to 0^+} y\ln(y)\right) + 1 = 2 \cdot 0 + 1 = \boxed{1}.
\]

\medskip
\textbf{3) UCLouvain -- Septembre 2025, Question 1a (et Septembre 2026, Question 1a) :}\\
\emph{Énoncé :} Pour $f(x) = x\ln^2(x)$ définie sur $\D = \,]0, +\infty[$, calculer $\lim_{x \to 0^+} x\ln^2(x)$ à l'aide de $\lim_{y \to 0^+} y^p\ln(y) = 0$ pour $p > 0$.\\
\emph{Résolution :} C'est une forme indéterminée $[0 \cdot (+\infty)]$. Pour faire apparaître un carré parfait du type $\bigl(x^p\ln(x)\bigr)^2$, écrivons $x = (\sqrt{x})^2 = (x^{1/2})^2$ :
\[
    \lim_{x \to 0^+} x\ln^2(x) = \lim_{x \to 0^+} \Bigl(x^{1/2}\ln(x)\Bigr)^2 = \left(\lim_{x \to 0^+} x^{1/2}\ln(x)\right)^2 = 0^2 = \boxed{0} \quad \left(\text{avec } p = \frac{1}{2} > 0\right).
\]
\emph{De même en Septembre 2026 (Q1a)} pour $f(x) = \frac{x^2}{2}\left(\ln^2(x) - \ln(x) + \frac{1}{2}\right)$ : en distribuant $x^2$, on obtient $\frac{1}{2}\bigl(x\ln(x)\bigr)^2 - \frac{1}{2}\bigl(x^2\ln(x)\bigr) + \frac{x^2}{4}$, dont chaque terme tend vers $0$ en $x \to 0^+$ (en appliquant l'aide avec $p=1$ et $p=2$).
\end{exemple}

% ----------------------------------------------------------------------------
\section{Continuité, prolongement et théorème des valeurs intermédiaires}
% ----------------------------------------------------------------------------

\subsection{Continuité en un point et sur un intervalle}

\begin{definition}[Fonction continue en un point $a \in \D$ et sur un intervalle]
Soit $f : \D \subset \R \to \R$ et soit $a \in \D$ :
\begin{enumerate}
    \item On dit que $f$ est \textbf{continue en $a$} si $f$ est définie en $a$ et si sa limite en $a$ existe et est égale à $f(a)$ :
    \[
        \boxed{f \text{ est continue en } a \iff \lim_{x \to a} f(x) = f(a) \iff \lim_{x \to a^-} f(x) = \lim_{x \to a^+} f(x) = f(a)}.
    \]
    \item On dit que $f$ est \textbf{continue sur un intervalle $I \subset \D$} si elle est continue en tout point $a \in I$ (aux bornes éventuelles d'un intervalle fermé $[a, b]$, la continuité s'entend à droite en $a$ et à gauche en $b$).
    \item \textbf{Prolongement par continuité :} Si $f$ n'est pas définie en un point $a \notin \D$, mais que $\lim_{x \to a} f(x) = L \in \R$ existe et est finie, on peut \textbf{prolonger $f$ par continuité en $a$} en posant $\tilde{f}(a) = L$ et $\tilde{f}(x) = f(x)$ pour $x \in \D$.
\end{enumerate}
\end{definition}

Toutes les fonctions usuelles (polynômes, fractions rationnelles, radicaux, $\sin$, $\cos$, $\tan$, $\arcsin$, $\arccos$, $\arctan$, $\ln$, $\exp$) ainsi que leurs sommes, produits, quotients et composées sont \textbf{continues sur tout intervalle inclus dans leur domaine de définition}.

\begin{exemple}[Continuité et discontinuités aux examens UCLouvain (Juillet 2022 Q6 et Juillet 2025 Q4)]
\textbf{1) Paramètre de continuité (UCLouvain -- Juillet 2022, Question 6) :}\\
\emph{Énoncé :} Soit $f : \R \to \R$ définie par $f(2\pi) = a$ et $f(x) = \dfrac{1 - \cos^2(x)}{(e^x - e^{2\pi})^2}$ si $x \neq 2\pi$. Sachant que $f$ est continue sur $\R$, déterminer la valeur de $a$.\\
\emph{Résolution :} Par définition de la continuité en $x = 2\pi$, il faut et il suffit que $a = f(2\pi) = \lim_{x \to 2\pi} f(x)$. Utilisons $\sin^2(x) = 1 - \cos^2(x)$ pour réécrire $f(x)$ comme le carré d'un quotient :
\[
    a = \lim_{x \to 2\pi} \frac{\sin^2(x)}{(e^x - e^{2\pi})^2} = \left(\lim_{x \to 2\pi} \frac{\sin(x)}{e^x - e^{2\pi}}\right)^2.
\]
Posons $u = x - 2\pi \to 0$ (ou appliquons la règle de L'Hospital au quotient intérieur de forme $\frac{0}{0}$) : comme $\sin(2\pi + u) = \sin(u)$ et $e^{2\pi+u} - e^{2\pi} = e^{2\pi}(e^u - 1)$, il vient :
\[
    \lim_{x \to 2\pi} \frac{\sin(x)}{e^x - e^{2\pi}} = \lim_{u \to 0} \left(\frac{1}{e^{2\pi}} \cdot \frac{\sin(u)}{u} \cdot \frac{u}{e^u - 1}\right) = \frac{1}{e^{2\pi}} \cdot 1 \cdot 1 = e^{-2\pi} \implies \boxed{a = \bigl(e^{-2\pi}\bigr)^2 = e^{-4\pi}}.
\]

\medskip
\textbf{2) Discontinuités par « saut » (UCLouvain -- Juillet 2025, Question 4) :}\\
\emph{Énoncé :} Soit $f : \R \to \R$ la fonction qui prend la représentation décimale de $x$, retire $1$ à la première décimale après la virgule si celle-ci est non nulle, et ne la change pas si elle est nulle (par exemple $f(3,14159) = 3,04159$, $f(1,01234) = 1,01234$, $f(-1,25) = -1,15$). Quelles sont les valeurs de $x \in \,]-2, 2[$ où $f$ n'est pas continue ? Sur $]-2, 2[$, $f$ est-elle croissante ?\\
\emph{Résolution :} Analysons l'expression explicite de $f(x)$ sur les sous-intervalles de $]0, 2[$ :
\begin{itemize}
    \item Sur $[0; 0,1[$, la première décimale est $0$, donc elle ne change pas : $f(x) = x$. Lorsque $x \to 0,1^-$, on a $\lim_{x \to 0,1^-} f(x) = 0,1$.
    \item En revanche, sur $[0,1; 1[$, la première décimale est dans $\{1, \dots, 9\}$, donc on lui retire $1$, ce qui revient à soustraire $0,1$ : $f(x) = x - 0,1$. En particulier, $f(0,1) = 0,0 \neq \lim_{x \to 0,1^-} f(x) = 0,1$ ! Il y a donc un \textbf{saut de discontinuité} en $x = 0,1$.
    \item En $x = 1,0$ : pour $x \to 1,0^-$ (par ex. $0,999\dots$), $f(x) = x - 0,1 \to 0,9$, tandis que sur $[1,0; 1,1[$ la première décimale redevient $0$, donc $f(x) = x$ et $f(1,0) = 1,0 \neq 0,9$ : discontinuité en $x = 1,0$, puis de nouveau en $x = 1,1$ (où $f(1,1) = 1,0$ alors que $\lim_{x \to 1,1^-} f(x) = 1,1$).
\end{itemize}
Par symétrie impaire pour les nombres négatifs ($f(-x) = -f(x)$), les $6$ points de discontinuité dans $]-2, 2[$ sont :
\[
    \boxed{x \in \bigl\{-1,1 \;;\; -1,0 \;;\; -0,1 \;;\; 0,1 \;;\; 1,0 \;;\; 1,1\bigr\}}.
\]
De plus, $f$ \textbf{n'est pas croissante} sur $]-2, 2[$ à cause de la chute brutale en $x = 0,1$ : par exemple, pour $x_1 = 0,09 < x_2 = 0,1$, on a $f(0,09) = 0,09 > f(0,1) = 0,0$, ce qui fournit un contre-exemple immédiat.
\end{exemple}

\subsection{Le théorème des valeurs intermédiaires (TVI / Bolzano)}

\begin{theoreme}[Théorème des valeurs intermédiaires (TVI) et théorème de Bolzano]
Soit $f : [a, b] \to \R$ une fonction \textbf{continue} sur l'intervalle fermé borné $[a, b]$.
\begin{enumerate}
    \item \textbf{Énoncé général (TVI) :} Pour tout nombre réel $k$ compris entre $f(a)$ et $f(b)$, il existe \textbf{au moins un} réel $c \in [a, b]$ tel que $f(c) = k$.
    \item \textbf{Cas particulier des racines (Théorème de Bolzano) :} Si $f$ est continue sur $[a, b]$ et si $f(a)$ et $f(b)$ sont de \textbf{signes strictement opposés} ($f(a) \cdot f(b) < 0$), alors l'équation $f(x) = 0$ possède \textbf{au moins une racine} $c \in \,]a, b[$.
    \item \textbf{Unicité (Corollaire de la bijection) :} Si de plus $f$ est \textbf{strictement monotone} (strictement croissante ou strictement décroissante) sur $[a, b]$, alors cette racine $c \in \,]a, b[$ est \textbf{unique}.
\end{enumerate}
\end{theoreme}

\begin{center}
\begin{tikzpicture}[scale=1.05, >=Stealth]
    \draw[->, thick, gray] (-0.5,0) -- (6.2,0) node[right, black] {$x$};
    \draw[->, thick, gray] (0,-1.6) -- (0,2.4) node[above, black] {$y$};

    % Courbe continue changeant de signe
    \draw[very thick, uclblue, domain=0.7:5.3, samples=100] plot (\x, {0.16*(\x-3)^3 + 0.35*(\x-3) + 0.25}) node[right] {$y = f(x)$};

    % Point a (f(a) < 0)
    \coordinate (A) at (1.0, -1.73);
    \draw[dashed, warnred] (1.0, 0) node[above, black] {$a$} -- (1.0, -1.73) -- (0, -1.73) node[left] {$f(a) < 0$};
    \fill[warnred] (1.0, -1.73) circle (2pt);

    % Point b (f(b) > 0)
    \coordinate (B) at (5.0, 2.23);
    \draw[dashed, defgreen] (5.0, 0) node[below, black] {$b$} -- (5.0, 2.23) -- (0, 2.23) node[left] {$f(b) > 0$};
    \fill[defgreen] (5.0, 2.23) circle (2pt);

    % Racine c
    \fill[exorange] (2.435, 0) circle (2.4pt) node[above left] {$c \in \,]a,b[$};
    \node[below right, exorange, font=\small\bfseries] at (2.435, -0.05) {$f(c) = 0$};
\end{tikzpicture}
\end{center}

\begin{astuce}[Quand utiliser le TVI à l'examen d'entrée ?]
Le théorème des valeurs intermédiaires est l'outil idéal dans trois situations récurrentes à l'examen :
\begin{enumerate}
    \item \textbf{En Question 1 (Étude de fonction) :} Lorsqu'on ne peut pas résoudre algébriquement $f(x) = 0$ ou $f'(x) = 0$ (par exemple pour $f'(x) = \frac{1 + 2x\tan x}{\cos^2 x}$ en Septembre 2023 Q1c, ou $g(x) = 2(\tan(x)(1-x^2)-x)$ en Septembre 2024 Q1c, ou $f(x) = (1-x^2)\ln(1-x)+x$ en Juillet 2026 Q1e), le constat que $f$ (ou $f'$) est continue sur un intervalle $[\alpha, \beta]$ sans asymptote et \textbf{change de signe} entre $\alpha$ et $\beta$ permet d'affirmer rigoureusement l'existence d'une racine ou d'un extremum dans $]\alpha, \beta[$.
    \item \textbf{Appliqué à la dérivée $f'$ :} Pour prouver que $f'$ possède au moins deux racines distinctes (comme en Septembre 2022 Q6 ou Septembre 2025 Q4), on peut calculer $f'$ en trois points $a < b < c$ et montrer que $f'(a) > 0$, $f'(b) < 0$ et $f'(c) > 0$ : par continuité de $f'$, le TVI garantit une première racine de $f'$ dans $]a, b[$ et une seconde dans $]b, c[$.
    \item \textbf{Pour prouver qu'une fonction continue qui ne s'annule jamais garde un signe constant} (Septembre 2024 Q5b).
\end{enumerate}
\end{astuce}

% ----------------------------------------------------------------------------
\section{Les asymptotes au graphe d'une fonction (AV, AH, AO)}
% ----------------------------------------------------------------------------

Une \textbf{asymptote} est une droite dont la courbe représentative de $f$ s'approche indéfiniment lorsque $x$ tend vers un bord fini de son domaine ou lorsque $x$ tend vers $\pm\infty$. On distingue trois types d'asymptotes.

\begin{theoreme}[Recherche systématique des asymptotes (AV, AH, AO)]
Soit $f : \D \subset \R \to \R$ une fonction réelle :
\begin{enumerate}
    \item \textbf{Asymptote verticale (AV) d'équation $\boxed{x = a}$ :}
    On cherche les asymptotes verticales aux \textbf{points finis $a$ exclus du domaine $\D$} (points d'annulation d'un dénominateur ou d'un logarithme) ou aux bornes ouvertes de $\D$. La droite verticale d'équation $x = a$ est une asymptote verticale au graphe de $f$ si \textbf{au moins l'une} des limites à gauche ou à droite de $f$ en $a$ est infinie :
    \[
        \lim_{x \to a^-} f(x) = \pm\infty \qquad \text{ou} \qquad \lim_{x \to a^+} f(x) = \pm\infty.
    \]
    \item \textbf{Asymptote horizontale (AH) d'équation $\boxed{y = b}$ en $+\infty$ (et/ou en $-\infty$) :}
    Lorsque le domaine $\D$ s'étend jusqu'à $+\infty$ (ou $-\infty$), la droite horizontale d'équation $y = b$ ($b \in \R$) est asymptote horizontale au graphe de $f$ en $+\infty$ (resp. en $-\infty$) si :
    \[
        \lim_{x \to +\infty} f(x) = b \in \R \qquad \bigl(\text{resp. } \lim_{x \to -\infty} f(x) = b \in \R\bigr).
    \]
    \item \textbf{Asymptote oblique (AO) d'équation $\boxed{y = mx + p}$ (avec $m \neq 0$) en $+\infty$ (et/ou en $-\infty$) :}
    Si $\lim_{x \to +\infty} f(x) = \pm\infty$ (il n'y a donc pas d'AH en $+\infty$), on cherche une éventuelle asymptote oblique $y = mx + p$ en $+\infty$ grâce aux \textbf{formules de Cauchy} en deux étapes :
    \[
        \boxed{\text{Étape 1 : } m = \lim_{x \to +\infty} \frac{f(x)}{x}} \; (m \in \R_0), \quad \boxed{\text{Étape 2 : } p = \lim_{x \to +\infty} \bigl(f(x) - mx\bigr)} \; (p \in \R).
    \]
    (On procède de même séparément en $x \to -\infty$ si le domaine s'étend vers $-\infty$).
\end{enumerate}
\end{theoreme}

\begin{attention}[Trois remarques cruciales sur les asymptotes à l'examen]
\begin{enumerate}
    \item \textbf{Une fonction périodique non constante n'admet JAMAIS d'asymptote horizontale ni d'asymptote oblique !} (UCLouvain -- Juillet 2021, Q1e). En effet, ses oscillations se répètent à l'identique à chaque période $T$, de sorte qu'elle ne peut jamais converger vers une droite en $\pm\infty$. En revanche, elle peut posséder \textbf{une infinité d'asymptotes verticales} périodiquement espacées (comme $x = \frac{\pi}{2} + k\pi$, $k \in \Z$).
    \item \textbf{Incompatibilité AH / AO d'un même côté :} En $+\infty$, le graphe d'une fonction ne peut pas avoir à la fois une AH et une AO (car l'AH correspond au cas $m=0$ dans $y = mx+p$). Par contre, une fonction peut très bien avoir une AH en $-\infty$ et une AO en $+\infty$ (par exemple $f(x) = x + \sqrt{x^2+1}$).
    \item \textbf{N'écrivez jamais « $\text{AV} = 0$ » ou « $\text{AO} = x$ » !} Une asymptote est une \textbf{droite} : sa réponse doit toujours être donnée sous la forme d'une \textbf{équation cartésienne de droite} : $\text{AV} \equiv x = a$, $\text{AH} \equiv y = b$, $\text{AO} \equiv y = mx + p$.
\end{enumerate}
\end{attention}

\begin{exemple}[Recherche complète des asymptotes : UCLouvain -- Septembre 2021, Question 1d]
\emph{Énoncé :} Soit la fonction $f(x) = x + 2\ln\left(1 + \dfrac{1}{x}\right)$, dont le domaine de définition est $\D = \,]-\infty, -1[ \;\cup\; ]0, +\infty[$. Donner l'équation de chacune des asymptotes éventuelles de $f$.

\emph{Résolution :}
\begin{enumerate}
    \item \textbf{Asymptotes verticales (aux bornes finies ouvertes $x = -1$ et $x = 0$ de $\D$) :}
    \begin{itemize}
        \item En $x \to 0^+$ (à droite de $0$, puisque $f$ n'est pas définie sur $[-1, 0]$) : $\frac{1}{x} \to +\infty$, donc $1 + \frac{1}{x} \to +\infty$ et $\ln\left(1+\frac{1}{x}\right) \to +\infty$, d'où :
        \[
            \lim_{x \to 0^+} \left(x + 2\ln\left(1 + \frac{1}{x}\right)\right) = 0 + 2(+\infty) = +\infty \implies \boxed{\text{AV}_1 \equiv x = 0}.
        \]
        \item En $x \to -1^-$ (à gauche de $-1$, c'est-à-dire $x < -1$, donc $\frac{1}{x} > -1$ et $1 + \frac{1}{x} \to 0^+$) : $\ln\left(1+\frac{1}{x}\right) \to \ln(0^+) = -\infty$, d'où :
        \[
            \lim_{x \to -1^-} \left(x + 2\ln\left(1 + \frac{1}{x}\right)\right) = -1 + 2(-\infty) = -\infty \implies \boxed{\text{AV}_2 \equiv x = -1}.
        \]
    \end{itemize}
    \item \textbf{Asymptotes horizontales ou obliques en $x \to +\infty$ et en $x \to -\infty$ :}
    Lorsque $x \to \pm\infty$, on a $\frac{1}{x} \to 0$, donc $\ln\left(1 + \frac{1}{x}\right) \to \ln(1) = 0$. Dès lors :
    \begin{itemize}
        \item $\lim_{x \to \pm\infty} f(x) = \lim_{x \to \pm\infty} (x + 0) = \pm\infty$ : il n'y a \textbf{pas d'asymptote horizontale}.
        \item \emph{Astuce directe pour l'asymptote oblique :} Comme $f(x) = x + 2\ln\left(1+\frac{1}{x}\right)$ est déjà écrite sous la forme $\underbrace{m x + p}_{= x} + \underbrace{\varepsilon(x)}_{= 2\ln(1+1/x)}$ où $\lim_{x \to \pm\infty} \varepsilon(x) = 2\ln(1) = 0$, on lit instantanément l'équation de l'asymptote oblique $y = x$ !
        \item \emph{Vérification par les formules de Cauchy :}
        \[
            m = \lim_{x \to \pm\infty} \frac{f(x)}{x} = \lim_{x \to \pm\infty} \left(1 + \frac{2\ln(1 + 1/x)}{x}\right) = 1 + \frac{0}{\pm\infty} = 1,
        \]
        \[
            p = \lim_{x \to \pm\infty} \bigl(f(x) - 1\cdot x\bigr) = \lim_{x \to \pm\infty} 2\ln\left(1 + \frac{1}{x}\right) = 2\ln(1) = 0.
        \]
        La fonction admet donc la même asymptote oblique en $+\infty$ et en $-\infty$ : $\boxed{\text{AO} \equiv y = x \quad (\text{en } +\infty \text{ et en } -\infty)}$.
    \end{itemize}
\end{enumerate}
\end{exemple}

% ----------------------------------------------------------------------------
\section{Exercices du Chapitre 2}
% ----------------------------------------------------------------------------

\emph{Les corrigés détaillés de tous ces exercices se trouvent au Chapitre 9, en fin d'ouvrage.}

\begin{enumerate}[label=\textbf{Exercice 2.\arabic* --}]
    \item \textbf{(Calculs de limites sans la règle de L'Hospital)} Calculer les limites suivantes en utilisant les factorisations, binômes conjugués, termes dominants ou limites fondamentales :
    \begin{enumerate}
        \item $\displaystyle \lim_{x \to 2} \frac{x^3 - 8}{x^2 - 3x + 2}$
        \item $\displaystyle \lim_{x \to 0} \frac{\sqrt{1 + 2x} - \sqrt{1 - 2x}}{x}$
        \item $\displaystyle \lim_{x \to +\infty} \left(\sqrt{x^2 + 4x + 1} - x\right)$
        \item $\displaystyle \lim_{x \to +\infty} \frac{3\ln(x^3) - 2\cos(x^2)}{2\ln(x) + 5\sin(x)}$
    \end{enumerate}

    \item \textbf{(Limites trigonométriques et exponentielles -- Annales UCLouvain)} Calculer les limites suivantes :
    \begin{enumerate}
        \item $\displaystyle \lim_{x \to 0} \frac{x\bigl(\exp(\sin(x)) - 1\bigr)}{\sin^2(x)}$ \quad \textbf{(Juillet 2021, Q3e)}
        \item $\displaystyle \lim_{x \to 0} \frac{\sin(x^2)}{x\tan(x)}$ \quad \textbf{(Juillet 2025, Q5)}
        \item $\displaystyle \lim_{x \to 0^+} x\,\ln\bigl(\sin(x)\bigr)$
    \end{enumerate}

    \item \textbf{(Définition formelle de la limite $\varepsilon$-$M$)} En utilisant la définition formelle $\forall \varepsilon > 0, \exists M > 0$ tel que $x > M \implies |f(x) - c| \le \varepsilon$, démontrer que :
    \begin{enumerate}
        \item $\displaystyle \lim_{x \to +\infty} \frac{2x + \cos(x)}{x} = 2$
        \item $\displaystyle \lim_{x \to +\infty} \frac{\ln(x) + 3}{\ln(x)} = 1$
    \end{enumerate}

    \item \textbf{(Continuité et prolongement par continuité)} Déterminer la valeur des paramètres réels $a$ et $b$ pour que la fonction $f : \R \to \R$ définie par :
    \[
        f(x) = \begin{cases}
            \dfrac{\sin(3x)}{x} & \text{si } x < 0, \\
            a & \text{si } x = 0, \\
            \sqrt{1 + bx} + 2 & \text{si } x > 0,
        \end{cases}
    \]
    soit continue en $x = 0$.

    \item \textbf{(Recherche complète d'asymptotes)} Déterminer le domaine de définition et les équations de toutes les asymptotes (AV, AH, AO) des fonctions suivantes :
    \begin{enumerate}
        \item $f_1(x) = \dfrac{2x^2 - 3x + 5}{x - 2}$
        \item $f_2(x) = \sqrt{x^2 + 2x}$
        \item $f_3(x) = \dfrac{e^x - 1}{e^x + 1}$
    \end{enumerate}

    \item \textbf{(Application du théorème des valeurs intermédiaires)} Démontrer que l'équation $\cos(x) = x$ possède une et une seule solution réelle $x^* \in \R$, et que cette solution appartient à l'intervalle $\left]0, \dfrac{\pi}{2}\right[$.
\end{enumerate}

% ============================================================================
% CHAPITRE 3 : DÉRIVÉES ET GRANDS THÉORÈMES D'ANALYSE
% ============================================================================
\chapter{Dérivées et Grands Théorèmes d'Analyse}

La \textbf{dérivée} est l'outil central du calcul différentiel : géométriquement, elle donne la pente de la droite tangente à une courbe en chaque point ; physiquement, elle mesure la vitesse instantanée de variation d'une grandeur ; analytiquement, son signe gouverne la croissance, la décroissance et les extrema d'une fonction. Ce chapitre couvre la définition du nombre dérivé, la dérivabilité en un point, toutes les règles de dérivation (y compris pour les fonctions composées et réciproques comme la fonction de Lambert $W_0$), la règle de L'Hospital, ainsi que les deux grands théorèmes théoriques systématiquement présents aux examens de l'UCLouvain : le \textbf{théorème de Rolle} et le \textbf{théorème des accroissements finis (TAF)}.

% ----------------------------------------------------------------------------
\section{Nombre dérivé, tangente et dérivabilité en un point}
% ----------------------------------------------------------------------------

\subsection{Définition par la limite du taux d'accroissement}

\begin{definition}[Nombre dérivé en un point $a$ et fonction dérivée $f'$]
Soit $f : \D \subset \R \to \R$ une fonction définie au voisinage d'un point $a \in \D$. On dit que $f$ est \textbf{dérivable en $a$} si la limite du \textbf{taux d'accroissement} de $f$ entre $a$ et $x$ existe et est finie lorsque $x \to a$ :
\[
    \boxed{f'(a) = \lim_{x \to a} \frac{f(x) - f(a)}{x - a} = \lim_{h \to 0} \frac{f(a+h) - f(a)}{h} \in \R}.
\]
\begin{itemize}
    \item Ce nombre réel $f'(a)$, noté aussi $\dfrac{\mathrm{d}f}{\mathrm{d}x}(a)$ ou $\left.\dfrac{\mathrm{d}f}{\mathrm{d}x}\right|_{x=a}$, est appelé le \textbf{nombre dérivé} de $f$ en $a$.
    \item La fonction $f' : \D' \subset \D \to \R, x \mapsto f'(x)$ qui à tout point $x$ où $f$ est dérivable associe $f'(x)$ est la \textbf{fonction dérivée première} de $f$.
    \item Si $f'$ est elle-même dérivable, sa dérivée $(f')' = f'' = \dfrac{\mathrm{d}^2 f}{\mathrm{d}x^2}$ est la \textbf{dérivée seconde} de $f$. Une fonction dérivable à tout ordre (et dont toutes les dérivées sont continues) est dite \textbf{de classe $C^\infty$}, tandis qu'une fonction $k$ fois continûment dérivable est dite \textbf{de classe $C^k$} (par exemple $C^1$ si $f'$ est continue, $C^2$ si $f''$ est continue).
\end{itemize}
\end{definition}

\begin{theoreme}[Interprétation géométrique : équation de la droite tangente]
Si $f$ est dérivable en $x_0 \in \D$, alors le graphe de $f$ admet au point $P\bigl(x_0, f(x_0)\bigr)$ une \textbf{droite tangente} non verticale $T_{x_0}$ dont le coefficient directeur (la pente) vaut $m = f'(x_0)$ et dont l'équation cartésienne est :
\[
    \boxed{T_{x_0} \equiv y = f(x_0) + f'(x_0)\,(x - x_0)}.
\]
\begin{itemize}
    \item Si $f'(x_0) = 0$, la tangente en $x_0$ est \textbf{horizontale} ($y = f(x_0)$).
    \item Si $\lim_{x \to x_0} \frac{f(x)-f(x_0)}{x-x_0} = \pm\infty$ (avec $f$ continue en $x_0$), la fonction n'est \textbf{pas dérivable} en $x_0$, mais son graphe admet une \textbf{tangente verticale} d'équation $x = x_0$ (c'est le cas de $\sqrt{x}$ en $x=0$, de $\sqrt{1+x}-x$ en $x=-1$ en Juillet 2024 Q1e, ou de la fonction de Lambert $W_0(y)$ en $y = -e^{-1}$ en Septembre 2026 Q4).
    \item Si la limite à gauche et la limite à droite du taux d'accroissement en $x_0$ sont finies mais différentes ($f'_g(x_0) \neq f'_d(x_0)$), la fonction n'est \textbf{pas dérivable} en $x_0$ et le graphe présente un \textbf{point anguleux} (demi-tangentes de pentes différentes, comme pour $|x|$ en $x=0$).
\end{itemize}
\end{theoreme}

\begin{attention}[Lien fondamental entre dérivabilité et continuité]
\textbf{Toute fonction dérivable en un point $a$ est nécessairement continue en $a$ :}
\[
    \boxed{f \text{ dérivable en } a \implies f \text{ continue en } a} \qquad \text{et par contraposée :} \qquad \boxed{f \text{ non continue en } a \implies f \text{ non dérivable en } a}.
\]
\textbf{Attention : la réciproque est FAUSSE !} Une fonction peut être parfaitement continue en un point $a$ sans y être dérivable. Les deux contre-exemples classiques à connaître par cœur (notamment pour tester les hypothèses du théorème de Rolle, comme en \textbf{Juillet 2023 -- Série 1, Question 5}) sont :
\begin{itemize}
    \item $x \mapsto |x|$ en $x = 0$ (continue en $0$, mais point anguleux : pente $-1$ à gauche et $+1$ à droite).
    \item $x \mapsto \sqrt{x}$ en $x = 0^+$, ou $x \mapsto \sqrt[3]{x}$ en $x = 0$, ou $x \mapsto 1 - \sqrt{|x|}$ en $x = 0$ (continues en $0$, mais pente infinie / tangente verticale en $0$).
\end{itemize}
\end{attention}

\begin{exemple}[Tangente et dérivabilité en un point aux examens UCLouvain (Septembre 2021, Q3a et Q3d)]
\textbf{1) Identification de paramètres par la tangente (UCLouvain -- Septembre 2021, Question 3a) :}\\
\emph{Énoncé :} Soit la fonction $f : \R \to \R$ définie par $f(x) = \dfrac{a}{1 + e^{-bx}}$. Déterminer les valeurs de $a$ et $b$ pour que $f(0) = \dfrac{1}{2}$ et que la tangente à $f$ en $x = 0$ passe par le point $(x, y) = (2, 1)$.\\
\emph{Résolution :}
\begin{itemize}
    \item En $x = 0$, comme $e^0 = 1$, on a $f(0) = \frac{a}{1+1} = \frac{a}{2}$. La condition $f(0) = \frac{1}{2}$ impose immédiatement $\boxed{a = 1}$.
    \item L'équation de la tangente au graphe de $f$ en $x_0 = 0$ est $y = f(0) + f'(0)(x - 0) = \frac{1}{2} + f'(0)\,x$. Dire que cette droite passe par le point $(2, 1)$ signifie que $1 = \frac{1}{2} + 2f'(0) \iff f'(0) = \frac{1}{4}$.
    \item Calculons la dérivée de $f(x) = (1 + e^{-bx})^{-1}$ pour $a = 1$ :
    \[
        f'(x) = -\frac{(-b\,e^{-bx})}{(1 + e^{-bx})^2} = \frac{b\,e^{-bx}}{(1 + e^{-bx})^2} \implies f'(0) = \frac{b \cdot 1}{(1+1)^2} = \frac{b}{4}.
    \]
    L'égalité $f'(0) = \frac{b}{4} = \frac{1}{4}$ donne donc $\boxed{b = 1}$.
\end{itemize}

\medskip
\textbf{2) Étude complète de la régularité en $0$ d'une fonction prolongée (UCLouvain -- Septembre 2021, Question 3d) :}\\
\emph{Énoncé :} Soit $f : \R \to \R$ définie par $f(x) = \dfrac{\sin(x)}{x}$ si $x \neq 0$ et $f(0) = 1$. Démontrer que : (i) $f$ est continue en $0$, (ii) $f$ est dérivable en $0$, et (iii) la fonction dérivée $f'$ est continue en $0$.\\
\emph{Résolution :}
\begin{enumerate}[label=(\roman*)]
    \item \textbf{Continuité de $f$ en $0$ :} Par la limite trigonométrique fondamentale, $\lim_{x \to 0} f(x) = \lim_{x \to 0} \frac{\sin(x)}{x} = 1 = f(0)$. Donc $f$ est continue en $0$.
    \item \textbf{Dérivabilité de $f$ en $0$ :} Pour savoir si une fonction définie par morceaux est dérivable au point de raccord $x = 0$, on \textbf{revient à la définition du nombre dérivé} par la limite du taux d'accroissement :
    \[
        f'(0) = \lim_{x \to 0} \frac{f(x) - f(0)}{x - 0} = \lim_{x \to 0} \frac{\dfrac{\sin(x)}{x} - 1}{x} = \lim_{x \to 0} \frac{\sin(x) - x}{x^2} \quad \left[\text{F.I. } \frac{0}{0}\right].
    \]
    En appliquant deux fois la règle de L'Hospital (voir Section 3.3 ci-dessous) :
    \[
        \lim_{x \to 0} \frac{\sin(x) - x}{x^2} \stackrel{\text{H}}{=} \lim_{x \to 0} \frac{\cos(x) - 1}{2x} \quad \left[\frac{0}{0}\right] \stackrel{\text{H}}{=} \lim_{x \to 0} \frac{-\sin(x)}{2} = \frac{0}{2} = 0.
    \]
    La limite existe et est finie, donc $f$ est dérivable en $0$ et $\boxed{f'(0) = 0}$.
    \item \textbf{Continuité de $f'$ en $0$ :} Pour tout $x \neq 0$, la règle du quotient donne :
    \[
        f'(x) = \frac{x\cos(x) - \sin(x)}{x^2}.
    \]
    Pour que $f'$ soit continue en $0$, il faut que $\lim_{x \to 0} f'(x) = f'(0) = 0$. Calculons cette limite (de forme $\frac{0}{0}$) par la règle de L'Hospital :
    \[
        \lim_{x \to 0} f'(x) = \lim_{x \to 0} \frac{x\cos(x) - \sin(x)}{x^2} \stackrel{\text{H}}{=} \lim_{x \to 0} \frac{\bigl(\cos(x) - x\sin(x)\bigr) - \cos(x)}{2x} = \lim_{x \to 0} \frac{-\cancel{x}\sin(x)}{2\cancel{x}} = \lim_{x \to 0} \frac{-\sin(x)}{2} = 0 = f'(0).
    \]
    La dérivée $f'$ est donc bien continue en $0$ (on dit que $f$ est au moins de classe $C^1$ sur $\R$).
\end{enumerate}
\end{exemple}

% ----------------------------------------------------------------------------
\section{Formulaire et règles de dérivation}
% ----------------------------------------------------------------------------

\subsection{Règles opératoires, fonctions composées et fonctions réciproques}

\begin{theoreme}[Les 5 règles fondamentales de dérivation]
Soient $u$ et $v$ deux fonctions dérivables et $c \in \R$ une constante :
\begin{enumerate}
    \item \textbf{Linéarité :} $\bigl(c\cdot u(x)\bigr)' = c\cdot u'(x)$ \quad et \quad $\bigl(u(x) \pm v(x)\bigr)' = u'(x) \pm v'(x)$.
    \item \textbf{Dérivée d'un produit :}
    \[
        \boxed{\bigl(u(x) \cdot v(x)\bigr)' = u'(x)\,v(x) + u(x)\,v'(x)}.
    \]
    \item \textbf{Dérivée d'un inverse et d'un quotient (où $v(x) \neq 0$) :}
    \[
        \boxed{\left(\frac{1}{v(x)}\right)' = -\frac{v'(x)}{\bigl(v(x)\bigr)^2}} \qquad \text{et} \qquad \boxed{\left(\frac{u(x)}{v(x)}\right)' = \frac{u'(x)\,v(x) - u(x)\,v'(x)}{\bigl(v(x)\bigr)^2}}.
    \]
    \item \textbf{Dérivée d'une fonction composée (règle de la chaîne / \emph{Chain Rule}) :}
    \[
        \boxed{\Bigl(f\bigl(u(x)\bigr)\Bigr)' = f'\bigl(u(x)\bigr) \cdot u'(x)}.
    \]
    \item \textbf{Dérivée d'une fonction réciproque $g = f^{-1}$ :} Si $f$ est bijective et dérivable en $x = g(y)$ avec $f'(g(y)) \neq 0$, alors en dérivant l'identité $f(g(y)) = y$ par rapport à $y$, on obtient $f'(g(y))\cdot g'(y) = 1$, soit :
    \[
        \boxed{\bigl(f^{-1}\bigr)'(y) = g'(y) = \frac{1}{f'\bigl(g(y)\bigr)} = \frac{1}{f'\bigl(f^{-1}(y)\bigr)}}.
    \]
\end{enumerate}
\end{theoreme}

\subsection{Tableau complet des dérivées usuelles et composées}

\begin{theoreme}[Tableau des dérivées des fonctions élémentaires et composées]
\begin{center}
\renewcommand{\arraystretch}{1.6}
\begin{tabular}{@{}lll@{}}
\toprule
\textbf{Fonction simple $f(x)$} & \textbf{Dérivée $f'(x)$} & \textbf{Version composée $\bigl(f(u(x))\bigr)'$} \\
\midrule
$c$ (constante) & $0$ & $0$ \\
$x^n \quad (n \in \R)$ & $n\,x^{n-1}$ & $\bigl(u(x)^n\bigr)' = n\,\bigl(u(x)\bigr)^{n-1} \cdot u'(x)$ \\
$\sqrt{x} = x^{1/2} \quad (x > 0)$ & $\dfrac{1}{2\sqrt{x}}$ & $\bigl(\sqrt{u(x)}\bigr)' = \dfrac{u'(x)}{2\sqrt{u(x)}}$ \\
$\dfrac{1}{x} = x^{-1} \quad (x \neq 0)$ & $-\dfrac{1}{x^2}$ & $\left(\dfrac{1}{u(x)}\right)' = -\dfrac{u'(x)}{\bigl(u(x)\bigr)^2}$ \\
$e^x$ & $e^x$ & $\bigl(e^{u(x)}\bigr)' = u'(x)\,e^{u(x)}$ \\
$a^x = e^{x\ln a} \quad (a > 0)$ & $(\ln a)\,a^x$ & $\bigl(a^{u(x)}\bigr)' = (\ln a)\,u'(x)\,a^{u(x)}$ \\
$\ln(x)$ (ou $\ln(|x|)$) & $\dfrac{1}{x}$ & $\Bigl(\ln\bigl(|u(x)|\bigr)\Bigr)' = \dfrac{u'(x)}{u(x)}$ \\
$\sin(x)$ & $\cos(x)$ & $\bigl(\sin(u(x))\bigr)' = u'(x)\,\cos\bigl(u(x)\bigr)$ \\
$\cos(x)$ & $-\sin(x)$ & $\bigl(\cos(u(x))\bigr)' = -u'(x)\,\sin\bigl(u(x)\bigr)$ \\
$\tan(x)$ & $\dfrac{1}{\cos^2(x)} = 1 + \tan^2(x)$ & $\bigl(\tan(u(x))\bigr)' = \dfrac{u'(x)}{\cos^2(u(x))} = u'(x)\bigl(1+\tan^2(u(x))\bigr)$ \\
$\cot(x)$ & $-\dfrac{1}{\sin^2(x)} = -\bigl(1 + \cot^2(x)\bigr)$ & $\bigl(\cot(u(x))\bigr)' = -\dfrac{u'(x)}{\sin^2(u(x))}$ \\
$\arcsin(x) \quad (|x| < 1)$ & $\dfrac{1}{\sqrt{1 - x^2}}$ & $\bigl(\arcsin(u(x))\bigr)' = \dfrac{u'(x)}{\sqrt{1 - u(x)^2}}$ \\
$\arccos(x) \quad (|x| < 1)$ & $-\dfrac{1}{\sqrt{1 - x^2}}$ & $\bigl(\arccos(u(x))\bigr)' = -\dfrac{u'(x)}{\sqrt{1 - u(x)^2}}$ \\
$\arctan(x) \quad (x \in \R)$ & $\dfrac{1}{1 + x^2}$ & $\bigl(\arctan(u(x))\bigr)' = \dfrac{u'(x)}{1 + u(x)^2}$ \\
\bottomrule
\end{tabular}
\end{center}
\end{theoreme}

\begin{astuce}[Comment retrouver en 10 secondes les dérivées de $\arcsin$, $\arccos$ et $\arctan$ ?]
Si vous hésitez le jour de l'examen sur la formule de $(\arcsin x)'$ ou $(\arctan x)'$, appliquez la méthode de dérivation implicite :
\begin{itemize}
    \item Pour $y(x) = \arcsin(x) \in \left]-\frac{\pi}{2}, \frac{\pi}{2}\right[$ : on a $\sin(y(x)) = x$. Dérivons les deux membres par rapport à $x$ :
    \[
        \cos\bigl(y(x)\bigr) \cdot y'(x) = 1 \implies y'(x) = \frac{1}{\cos(\arcsin x)} = \frac{1}{\sqrt{1 - \sin^2(\arcsin x)}} = \frac{1}{\sqrt{1 - x^2}}.
    \]
    \item Pour $y(x) = \arctan(x)$ : on a $\tan(y(x)) = x$. En dérivant par rapport à $x$ :
    \[
        \bigl(1 + \tan^2(y(x))\bigr)\cdot y'(x) = 1 \implies y'(x) = \frac{1}{1 + \tan^2(\arctan x)} = \frac{1}{1 + x^2}.
    \]
\end{itemize}
C'est exactement cette méthode qui a été demandée en \textbf{Septembre 2026 (Question 4)} pour calculer la dérivée de la fonction de Lambert $W_0(y)$ !
\end{astuce}

\begin{exemple}[Dérivation de fonctions réciproques et d'équations fonctionnelles aux examens UCLouvain]
\textbf{1) Dérivée de la fonction de Lambert $W_0(y)$ (UCLouvain -- Septembre 2026, Question 4) :}\\
\emph{Énoncé :} La fonction de Lambert $W_0(y)$, définie sur $D = [-e^{-1}, +\infty[$, est la réciproque de $f(x) = x\,e^x$ sur $I = [-1, +\infty[$, et vérifie donc l'identité $W_0(y)\,e^{W_0(y)} = y$ (avec $W_0(0) = 0$, $W_0(e) = 1$ et $W_0(-e^{-1}) = -1$). Obtenir la dérivée $W'_0(y)$. Que valent $W'_0(0)$ et $W'_0(e)$ ? La dérivée de $W_0(y)$ existe-t-elle en $y = -e^{-1}$ ? Pourquoi ?\\
\emph{Résolution :} Présentons les deux méthodes officielles du solutionnaire :
\begin{itemize}
    \item \emph{Méthode 1 (Dérivation implicite de l'identité $W_0(y)\,e^{W_0(y)} = y$) :} En dérivant les deux membres par rapport à $y$ (par la règle du produit et de la composée) :
    \[
        W'_0(y)\,e^{W_0(y)} + W_0(y)\cdot W'_0(y)\,e^{W_0(y)} = 1 \iff W'_0(y)\,e^{W_0(y)}\bigl(1 + W_0(y)\bigr) = 1,
    \]
    d'où, pour tout $y > -e^{-1}$ (c'est-à-dire $W_0(y) > -1$) :
    \[
        \boxed{W'_0(y) = \frac{e^{-W_0(y)}}{1 + W_0(y)} = \bigl(1 + W_0(y)\bigr)^{-1} e^{-W_0(y)}} \quad \left(\text{ou encore } \frac{1}{e^{W_0(y)} + y} = \frac{W_0(y)}{y(1+W_0(y))} \text{ pour } y \neq 0\right).
    \]
    \item \emph{Méthode 2 (Formule de la dérivée d'une fonction réciproque) :} Comme $f(x) = x\,e^x$, on a $f'(x) = (1+x)e^x$. La formule $g'(y) = \frac{1}{f'(g(y))}$ appliquée à $g = W_0$ donne directement :
    \[
        W'_0(y) = \frac{1}{f'\bigl(W_0(y)\bigr)} = \frac{1}{\bigl(1 + W_0(y)\bigr)\,e^{W_0(y)}} = \bigl(1 + W_0(y)\bigr)^{-1} e^{-W_0(y)}.
    \]
    \item \emph{Évaluations et non-dérivabilité au bord $y = -e^{-1}$ :}
    \begin{itemize}
        \item En $y = 0$, $W_0(0) = 0 \implies W'_0(0) = (1+0)^{-1}e^{-0} = \boxed{1}$.
        \item En $y = e$, $W_0(e) = 1 \implies W'_0(e) = (1+1)^{-1}e^{-1} = \boxed{\dfrac{1}{2e}}$.
        \item En $y = -e^{-1}$, on a $W_0(-e^{-1}) = -1$, ce qui \textbf{annule le dénominateur} $1 + W_0(y) = 0$ (géométriquement, la fonction directe $f(x) = x\,e^x$ a une tangente horizontale $f'(-1) = 0$ en $x = -1$, donc par symétrie par rapport à $y = x$, la fonction réciproque $W_0(y)$ a une \textbf{tangente verticale} de pente infinie en $y = -e^{-1}$). La dérivée $W'_0(-e^{-1})$ \textbf{n'existe donc pas}.
    \end{itemize}
\end{itemize}

\medskip
\textbf{2) Fonctions liées et équations différentielles par dérivation composée (UCLouvain -- Juillet 2022 Q5, Septembre 2023 Q4 et Juillet 2024 Q4) :}\\
Ces trois superbes questions d'examen reposent toutes sur le même principe : \textbf{reconnaître la dérivée d'une fonction composée pour montrer qu'une combinaison est constante ou linéaire}.
\begin{itemize}
    \item \textbf{Juillet 2022, Question 5 :} Soit $f : \R^+ \to \R_0^+$ dérivable telle que $f'(x) = -\bigl(f(x)\bigr)^2$. Que vaut $\left(\dfrac{1}{f(x)}\right)'$ ? En déduire $f(x)$.\\
    \emph{Résolution :} Par la règle de dérivation d'un inverse :
    \[
        \left(\frac{1}{f(x)}\right)' = -\frac{f'(x)}{\bigl(f(x)\bigr)^2} = -\frac{-\bigl(f(x)\bigr)^2}{\bigl(f(x)\bigr)^2} = 1.
    \]
    Puisque la dérivée de $\frac{1}{f(x)}$ vaut la constante $1$, on a en intégrant $\frac{1}{f(x)} = x + C$, soit $\boxed{f(x) = \dfrac{1}{x + C}}$ (avec $C > 0$ pour que $f$ soit définie et strictement positive sur tout $\R^+$).
    \item \textbf{Septembre 2023, Question 4 :} Soit $f : I \to \R$ deux fois continûment dérivable telle que (i) $f''(x) = -f(x)$ et (ii) $f'(x) \ge 0$.
    \begin{enumerate}[label=\alph*)]
        \item Prouver que $f^2(x) + f'^2(x)$ est une constante positive notée $k^2$.\\
        \emph{Résolution :} Dérivons la somme de carrés $f^2(x) + f'^2(x)$ par la règle de la chaîne $\bigl(u^2(x)\bigr)' = 2u(x)u'(x)$ :
        \[
            \Bigl(f^2(x) + f'^2(x)\Bigr)' = 2f(x)f'(x) + 2f'(x)f''(x) = 2f(x)f'(x) + 2f'(x)\bigl(-f(x)\bigr) = 0.
        \]
        Sa dérivée étant identiquement nulle sur l'intervalle $I$, la fonction $f^2(x) + f'^2(x)$ est constante sur $I$, et cette constante est positive ou nulle comme somme de deux carrés : $f^2(x) + f'^2(x) = k^2$.
        \item On pose $g(x) = \arcsin\left(\dfrac{f(x)}{k}\right)$ (pour $k > 0$). Calculer et simplifier $g'(x)$, puis en déduire $f(x)$.\\
        \emph{Résolution :} Par la dérivée composée de $\arcsin(u(x))$ :
        \[
            g'(x) = \frac{1}{\sqrt{1 - \left(\dfrac{f(x)}{k}\right)^2}} \cdot \frac{f'(x)}{k} = \frac{f'(x)}{\sqrt{k^2 - f^2(x)}} = \frac{f'(x)}{\sqrt{f'^2(x)}} = \frac{f'(x)}{f'(x)} = \boxed{1} \quad \bigl(\text{car } f'(x) \ge 0 \implies \sqrt{f'^2(x)} = f'(x)\bigr).
        \]
        Dès lors, $g'(x) = 1 \implies g(x) = \arcsin\left(\frac{f(x)}{k}\right) = x + C \implies \boxed{f(x) = k\sin(x + C)}$.
    \end{enumerate}
    \item \textbf{Juillet 2024, Question 4 :} Soit $f(x) > 0$ dérivable telle que $f(0) = 1$, et soit $g(x)$ vérifiant $\dfrac{g'(x)}{f(x)} = g(x)f'(x)$ avec $g(0) = 1$. Exprimer $g(x)$ en fonction de $f(x)$ et montrer que $g(x) > C$ pour une constante $C > 0$.\\
    \emph{Résolution :} En multipliant par $f(x)$ et en divisant par $g(x)$ (qui est $>0$ au voisinage de $0$ car $g(0) = 1$), on sépare les fonctions $g$ et $f$ :
    \[
        \frac{g'(x)}{g(x)} = f(x)f'(x) \iff \Bigl(\ln\bigl(g(x)\bigr)\Bigr)' = \left(\frac{1}{2}f^2(x)\right)'.
    \]
    En intégrant entre $0$ et $x$ avec $f(0) = 1$ et $g(0) = 1$ ($\ln(g(0)) = \ln(1) = 0$), il vient :
    \[
        \ln\bigl(g(x)\bigr) - \ln\bigl(g(0)\bigr) = \frac{1}{2}\bigl(f^2(x) - f^2(0)\bigr) \iff \ln\bigl(g(x)\bigr) = \frac{1}{2}\bigl(f^2(x) - 1\bigr) \implies \boxed{g(x) = e^{\frac{1}{2}(f^2(x) - 1)}}.
    \]
    Comme $f(x) > 0 \implies f^2(x) > 0$, la croissance stricte de l'exponentielle donne $g(x) > e^{\frac{1}{2}(0-1)} = \boxed{e^{-1/2} = \dfrac{1}{\sqrt{e}}}$ (donc $C = e^{-1/2}$).
\end{itemize}
\end{exemple}

% ----------------------------------------------------------------------------
\section{La règle de L'Hospital pour les limites indéterminées}
% ----------------------------------------------------------------------------

La \textbf{règle de L'Hospital} (ou de Bernoulli-L'Hospital) est l'outil privilégié à l'examen d'entrée pour lever rapidement les formes indéterminées de type quotient $\left[\dfrac{0}{0}\right]$ ou $\left[\dfrac{\infty}{\infty}\right]$.

\begin{theoreme}[Théorème de Bernoulli -- Règle de L'Hospital]
Soient $f$ et $g$ deux fonctions dérivables sur un voisinage épointé de $a$ (où $a \in \R$ ou $a = \pm\infty$), telles que $g'(x) \neq 0$ sur ce voisinage. Si l'on est en présence d'une \textbf{forme indéterminée} de type :
\[
    \lim_{x \to a} f(x) = \lim_{x \to a} g(x) = 0 \quad \left[\text{F.I. } \frac{0}{0}\right] \qquad \text{ou} \qquad \lim_{x \to a} |f(x)| = \lim_{x \to a} |g(x)| = +\infty \quad \left[\text{F.I. } \frac{\infty}{\infty}\right],
\]
et si la limite du quotient des dérivées $\lim_{x \to a} \dfrac{f'(x)}{g'(x)}$ existe (finie ou égale à $\pm\infty$), alors :
\[
    \boxed{\lim_{x \to a} \frac{f(x)}{g(x)} \stackrel{\text{H}}{=} \lim_{x \to a} \frac{f'(x)}{g'(x)}}.
\]
\emph{Itération :} Si le nouveau quotient $\frac{f'(x)}{g'(x)}$ présente à son tour une forme indéterminée $\left[\frac{0}{0}\right]$ ou $\left[\frac{\infty}{\infty}\right]$, on peut réappliquer la règle de L'Hospital aux dérivées secondes $\lim_{x \to a} \frac{f''(x)}{g''(x)}$, et ainsi de suite.
\end{theoreme}

\begin{attention}[Les 4 règles de prudence avec la règle de L'Hospital]
\begin{enumerate}
    \item \textbf{Vérifiez TOUJOURS qu'il s'agit bien d'une forme $\frac{0}{0}$ ou $\frac{\infty}{\infty}$ avant d'appliquer L'Hospital !} Appliquer L'Hospital à un quotient qui n'est pas indéterminé donne un résultat totalement faux.
    \item \textbf{On dérive le numérateur ET le dénominateur séparément :} C'est $\dfrac{f'(x)}{g'(x)}$, et \textbf{surtout pas} la dérivée du quotient $\left(\frac{f}{g}\right)' = \frac{f'g-fg'}{g^2}$ !
    \item \textbf{Simplifiez algébriquement entre deux applications de L'Hospital !} Avant de réappliquer L'Hospital une deuxième fois, nettoyez l'expression $\frac{f'(x)}{g'(x)}$ (mettez au même dénominateur, remplacez $\lim \cos(x) = 1$, utilisez les propriétés de $\ln$...) : cela évite une explosion de dérivées de produits ou de quotients.
    \item \textbf{Transformez un produit $[0 \cdot \infty]$ en quotient avant L'Hospital :} Pour calculer $\lim_{x \to 0^+} x\ln^2(x)$ (Septembre 2025 Q1a) par L'Hospital, on fait passer $x$ au dénominateur sous la forme $1/x$ (jamais $\ln(x)$ au dénominateur !) :
    \[
        \lim_{x \to 0^+} x\ln^2(x) = \lim_{x \to 0^+} \frac{\ln^2(x)}{1/x} \quad \left[\frac{\infty}{\infty}\right] \stackrel{\text{H}}{=} \lim_{x \to 0^+} \frac{2\ln(x)\cdot \frac{1}{x}}{-1/x^2} = \lim_{x \to 0^+} \frac{-2\ln(x)}{1/x} \quad \left[\frac{\infty}{\infty}\right] \stackrel{\text{H}}{=} \lim_{x \to 0^+} \frac{-2/x}{-1/x^2} = \lim_{x \to 0^+} 2x = 0.
    \]
\end{enumerate}
\end{attention}

\begin{exemple}[Trois limites d'examen résolues par la règle de L'Hospital (Juillet 2023 Q3, Septembre 2024 Q4b, Juillet 2025 Q5)]
\textbf{1) UCLouvain -- Juillet 2023 (Série 1), Question 3 :}\\
\emph{Énoncé :} Calculer la limite suivante :
\[
    L = \lim_{x \to 0} \frac{\ln(1+x) + \ln(1-x)}{1 - \cos(x)}.
\]
\emph{Résolution :} En $x = 0$, le numérateur vaut $\ln(1) + \ln(1) = 0$ et le dénominateur vaut $1 - \cos(0) = 0$ : c'est une forme indéterminée $\left[\frac{0}{0}\right]$.
Simplifions d'abord le numérateur par la propriété du logarithme d'un produit : $\ln(1+x) + \ln(1-x) = \ln\bigl((1+x)(1-x)\bigr) = \ln(1-x^2)$. Appliquons alors une première fois la règle de L'Hospital :
\[
    L = \lim_{x \to 0} \frac{\ln(1-x^2)}{1 - \cos(x)} \stackrel{\text{H}}{=} \lim_{x \to 0} \frac{\dfrac{-2x}{1-x^2}}{\sin(x)} = \lim_{x \to 0} \left(\frac{-2}{1-x^2} \cdot \frac{x}{\sin(x)}\right).
\]
À ce stade, soit on conclut immédiatement grâce à la limite fondamentale $\lim_{x \to 0} \frac{x}{\sin(x)} = 1$, soit on réapplique une seconde fois la règle de L'Hospital à $\lim_{x \to 0} \frac{-2x}{(1-x^2)\sin(x)}$ (forme $\frac{0}{0}$) :
\[
    L \stackrel{\text{H}}{=} \lim_{x \to 0} \frac{-2}{-2x\sin(x) + (1-x^2)\cos(x)} = \frac{-2}{0 + 1 \cdot 1} = \boxed{-2}.
\]

\medskip
\textbf{2) UCLouvain -- Septembre 2024, Question 4b (Limite en une racine implicite $x_0$) :}\\
\emph{Énoncé :} Soit $x_0$ la racine de $f(x) = \sqrt{1+x} - x$ (donc $\sqrt{1+x_0} = x_0 > 0$). On a montré à l'Exercice 1.3 que $x_0$ est aussi racine de $g(x) = \sqrt{1 + \sqrt{1+x}} - x$. Soit $h(x) = \dfrac{f(x)}{g(x)}$. Obtenir $\lim_{x \to x_0} h(x)$ exprimée uniquement en fonction de $x_0$ (sans aucune racine carrée) et simplifiée au maximum.\\
\emph{Résolution :} Puisque $f(x_0) = 0$ et $g(x_0) = 0$, la limite $\lim_{x \to x_0} \frac{f(x)}{g(x)}$ est une forme indéterminée $\left[\frac{0}{0}\right]$. Par la règle de L'Hospital :
\[
    \lim_{x \to x_0} h(x) = \lim_{x \to x_0} \frac{f(x)}{g(x)} \stackrel{\text{H}}{=} \frac{f'(x_0)}{g'(x_0)}.
\]
Calculons les dérivées $f'(x)$ et $g'(x)$ pour tout $x > -1$ :
\[
    f'(x) = \frac{1}{2\sqrt{1+x}} - 1 \qquad \text{et} \qquad g'(x) = \frac{1}{2\sqrt{1 + \sqrt{1+x}}} \cdot \frac{1}{2\sqrt{1+x}} - 1 = \frac{1}{4\sqrt{1+\sqrt{1+x}}\,\sqrt{1+x}} - 1.
\]
Or, en $x = x_0$, par définition de la racine $x_0$, on a $\sqrt{1+x_0} = x_0$ et $\sqrt{1 + \sqrt{1+x_0}} = \sqrt{1+x_0} = x_0$ ! En remplaçant chaque radical en $x_0$ par $x_0$, les deux dérivées se simplifient merveilleusement :
\[
    f'(x_0) = \frac{1}{2x_0} - 1 = \frac{1 - 2x_0}{2x_0} \qquad \text{et} \qquad g'(x_0) = \frac{1}{4x_0^2} - 1 = \frac{1 - 4x_0^2}{4x_0^2} = \frac{(1 - 2x_0)(1 + 2x_0)}{4x_0^2}.
\]
En effectuant le quotient et en simplifiant par $(1 - 2x_0) \neq 0$ et par $2x_0 \neq 0$ :
\[
    \lim_{x \to x_0} h(x) = \frac{f'(x_0)}{g'(x_0)} = \frac{1 - 2x_0}{2x_0} \cdot \frac{4x_0^2}{(1 - 2x_0)(1 + 2x_0)} = \boxed{\frac{2x_0}{1 + 2x_0}}.
\]

\medskip
\textbf{3) UCLouvain -- Juillet 2025, Question 5 (Comparaison L'Hospital vs décomposition) :}\\
\emph{Énoncé :} Calculer $\displaystyle \lim_{x \to 0} \frac{\sin(x^2)}{x\tan(x)}$.\\
\emph{Résolution :}
\begin{itemize}
    \item \emph{Méthode 1 (par la règle de L'Hospital deux fois) :} C'est une forme $\left[\frac{0}{0}\right]$.
    \[
        \lim_{x \to 0} \frac{\sin(x^2)}{x\tan(x)} \stackrel{\text{H}}{=} \lim_{x \to 0} \frac{2x\cos(x^2)}{\tan(x) + \dfrac{x}{\cos^2(x)}} \quad \left[\frac{0}{0}\right] \stackrel{\text{H}}{=} \lim_{x \to 0} \frac{2\cos(x^2) - 4x^2\sin(x^2)}{\dfrac{2}{\cos^2(x)} + \dfrac{2x\sin(x)}{\cos^3(x)}} = \frac{2 \cdot 1 - 0}{\dfrac{2}{1} + 0} = \boxed{1}.
    \]
    \item \emph{Méthode 2 (beaucoup plus rapide : décomposition en limites fondamentales ou L'Hospital séparé) :} En multipliant et divisant par $x$ :
    \[
        \lim_{x \to 0} \frac{\sin(x^2)}{x\tan(x)} = \left(\lim_{x \to 0} \frac{\sin(x^2)}{x^2}\right) \cdot \left(\lim_{x \to 0} \frac{x}{\tan(x)}\right) = 1 \cdot 1 = \boxed{1}.
    \]
\end{itemize}
\end{exemple}

% ----------------------------------------------------------------------------
\section{Théorème de Rolle et Théorème des Accroissements Finis (TAF)}
% ----------------------------------------------------------------------------

Presque chaque session de l'examen d'admission d'Analyse à l'UCLouvain comporte une question dédiée au \textbf{théorème de Rolle} ou au \textbf{théorème des accroissements finis (TAF)} (Septembre 2021 Q3f, Septembre 2022 Q6, Juillet 2023 Série 1 Q5, Juillet 2024 Q5, Septembre 2025 Q4). Ce sont des questions très rentables dès lors qu'on maîtrise leurs trois hypothèses et leurs schémas d'application.

\subsection{Le théorème de Rolle}

\begin{theoreme}[Théorème de Rolle]
Soient $a$ et $b$ deux nombres réels tels que $a < b$, et soit $f : [a, b] \to \R$ une fonction vérifiant les \textbf{trois hypothèses} suivantes :
\begin{enumerate}
    \item $f$ est \textbf{continue sur l'intervalle fermé $[a, b]$} ;
    \item $f$ est \textbf{dérivable sur l'intervalle ouvert $]a, b[$} ;
    \item $f$ prend la même valeur aux deux extrémités : $\boxed{f(a) = f(b)}$.
\end{enumerate}
Alors il existe \textbf{au moins un} réel $c \in \,]a, b[$ (strictement compris entre $a$ et $b$) tel que :
\[
    \boxed{f'(c) = 0}.
\]
\emph{Interprétation géométrique :} Si une courbe continue et lisse (sans point anguleux ni rupture) part d'une certaine hauteur $f(a)$ et revient à la même hauteur $f(b) = f(a)$, elle possède nécessairement au moins un point intérieur $c \in \,]a, b[$ où la tangente est horizontale.
\end{theoreme}

\subsection{Le théorème des accroissements finis (TAF / Lagrange)}

En inclinant la figure du théorème de Rolle lorsque $f(a) \neq f(b)$, la corde reliant $A\bigl(a, f(a)\bigr)$ à $B\bigl(b, f(b)\bigr)$ a pour pente le taux d'accroissement moyen $\dfrac{f(b) - f(a)}{b - a}$. Le théorème des accroissements finis affirme qu'il existe au moins un point intérieur $c \in \,]a, b[$ où la tangente est parallèle à la corde $[AB]$.

\begin{theoreme}[Théorème des Accroissements Finis (TAF)]
Soient $a$ et $b$ deux nombres réels tels que $a < b$, et soit $f : [a, b] \to \R$ une fonction vérifiant les \textbf{deux hypothèses} suivantes :
\begin{enumerate}
    \item $f$ est \textbf{continue sur l'intervalle fermé $[a, b]$} ;
    \item $f$ est \textbf{dérivable sur l'intervalle ouvert $]a, b[$}.
\end{enumerate}
Alors il existe \textbf{au moins un} réel $c \in \,]a, b[$ tel que :
\[
    \boxed{f'(c) = \frac{f(b) - f(a)}{b - a}} \qquad \Longleftrightarrow \qquad \boxed{f(b) - f(a) = f'(c)\,(b - a)}.
\]
\emph{Interprétation cinématique :} Si une voiture parcourt $120\text{ km}$ en $1\text{ heure}$ (vitesse moyenne de $120\text{ km/h}$), son compteur de vitesse instantanée $v(t) = x'(t)$ a nécessairement affiché exactement $120\text{ km/h}$ à au moins un instant $c$ durant le trajet !
\end{theoreme}

\begin{center}
\begin{tikzpicture}[scale=1.0, >=Stealth]
    % Théorème de Rolle
    \begin{scope}[shift={(0,0)}]
        \draw[->, thick, gray] (-0.4,0) -- (5.2,0) node[right, black] {$x$};
        \draw[->, thick, gray] (0,-0.4) -- (0,3.2) node[above, black] {$y$};
        \draw[very thick, uclblue, domain=0.8:4.4, samples=80] plot (\x, {0.8 + 1.8*sin(deg((\x-0.8)*3.14159/3.6))});
        \draw[dashed, darkgray] (0,0.8) node[left] {\small $f(a)=f(b)$} -- (4.4,0.8);
        \draw[dashed, gray] (0.8,0) node[below, black] {$a$} -- (0.8,0.8);
        \draw[dashed, gray] (4.4,0) node[below, black] {$b$} -- (4.4,0.8);
        \draw[dashed, warnred] (2.6,0) node[below, warnred] {$c$} -- (2.6,2.6);
        \draw[thick, warnred, <->] (1.6,2.6) -- (3.6,2.6) node[right] {\small $f'(c)=0$};
        \fill[uclblue] (0.8,0.8) circle (2pt) node[above left] {\small $A$};
        \fill[uclblue] (4.4,0.8) circle (2pt) node[above right] {\small $B$};
        \fill[warnred] (2.6,2.6) circle (2pt);
        \node[above, uclblue, font=\bfseries\small] at (2.6,3.2) {Théorème de Rolle};
    \end{scope}

    % Théorème des accroissements finis
    \begin{scope}[shift={(7.2,0)}]
        \draw[->, thick, gray] (-0.4,0) -- (5.2,0) node[right, black] {$x$};
        \draw[->, thick, gray] (0,-0.4) -- (0,3.2) node[above, black] {$y$};
        \draw[very thick, uclblue, domain=0.8:4.4, samples=80] plot (\x, {0.5 + 0.45*(\x-0.8) + 1.2*sin(deg((\x-0.8)*3.14159/3.6))});
        % Corde AB
        \draw[thick, defgreen, dashed] (0.8,0.5) -- (4.4,2.12) node[pos=0.55, below right=-2pt] {\footnotesize Corde $[AB]$};
        % Tangente parallèle en c=2.6
        \draw[thick, warnred, <->] (1.5, 2.015) -- (3.7, 3.005) node[right] {\footnotesize $f'(c) = \frac{f(b)-f(a)}{b-a}$};
        \draw[dashed, gray] (0.8,0) node[below, black] {$a$} -- (0.8,0.5);
        \draw[dashed, gray] (4.4,0) node[below, black] {$b$} -- (4.4,2.12);
        \draw[dashed, warnred] (2.6,0) node[below, warnred] {$c$} -- (2.6,2.51);
        \fill[defgreen] (0.8,0.5) circle (2pt) node[above left] {\small $A$};
        \fill[defgreen] (4.4,2.12) circle (2pt) node[below right] {\small $B$};
        \fill[warnred] (2.6,2.51) circle (2pt);
        \node[above, uclblue, font=\bfseries\small] at (2.6,3.2) {Théorème des Accroissements Finis (TAF)};
    \end{scope}
\end{tikzpicture}
\end{center}

\begin{astuce}[Les 3 stratégies récurrentes avec Rolle et le TAF aux examens UCLouvain]
\begin{enumerate}
    \item \textbf{Vérifier minutieusement les 3 hypothèses de Rolle (Juillet 2023 -- Série 1, Q5) :} Si la fonction n'est pas définie en un seul point de $[a, b]$ (ex: $g(x) = x - \frac{1}{x}$ sur $[-1, 1]$, non définie en $0$), ou n'est pas dérivable en un seul point de $]a, b[$ (ex: $f(x) = 1 - \sqrt{|x|}$ sur $[-1, 1]$, non dérivable en $0$), ou si $f(a) \neq f(b)$, \textbf{le théorème de Rolle ne s'applique pas} !
    \item \textbf{Découper un intervalle en sous-intervalles disjoints $[x_i, x_{i+1}]$ :} Pour prouver que $f'$ s'annule (ou vaut une constante $k$) en \textbf{plusieurs points distincts} $c_1, c_2, \dots$, on applique Rolle ou le TAF sur des intervalles consécutifs $[a, b]$ et $[b, c]$ dont les intérieurs $]a, b[$ et $]b, c[$ sont \textbf{disjoints}, ce qui garantit automatiquement que $c_1 \in \,]a, b[$ et $c_2 \in \,]b, c[$ vérifient $c_1 < b < c_2$, donc $c_1 \neq c_2$ !
    \item \textbf{Enchaîner en cascade Rolle sur $f$, puis Rolle sur $f'$ pour atteindre $f''(d) = 0$ :} Dès que l'on a trouvé deux points distincts $c_1 < c_2$ où la dérivée première prend la même valeur ($f'(c_1) = f'(c_2)$), on applique une seconde fois le théorème de Rolle (ou le TAF) \textbf{à la fonction dérivée $f'$} sur l'intervalle $[c_1, c_2]$ pour conclure qu'il existe un point $c_4 \in \,]c_1, c_2[$ tel que $(f')'(c_4) = f''(c_4) = 0$ !
\end{enumerate}
\end{astuce}

\begin{exemple}[Quatre questions complètes d'examen UCLouvain sur Rolle et le TAF (2021--2025)]
\textbf{1) UCLouvain -- Septembre 2021, Question 3f (Cascade du TAF sur $f$ puis sur $f'$) :}\\
\emph{Énoncé :} Soit $f : \R^+ \to \R^+$ une fonction de classe $C^\infty$ telle que $f(n) = n$ pour tout entier $n \ge 0$. Démontrer qu'alors, pour tout entier $p \ge 2$, l'intervalle $[0, p]$ contient : (i) au moins $p$ points $x_i$ distincts pour lesquels $f'(x_i) = 1$, et (ii) au moins $p - 1$ points $y_i$ distincts pour lesquels $f''(y_i) = 0$.\\
\emph{Résolution :}
\begin{enumerate}[label=(\roman*)]
    \item Subdivisons $[0, p]$ en $p$ sous-intervalles de longueur $1$ : $[i, i+1]$ pour $i \in \{0, 1, \dots, p-1\}$. Sur chaque intervalle $[i, i+1]$, $f$ est continue et dérivable (car de classe $C^\infty$). Le théorème des accroissements finis appliqué à $f$ sur $[i, i+1]$ fournit l'existence d'un réel $x_i \in \,]i, i+1[$ tel que :
    \[
        f'(x_i) = \frac{f(i+1) - f(i)}{(i+1) - i} = \frac{(i+1) - i}{1} = 1.
    \]
    Comme les $p$ intervalles ouverts $]0, 1[,\; ]1, 2[,\; \dots,\; ]p-1, p[$ sont deux à deux disjoints, les $p$ points obtenus vérifient $0 < x_0 < 1 < x_1 < 2 < \dots < x_{p-1} < p$ et sont donc bien tous distincts !
    \item Considérons maintenant les $p-1$ intervalles $[x_i, x_{i+1}]$ pour $i \in \{0, 1, \dots, p-2\}$. La fonction dérivée $f'$ est elle-même continue sur $[x_i, x_{i+1}]$, dérivable sur $]x_i, x_{i+1}[$, et prend la même valeur aux bornes : $f'(x_i) = f'(x_{i+1}) = 1$. En appliquant le théorème de Rolle (ou le TAF) \textbf{à la fonction $f'$} sur chaque $[x_i, x_{i+1}]$, il existe un point $y_i \in \,]x_i, x_{i+1}[$ tel que :
    \[
        f''(y_i) = \frac{f'(x_{i+1}) - f'(x_i)}{x_{i+1} - x_i} = \frac{1 - 1}{x_{i+1} - x_i} = 0.
    \]
    Les $p-1$ intervalles ouverts $]x_i, x_{i+1}[$ étant disjoints, les $p-1$ points $y_0 < y_1 < \dots < y_{p-2}$ sont tous distincts dans $[0, p]$.
\end{enumerate}

\medskip
\textbf{2) UCLouvain -- Septembre 2022, Question 6 (Trois méthodes pour les racines de $f'$) :}\\
\emph{Énoncé :} Soient $a < b < c$ trois nombres réels, et soit $f : \R \to \R$ définie par $f(x) = (x-a)(x-b)(x-c)$. Prouver que la fonction dérivée $f'$ possède au moins deux racines distinctes.\\
\emph{Résolution :}
\begin{itemize}
    \item \emph{Méthode 1 (la plus élégante : théorème de Rolle) :} $f$ est un polynôme, donc continue et dérivable sur tout $\R$, et elle s'annule en $a$, $b$ et $c$ : $f(a) = f(b) = f(c) = 0$.
    \begin{itemize}
        \item Sur $[a, b]$, comme $f(a) = f(b) = 0$, le théorème de Rolle garantit l'existence de $r_1 \in \,]a, b[$ tel que $f'(r_1) = 0$.
        \item Sur $[b, c]$, comme $f(b) = f(c) = 0$, le théorème de Rolle garantit l'existence de $r_2 \in \,]b, c[$ tel que $f'(r_2) = 0$.
    \end{itemize}
    Puisque $a < r_1 < b < r_2 < c$, on a $r_1 \neq r_2$ : $f'$ possède bien deux racines distinctes.
    \item \emph{Méthode 2 (par le TVI sur $f'$) :} Par la règle du produit, $f'(x) = (x-b)(x-c) + (x-a)(x-c) + (x-a)(x-b)$. Évaluons le signe de $f'$ en $a, b, c$ (sachant que $a < b < c$) :
    \[
        f'(a) = \underbrace{(a-b)}_{<0}\underbrace{(a-c)}_{<0} > 0, \qquad f'(b) = \underbrace{(b-a)}_{>0}\underbrace{(b-c)}_{<0} < 0, \qquad f'(c) = \underbrace{(c-a)}_{>0}\underbrace{(c-b)}_{>0} > 0.
    \]
    La fonction continue $f'$ change de signe sur $[a, b]$ et sur $[b, c]$ : par le TVI, elle s'annule dans $]a, b[$ et dans $]b, c[$.
    \item \emph{Méthode 3 (algébrique par le discriminant $\Delta$) :} En développant $f'(x) = 3x^2 - 2(a+b+c)x + (ab+bc+ac)$, son discriminant vaut $\Delta = 4(a+b+c)^2 - 12(ab+bc+ac) = 2\bigl[(a-b)^2 + (b-c)^2 + (a-c)^2\bigr] > 0$ car $a < b < c$.
\end{itemize}

\medskip
\textbf{3) UCLouvain -- Juillet 2024, Question 5 (Égalité des carrés $f^2(a) = f^2(b) = f^2(c)$) :}\\
\emph{Énoncé :} Soient $a < b < c$ trois nombres réels et $f : [a, c] \to \R$ une fonction continue sur $[a, c]$ et dérivable sur $]a, c[$ telle que $f^2(a) = f^2(b) = f^2(c)$. En partant du théorème de Rolle, prouver qu'il existe au moins un réel $d \in \,]a, c[$ tel que $f'(d) = 0$.\\
\emph{Résolution (2 démonstrations officielles) :}
\begin{itemize}
    \item \emph{Démonstration 1 (Principe des tiroirs sur les signes) :} Puisque $f^2(a) = f^2(b) = f^2(c)$, les trois nombres $f(a), f(b), f(c)$ ont la même valeur absolue $|f(a)| = |f(b)| = |f(c)|$. Comme il n'y a que deux signes possibles ($+$ ou $-$), parmi les trois nombres $f(a), f(b), f(c)$, il y en a nécessairement \textbf{au moins deux qui ont le même signe} et qui sont donc \textbf{égaux} :
    \begin{itemize}
        \item Si $f(a) = f(b)$, Rolle appliqué à $f$ sur $[a, b]$ donne $d \in \,]a, b[ \;\subset \,]a, c[$ tel que $f'(d) = 0$.
        \item Si $f(b) = f(c)$, Rolle appliqué à $f$ sur $[b, c]$ donne $d \in \,]b, c[ \;\subset \,]a, c[$ tel que $f'(d) = 0$.
        \item Si $f(a) = f(c)$, Rolle appliqué à $f$ sur $[a, c]$ donne directement $d \in \,]a, c[$ tel que $f'(d) = 0$.
    \end{itemize}
    \item \emph{Démonstration 2 (Rolle appliqué à la fonction $f^2$) :} La fonction $f^2$ vérifie les hypothèses de Rolle sur $[a, b]$ et sur $[b, c]$. Il existe donc $d_1 \in \,]a, b[$ et $d_2 \in \,]b, c[$ tels que $(f^2)'(d_1) = 2f(d_1)f'(d_1) = 0$ et $(f^2)'(d_2) = 2f(d_2)f'(d_2) = 0$. Si $f'(d_1) = 0$ ou $f'(d_2) = 0$, la thèse est prouvée (avec $d = d_1$ ou $d = d_2$). Sinon, on a simultanément $f(d_1) = 0$ et $f(d_2) = 0$ : on applique alors le théorème de Rolle à $f$ sur $[d_1, d_2] \subset [a, c]$, ce qui fournit $d_3 \in \,]d_1, d_2[ \;\subset \,]a, c[$ tel que $f'(d_3) = 0$ !
\end{itemize}

\medskip
\textbf{4) UCLouvain -- Septembre 2025, Question 4 (TAF et Rolle combinés sur une fraction rationnelle) :}\\
\emph{Énoncé :} Soit $f(x) = \dfrac{x^2 + x + 2}{x^3 + 1}$. Montrer qu'il existe dans l'intervalle $\left]-\dfrac{1}{2}, 2\right[$ :
(1) un nombre $c_1$ où $f'(c_1) = -\dfrac{4}{9}$ ;
(2) et (3) deux nombres distincts $c_2$ et $c_3$ où $f'(c_2) = f'(c_3) = 0$ ;
(4) un nombre $c_4$ où $f''(c_4) = 0$.
\emph{Aide :} examiner la valeur de $f$ pour certains entiers de l'intervalle.\\
\emph{Résolution :} Le dénominateur $x^3 + 1$ ne s'annule qu'en $x = -1 \notin \left[-\frac{1}{2}, 2\right]$. La fraction rationnelle $f$ est donc de classe $C^\infty$ (continue et indéfiniment dérivable) sur $\left[-\frac{1}{2}, 2\right]$. Calculons les valeurs de $f$ aux bornes $-\frac{1}{2}$ et $2$, ainsi qu'aux entiers $0$ et $1$ de l'intervalle :
\[
    f\left(-\frac{1}{2}\right) = \frac{\frac{1}{4} - \frac{1}{2} + 2}{-\frac{1}{8} + 1} = \frac{7/4}{7/8} = 2, \qquad f(0) = \frac{2}{1} = 2, \qquad f(1) = \frac{1+1+2}{1+1} = \frac{4}{2} = 2, \qquad f(2) = \frac{4+2+2}{8+1} = \frac{8}{9}.
\]
\begin{itemize}
    \item \textbf{(1)} Appliquons le \textbf{théorème des accroissements finis} à $f$ sur l'intervalle complet $\left[-\frac{1}{2}, 2\right]$ : il existe $c_1 \in \left]-\frac{1}{2}, 2\right[$ tel que :
    \[
        f'(c_1) = \frac{f(2) - f(-1/2)}{2 - (-1/2)} = \frac{\frac{8}{9} - 2}{\frac{5}{2}} = \frac{-10/9}{5/2} = -\frac{4}{9}.
    \]
    \item \textbf{(2) et (3)} Comme $f\left(-\frac{1}{2}\right) = f(0) = f(1) = 2$, appliquons le \textbf{théorème de Rolle} (ou le TAF) à $f$ sur les deux sous-intervalles $\left[-\frac{1}{2}, 0\right]$ et $[0, 1]$ : il existe $c_2 \in \left]-\frac{1}{2}, 0\right[$ tel que $f'(c_2) = 0$ et $c_3 \in \,]0, 1[$ tel que $f'(c_3) = 0$. Puisque $-\frac{1}{2} < c_2 < 0 < c_3 < 1 < 2$, $c_2$ et $c_3$ sont bien distincts et appartiennent à $\left]-\frac{1}{2}, 2\right[$.
    \item \textbf{(4)} Appliquons enfin le \textbf{théorème de Rolle} (ou le TAF) \textbf{à la fonction dérivée $f'$} sur l'intervalle $[c_2, c_3]$ (puisque $f'(c_2) = f'(c_3) = 0$) : il existe $c_4 \in \,]c_2, c_3[ \;\subset \left]-\frac{1}{2}, 2\right[$ tel que $(f')'(c_4) = f''(c_4) = 0$.
\end{itemize}
\end{exemple}

% ----------------------------------------------------------------------------
\section{Exercices du Chapitre 3}
% ----------------------------------------------------------------------------

\emph{Les corrigés détaillés de tous ces exercices se trouvent au Chapitre 9, en fin d'ouvrage.}

\begin{enumerate}[label=\textbf{Exercice 3.\arabic* --}]
    \item \textbf{(Calculs de dérivées premières et secondes d'annales UCLouvain)} Calculer et factoriser au maximum la dérivée première $f'(x)$ (et la dérivée seconde $f''(x)$ lorsqu'elle est demandée) des fonctions suivantes :
    \begin{enumerate}
        \item $f_1(x) = x^2\sin(x^2) + \cos(x^2)$ \quad (calculer $f'_1$ et $f''_1$) \quad \textbf{(Juillet 2023 -- Série 1, Q1a-b)}
        \item $f_2(x) = 3\cos(x) - \cos^3(x)$ \quad (calculer $f'_2$ et $f''_2$) \quad \textbf{(Juillet 2023 -- Série 2, Q1a-b)}
        \item $f_3(x) = \dfrac{1 - x^2}{\cos^2(x)}$ \quad (écrire $f'_3(x)$ sous la forme $\dfrac{g(x)}{\cos^2(x)}$) \quad \textbf{(Septembre 2024, Q1b)}
        \item $f_4(x) = \dfrac{x^2}{2}\left(\ln^2(x) - \ln(x) + \dfrac{1}{2}\right)$ \quad (calculer $f'_4$ et $f''_4$) \quad \textbf{(Septembre 2026, Q1b-c)}
    \end{enumerate}

    \item \textbf{(Dérivées de fonctions réciproques et composées)} Calculer la dérivée des fonctions suivantes sur leur domaine de dérivabilité :
    \begin{enumerate}
        \item $g_1(x) = \arcsin(2x - 1)$
        \item $g_2(x) = \arctan\left(\dfrac{1+x}{1-x}\right)$ \quad (interpréter géométriquement le résultat simplifié !)
        \item $g_3(x) = x^x = e^{x\ln(x)} \quad (x > 0)$
    \end{enumerate}

    \item \textbf{(Limites par la règle de L'Hospital)} Calculer les limites suivantes en vérifiant à chaque étape les conditions d'application de la règle de L'Hospital :
    \begin{enumerate}
        \item $\displaystyle \lim_{x \to 0} \frac{e^x - e^{-x} - 2x}{x - \sin(x)}$
        \item $\displaystyle \lim_{x \to \pi/2} \frac{\ln(\sin x)}{(\pi - 2x)^2}$
        \item $\displaystyle \lim_{x \to 0^+} x^{\sin(x)}$
        \item $\displaystyle \lim_{x \to +\infty} \left(1 + \frac{y}{x}\right)^x \quad (y \in \R \text{ fixé})$
    \end{enumerate}

    \item \textbf{(UCLouvain -- Juillet 2023 -- Série 1, Q5 -- Applicabilité du théorème de Rolle)} Pour chacune des quatre fonctions suivantes, indiquer par \textbf{OUI} ou \textbf{NON} si l'on peut utiliser le théorème de Rolle sur l'intervalle $[-1, 1]$ pour prouver l'existence d'un point $c \in \,]-1, 1[$ auquel la dérivée s'annule, et justifier succinctement sur base des hypothèses du théorème :
    \begin{enumerate}
        \item $f(x) = 1 - \sqrt{|x|}$
        \item $g(x) = x - \dfrac{1}{x}$
        \item $h(x) = \ln(1 + x^2)$
        \item $w(x) = x - x^5$
    \end{enumerate}

    \item \textbf{(Application du théorème de Rolle aux polynômes)} Soit $P(x) = x^4 - 4x + 1$.
    \begin{enumerate}
        \item En utilisant le théorème de Rolle (ou son raisonnement par l'absurde), démontrer que l'équation $x^4 - 4x + 1 = 0$ ne peut pas avoir plus de deux racines réelles distinctes.
        \item En utilisant le théorème des valeurs intermédiaires, montrer qu'elle possède exactement deux racines réelles.
    \end{enumerate}

    \item \textbf{(Inégalités par le théorème des accroissements finis)} En appliquant le théorème des accroissements finis à une fonction judicieusement choisie sur un intervalle $[a, b]$, démontrer que :
    \begin{enumerate}
        \item Pour tous réels $a, b \in \R$, $\bigl|\sin(b) - \sin(a)\bigr| \le |b - a|$.
        \item Pour tout $x > 0$, $\dfrac{x}{1 + x^2} < \arctan(x) < x$.
    \end{enumerate}
\end{enumerate}

% ============================================================================
% CHAPITRE 4 : ÉTUDE COMPLÈTE DE FONCTIONS
% ============================================================================
\chapter{Étude Complète de Fonctions}

La \textbf{Question 1} de l'examen d'admission d'Analyse à l'UCLouvain (qui représente typiquement $15$ points sur $50$, soit $30\%$ de la note totale !) est systématiquement consacrée à l'\textbf{étude complète d'une fonction} $f : \D \subset \R \to \R$ aboutissant à son tableau de variations et à l'esquisse précise de son graphe sur un repère quadrillé. Ce chapitre expose la théorie reliant les dérivées $f'$ et $f''$ à la géométrie de la courbe (croissance, extrema locaux, concavité, points d'inflexion), puis détaille pas à pas la méthodologie attendue par les correcteurs à travers cinq études complètes issues des annales 2021 à 2026.

% ----------------------------------------------------------------------------
\section{Croissance, extrema locaux, concavité et points d'inflexion}
% ----------------------------------------------------------------------------

\subsection{Signe de la dérivée première $f'$ et sens de variation}

En appliquant le théorème des accroissements finis $f(x_2) - f(x_1) = f'(c)(x_2 - x_1)$ entre deux points quelconques $x_1 < x_2$ d'un intervalle $I$, on voit immédiatement que le signe de $f(x_2) - f(x_1)$ est le même que celui de $f'(c)$.

\begin{theoreme}[Théorème fondamental des variations sur un intervalle $I$]
Soit $f$ une fonction continue sur un intervalle $I$ et dérivable à l'intérieur de $I$ :
\begin{enumerate}
    \item $f'(x) \ge 0$ sur $I \iff f$ est \textbf{croissante} sur $I$ ($\nearrow$).\\
    Si $f'(x) > 0$ sur $I$ (sauf éventuellement en un nombre fini de points isolés où $f'(x) = 0$), alors $f$ est \textbf{strictement croissante} sur $I$.
    \item $f'(x) \le 0$ sur $I \iff f$ est \textbf{décroissante} sur $I$ ($\searrow$).\\
    Si $f'(x) < 0$ sur $I$ (sauf éventuellement en des points isolés), alors $f$ est \textbf{strictement décroissante} sur $I$.
    \item $f'(x) = 0$ pour tout $x \in I \iff f$ est \textbf{constante} sur l'intervalle $I$.
\end{enumerate}
\end{theoreme}

\subsection{Extrema locaux (minima et maxima) : conditions du 1\ier{} et du 2\ieme{} ordre}

\begin{definition}[Maximum local, minimum local et extremum global]
Soit $f : \D \subset \R \to \R$ et soit $x^* \in \D$ :
\begin{itemize}
    \item On dit que $f$ admet un \textbf{maximum local} (ou relatif) en $x^*$ s'il existe un intervalle ouvert $]x^* - \delta, x^* + \delta[$ ($\delta > 0$) tel que $\forall x \in \D \cap \,]x^* - \delta, x^* + \delta[,\; f(x) \le f(x^*)$. Si cette inégalité est vraie sur tout le domaine $\D$, c'est un \textbf{maximum global} (ou absolu).
    \item On dit que $f$ admet un \textbf{minimum local} (ou relatif) en $x^*$ s'il existe $\delta > 0$ tel que $\forall x \in \D \cap \,]x^* - \delta, x^* + \delta[,\; f(x) \ge f(x^*)$ (et \textbf{minimum global} si c'est vrai sur tout $\D$).
    \item Un \textbf{point critique} (ou point stationnaire) de $f$ est un point $x^* \in \D$ où la dérivée s'annule : $\boxed{f'(x^*) = 0}$.
\end{itemize}
\end{definition}

\begin{theoreme}[Caractérisation des extrema locaux d'une fonction dérivable]
Soit $f$ une fonction dérivable sur un intervalle ouvert $I \subset \D$ et soit $x^* \in I$ un point intérieur à $I$ :
\begin{enumerate}
    \item \textbf{Condition nécessaire du 1\ier{} ordre (Théorème de Fermat) :}
    \[
        \boxed{\text{Si } f \text{ a un extremum local en un point intérieur } x^* \in I, \text{ alors } f'(x^*) = 0}.
    \]
    \item \textbf{Condition suffisante du 1\ier{} ordre (changement de signe de $f'$) :} Soit $x^* \in I$ tel que $f'(x^*) = 0$ :
    \begin{itemize}
        \item Si $f'(x)$ passe du signe \textbf{$+$ à gauche au signe $-$ à droite} de $x^*$ ($\nearrow \;\text{max}\; \searrow$), alors $f$ admet un \textbf{maximum local} en $x^*$.
        \item Si $f'(x)$ passe du signe \textbf{$-$ à gauche au signe $+$ à droite} de $x^*$ ($\searrow \;\text{min}\; \nearrow$), alors $f$ admet un \textbf{minimum local} en $x^*$.
        \item Si $f'(x)$ \textbf{garde le même signe} à gauche et à droite de $x^*$ (tout en s'annulant en $x^*$, comme $x^3$ en $0$ ou la fonction de \textbf{Septembre 2026 Q1} en $x = 1$), alors $x^*$ \textbf{n'est PAS un extremum local} : c'est un \textbf{point d'inflexion à tangente horizontale} !
    \end{itemize}
    \item \textbf{Condition suffisante du 2\ieme{} ordre (signe de la dérivée seconde $f''(x^*)$) :} Si $f$ est deux fois dérivable en un point critique $x^* \in I$ ($f'(x^*) = 0$) :
    \[
        \boxed{\begin{cases}
            f'(x^*) = 0 \;\text{ et }\; f''(x^*) < 0 &\implies f \text{ admet un \textbf{maximum local strict} en } x^* \;\; (\cap) \\
            f'(x^*) = 0 \;\text{ et }\; f''(x^*) > 0 &\implies f \text{ admet un \textbf{minimum local strict} en } x^* \;\; (\cup)
        \end{cases}}
    \]
    \emph{Attention :} Si $f'(x^*) = 0$ et $f''(x^*) = 0$, le test de la dérivée seconde seul ne permet pas de conclure : il faut regarder si $f'(x)$ change de signe de part et d'autre de $x^*$ (par exemple, pour $x^4$ ou $f(x) = x^2\sin(x^2)+\cos(x^2)$ en $x=0$ en Juillet 2023 Série 1 Q1, on a $f'(0) = f''(0) = 0$ et c'est un minimum très plat ; tandis que pour $x^3$ en $0$, $f'(0) = f''(0) = 0$ et c'est un point d'inflexion !).
\end{enumerate}
\end{theoreme}

\subsection{Signe de la dérivée seconde $f''$, concavité et points d'inflexion}

\begin{definition}[Concavité tournée vers le haut / vers le bas et point d'inflexion]
Soit $f$ une fonction deux fois dérivable sur un intervalle $I$ :
\begin{enumerate}
    \item \textbf{Concavité positive (tournée vers le haut $\cup$, fonction convexe) :}
    \[
        \boxed{f''(x) > 0 \text{ sur } I \implies f' \text{ croissante} \implies \text{concavité tournée vers le haut } (\cup)}.
    \]
    Géométriquement, la courbe est située \textbf{au-dessus de toutes ses tangentes} sur $I$.
    \item \textbf{Concavité négative (tournée vers le bas $\cap$, fonction concave) :}
    \[
        \boxed{f''(x) < 0 \text{ sur } I \implies f' \text{ décroissante} \implies \text{concavité tournée vers le bas } (\cap)}.
    \]
    Géométriquement, la courbe est située \textbf{en dessous de toutes ses tangentes} sur $I$.
    \item \textbf{Point d'inflexion (PI) :} Un point $x_0 \in \D$ (où $f$ est continue) est un \textbf{point d'inflexion} du graphe de $f$ si la fonction \textbf{change de sens de concavité} en passant par $x_0$ (c'est-à-dire si la dérivée seconde $f''(x)$ \textbf{s'annule en $x_0$ en changeant de signe} de part et d'autre de $x_0$). En un point d'inflexion, la courbe \textbf{traverse sa droite tangente}.
\end{enumerate}
\end{definition}

\begin{center}
\begin{tikzpicture}[scale=0.88, >=Stealth]
    % Concavité vers le haut (minimum)
    \begin{scope}[shift={(0,0)}]
        \draw[very thick, uclblue, domain=-1.5:1.5, samples=60] plot (\x, {0.65*(\x)^2 + 0.4});
        \draw[thick, defgreen] (-1.2,0.4) -- (1.2,0.4);
        \fill[defgreen] (0,0.4) circle (2pt) node[below] {\footnotesize Min ($f'=0, f''>0$)};
        \node[above, uclblue, font=\bfseries\footnotesize] at (0,2.1) {Concavité positive ($f'' > 0, \cup$)};
        \node[below, darkgray, font=\footnotesize] at (0,-0.25) {Au-dessus des tangentes};
    \end{scope}

    % Concavité vers le bas (maximum)
    \begin{scope}[shift={(5.0,0)}]
        \draw[very thick, warnred, domain=-1.5:1.5, samples=60] plot (\x, {1.85 - 0.65*(\x)^2});
        \draw[thick, defgreen] (-1.2,1.85) -- (1.2,1.85);
        \fill[defgreen] (0,1.85) circle (2pt) node[above] {\footnotesize Max ($f'=0, f''<0$)};
        \node[above, warnred, font=\bfseries\footnotesize] at (0,2.4) {Concavité négative ($f'' < 0, \cap$)};
        \node[below, darkgray, font=\footnotesize] at (0,-0.25) {En dessous des tangentes};
    \end{scope}

    % Point d'inflexion
    \begin{scope}[shift={(10.0,0.9)}]
        \draw[very thick, exorange, domain=-1.45:1.45, samples=60] plot (\x, {0.42*(\x)^3 + 0.15*\x});
        \draw[thick, dashed, darkgray] (-1.3,-0.195) -- (1.3,0.195);
        \fill[exorange] (0,0) circle (2.2pt) node[below right] {\footnotesize PI ($f''=0$)};
        \node[above, exorange, font=\bfseries\footnotesize] at (0,1.5) {Point d'inflexion ($\cap \to \cup$)};
        \node[below, darkgray, font=\footnotesize] at (0,-1.15) {Traverse sa tangente};
    \end{scope}
\end{tikzpicture}
\end{center}

\begin{exemple}[Extrema locaux en fonction de paramètres : UCLouvain -- Juillet 2021, Question 3a]
\emph{Énoncé :} Soient $a, b, c$ des constantes non nulles. On considère la fonction $f$ définie sur $\R_0$ par :
\[
    f(x) = \frac{ax^2 + bx + c}{x^2}.
\]
En quel(s) point(s) $x$ la fonction $f$ admet-elle un maximum et/ou un minimum local ? S'agit-il d'un maximum ou d'un minimum (la réponse peut dépendre des valeurs de $a, b, c$) ?

\emph{Résolution :}
\begin{itemize}
    \item \emph{Astuce de réécriture avant de dériver :} Au lieu d'appliquer la lourde formule de dérivation d'un quotient $\frac{u}{v}$, divisons chaque terme du numérateur par $x^2$ :
    \[
        f(x) = a + \frac{b}{x} + \frac{c}{x^2} = a + b\,x^{-1} + c\,x^{-2} \quad (\forall x \neq 0).
    \]
    \item \emph{Dérivée première et point critique :} En dérivant terme à terme :
    \[
        f'(x) = -\frac{b}{x^2} - \frac{2c}{x^3} = -\frac{bx + 2c}{x^3}.
    \]
    Comme $b \neq 0$ et $c \neq 0$, la dérivée s'annule en un unique point critique $x^* \neq 0$ :
    \[
        f'(x^*) = 0 \iff bx^* + 2c = 0 \iff \boxed{x^* = -\frac{2c}{b}}.
    \]
    \item \emph{Nature de l'extremum par la dérivée seconde $f''(x^*)$ :} Dérivons une seconde fois $f'(x) = -b\,x^{-2} - 2c\,x^{-3}$ :
    \[
        f''(x) = \frac{2b}{x^3} + \frac{6c}{x^4} = \frac{2bx + 6c}{x^4}.
    \]
    Évaluons $f''$ en $x^* = -\frac{2c}{b}$ (où $2bx^* + 6c = 2b\left(-\frac{2c}{b}\right) + 6c = -4c + 6c = 2c$) :
    \[
        f''(x^*) = \frac{2c}{(x^*)^4} = \frac{2c}{\left(-\frac{2c}{b}\right)^4} = \frac{2c\,b^4}{16c^4} = \frac{b^4}{8c^3}.
    \]
    Puisque $b^4 > 0$ pour tout $b \neq 0$, le signe de $f''(x^*)$ ne dépend que du signe de $c$ (ni $a$ ni $b$ n'influencent la nature de l'extremum !) :
    \begin{itemize}
        \item \textbf{Si $c > 0$ :} $f''(x^*) > 0$, donc $f$ admet un \textbf{minimum local} en $x^* = -\dfrac{2c}{b}$.
        \item \textbf{Si $c < 0$ :} $f''(x^*) < 0$, donc $f$ admet un \textbf{maximum local} en $x^* = -\dfrac{2c}{b}$.
    \end{itemize}
    \emph{Remarque élégante du solutionnaire :} En posant $y = \frac{1}{x}$, on a $f(x) = a + by + cy^2$, qui est une simple parabole en $y$ dont le sommet est en $y^* = -\frac{b}{2c} \iff x^* = -\frac{2c}{b}$, tournée vers le haut (minimum) si $c > 0$ et vers le bas (maximum) si $c < 0$ !
\end{itemize}
\end{exemple}

% ----------------------------------------------------------------------------
\section{Méthodologie en 6 étapes pour réussir la Question 1 de l'examen}
% ----------------------------------------------------------------------------

À l'examen d'admission de l'UCLouvain, la Question 1 est guidée par des sous-questions a), b), c), d), e) (et parfois f)). Voici la feuille de route méthodologique qui répond exactement aux attentes du barème officiel :

\begin{astuce}[Les 6 étapes du succès à la Question 1 (Étude de fonction)]
\begin{enumerate}
    \item \textbf{Domaine $\D$, symétries, racines et comportement aux bords :}
    \begin{itemize}
        \item Écrire les conditions d'existence et exprimer clairement $\D$.
        \item Repérer si $f$ est paire, impaire ou périodique.
        \item Chercher les racines $f(x) = 0$ si l'équation est résoluble, ainsi que les limites aux bords ouverts de $\D$ (asymptotes verticales AV, horizontales AH ou obliques AO).
    \end{itemize}
    \item \textbf{Dérivée première $f'(x)$ et ses racines :}
    \begin{itemize}
        \item Calculer $f'(x)$ proprement et \textbf{la mettre sous forme factorisée / sur un même dénominateur} (sans factorisation, il est impossible d'étudier son signe !).
        \item Résoudre $f'(x) = 0$ sur $\D$ pour identifier les candidats extrema locaux.
    \end{itemize}
    \item \textbf{Dérivée seconde $f''(x)$ et ses racines (si demandée) :}
    \begin{itemize}
        \item Partir de la forme la plus simple de $f'(x)$, calculer $f''(x)$ et la factoriser pour trouver ses racines (candidats points d'inflexion) et son signe.
    \end{itemize}
    \item \textbf{Calcul des valeurs exactes ET numériques approchées aux points remarquables :}
    \begin{itemize}
        \item Depuis 2024, l'énoncé fournit une \textbf{Aide numérique} (par ex. $\sqrt{2} \simeq 1,414$, $\sqrt{5} \simeq 2,236$, $e^{-1} \simeq 0,368$, $e^{-2} \simeq 0,135$, $\ln(2) \simeq 0,693$, $\pi^2 \simeq 9,87$).
        \item Calculez au brouillon, puis reportez dans le tableau la \textbf{valeur mathématique exacte ET la valeur décimale approchée} de $f(x)$ (et de $f'(x)$ si demandé) en chaque point remarquable : c'est indispensable pour placer les points au bon endroit sur le quadrillage !
    \end{itemize}
    \item \textbf{Tableau de variations dans l'ordre logique officiel :}
    \begin{itemize}
        \item Les énoncés récents (2025--2026) demandent explicitement de présenter les lignes du tableau dans l'\textbf{ordre logique} :
        \[
            \mathbf{1.\; x} \;\to\; \mathbf{2.\; f'(x)} \;\text{(signe)} \;\to\; \mathbf{3.\; f''(x)} \;\text{(signe)} \;\to\; \mathbf{4.\; f(x)} \;\text{(valeurs, } \nearrow\searrow, \cup\cap\text{)}.
        \]
        \item Indiquez dans la dernière ligne les mentions \textbf{max}, \textbf{min}, \textbf{inflexion (PI)}, \textbf{racine} et \textbf{AV}.
    \end{itemize}
    \item \textbf{Esquisse graphique sur le quadrillage avec axes chiffrés :}
    \begin{itemize}
        \item \textbf{Graduez et chiffrez vos deux axes $x$ et $y$} en choisissant une échelle adaptée aux valeurs de votre tableau.
        \item Tracez d'abord les \textbf{asymptotes en pointillés} (AV, AH, AO), puis placez tous les \textbf{points remarquables} (racines, extrema avec un petit segment horizontal pour la tangente horizontale, points d'inflexion, bords avec tangente verticale éventuelle).
        \item Reliez les points par une courbe lisse respectant scrupuleusement le sens de concavité ($\cup$ ou $\cap$) et annotez clairement sur le dessin : \emph{racine}, \emph{max}, \emph{min}, \emph{point d'inflexion}, \emph{AV}.
    \end{itemize}
\end{enumerate}
\end{astuce}

% ----------------------------------------------------------------------------
\section{Cinq études complètes de fonctions issues des annales UCLouvain}
% ----------------------------------------------------------------------------

Étudions maintenant de A à Z cinq sujets représentatifs couvrant toute la diversité des fonctions posées à l'examen d'admission entre 2022 et 2026.

\begin{exemple}[Étude 1 : Fonction trigonométrique périodique -- UCLouvain, Septembre 2022, Question 1]
\emph{Énoncé :} On considère la fonction $f : \D \subset \R \to \R$ définie par :
\[
    f(x) = \frac{\cos^2(x)}{1 + \cos(x)}.
\]
\begin{enumerate}[label=\alph*)]
    \item Quel est le domaine de définition $\D$ de $f$ ?
    \item Calculez la fonction dérivée $f'$.
    \item Quels sont les minima et maxima locaux éventuels de $f$ ? \emph{(Note : il ne vous est pas demandé de calculer la dérivée seconde).}
    \item Sur base de vos résultats précédents, esquissez le graphique de la fonction $f$ pour $x \in [-\pi, \pi]$. Veillez à aussi y indiquer la position des racines, s'il y en a.
\end{enumerate}

\emph{Résolution complète :}
\begin{enumerate}[label=\alph*)]
    \item \textbf{Domaine, parité, périodicité et racines :}
    \begin{itemize}
        \item \textbf{CE :} $1 + \cos(x) \neq 0 \iff \cos(x) \neq -1 \iff x \neq \pi + 2k\pi$ ($k \in \Z$). D'où $\D = \R \setminus \{\pi + 2k\pi \mid k \in \Z\}$.
        \item La fonction $f$ est \textbf{paire} ($f(-x) = f(x)$) et \textbf{périodique de période $2\pi$}. Sur $[-\pi, \pi]$, le domaine est $]-\pi, \pi[$.
        \item \textbf{Racines :} $f(x) = 0 \iff \cos^2(x) = 0$ (avec $\cos(x) \neq -1$) $\iff \cos(x) = 0 \iff x = \frac{\pi}{2} + k\pi$. Sur $]-\pi, \pi[$, il y a deux racines : $x = -\frac{\pi}{2}$ et $x = \frac{\pi}{2}$.
        \item \textbf{Asymptotes verticales :} Pour $x \to \pm\pi$ (à l'intérieur de $]-\pi, \pi[$), $\cos(x) \to -1^+$ donc $1 + \cos(x) \to 0^+$ et $\cos^2(x) \to (-1)^2 = 1 > 0$. Ainsi $\lim_{x \to \pi^-} f(x) = \lim_{x \to -\pi^+} f(x) = \frac{1}{0^+} = +\infty$, d'où les asymptotes verticales $\text{AV} \equiv x = \pi + 2k\pi$.
    \end{itemize}
    \item \textbf{Dérivée première $f'(x)$ :} Par la règle du quotient avec $(\cos^2 x)' = -2\cos(x)\sin(x)$ :
    \begin{align*}
        f'(x) &= \frac{-2\cos(x)\sin(x)\bigl(1+\cos(x)\bigr) - \cos^2(x)\bigl(-\sin(x)\bigr)}{\bigl(1+\cos(x)\bigr)^2} \\
        &= \boxed{-\frac{\cos(x)\sin(x)\bigl(2 + \cos(x)\bigr)}{\bigl(1+\cos(x)\bigr)^2}} = -\frac{\sin(2x)\bigl(2+\cos(x)\bigr)}{2\bigl(1+\cos(x)\bigr)^2}.
    \end{align*}
    \item \textbf{Signe de $f'(x)$ et extrema locaux :} Sur tout le domaine $\D$, on a $\bigl(1+\cos(x)\bigr)^2 > 0$ et $2 + \cos(x) \ge 2 - 1 = 1 > 0$. Le signe de $f'(x)$ est donc exactement le signe de $-\cos(x)\sin(x) = -\frac{1}{2}\sin(2x)$.
    \[
        f'(x) = 0 \iff \sin(x) = 0 \text{ ou } \cos(x) = 0 \iff x = 2k\pi \text{ ou } x = \frac{\pi}{2} + k\pi \;\; (k \in \Z).
    \]
    Dressons le tableau de signes et de variations sur $[-\pi, \pi]$ :
    \begin{center}
    {\small
    \renewcommand{\arraystretch}{1.35}
    \setlength{\tabcolsep}{4pt}
    \begin{tabular}{|c||c|c|c|c|c|c|c|c|c|}
    \hline
    $x$ & $-\pi$ & & $-\dfrac{\pi}{2}$ & & $0$ & & $\dfrac{\pi}{2}$ & & $\pi$ \\ \hline\hline
    $\sin(2x)$ & $0$ & $+$ & $0$ & $-$ & $0$ & $+$ & $0$ & $-$ & $0$ \\ \hline
    $1 + \cos(x)$ & $0$ & $+$ & $1$ & $+$ & $2$ & $+$ & $1$ & $+$ & $0$ \\ \hline
    $f'(x)$ & AV & $-$ & $0$ & $+$ & $0$ & $-$ & $0$ & $+$ & AV \\ \hline
    $f(x)$ & $\text{AV}\;(+\infty)$ & $\searrow$ & $\mathbf{0}$ (min) & $\nearrow$ & $\mathbf{\dfrac{1}{2}}$ (max) & $\searrow$ & $\mathbf{0}$ (min) & $\nearrow$ & $(+\infty)\;\text{AV}$ \\ \hline
    \end{tabular}
    }
    \end{center}
    La fonction admet donc :
    \begin{itemize}
        \item des \textbf{minima locaux} (et globaux, qui sont aussi les racines de $f$, de valeur $0$) en $\boxed{x = \dfrac{\pi}{2} + k\pi}$ ($k \in \Z$) ;
        \item des \textbf{maxima locaux} (de valeur $f(0) = \dfrac{1}{2}$) en $\boxed{x = 2k\pi}$ ($k \in \Z$).
    \end{itemize}
    \item \textbf{Esquisse graphique sur $[-\pi, \pi]$ :}
    \begin{center}
    \begin{tikzpicture}[scale=1.05, >=Stealth]
        \begin{axis}[
            width=10.5cm, height=6.5cm,
            axis lines=middle,
            xmin=-3.6, xmax=3.6,
            ymin=-0.3, ymax=3.2,
            xtick={-3.1416, -1.5708, 0, 1.5708, 3.1416},
            xticklabels={$-\pi$, $-\frac{\pi}{2}$, $0$, $\frac{\pi}{2}$, $\pi$},
            ytick={0.5, 1, 2, 3},
            yticklabels={$\frac{1}{2}$, $1$, $2$, $3$},
            xlabel={$x$},
            ylabel={$f(x)$},
            grid=both,
            grid style={gray!20},
            samples=120
        ]
            \addplot[very thick, uclblue, domain=-2.72:2.72] {(cos(deg(x)))^2 / (1 + cos(deg(x)))};
            \draw[dashed, warnred, thick] (axis cs:-3.1416, -0.2) -- (axis cs:-3.1416, 3.1) node[pos=0.85, right] {\small AV};
            \draw[dashed, warnred, thick] (axis cs:3.1416, -0.2) -- (axis cs:3.1416, 3.1) node[pos=0.85, left] {\small AV};
            \fill[defgreen] (axis cs:-1.5708, 0) circle (2.2pt) node[below] {\footnotesize min \& racine};
            \fill[defgreen] (axis cs:1.5708, 0) circle (2.2pt) node[below] {\footnotesize min \& racine};
            \fill[exorange] (axis cs:0, 0.5) circle (2.2pt) node[above right] {\footnotesize max $(0, 1/2)$};
        \end{axis}
    \end{tikzpicture}
    \end{center}
\end{enumerate}
\end{exemple}

\begin{exemple}[Étude 2 : Fonction algébrique avec radical -- UCLouvain, Juillet 2024, Question 1]
\emph{Énoncé :} On considère la fonction $f : \D \subset \R \to \R$ définie par :
\[
    f(x) = \sqrt{1 + x} - x.
\]
\begin{enumerate}[label=\alph*)]
    \item Quel est le domaine de définition $\D$ de $f$ ? Obtenez aussi la(les) racine(s) de $f$.
    \item Obtenez la fonction dérivée première $f'$.
    \item Obtenez la fonction dérivée seconde $f''$.
    \item Réalisez un tableau avec les points $x = -1, x = 0, x = 1, x = 2$, la(les) racine(s) de $f$ et le(s) point(s) où $f'$ s'annule. Reportez-y clairement toutes les valeurs de $f$ et $f'$, ainsi que le signe de $f''$.
    \item Esquissez proprement le graphe de $f$ dans l'intervalle $[-1, 2]$. \textbf{Aide :} $\sqrt{2} \simeq 1,414$, $\sqrt{3} \simeq 1,732$, $\sqrt{5} \simeq 2,236$.
\end{enumerate}

\emph{Résolution complète :}
\begin{enumerate}[label=\alph*)]
    \item \textbf{Domaine et racine :} Comme démontré à l'Exemple du Chapitre 1, $1 + x \ge 0 \iff x \ge -1$, donc $\boxed{\D = [-1, +\infty[}$. L'équation $\sqrt{1+x} = x$ (qui impose $x \ge 0$ et $x^2 - x - 1 = 0$) admet \textbf{une unique racine} dans $\D$ :
    \[
        \boxed{x_1 = \frac{1 + \sqrt{5}}{2} \simeq \frac{1 + 2,236}{2} = 1,618}.
    \]
    \item \textbf{Dérivée première $f'(x)$ et point critique :} Pour tout $x > -1$ :
    \[
        \boxed{f'(x) = \frac{1}{2\sqrt{1+x}} - 1}.
    \]
    (En $x \to -1^+$, $\lim_{x \to -1^+} f'(x) = +\infty$ : demi-tangente verticale au bord gauche $x = -1$). Cherchons où $f'$ s'annule :
    \[
        f'(x) = 0 \iff \sqrt{1+x} = \frac{1}{2} \iff 1 + x = \frac{1}{4} \iff \boxed{x = -\frac{3}{4} = -0,75}.
    \]
    En ce point $x = -\frac{3}{4}$, la fonction vaut $f\left(-\frac{3}{4}\right) = \sqrt{\frac{1}{4}} - \left(-\frac{3}{4}\right) = \frac{1}{2} + \frac{3}{4} = \dfrac{5}{4} = 1,25$.
    \item \textbf{Dérivée seconde $f''(x)$ :} En dérivant $f'(x) = \frac{1}{2}(1+x)^{-1/2} - 1$ pour $x > -1$ :
    \[
        \boxed{f''(x) = -\frac{1}{4}(1+x)^{-3/2} = -\frac{1}{4(1+x)\sqrt{1+x}} < 0 \quad (\forall x > -1)}.
    \]
    La dérivée seconde est \textbf{strictement négative partout} sur $]-1, +\infty[$ : la concavité est toujours tournée vers le bas ($\cap$), il n'y a \textbf{aucun point d'inflexion}, et le point critique $x = -\frac{3}{4}$ est un \textbf{maximum global}.
    \item \textbf{Tableau de variations et de valeurs sur $[-1, 2]$ :}
    \begin{center}
    {\scriptsize
    \renewcommand{\arraystretch}{1.45}
    \setlength{\tabcolsep}{1.4pt}
    \begin{tabular}{|c||c|c|c|c|c|c|c|c|c|c|c|}
    \hline
    $x$ & $-1$ & & $-\dfrac{3}{4}$ & & $0$ & & $1$ & & $x_1 = \dfrac{1+\sqrt{5}}{2}$ & & $2$ \\ \hline\hline
    $f'(x)$ & $+\infty$ & $+$ & $0$ & $-$ & $-\dfrac{1}{2}$ & $-$ & $\dfrac{1}{2\sqrt{2}}-1 \simeq -0,65$ & $-$ & $\dfrac{1}{2x_1}-1 < 0$ & $-$ & $\dfrac{1}{2\sqrt{3}}-1 \simeq -0,71$ \\ \hline
    $f''(x)$ & $-\infty$ & $-$ & $-$ & $-$ & $-$ & $-$ & $-$ & $-$ & $-$ & $-$ & $-$ \\ \hline
    $f(x)$ & $1$ & $\nearrow$ & $\mathbf{\dfrac{5}{4} = 1,25}$ (max) & $\searrow$ & $1$ & $\searrow$ & $\sqrt{2}-1 \simeq 0,41$ & $\searrow$ & $\mathbf{0}$ (racine) & $\searrow$ & $\sqrt{3}-2 \simeq -0,27$ \\ \hline
    \end{tabular}
    }
    \end{center}
    \item \textbf{Esquisse graphique sur $[-1, 2]$ :}
    \begin{center}
    \begin{tikzpicture}[scale=1.05, >=Stealth]
        \begin{axis}[
            width=10.5cm, height=6.5cm,
            axis lines=middle,
            xmin=-1.25, xmax=2.25,
            ymin=-0.45, ymax=1.45,
            xtick={-1, -0.75, 0, 1, 1.618, 2},
            xticklabels={$-1$, $-\frac{3}{4}$, $0$, $1$, $x_1 \simeq 1{,}62$, $2$},
            ytick={-0.268, 0.414, 1, 1.25},
            yticklabels={$-0{,}27$, $0{,}41$, $1$, $\frac{5}{4}$},
            xlabel={$x$},
            ylabel={$f(x)$},
            grid=both,
            grid style={gray!20},
            samples=150
        ]
            \addplot[very thick, uclblue, domain=-1:2.05] {sqrt(1+x) - x};
            \draw[thick, warnred, <->] (axis cs:-0.95, 1.25) -- (axis cs:-0.55, 1.25);
            \draw[thick, exorange, ->] (axis cs:-1, 1) -- (axis cs:-1, 1.32) node[above right] {\footnotesize pente $\infty$};
            \fill[exorange] (axis cs:-1, 1) circle (2pt);
            \fill[warnred] (axis cs:-0.75, 1.25) circle (2.2pt) node[above] {\footnotesize max $(-3/4, 5/4)$};
            \fill[uclblue] (axis cs:0, 1) circle (2pt);
            \fill[uclblue] (axis cs:1, 0.414) circle (2pt);
            \fill[defgreen] (axis cs:1.618, 0) circle (2.2pt) node[above right] {\footnotesize racine $x_1$};
            \fill[uclblue] (axis cs:2, -0.268) circle (2pt);
        \end{axis}
    \end{tikzpicture}
    \end{center}
\end{enumerate}
\end{exemple}

\begin{exemple}[Étude 3 : Fonction logarithmique et trigonométrique -- UCLouvain, Juillet 2025, Question 1]
\emph{Énoncé :} On considère la fonction $f : \D \subset \R \to \R$ définie par :
\[
    f(x) = x + \ln\bigl(\sin(x) + \cos(x)\bigr).
\]
\begin{enumerate}[label=\alph*)]
    \item Obtenez le domaine de définition $\D = \bigcup_{k\in\Z} \,]a_k, b_k[$.
    \item Obtenez la fonction dérivée première $f'$.
    \item Obtenez la fonction dérivée seconde $f''$.
    \item Réalisez le tableau de variation de $f$ pour la partie de $\D$ comprise dans $[-\pi, \pi]$, dans l'ordre logique $f'$, puis $f''$, puis $f$.
    \item Esquissez proprement le graphe de $f$ sur cette même partie de $\D$.
\end{enumerate}

\emph{Résolution complète :}
\begin{enumerate}[label=\alph*)]
    \item \textbf{Domaine :} Vu à l'Exemple du Chapitre 1 : $\sin(x) + \cos(x) = \sqrt{2}\cos\left(x - \frac{\pi}{4}\right) > 0 \iff x \in \,]-\frac{\pi}{4} + 2k\pi,\; \frac{3\pi}{4} + 2k\pi[$. Dans $[-\pi, \pi]$, la partie du domaine est l'intervalle ouvert $\left]-\dfrac{\pi}{4}, \dfrac{3\pi}{4}\right[$.
    \item \textbf{Dérivée première $f'(x)$ :} En dérivant terme à terme et en mettant au même dénominateur :
    \[
        f'(x) = 1 + \frac{\cos(x) - \sin(x)}{\sin(x) + \cos(x)} = \frac{\sin(x) + \cos(x) + \cos(x) - \sin(x)}{\sin(x) + \cos(x)} = \boxed{\frac{2\cos(x)}{\sin(x) + \cos(x)}}.
    \]
    Sur $\left]-\frac{\pi}{4}, \frac{3\pi}{4}\right[$, le dénominateur $\sin(x) + \cos(x)$ est strictement positif. Donc $f'(x)$ a le même signe que $\cos(x)$, et s'annule uniquement en $\cos(x) = 0 \iff \boxed{x = \dfrac{\pi}{2}}$ (avec $f'(x) > 0$ sur $\left]-\frac{\pi}{4}, \frac{\pi}{2}\right[$ et $f'(x) < 0$ sur $\left]\frac{\pi}{2}, \frac{3\pi}{4}\right[$).
    \item \textbf{Dérivée seconde $f''(x)$ :} Dérivons le quotient $f'(x) = \dfrac{2\cos(x)}{\sin(x)+\cos(x)}$ :
    \begin{align*}
        f''(x) &= 2\,\frac{-\sin(x)\bigl(\sin(x)+\cos(x)\bigr) - \cos(x)\bigl(\cos(x)-\sin(x)\bigr)}{\bigl(\sin(x)+\cos(x)\bigr)^2} \\
        &= \frac{2\bigl(-\sin^2(x) - \cos^2(x)\bigr)}{\bigl(\sin(x)+\cos(x)\bigr)^2} = \boxed{-\frac{2}{\bigl(\sin(x)+\cos(x)\bigr)^2} < 0}.
    \end{align*}
    (On peut aussi l'écrire $-\dfrac{2}{1 + \sin(2x)}$ car $(\sin x + \cos x)^2 = \sin^2 x + \cos^2 x + 2\sin x\cos x = 1 + \sin(2x)$).
    La dérivée seconde est \textbf{strictement négative partout} : la concavité est négative ($\cap$) sur tout l'intervalle, sans point d'inflexion.
    \item \textbf{Comportement aux bornes, racines et tableau de variations sur $\left]-\frac{\pi}{4}, \frac{3\pi}{4}\right[$ :}
    \begin{itemize}
        \item En $x \to -\frac{\pi}{4}^+$ et en $x \to \frac{3\pi}{4}^-$, $\sin(x) + \cos(x) \to 0^+ \implies \ln(\sin x + \cos x) \to -\infty$, d'où deux \textbf{asymptotes verticales} $\text{AV}_1 \equiv x = -\dfrac{\pi}{4}$ et $\text{AV}_2 \equiv x = \dfrac{3\pi}{4}$.
        \item Au maximum $x = \frac{\pi}{2}$, on a $f\left(\frac{\pi}{2}\right) = \frac{\pi}{2} + \ln(1 + 0) = \dfrac{\pi}{2} \simeq 1,57$.
        \item En $x = 0$, $f(0) = 0 + \ln(0 + 1) = 0$ : $x = 0$ est une première \textbf{racine} évidente (avec pente $f'(0) = \frac{2}{1} = 2$). Comme $f$ redescend continûment de $f\left(\frac{\pi}{2}\right) = \frac{\pi}{2} > 0$ vers $-\infty$ en $\frac{3\pi}{4}^-$, le TVI garantit une \textbf{seconde racine} dans $\left]\frac{\pi}{2}, \frac{3\pi}{4}\right[$.
    \end{itemize}
    \begin{center}
    \renewcommand{\arraystretch}{1.4}
    \begin{tabular}{|c||c|c|c|c|c|c|c|}
    \hline
    $x$ & $-\dfrac{\pi}{4}$ & & $0$ & & $\dfrac{\pi}{2}$ & & $\dfrac{3\pi}{4}$ \\ \hline\hline
    $f'(x)$ & AV & $+$ & $2$ & $+$ & $0$ & $-$ & AV \\ \hline
    $f''(x)$ & AV & $-$ & $-2$ & $-$ & $-2$ & $-$ & AV \\ \hline
    $f(x)$ & $\text{AV}\;(-\infty)$ & $\nearrow\;\cap$ & $\mathbf{0}$ (racine) & $\nearrow\;\cap$ & $\mathbf{\dfrac{\pi}{2}}$ (max) & $\searrow\;\cap$ & $(-\infty)\;\text{AV}$ \\ \hline
    \end{tabular}
    \end{center}
    \item \textbf{Esquisse graphique sur $\left]-\frac{\pi}{4}, \frac{3\pi}{4}\right[$ :}
    \begin{center}
    \begin{tikzpicture}[scale=1.05, >=Stealth]
        \begin{axis}[
            width=10.0cm, height=6.5cm,
            axis lines=middle,
            xmin=-1.1, xmax=2.65,
            ymin=-2.4, ymax=1.95,
            xtick={-0.7854, 0, 0.7854, 1.5708, 2.3562},
            xticklabels={$-\frac{\pi}{4}$, $0$, $\frac{\pi}{4}$, $\frac{\pi}{2}$, $\frac{3\pi}{4}$},
            ytick={-2, -1, 0.7854, 1.5708},
            yticklabels={$-2$, $-1$, $\frac{\pi}{4}$, $\frac{\pi}{2}$},
            xlabel={$x$},
            ylabel={$f(x)$},
            grid=both,
            grid style={gray!20},
            samples=140
        ]
            \addplot[very thick, uclblue, domain=-0.75:2.335] {x + ln(sin(deg(x)) + cos(deg(x)))};
            \draw[dashed, warnred, thick] (axis cs:-0.7854, -2.35) -- (axis cs:-0.7854, 1.85) node[pos=0.15, right] {\small AV};
            \draw[dashed, warnred, thick] (axis cs:2.3562, -2.35) -- (axis cs:2.3562, 1.85) node[pos=0.15, left] {\small AV};
            \draw[thick, warnred, <->] (axis cs:1.25, 1.5708) -- (axis cs:1.89, 1.5708);
            \fill[warnred] (axis cs:1.5708, 1.5708) circle (2.2pt) node[above] {\footnotesize max $(\pi/2, \pi/2)$};
            \fill[defgreen] (axis cs:0, 0) circle (2.2pt) node[below right] {\footnotesize racine $(0,0)$};
            \fill[defgreen] (axis cs:2.253, 0) circle (2.2pt) node[above left] {\footnotesize 2\ieme{} racine};
        \end{axis}
    \end{tikzpicture}
    \end{center}
\end{enumerate}
\end{exemple}

\begin{exemple}[Étude 4 : Logarithmes, extremum et point d'inflexion -- UCLouvain, Septembre 2025, Question 1]
\emph{Énoncé :} On considère la fonction $f : \D \subset \R \to \R$ définie par :
\[
    f(x) = x\,\ln^2(x).
\]
\textbf{Aide :} $\lim_{y \to 0^+} y^p\ln(y) = 0$ pour $p > 0$, $e \simeq 2,72$, $\dfrac{1}{e} \simeq 0,368$ et $\ln(2) \simeq 0,693$.
\begin{enumerate}[label=\alph*)]
    \item Obtenez le domaine de définition $\D$ et la/les racine(s) de $f$.
    \item Obtenez la dérivée première $f'$ et ses racines.
    \item Obtenez la dérivée seconde $f''$ et ses racines.
    \item Réalisez le tableau de variation de $f$ (ordre logique $f'$, $f''$, $f$).
    \item Esquissez proprement le graphe de $f$ pour $x \le 2$ en indiquant les valeurs numériques aux points remarquables.
\end{enumerate}

\emph{Résolution complète :}
\begin{enumerate}[label=\alph*)]
    \item \textbf{Domaine, comportement en $0^+$ et racines :}
    \begin{itemize}
        \item \textbf{CE :} $x > 0 \implies \boxed{\D = \,]0, +\infty[}$.
        \item \textbf{Limite en $0^+$ :} $\lim_{x \to 0^+} x\ln^2(x) = \lim_{x \to 0^+} \bigl(x^{1/2}\ln(x)\bigr)^2 = 0^2 = 0$ (pas d'asymptote verticale en $0$, la courbe part du point limite $(0, 0)$).
        \item \textbf{Racine sur $\D = \,]0, +\infty[$ :} $f(x) = 0 \iff \ln(x) = 0 \iff \boxed{x = 1}$ (et $f(x) > 0$ pour tout $x \in \D \setminus \{1\}$).
    \end{itemize}
    \item \textbf{Dérivée première $f'(x)$ et ses racines :} Par la règle du produit et de la chaîne $\bigl(\ln^2(x)\bigr)' = 2\ln(x)\cdot\frac{1}{x}$ :
    \[
        f'(x) = 1\cdot\ln^2(x) + x \cdot 2\ln(x)\frac{1}{x} = \ln^2(x) + 2\ln(x) = \boxed{\ln(x)\bigl(\ln(x) + 2\bigr)}.
    \]
    (Notons qu'en $x \to 0^+$, $\lim_{x \to 0^+} f'(x) = (-\infty)(-\infty) = +\infty$ : demi-tangente verticale en $0^+$).
    \[
        f'(x) = 0 \iff \ln(x) = 0 \;\text{ ou }\; \ln(x) = -2 \iff \boxed{x = 1 \;\text{ ou }\; x = e^{-2} = \frac{1}{e^2} \simeq 0,135}.
    \]
    \item \textbf{Dérivée seconde $f''(x)$ et ses racines :} En dérivant $f'(x) = \ln^2(x) + 2\ln(x)$ :
    \[
        f''(x) = \frac{2\ln(x)}{x} + \frac{2}{x} = \boxed{\frac{2}{x}\bigl(\ln(x) + 1\bigr)}.
    \]
    Comme $x > 0$, $f''(x)$ s'annule en changeant de signe uniquement en :
    \[
        f''(x) = 0 \iff \ln(x) = -1 \iff \boxed{x = e^{-1} = \frac{1}{e} \simeq 0,368} \quad (\textbf{point d'inflexion}).
    \]
    \item \textbf{Valeurs remarquables et tableau de variations :}
    \begin{itemize}
        \item En $x = e^{-2} \simeq 0,135$ (maximum local) : $f(e^{-2}) = e^{-2}(-2)^2 = \dfrac{4}{e^2} \simeq 4 \times 0,135 = \mathbf{0,542}$.
        \item En $x = e^{-1} \simeq 0,368$ (point d'inflexion) : $f(e^{-1}) = e^{-1}(-1)^2 = \dfrac{1}{e} \simeq \mathbf{0,368}$.
        \item En $x = 1$ (minimum local et racine) : $f(1) = \mathbf{0}$.
        \item En $x = 2$ (bord droit du tracé) : $f(2) = 2\ln^2(2) \simeq 2 \times (0,693)^2 \simeq 2 \times 0,480 = \mathbf{0,960}$.
    \end{itemize}
    \begin{center}
    {\footnotesize
    \renewcommand{\arraystretch}{1.4}
    \setlength{\tabcolsep}{2.8pt}
    \begin{tabular}{|c||c|c|c|c|c|c|c|c|}
    \hline
    $x$ & $0$ & & $e^{-2} \simeq 0,135$ & & $e^{-1} \simeq 0,368$ & & $1$ & \\ \hline\hline
    $f'(x)$ & $+\infty$ & $+$ & $0$ & $-$ & $-1$ & $-$ & $0$ & $+$ \\ \hline
    $f''(x)$ & $-\infty$ & $-$ & $-2e^2 < 0$ & $-$ & $0$ & $+$ & $2 > 0$ & $+$ \\ \hline
    $f(x)$ & $\lim = 0$ & $\nearrow\;\cap$ & $\mathbf{\dfrac{4}{e^2} \simeq 0,542}$ (max) & $\searrow\;\cap$ & $\mathbf{\dfrac{1}{e} \simeq 0,368}$ (PI) & $\searrow\;\cup$ & $\mathbf{0}$ (min \& racine) & $\nearrow\;\cup$ \\ \hline
    \end{tabular}
    }
    \end{center}
    \item \textbf{Esquisse graphique sur $]0, 2]$ :}
    \begin{center}
    \begin{tikzpicture}[scale=1.05, >=Stealth]
        \begin{axis}[
            width=10.0cm, height=6.5cm,
            axis lines=middle,
            xmin=-0.1, xmax=2.15,
            ymin=-0.12, ymax=1.08,
            xtick={0.135, 0.368, 1, 2},
            xticklabels={$0{,}135$, $0{,}368$, $1$, $2$},
            ytick={0.368, 0.541, 0.960},
            yticklabels={$0{,}368$, $0{,}542$, $0{,}960$},
            xlabel={$x$},
            ylabel={$f(x)$},
            grid=both,
            grid style={gray!20},
            samples=200
        ]
            \addplot[very thick, uclblue, domain=0.001:2.03] {x * (ln(x))^2};
            \draw[uclblue, thick] (axis cs:0,0) circle (2pt);
            \fill[warnred] (axis cs:0.1353, 0.5413) circle (2.2pt) node[above right] {\footnotesize max $(e^{-2}, 4/e^2)$};
            \fill[exorange] (axis cs:0.3679, 0.3679) circle (2.2pt) node[above right] {\footnotesize inflexion $(1/e, 1/e)$};
            \fill[defgreen] (axis cs:1, 0) circle (2.2pt) node[above right] {\footnotesize min \& racine $(1,0)$};
            \fill[uclblue] (axis cs:2, 0.961) circle (2pt) node[left] {\footnotesize $(2;\; 0{,}96)$};
        \end{axis}
    \end{tikzpicture}
    \end{center}
\end{enumerate}
\end{exemple}

\begin{exemple}[Étude 5 : Le cas d'un point d'inflexion à tangente horizontale -- UCLouvain, Septembre 2026, Question 1]
\emph{Énoncé :} On considère la fonction $f : \,]0, +\infty[ \;\to \R$ définie par :
\[
    f(x) = \frac{x^2}{2}\left(\ln^2(x) - \ln(x) + \frac{1}{2}\right).
\]
\textbf{Aide :} $\lim_{y \to 0^+} y^p\ln(y) = 0$ pour $p > 0$ ; $e^{-2} \simeq 0,135$ ; $f(2) \simeq 0,575$.
\begin{enumerate}[label=\alph*)]
    \item Obtenez $\lim_{x \to 0} f(x)$ en détail.
    \item Obtenez la dérivée première $f'$, sa/ses racine(s) ainsi que $\lim_{x \to 0} f'(x)$.
    \item Obtenez la dérivée seconde $f''$, sa/ses racine(s) ainsi que $\lim_{x \to 0} f''(x)$.
    \item Réalisez le tableau de variation de $f$ (avec valeurs numériques approchées).
    \item Esquissez proprement le graphe de $f$ pour $x \le 2$.
\end{enumerate}

\emph{Résolution complète :}
\begin{enumerate}[label=\alph*)]
    \item \textbf{Limite en $0^+$ :} En développant $f(x)$ :
    \[
        \lim_{x \to 0^+} f(x) = \frac{1}{2}\lim_{x \to 0^+} \bigl(x\ln(x)\bigr)^2 - \frac{1}{2}\lim_{x \to 0^+} \bigl(x^2\ln(x)\bigr) + \frac{1}{4}\lim_{x \to 0^+} x^2 = 0 - 0 + 0 = \boxed{0}.
    \]
    (De plus, comme le trinôme $u^2 - u + \frac{1}{2} = \left(u - \frac{1}{2}\right)^2 + \frac{1}{4} \ge \frac{1}{4} > 0$ ne s'annule jamais, $f$ ne possède \textbf{aucune racine} sur $]0, +\infty[$ !).
    \item \textbf{Dérivée première $f'(x)$ (simplification remarquable !) :} Par la règle du produit :
    \begin{align*}
        f'(x) &= x\left(\ln^2(x) - \ln(x) + \frac{1}{2}\right) + \frac{x^2}{2}\left(2\ln(x)\cdot\frac{1}{x} - \frac{1}{x}\right) \\
        &= x\ln^2(x) - \cancel{x\ln(x)} + \cancel{\frac{x}{2}} + \cancel{x\ln(x)} - \cancel{\frac{x}{2}} = \boxed{x\,\ln^2(x)}.
    \end{align*}
    \emph{Remarque :} La fonction $f(x)$ de Septembre 2026 est tout simplement la primitive (trouvée en Septembre 2025 Q3 !) de la fonction $x\ln^2(x)$ étudiée en Septembre 2025 Q1 !
    \begin{itemize}
        \item \textbf{Racine de $f'$ :} $f'(x) = x\ln^2(x) = 0 \iff \ln(x) = 0 \iff \boxed{x = 1}$.
        \item \textbf{Attention au signe de $f'(x)$ :} Sur $]0, +\infty[$, on a $x > 0$ et $\ln^2(x) \ge 0$ (c'est un carré !). Donc $f'(x) \ge 0$ \textbf{partout} : $f'(x)$ s'annule en $x = 1$ \textbf{sans changer de signe} ! La fonction $f$ est donc \textbf{strictement croissante sur tout $]0, +\infty[$} et n'admet \textbf{aucun extremum local} !
        \item \textbf{Limite de $f'$ en $0^+$ :} $\lim_{x \to 0^+} f'(x) = \lim_{x \to 0^+} \bigl(x^{1/2}\ln(x)\bigr)^2 = 0^2 = \boxed{0}$ (demi-tangente horizontale en $(0,0)$ !).
    \end{itemize}
    \item \textbf{Dérivée seconde $f''(x)$ :} Comme $f'(x) = x\ln^2(x)$, on retrouve le calcul de l'Étude 4 :
    \[
        \boxed{f''(x) = \ln^2(x) + 2\ln(x) = \ln(x)\bigl(\ln(x) + 2\bigr)}.
    \]
    \begin{itemize}
        \item \textbf{Racines de $f''$ (deux points d'inflexion !) :} $f''(x) = 0 \iff \ln(x) = -2$ ou $\ln(x) = 0 \iff \boxed{x = e^{-2} \simeq 0,135}$ ou $\boxed{x = 1}$. Comme $f''(x)$ change de signe en chacun de ces deux points ($+$ avant $e^{-2}$, $-$ entre $e^{-2}$ et $1$, $+$ après $1$), la courbe possède \textbf{deux points d'inflexion}, dont celui en $x = 1$ est un \textbf{point d'inflexion à tangente horizontale} ($f'(1) = f''(1) = 0$).
        \item \textbf{Limite de $f''$ en $0^+$ :} $\lim_{x \to 0^+} \ln(x)\bigl(\ln(x)+2\bigr) = (-\infty)(-\infty) = \boxed{+\infty}$.
    \end{itemize}
    \item \textbf{Valeurs remarquables et tableau de variations :}
    \begin{itemize}
        \item En $x = e^{-2} \simeq 0,135$ : $f'(e^{-2}) = 4e^{-2} \simeq 0,540$ et $f(e^{-2}) = \frac{e^{-4}}{2}\left(4 - (-2) + \frac{1}{2}\right) = \frac{13}{4}e^{-4} \simeq 3,25 \times (0,135)^2 \simeq \mathbf{0,059}$.
        \item En $x = 1$ : $f'(1) = 0$, $f''(1) = 0$ et $f(1) = \frac{1}{2}\left(0 - 0 + \frac{1}{2}\right) = \dfrac{1}{4} = \mathbf{0,250}$.
    \end{itemize}
    \begin{center}
    {\small
    \renewcommand{\arraystretch}{1.4}
    \setlength{\tabcolsep}{4pt}
    \begin{tabular}{|c||c|c|c|c|c|c|}
    \hline
    $x$ & $\text{limite } 0$ & & $e^{-2} \simeq 0,135$ & & $1$ & \\ \hline\hline
    $f'(x)$ & $0$ & $+$ & $4e^{-2} \simeq 0,540$ & $+$ & $0$ & $+$ \\ \hline
    $f''(x)$ & $+\infty$ & $+$ & $0$ & $-$ & $0$ & $+$ \\ \hline
    $f(x)$ & $0$ & $\nearrow\;\cup$ & $\mathbf{\dfrac{13}{4}e^{-4} \simeq 0,059}$ (PI) & $\nearrow\;\cap$ & $\mathbf{\dfrac{1}{4} = 0,250}$ (PI tang. horiz.) & $\nearrow\;\cup$ \\ \hline
    \end{tabular}
    }
    \end{center}
    \item \textbf{Esquisse graphique sur $]0, 2]$ :}
    \begin{center}
    \begin{tikzpicture}[scale=1.05, >=Stealth]
        \begin{axis}[
            width=10.0cm, height=6.5cm,
            axis lines=middle,
            xmin=-0.1, xmax=2.15,
            ymin=-0.05, ymax=0.65,
            xtick={0.135, 1, 2},
            xticklabels={$0{,}135$, $1$, $2$},
            ytick={0.059, 0.250, 0.500, 0.575},
            yticklabels={$0{,}059$, $0{,}250$, $0{,}500$, $0{,}575$},
            xlabel={$x$},
            ylabel={$f(x)$},
            grid=both,
            grid style={gray!20},
            samples=200
        ]
            \addplot[very thick, uclblue, domain=0.001:2.03] {0.5 * x^2 * ((ln(x))^2 - ln(x) + 0.5)};
            \draw[thick, warnred, <->] (axis cs:0.65, 0.25) -- (axis cs:1.35, 0.25);
            \fill[exorange] (axis cs:0.1353, 0.0595) circle (2.2pt) node[below right] {\footnotesize PI $(e^{-2};\; 0{,}059)$};
            \fill[warnred] (axis cs:1, 0.25) circle (2.2pt) node[above left] {\footnotesize PI à tangente horizontale $(1;\; 0{,}25)$};
            \fill[uclblue] (axis cs:2, 0.575) circle (2pt) node[left] {\footnotesize $(2;\; 0{,}575)$};
        \end{axis}
    \end{tikzpicture}
    \end{center}
\end{enumerate}
\end{exemple}

% ----------------------------------------------------------------------------
\section{Exercices du Chapitre 4}
% ----------------------------------------------------------------------------

\emph{Les corrigés détaillés de tous ces exercices se trouvent au Chapitre 9, en fin d'ouvrage.}

\begin{enumerate}[label=\textbf{Exercice 4.\arabic* --}]
    \item \textbf{(UCLouvain -- Juillet 2021, Q1 -- Étude complète)} On considère la fonction $f : \D \subset \R \to \R$ définie par :
    \[
        f(x) = \ln\left(\frac{1}{\cos(x)}\right).
    \]
    \begin{enumerate}
        \item Déterminer le domaine de définition $\D$ de $f$.
        \item Préciser le domaine de dérivabilité et calculer $f'(x)$ ainsi que la dérivée seconde $f''(x)$.
        \item Déterminer les minima et maxima locaux éventuels de $f$ ainsi que ses asymptotes.
        \item Dresser le tableau de variations et de concavité sur $[-\pi, \pi[$, puis expliquer comment en déduire le graphe sur $[-4\pi, 4\pi]$ de $f(x)$ et de $\bar{f}(x) = \ln\left(\dfrac{1}{\cos^2(x)}\right)$.
    \end{enumerate}

    \item \textbf{(UCLouvain -- Juillet 2022, Q1 -- Étude complète)} Réaliser l'étude complète (domaine $\D$, dérivée première $f'$, dérivée seconde $f''$, extrema locaux, asymptotes, racines et tableau de variations sur $[-\pi, \pi]$) de la fonction :
    \[
        f(x) = \ln\left(\frac{1}{1 + \cos(x)}\right).
    \]

    \item \textbf{(UCLouvain -- Juillet 2023 -- Série 1, Q1 -- Étude complète)} On considère la fonction $f : \R \to \R$ définie par :
    \[
        f(x) = x^2\sin(x^2) + \cos(x^2).
    \]
    \begin{enumerate}
        \item Obtenir la dérivée première $f'(x)$ et identifier tous les points qui pourraient être un minimum ou un maximum local.
        \item Obtenir la dérivée seconde $f''(x)$.
        \item Réaliser un tableau reprenant les zones de croissance et de décroissance ainsi que les extrema de $f$ sur $[0, 2\sqrt{\pi}]$.
    \end{enumerate}

    \item \textbf{(UCLouvain -- Juillet 2023 -- Série 2, Q1 -- Étude complète)} On considère la fonction $f : \R \to \R$ définie par :
    \[
        f(x) = 3\cos(x) - \cos^3(x).
    \]
    \begin{enumerate}
        \item Obtenir $f'(x)$, identifier les candidats extrema locaux et trouver les racines de $f$.
        \item Obtenir $f''(x)$ et ses racines.
        \item Réaliser le tableau complet de variations et de concavité de $f$ sur $[0, 3\pi]$.
    \end{enumerate}

    \item \textbf{(UCLouvain -- Septembre 2023, Q1 -- Étude complète)} On considère la fonction $f : \D \subset \R \to \R$ définie par :
    \[
        f(x) = \frac{x}{\cos^2(x)}.
    \]
    Déterminer son domaine $\D$, sa dérivée première $f'(x)$, le tableau des valeurs de $f$ et $f'$ en $x \in \left\{0, \frac{\pi}{4}, \frac{\pi}{2}, \frac{3\pi}{4}, \pi\right\}$, et justifier pourquoi $f$ possède un minimum local sur $\left]\frac{3\pi}{4}, \pi\right[$.

    \item \textbf{(UCLouvain -- Juillet 2026, Q1 -- Étude complète)} On considère la fonction $f : \D \subset \R \to \R$ définie par :
    \[
        f(x) = (1 - x^2)\ln(1 - x) + x.
    \]
    \textbf{Aide :} $\lim_{y \to 0^+} y\ln(y) = 0$, $\ln(2) \simeq 0,6931$ et $e^{-1/2} \simeq 0,6065$.
    Déterminer $\D$, $\lim_{x \to 1} f(x)$, la dérivée première $f'(x)$ et ses deux racines $x_1, x_2$, la dérivée seconde $f''(x)$, puis dresser le tableau complet de variations et de concavité sur $[-1, 1[$.
\end{enumerate}

% ============================================================================
% CHAPITRE 5 : PROBLÈMES D'OPTIMISATION APPLIQUÉE
% ============================================================================
\chapter{Problèmes d'Optimisation Appliquée}

La \textbf{Question 2} de l'examen d'admission d'Analyse à l'UCLouvain (qui vaut elle aussi typiquement $15$ points sur $50$) est un \textbf{problème concret d'optimisation} inspiré du métier d'ingénieur : concevoir l'épaisseur optimale d'un isolant thermique, dimensionner une boîte métallique ou un enclos agricole, choisir la vitesse économique d'un drone-lanceur ou d'un navire marchand, optimiser un panneau solaire photovoltaïque, un accélérateur de particules, un propulseur spatial ou l'heure de prise d'un café avant l'examen ! Ce chapitre vous donne une méthode universelle en 5 étapes pour traduire n'importe quel énoncé en équations et le résoudre sans erreur.

% ----------------------------------------------------------------------------
\section{Méthodologie générale de résolution en 5 étapes}
% ----------------------------------------------------------------------------

Face à un long texte de problème (parfois une demi-page décrivant un système physique ou économique), la difficulté principale n'est pas le calcul de la dérivée, mais la \textbf{mise en équation} et l'\textbf{identification de la variable par rapport à laquelle il faut dériver}.

\begin{theoreme}[La méthode universelle en 5 étapes pour les problèmes d'optimisation]
\begin{enumerate}
    \item \textbf{Étape 1 -- Identifier la grandeur à optimiser (fonction objectif) et les variables :}
    Repérez dans la question finale ce qu'il faut \textbf{minimiser} ou \textbf{maximiser} : est-ce un coût $C$, une énergie $E$, une masse de métal $M$, un temps de trajet $T$, un volume $V$, une aire $A$, une vitesse $v_e$, une accélération $A$ ? Donnez un nom clair à cette grandeur et listez les variables dont elle dépend a priori (par ex. largeur $l$ et longueur $L$, ou côté $c$ et hauteur $h$), en les distinguant bien des \textbf{constantes fixées par l'énoncé} ($V_0, H, \alpha, \beta, f_0, K\dots$).
    \item \textbf{Étape 2 -- Écrire la ou les équation(s) de contrainte :}
    Si votre fonction objectif dépend de deux (ou trois) variables, cherchez dans l'énoncé la contrainte géométrique, physique ou budgétaire qui les relie (par ex. volume imposé $V = c^2 h$, longueur de clôture $C = L + 2l$, budget $A + V = 560$, ou conditions aux bords $f(1) = f(-1) = 1$).
    \item \textbf{Étape 3 -- Éliminer les variables liées pour obtenir une fonction d'UNE SEULE variable :}
    Isolez l'une des variables dans l'équation de contrainte (par ex. $h = \frac{V}{c^2}$ ou $L = C - 2l$) et remplacez-la dans la fonction objectif afin d'obtenir une fonction $\Phi(x)$ ne dépendant plus que d'\textbf{une seule variable indépendante $x$}, en précisant son intervalle de validité physique (souvent $x > 0$ ou $x \in [a, b]$).
    \item \textbf{Étape 4 -- Dériver par rapport à cette variable et résoudre $\Phi'(x^*) = 0$ :}
    Calculez $\Phi'(x) = \dfrac{\mathrm{d}\Phi}{\mathrm{d}x}$, factorisez-la et résolvez $\Phi'(x^*) = 0$ sur l'intervalle physique.
    \item \textbf{Étape 5 -- Vérifier la nature de l'extremum et répondre à TOUTES les sous-questions :}
    \begin{itemize}
        \item Justifiez brièvement que $x^*$ est bien un minimum (resp. un maximum), soit par le signe de la dérivée seconde $\Phi''(x^*) > 0$ (resp. $< 0$), soit par un petit tableau de signes de $\Phi'(x)$ autour de $x^*$.
        \item Calculez la valeur optimale $\Phi(x^*)$ ainsi que les autres dimensions demandées, et vérifiez les \textbf{unités} (mètres, hectares, heures/minutes, kg/min...).
    \end{itemize}
\end{enumerate}
\end{theoreme}

\begin{astuce}[Quatre astuces d'ingénieur qui simplifient radicalement les calculs]
\begin{enumerate}
    \item \textbf{Gardez les constantes sous forme de lettres jusqu'à la fin !} Comme le rappelle avec insistance le solutionnaire de \textbf{Juillet 2026 (Question 2)}, n'injectez jamais les grands nombres ($2300$ milles nautiques, $5$ millions d'euros, $f_0 = 1,146\cdot 10^{14}\text{ Hz}$...) au début du calcul : en travaillant avec des lettres ($C_{\text{total}}(V) = aV^2 + \frac{d}{V}$), la dérivée tient en une demi-ligne et les constantes inutiles s'éliminent d'elles-mêmes !
    \item \textbf{Optimiser le carré d'une fonction positive $A(x) \ge 0$ pour éliminer une racine carrée :} (Septembre 2026, Q2 -- Méthode 3). Si vous devez maximiser une grandeur positive contenant une racine carrée, telle que l'aire d'un triangle $A(b) = \frac{1}{2}b\sqrt{(C/2)^2 - (b/2)^2}$, remarquez que la fonction carré $u \mapsto u^2$ est strictement croissante sur $\R^+$ : \textbf{$A(b)$ atteint son maximum exactement au même point $b^*$ que son carré $A^2(b) = \frac{1}{16}(b^2 C^2 - b^4)$}, qui est un simple polynôme sans aucune racine carrée à dériver !
    \item \textbf{Penser au paramétrage trigonométrique dans les problèmes géométriques :} ( Juillet 2023 Série 1 Q2 et Septembre 2026 Q2 -- Méthodes 5 et 6). Dans un problème de triangle ou de trajet oblique, choisir l'\textbf{angle $\theta$} comme variable plutôt qu'une longueur $x$ transforme souvent les racines carrées de Pythagore en simples $\sin\theta$ et $\cos\theta$ !
    \item \textbf{Quand la fonction objectif est définie par une intégrale :} (Juillet 2022 Q2, Septembre 2022 Q2, Septembre 2023 Q2, Juillet 2024 Q2). Calculez d'abord l'intégrale définie par rapport au temps $t$ (ou à l'espace $x$) pour obtenir une expression explicite en fonction du paramètre à optimiser ($v$, $a$, $f_s$), puis dérivez ce résultat par rapport au paramètre !
\end{enumerate}
\end{astuce}

% ----------------------------------------------------------------------------
\section{Optimisation de coûts, d'énergie et de vitesse (Modèles $a\,x^p + b\,x^{-q}$)}
% ----------------------------------------------------------------------------

Une grande famille de problèmes d'ingénierie consiste à arbitrer entre deux effets antagonistes lorsque l'on fait varier un paramètre $x > 0$ (épaisseur d'isolant, vitesse d'un navire ou d'un drone) : un terme qui croît avec $x$ ($a\,x^p$) et un terme qui décroît avec $x$ ($\frac{b}{x^q}$).

\begin{exemple}[compromis coût de production / déperdition ou temps : Juillet 2021 Q2, Juillet 2022 Q2 et Juillet 2026 Q2]
\textbf{1) Épaisseur optimale d'un isolant thermique (UCLouvain -- Juillet 2021, Question 2) :}\\
\emph{Énoncé résumé :} On souhaite isoler un bâtiment pour une durée de $20\text{ ans}$ afin de minimiser la dépense énergétique totale par $\text{m}^2$ de mur (somme de la déperdition de chauffage et de l'énergie de fabrication de l'isolant) :
\begin{itemize}
    \item La déperdition annuelle par $\text{m}^2$ est inversement proportionnelle à l'épaisseur $e$ (en m) de l'isolant (soit de la forme $\frac{k}{e}$) et vaut $5\text{ kWh}/(\text{m}^2\cdot\text{an})$ pour $e = 0,1\text{ m}$.
    \item La fabrication de l'isolant demande $1000\text{ kWh}$ par $\text{m}^3$ d'isolant.
\end{itemize}
Quelle épaisseur $e$ permet de minimiser l'énergie totale sur $20\text{ ans}$ ?\\
\emph{Résolution :}
\begin{itemize}
    \item \textbf{Mise en équation pour $1\text{ m}^2$ de mur sur $20\text{ ans}$ en fonction de l'épaisseur $e > 0$ :}
    \begin{itemize}
        \item La déperdition annuelle vaut $\frac{k}{e}$. Comme elle vaut $5\text{ kWh}$ pour $e = 0,1\text{ m}$, on a $k = 5 \times 0,1 = 0,5\text{ kWh}\cdot\text{m/an}$. Sur $20\text{ ans}$, la déperdition totale par $\text{m}^2$ vaut donc $20 \times \frac{0,5}{e} = \dfrac{10}{e}\text{ kWh}$.
        \item Pour $1\text{ m}^2$ de mur et une épaisseur $e$ (en m), le volume d'isolant est $1 \times e = e\text{ m}^3$, ce qui consomme à la fabrication $1000\,e\text{ kWh}$.
    \end{itemize}
    La dépense énergétique totale par $\text{m}^2$ sur $20\text{ ans}$ est donc la fonction d'une seule variable $e > 0$ :
    \[
        f(e) = 1000\,e + \frac{10}{e}.
    \]
    \item \textbf{Dérivation et optimisation :}
    \[
        f'(e) = 1000 - \frac{10}{e^2} = 0 \iff e^2 = \frac{10}{1000} = 0,01 \iff \boxed{e^* = 0,1\text{ m} = 10\text{ cm}} \; (e > 0).
    \]
    Comme $f''(e) = \frac{20}{e^3} > 0$ pour tout $e > 0$, il s'agit bien d'un minimum global.
\end{itemize}

\medskip
\textbf{2) Vitesse optimale d'un navire de transport maritime (UCLouvain -- Juillet 2026, Question 2) :}\\
\emph{Énoncé :} Une entreprise livre du matériel d'Anvers à Marseille (distance $D = 2300$ milles nautiques). Le coût total du voyage est $C_{\text{total}} = C_E + C_T$, où le coût énergétique augmente quadratiquement avec la vitesse $V$ du bateau ($C_E = aV^2$) et le coût salarial/amortissement est proportionnel à la durée $T$ du voyage ($C_T = bT$).
\begin{enumerate}[label=(\arabic*)]
    \item \textbf{Premier voyage :} Pour $V = 20\text{ nœuds}$, le voyage coûte $5$ millions d'euros avec $C_E = C_T / 2$. Cette vitesse $V = 20\text{ nœuds}$ était-elle optimale ? Quelle est $V_{\text{opt1}}$ ?
    \item \textbf{Second voyage :} Les paramètres $a$ et $b$ ont été réduits : le coût total vaut maintenant $4$ millions d'euros pour $V = 20\text{ nœuds}$, et vaudrait $\dfrac{17}{3}$ millions d'euros pour $V = 10\text{ nœuds}$. Quelle est la nouvelle vitesse optimale $V_{\text{opt2}}$ ?
\end{enumerate}
\emph{Résolution :}
\begin{itemize}
    \item \textbf{Modélisation générale pour tout voyage :} Comme la distance parcourue est $D = T \cdot V$, la durée du voyage est $T = \dfrac{D}{V}$. En posant la constante globale $d = bD$, le coût total s'exprime en fonction de la seule variable $V > 0$ :
    \[
        C_{\text{total}}(V) = aV^2 + b\,\frac{D}{V} = aV^2 + \frac{d}{V}.
    \]
    Dérivons par rapport à $V$ :
    \[
        \frac{\mathrm{d}C_{\text{total}}}{\mathrm{d}V} = 2aV - \frac{d}{V^2} = 0 \iff 2aV = \frac{d}{V^2} \iff \boxed{V_{\text{opt}}^3 = \frac{d}{2a}},
    \]
    soit $aV_{\text{opt}}^2 = \frac{1}{2}\frac{d}{V_{\text{opt}}}$, c'est-à-dire $C_E = \frac{1}{2}C_T$.
    La dérivée seconde $\frac{\mathrm{d}^2 C_{\text{total}}}{\mathrm{d}V^2} = 2a + \frac{2d}{V^3} > 0$ confirme que c'est un minimum.
    \item \textbf{Application au premier voyage :} La condition d'optimalité générale $aV^2 = \frac{d}{2V} \iff C_E = \frac{1}{2}C_T$ est \textbf{précisément celle qui est observée à $V = 20\text{ nœuds}$} dans l'énoncé ! Sans aucun calcul numérique, on conclut immédiatement que la vitesse de $20\text{ nœuds}$ était déjà optimale : $\boxed{V_{\text{opt1}} = 20\text{ nœuds}}$.
    \item \textbf{Application au second voyage :} Traduisons les deux coûts donnés pour $V = 20$ et $V = 10$ dans $C_{\text{total}}(V) = aV^2 + \frac{d}{V}$ :
    \[
        \begin{cases}
            400\,a + \dfrac{d}{20} = 4 \\[6pt]
            100\,a + \dfrac{d}{10} = \dfrac{17}{3}
        \end{cases}
    \]
    Multiplions la seconde équation par $4$ ($400\,a + \frac{4d}{10} = \frac{68}{3}$) et soustrayons la première pour éliminer $400a$ :
    \[
        \left(\frac{4}{10} - \frac{1}{20}\right)d = \frac{68}{3} - 4 \iff \frac{7}{20}d = \frac{56}{3} \iff d = \frac{56 \times 20}{3 \times 7} = \frac{160}{3}.
    \]
    En remplaçant dans la seconde équation : $100a = \frac{17}{3} - \frac{16}{3} = \frac{1}{3} \implies a = \frac{1}{300}$. On injecte dans la formule (5) de la vitesse optimale :
    \[
        V_{\text{opt2}}^3 = \frac{d}{2a} = \frac{160/3}{2/300} = \frac{160}{3} \times \frac{300}{2} = 80 \times 100 = 8000 = 20^3 \implies \boxed{V_{\text{opt2}} = 20\text{ nœuds}}.
    \]
\end{itemize}
\end{exemple}

% ----------------------------------------------------------------------------
\section{Optimisation géométrique et cinématique (Aires, volumes, trajets)}
% ----------------------------------------------------------------------------

\begin{exemple}[Trois problèmes géométriques et cinématiques d'annales : Septembre 2021 Q2, Juillet 2023 Q2, Septembre 2026 Q2]
\textbf{1) Dimensionnement d'une boîte métallique avec et sans couvercle (UCLouvain -- Septembre 2021, Question 2) :}\\
\emph{Énoncé résumé :} On conçoit une boîte métallique de volume fixé $V = 1\text{ dm}^3$ ayant la forme d'un parallélépipède rectangle à base carrée (côté de base $c$, hauteur $h$). La quantité de métal est $M = \alpha S$ où $S$ est la surface totale des parois métalliques.
\begin{enumerate}[label=\alph*)]
    \item Quelles dimensions $(c, h)$ minimisent la quantité de métal $M_1$ pour une boîte fermée (avec couvercle) ? Quelle est sa forme ?
    \item Comme les boîtes sont empilées, on supprime le couvercle supérieur. Si l'on garde les dimensions $(c, h)$ trouvées en a), quelle est la nouvelle quantité de métal $M_2$ ?
    \item Si l'on ré-optimise les dimensions $(c, h)$ en tenant compte de l'absence de couvercle, quelle quantité minimale de métal $M_3$ obtient-on ?
\end{enumerate}
\emph{Résolution :}
\begin{enumerate}[label=\alph*)]
    \item \textbf{Boîte fermée (2 bases carrées d'aire $c^2$ + 4 faces latérales d'aire $ch$) :}
    La contrainte de volume est $V = c^2 h \implies h = \dfrac{V}{c^2}$. En remplaçant $h$ dans $M(c, h) = \alpha(2c^2 + 4ch)$ :
    \begin{align*}
        M(c) &= \alpha\left(2c^2 + 4c\,\frac{V}{c^2}\right) = \alpha\left(2c^2 + \frac{4V}{c}\right) \\
        \implies M'(c) &= \alpha\left(4c - \frac{4V}{c^2}\right) = 0 \iff c^3 = V \iff \boxed{c = \sqrt[3]{V}}.
    \end{align*}
    On en déduit $h = \frac{V}{c^2} = \frac{V}{V^{2/3}} = \sqrt[3]{V} = c$ : la boîte optimale est un \textbf{cube} ! Et $M_1 = \alpha(2V^{2/3} + 4V^{2/3}) = \boxed{6\alpha V^{2/3}}$.
    \item \textbf{Suppression du couvercle à dimensions inchangées ($c = h = V^{1/3}$) :} Il n'y a plus qu'une seule base carrée ($S = c^2 + 4ch = 5c^2$), donc $\boxed{M_2 = 5\alpha V^{2/3}}$ (économie de $1/6 \approx 16,7\%$).
    \item \textbf{Ré-optimisation sans couvercle ($S = c^2 + 4ch$) :} Avec $h = \frac{V}{c^2}$, la nouvelle fonction à minimiser est :
    \[
        M(c) = \alpha\left(c^2 + \frac{4V}{c}\right) \implies M'(c) = \alpha\left(2c - \frac{4V}{c^2}\right) = 0 \iff \boxed{c^* = \sqrt[3]{2V} = 2^{1/3}V^{1/3}}.
    \]
    La nouvelle hauteur optimale vaut $h^* = \frac{V}{(c^*)^2} = \frac{V}{2^{2/3}V^{2/3}} = \frac{1}{2^{2/3}}V^{1/3} = \sqrt[3]{\frac{V}{4}} = \frac{c^*}{2}$ (un demi-cube !), et la quantité de métal minimale vaut :
    \[
        M_3 = \alpha\left(2^{2/3}V^{2/3} + \frac{4V}{2^{1/3}V^{1/3}}\right) = \alpha\bigl(2^{2/3} + 2\cdot 2^{2/3}\bigr)V^{2/3} = \boxed{3\sqrt[3]{2}\,\alpha V^{2/3} \approx 4,76\,\alpha V^{2/3}}.
    \]
\end{enumerate}

\medskip
\textbf{2) Trajet le plus rapide sable / route -- Loi de Snell-Descartes (UCLouvain -- Juillet 2023 -- Série 1, Question 2) :}\\
\emph{Énoncé résumé :} Vous êtes dans le sable à $1\text{ km}$ au nord d'une route est-ouest, et votre destination est également dans le sable à $1\text{ km}$ au nord de la route et à une distance $L$ à l'est de votre position. Vous marchez à $v_1 = 3\text{ km/h}$ dans le sable et à $v_2 = 5\text{ km/h}$ sur la route. Vous suivez une diagonale jusqu'à un point de la route décalé de $x\text{ km}$ vers l'est, puis la route sur une distance $L - 2x$, puis une diagonale symétrique jusqu'à la destination.
\begin{enumerate}[label=\alph*)]
    \item Quelle est la durée minimale de votre trajet ?
    \item Pour quelles distances $L$ aurait-il été plus rapide de marcher directement en ligne droite dans le sable ?
\end{enumerate}
\begin{center}
\begin{tikzpicture}[scale=1.1, >=Stealth]
    \draw[very thick, darkgray] (-0.5,0) -- (6.5,0) node[midway, below=3pt] {\small Route ($v_2 = 5\text{ km/h}$, distance parcourue $L - 2x$)};
    \coordinate (A) at (0.5, 1.8);
    \coordinate (B) at (5.5, 1.8);
    \coordinate (A0) at (0.5, 0);
    \coordinate (B0) at (5.5, 0);
    \coordinate (P1) at (1.85, 0);
    \coordinate (P2) at (4.15, 0);

    \draw[dashed, gray] (A) -- (A0) node[midway, left] {\small $1\text{ km}$};
    \draw[dashed, gray] (B) -- (B0) node[midway, right] {\small $1\text{ km}$};
    \draw[dotted, gray, <->] (0.5, 2.05) -- (5.5, 2.05) node[midway, above] {\small $L$};
    \draw[very thick, uclblue] (A) -- (P1) node[midway, above right=-2pt] {\footnotesize $\sqrt{1+x^2}$};
    \draw[very thick, uclblue] (P1) -- (P2);
    \draw[very thick, uclblue] (P2) -- (B) node[midway, above left=-2pt] {\footnotesize $\sqrt{1+x^2}$};

    \draw[<->, warnred] (0.5, -0.25) -- (1.85, -0.25) node[midway, below] {\small $x$};
    \draw[<->, warnred] (4.15, -0.25) -- (5.5, -0.25) node[midway, below] {\small $x$};

    \fill[uclblue] (A) circle (2.2pt) node[above left] {\small Départ};
    \fill[uclblue] (B) circle (2.2pt) node[above right] {\small Destination};
    \fill[warnred] (P1) circle (1.8pt);
    \fill[warnred] (P2) circle (1.8pt);
\end{tikzpicture}
\end{center}
\emph{Résolution :}
\begin{enumerate}[label=\alph*)]
    \item Par le théorème de Pythagore, chaque diagonale dans le sable mesure $\sqrt{1 + x^2}\text{ km}$. Parcourir les deux diagonales à $3\text{ km/h}$ prend un temps $\frac{2\sqrt{1+x^2}}{3}\text{ h}$, et parcourir le segment de route $L - 2x$ à $5\text{ km/h}$ prend $\frac{L-2x}{5}\text{ h}$. Le temps total $T(x)$ vaut donc :
    \[
        T(x) = \frac{1}{5}(L - 2x) + \frac{2}{3}\sqrt{1 + x^2} \implies T'(x) = -\frac{2}{5} + \frac{2}{3}\cdot\frac{\cancel{2}x}{\cancel{2}\sqrt{1+x^2}} = -\frac{2}{5} + \frac{2x}{3\sqrt{1+x^2}}.
    \]
    Annulons la dérivée (avec $x \ge 0$) :
    \[
        T'(x) = 0 \iff 5x = 3\sqrt{1+x^2} \iff 16x^2 = 9 \iff \boxed{x^* = \frac{3}{4}\text{ km}}.
    \]
    En remplaçant $x^* = \frac{3}{4}$ (notons le triplet pythagoricien $\sqrt{1 + (3/4)^2} = \sqrt{25/16} = \frac{5}{4}$) dans $T(x)$ :
    \[
        T(x^*) = \frac{1}{5}\left(L - \frac{3}{2}\right) + \frac{2}{3} \cdot \frac{5}{4} = \frac{L}{5} - \frac{3}{10} + \frac{5}{6} = \frac{L}{5} + \frac{-9 + 25}{30} = \boxed{\frac{1}{5}L + \frac{8}{15}\text{ heures}}.
    \]
    \emph{Remarque :} Si l'on paramètre par l'angle d'incidence $\theta$, la condition $T'(\theta) = 0$ donne immédiatement $\sin\theta = \frac{3}{5} = \frac{v_1}{v_2}$, qui n'est autre que la loi de la réfraction de Snell-Descartes en optique !
    \item Marcher en ligne droite dans le sable sur la distance $L$ à $3\text{ km/h}$ prend un temps $T_{\text{direct}} = \frac{1}{3}L$. Ce trajet direct est plus rapide si et seulement si :
    \[
        \frac{1}{3}L < \frac{1}{5}L + \frac{8}{15} \iff \left(\frac{5 - 3}{15}\right)L < \frac{8}{15} \iff \frac{2}{15}L < \frac{8}{15} \iff \boxed{L < 4\text{ km}}.
    \]
\end{enumerate}

\medskip
\textbf{3) Enclos agricole rectangulaire vs triangulaire adossé à un mur (UCLouvain -- Septembre 2026, Question 2) :}\\
\emph{Énoncé résumé :} Une agricultrice dispose de $C = 200\text{ m}$ de clôture pour délimiter un enclos adossé à un mur rectiligne (le mur ne nécessite pas de clôture).
\begin{enumerate}[label=(\arabic*)]
    \item \textbf{Cas rectangulaire :} Quelles doivent être la longueur $L$ (parallèle au mur) et la largeur $l$ (perpendiculaire au mur) pour maximiser l'aire $A_{\text{rectangle}}$ ? Que vaut $A_{\text{rectangle}}$ en $\text{m}^2$ et en hectares ($1\text{ ha} = 10\,000\text{ m}^2$) ?
    \item \textbf{Cas triangulaire :} Pour économiser un pieu, elle envisage un enclos en forme de triangle isocèle dont la base $b$ est le mur (les deux côtés clôturés mesurant chacun $C/2 = 100\text{ m}$). Quelles sont la base $b$ et la hauteur $h$ qui maximisent $A_{\text{triangle}}$, et que vaut $A_{\text{triangle}}$ ? Que doit-elle choisir ? (\textbf{Aide :} $\sqrt{1/2} \simeq 0,7071$).
\end{enumerate}
\emph{Résolution :}
\begin{enumerate}[label=(\arabic*)]
    \item \textbf{Rectangle :} La clôture couvre la longueur $L$ et les deux largeurs $l$, soit $L + 2l = C \implies L = C - 2l$. L'aire en fonction de $l$ est :
    \[
        A_{\text{rect}}(l) = L\cdot l = (C - 2l)l = Cl - 2l^2 \implies \frac{\mathrm{d}A_{\text{rect}}}{\mathrm{d}l} = C - 4l = 0 \iff \boxed{l = \frac{C}{4} = 50\text{ m}},
    \]
    (avec dérivée seconde $-4 < 0$), d'où $\boxed{L = C - 2l = \frac{C}{2} = 100\text{ m}}$ et $\boxed{A_{\text{rectangle}} = 100 \times 50 = \frac{C^2}{8} = 5000\text{ m}^2 = 0,5\text{ ha}}$.
    \item \textbf{Triangle isocèle (présentons 3 des 6 méthodes du solutionnaire officiel !) :}
    Chaque côté latéral clôturé mesure $\frac{C}{2} = 100\text{ m}$. Par Pythagore dans le demi-triangle rectangle de base $\frac{b}{2}$ et de hauteur $h$, on a $\left(\frac{b}{2}\right)^2 + h^2 = \left(\frac{C}{2}\right)^2$, et l'aire vaut $A_{\text{tri}} = \frac{1}{2}bh$.
    \begin{itemize}
        \item \emph{Méthode analytique standard (par la dérivée de $A_{\text{tri}}(b)$ ou de son carré $A_{\text{tri}}^2(b)$) :} Comme $h = \sqrt{(C/2)^2 - (b/2)^2}$, on a $A_{\text{tri}}^2(b) = \frac{1}{4}b^2\left(\frac{C^2}{4} - \frac{b^2}{4}\right) = \frac{1}{16}(C^2 b^2 - b^4)$. En dérivant ce polynôme par rapport à $b$ :
        \[
            \frac{\mathrm{d}(A_{\text{tri}}^2)}{\mathrm{d}b} = \frac{b}{8}(C^2 - 2b^2) = 0 \iff \boxed{b = \frac{C}{\sqrt{2}} = 100\sqrt{2} \simeq 141,42\text{ m}},
        \]
        d'où $\boxed{h = \sqrt{\frac{C^2}{4} - \frac{C^2}{8}} = \frac{C}{2\sqrt{2}} = 50\sqrt{2} \simeq 70,71\text{ m}}$ et $\boxed{A_{\text{triangle}} = \frac{1}{2}\cdot\frac{C}{\sqrt{2}}\cdot\frac{C}{2\sqrt{2}} = \frac{C^2}{8} = 5000\text{ m}^2 = 0,5\text{ ha}}$.
        \item \emph{Méthode géométrique fulgurante (Méthode 6 du solutionnaire) :} Notons $\beta$ l'angle au sommet entre les deux côtés clôturés de longueur fixe $\frac{C}{2} = 100\text{ m}$. L'aire d'un triangle de côtés $\frac{C}{2}$ et $\frac{C}{2}$ entourant l'angle $\beta$ vaut :
        \[
            A_{\text{triangle}} = \frac{1}{2}\left(\frac{C}{2}\right)\left(\frac{C}{2}\right)\sin\beta = \frac{C^2}{8}\sin\beta.
        \]
        Elle est immédiatement maximale lorsque $\sin\beta = 1 \iff \beta = 90\degree$ (triangle rectangle isocèle d'hypoténuse $b = \sqrt{2}\cdot\frac{C}{2} = \frac{C}{\sqrt{2}}$ et d'aire maximale $\frac{C^2}{8}$) !
    \end{itemize}
    \textbf{Conclusion :} Les deux aires maximales sont \textbf{strictement égales} ($5000\text{ m}^2 = 0,5\text{ ha}$), donc en tenant compte de l'économie d'un pieu, l'agricultrice a intérêt à choisir \textbf{le triangle} !
\end{enumerate}
\end{exemple}

% ----------------------------------------------------------------------------
\section{Optimisation de fonctions exponentielles, logarithmiques et intégrales}
% ----------------------------------------------------------------------------

\begin{exemple}[Quatre problèmes avancés d'annales : Septembre 2022 Q2, Juillet 2023 Série 2 Q2, Juillet 2024 Q2, Septembre 2025 Q2]
\textbf{1) Conception d'une pièce mécanique de volume minimal par intégration (UCLouvain -- Septembre 2022, Question 2) :}\\
\emph{Énoncé :} Vous concevez une pièce mécanique de longueur $2$ obtenue par révolution d'une fonction $f : [-1, 1] \to \R$ autour de l'axe des abscisses, et devez choisir $f$ de façon à minimiser le volume $V = \pi \int_{-1}^1 f^2(x)\dd x$ sous deux contraintes :
(i) le rayon aux extrémités doit valoir $1$, soit $f(1) = f(-1) = 1$ ;
(ii) $f$ doit être un polynôme du second degré $f(x) = ax^2 + bx + c$ (qui peut prendre des valeurs négatives).
Quelle fonction $f(x)$ utiliserez-vous et quel sera le volume minimal $V$ ?\\
\emph{Résolution :}
\begin{itemize}
    \item \textbf{Élimination de $b$ et $c$ grâce aux contraintes aux bords :}
    \[
        \begin{cases} f(1) = a + b + c = 1 \\ f(-1) = a - b + c = 1 \end{cases} \implies \begin{cases} b = 0 \\ c = 1 - a \end{cases} \implies f(x) = ax^2 + (1 - a).
    \]
    \item \textbf{Calcul du volume $V(a)$ en fonction du seul paramètre $a \in \R$ :} Par parité de $f^2(x)$ sur $[-1, 1]$ :
    \begin{align*}
        V(a) &= \pi \int_{-1}^1 \bigl(ax^2 + (1-a)\bigr)^2\dd x = \pi \int_{-1}^1 \bigl(a^2 x^4 + 2a(1-a)x^2 + (1-a)^2\bigr)\dd x \\
        &= \pi\left[a^2\cdot\frac{2}{5} + 2a(1-a)\cdot\frac{2}{3} + (1-a)^2\cdot 2\right] \\
        &= \frac{\pi}{15}\bigl(6a^2 + 20a(1-a) + 30(1 - 2a + a^2)\bigr) = \frac{\pi}{15}\bigl(16a^2 - 40a + 30\bigr).
    \end{align*}
    \item \textbf{Optimisation par rapport à $a$ :}
    \[
        V'(a) = \frac{\pi}{15}(32a - 40) = 0 \iff \boxed{a^* = \frac{40}{32} = \frac{5}{4}} \quad (\text{minimum car } V''(a) = \frac{32\pi}{15} > 0).
    \]
    On obtient $c^* = 1 - \frac{5}{4} = -\frac{1}{4}$, soit $\boxed{f(x) = \frac{5}{4}x^2 - \frac{1}{4}}$ et le volume minimal $V\left(\frac{5}{4}\right) = \frac{\pi}{15}\left(16\cdot\frac{25}{16} - 40\cdot\frac{5}{4} + 30\right) = \frac{\pi}{15}(25 - 50 + 30) = \boxed{\dfrac{\pi}{3}}$.
\end{itemize}

\medskip
\textbf{2) Accélérateur de particules et vitesse maximale (UCLouvain -- Juillet 2023 -- Série 2, Question 2) :}\\
\emph{Énoncé :} La position d'une particule au temps $t \in \R$ (en ns) est $x(t) = b\,e^{-\frac{1}{2}at^2}$ et sa vitesse est $v(t) = x'(t)$.
(i) À $t = 0$, la particule passe en $x = 1$.
(ii) Elle est éclairée par un laser $1\text{ ns}$ \emph{avant} d'arriver en $t = 0$ (donc à l'instant $t = -1$), à une position $x^* = x(-1)$.
(iii) Elle doit aller le plus vite possible au moment où elle est éclairée. Déterminer $a$, $b$ et $x^*$.\\
\emph{Résolution :}
\begin{itemize}
    \item Condition à $t = 0$ : $x(0) = b\,e^0 = b = 1 \implies \boxed{b = 1}$, donc $x(t) = e^{-\frac{1}{2}at^2}$ et $x^* = x(-1) = e^{-\frac{1}{2}a}$.
    \item La vitesse à l'instant $t$ est $v(t) = x'(t) = -at\,e^{-\frac{1}{2}at^2}$. À l'instant de l'éclairage $t = -1$, la vitesse de la particule vaut donc, en fonction du paramètre $a$ :
    \[
        v_e(a) = v(-1) = a\,e^{-\frac{1}{2}a}.
    \]
    \item Maximisons $v_e(a)$ en dérivant par rapport à $a$ (règle du produit) :
    \[
        \frac{\mathrm{d}v_e}{\mathrm{d}a} = 1\cdot e^{-\frac{1}{2}a} + a\left(-\frac{1}{2}\right)e^{-\frac{1}{2}a} = e^{-\frac{1}{2}a}\left(1 - \frac{1}{2}a\right) = 0 \iff \boxed{a = 2} \quad (\text{car } e^{-\frac{1}{2}a} > 0).
    \]
    Comme $\frac{\mathrm{d}v_e}{\mathrm{d}a}$ est $> 0$ pour $a < 2$ et $< 0$ pour $a > 2$, $a = 2$ est bien un maximum global, et la position du laser est $\boxed{x^* = e^{-\frac{1}{2}(2)} = e^{-1} = \dfrac{1}{e}}$.
\end{itemize}

\medskip
\textbf{3) Panneau solaire photovoltaïque et loi de Planck / Wien (UCLouvain -- Juillet 2024, Question 2) :}\\
\emph{Énoncé résumé :} L'énergie électrique totale générée par un panneau solaire de fréquence seuil $f_s > 0$ (dans l'approximation de Wien $K = 0$) est donnée par :
\[
    E(f_s) = S\,C_1\,C_2\,f_s \int_{f_s}^\infty f^{n-1}\,e^{-f/f_0}\dd f, \qquad \text{où } f_0 = 1,146\cdot 10^{14}\text{ Hz}.
\]
Pour $n = 2$, calculer l'intégrale (par changement de variable $x = f/f_0$ et intégration par parties), puis déterminer la fréquence seuil optimale $f_s$ qui maximise $E(f_s)$ (avec $\sqrt{5} \simeq 2,236$).\\
\emph{Résolution :}
\begin{itemize}
    \item \textbf{Calcul de l'intégrale pour $n = 2$ :} Posons $x = \frac{f}{f_0} \implies f = f_0 x$ et $\mathrm{d}f = f_0\dd x$. Lorsque $f$ varie de $f_s$ à $\infty$, $x$ varie de $x_s = \frac{f_s}{f_0}$ à $\infty$ :
    \[
        E(f_s) = S\,C_1\,C_2\,f_s\,f_0^2 \int_{f_s/f_0}^\infty x\,e^{-x}\dd x.
    \]
    Une intégration par parties ($u = x \implies \mathrm{d}u = \mathrm{d}x$ ; $\mathrm{d}v = e^{-x}\dd x \implies v = -e^{-x}$) donne la primitive $\int x\,e^{-x}\dd x = -x\,e^{-x} + \int e^{-x}\dd x = -(x+1)e^{-x}$. Comme $\lim_{x \to \infty} -(x+1)e^{-x} = 0$ (croissance comparée), il vient :
    \[
        E(f_s) = S\,C_1\,C_2\,f_0^3\,\left[\frac{f_s}{f_0}\left(\frac{f_s}{f_0} + 1\right)e^{-f_s/f_0}\right] = K_0\,\bigl(x_s^2 + x_s\bigr)e^{-x_s} \quad \text{où } x_s = \frac{f_s}{f_0} > 0.
    \]
    \item \textbf{Maximisation par rapport à $x_s = f_s / f_0$ :} Dérivons $\phi(x_s) = (x_s^2 + x_s)e^{-x_s}$ par rapport à $x_s$ :
    \[
        \phi'(x_s) = (2x_s + 1)e^{-x_s} - (x_s^2 + x_s)e^{-x_s} = \bigl(-x_s^2 + x_s + 1\bigr)e^{-x_s} = 0 \iff x_s^2 - x_s - 1 = 0.
    \]
    On retrouve l'équation du nombre d'or ! Comme $x_s > 0$, l'unique racine positive est :
    \[
        \frac{f_s}{f_0} = \frac{1 + \sqrt{5}}{2} \simeq 1,618 \implies \boxed{f_s = \frac{1+\sqrt{5}}{2}f_0 \simeq 1,618 \times 1,146\cdot 10^{14} \simeq 1,85\cdot 10^{14}\text{ Hz}}.
    \]
\end{itemize}
\end{exemple}

% ----------------------------------------------------------------------------
\section{Exercices du Chapitre 5}
% ----------------------------------------------------------------------------

\emph{Les corrigés détaillés de tous ces exercices se trouvent au Chapitre 9, en fin d'ouvrage.}

\begin{enumerate}[label=\textbf{Exercice 5.\arabic* --}]
    \item \textbf{(UCLouvain -- Juillet 2022, Q2 -- Consommation énergétique d'un drone-lanceur)} Un drone-lanceur monte verticalement du sol ($x=0$) jusqu'à l'altitude $H$ à vitesse constante $v$, de sorte que $x(t) = vt$ pour $t \in [0, T]$. Durant la montée, sa puissance consommée par unité de temps est $Q(t) = K + C_x(t)v^3$, où le coefficient de traînée diminue avec l'altitude selon $C_x(t) = \dfrac{C}{x(t) + x_0}$ ($K, C, x_0, H > 0$ sont des constantes).
    \begin{enumerate}
        \item À quel instant $T$ le lanceur atteint-il l'altitude $H$ ?
        \item En intégrant $Q(t)$ entre $0$ et $T$, calculer la consommation énergétique totale $E(v)$ de la montée en fonction de $v$ et des constantes $K, C, x_0, H$.
        \item Quelle vitesse $v^*$ faut-il choisir pour minimiser cette consommation totale $E(v)$ ? Quelle est alors la consommation minimale $E(v^*)$ ?
    \end{enumerate}

    \item \textbf{(UCLouvain -- Septembre 2023, Q2 -- Éclairage optimal dans un train de nuit)} L'intensité lumineuse dans un train de nuit au temps $t \in \left[-\frac{\pi}{2}, \frac{\pi}{2}\right]$ est $\ell(t) = a + b\bigl(1 + \sin(t)\bigr)$, où $b > 0$ est donné et $a > 0$ est l'éclairage d'appoint à choisir. Le degré d'inconfort instantané est modélisé par :
    \[
        c(t) = \ell(t) + \frac{|\ell'(t)|}{\ell(t)},
    \]
    et l'inconfort moyen sur $\left[-\frac{\pi}{2}, \frac{\pi}{2}\right]$ est $\bar{c}(a) = \dfrac{1}{\pi}\displaystyle\int_{-\pi/2}^{\pi/2} c(t)\dd t$.
    \begin{enumerate}
        \item En remarquant que $\ell'(t) = b\cos(t) \ge 0$ sur $\left[-\frac{\pi}{2}, \frac{\pi}{2}\right]$, calculer explicitement l'intégrale donnant $\bar{c}(a)$ en fonction de $a$ et $b$.
        \item Déterminer la valeur de $a > 0$ qui minimise l'inconfort moyen $\bar{c}(a)$.
        \item Même question si l'intensité vaut $\ell(t) = a + b\bigl(1 + \sin(\lambda t)\bigr)$ avec une fréquence $\lambda > 0$ sur l'intervalle $\left[-\frac{\pi}{2\lambda}, \frac{\pi}{2\lambda}\right]$.
    \end{enumerate}

    \item \textbf{(UCLouvain -- Septembre 2024, Q2 -- Le thermos de café avant l'examen)} Vous buvez un café à l'instant $t = 0$ pour que la concentration de caféine dans votre sang soit maximale exactement $60\text{ minutes}$ après le début de l'examen. La concentration à l'instant $t \ge 0$ (en heures) vaut :
    \[
        c(t) = \frac{D}{1 - \beta/\alpha}\Bigl((1 - e^{-\alpha t}) - (1 - e^{-\beta t})\Bigr),
    \]
    avec $\alpha = 1,0\text{ h}^{-1}$ et $\beta = 0,2\text{ h}^{-1}$. (\textbf{Aide :} $\ln(5) \simeq 1,61$).
    \begin{enumerate}
        \item Combien de minutes avant le début de l'examen devez-vous boire votre café ?
        \item Prouver par le calcul de la dérivée seconde $\left.\dfrac{\mathrm{d}^2 c}{\mathrm{d}t^2}\right|_{t_{\max}}$ qu'il s'agit bien d'un maximum (en justifiant le signe de la constante $D$).
    \end{enumerate}

    \item \textbf{(UCLouvain -- Juillet 2025, Q2 -- Le tournoi de la ligue des « legendgénieurs »)} Vous disposez d'un budget de $560\text{ bitcoins}$ et de $7$ points de vie initiaux. Vous répartissez votre budget entre l'achat d'une armure (coût $A$ bitcoins) et l'achat de points de vie supplémentaires (coût $V$ bitcoins, avec $A + V = 560$, sachant qu'un point de vie coûte $10$ bitcoins). L'efficacité de l'armure est proportionnelle au carré de son coût ($E = kA^2$) et vaut $E = 1$ si $A = 560$. Votre durée de vie en jours est le produit de votre nombre total de points de vie par l'efficacité de l'armure : $D = \left(7 + \dfrac{V}{10}\right)E$.
    Déterminer la répartition optimale $(A, V)$ qui maximise $D$, et exprimer la durée de vie maximale $D_{\max}$ sous forme de fraction, puis en jours $J$, heures $H$ et minutes $M$.

    \item \textbf{(UCLouvain -- Septembre 2025, Q2 -- Le moteur propulseur spatial et l'effet d'embouteillage)} Un moteur spatial émet un flux de particules $F$ (en $\text{g/s}$). La vitesse des particules décroît avec le flux selon $V(F) = V_{\max}\exp(-F/F_c)$, où $V_{\max} = 100\text{ km/s}$ et $F_c = 100\text{ g/s}$. L'accélération du vaisseau est $A(F) = k\,V(F)\,F$ ($k > 0$).
    \begin{enumerate}
        \item Quel flux optimal $F^*$ (en $\text{g/s}$ puis en $\text{kg/min}$) maximise l'accélération $A(F)$ ? Justifier qu'il s'agit d'un maximum.
        \item Quel devient le flux optimal $F^*$ (en $\text{g/s}$) si le facteur de réduction de vitesse est $\exp\bigl[-(F/F_c)^{1/2}\bigr]$ ?
    \end{enumerate}

    \item \textbf{(Volume maximal d'un cylindre inscrit dans une sphère ou un cône)} Parmi tous les cylindres circulaires droits de rayon $r$ et de hauteur $h$ inscrits dans une sphère de rayon donné $R > 0$, déterminer les dimensions $(r, h)$ et le volume maximal du cylindre.
\end{enumerate}

% ============================================================================
% CHAPITRE 6 : CALCUL INTÉGRAL ET PRIMITIVES
% ============================================================================
\chapter{Calcul Intégral et Primitives}

La \textbf{Question 3} de l'examen d'admission d'Analyse à l'UCLouvain (notée typiquement sur $10$ points) consiste à calculer une \textbf{primitive} $\int f(x)\dd x$ (en Juillet et Septembre 2021, il y avait deux ou trois primitives au sein de la Question 3 ; depuis 2022, une question entière est dédiée à une primitive riche que les solutionnaires officiels résolvent systématiquement de $2$ à $6$ façons différentes !). Ce chapitre détaille toutes les techniques de primitivation au programme : primitives immédiates, intégration par substitution (changement de variable), intégration par parties (IPP), fractions rationnelles et primitives trigonométriques.

% ----------------------------------------------------------------------------
\section{Définition d'une primitive et tableau des primitives immédiates}
% ----------------------------------------------------------------------------

\subsection{Définition et unicité à une constante près}

\begin{definition}[Primitive et intégrale indéfinie]
Soit $f : I \to \R$ une fonction définie sur un intervalle $I \subset \R$. Une fonction $F : I \to \R$ est appelée une \textbf{primitive} de $f$ sur $I$ si $F$ est dérivable sur $I$ et a pour dérivée $f$ :
\[
    \boxed{\forall x \in I, \quad F'(x) = f(x)}.
\]
Sur un intervalle $I$, si $F(x)$ est une primitive de $f(x)$, alors \textbf{toutes les primitives} de $f$ sur $I$ sont de la forme $F(x) + C$ où $C \in \R$ est une constante réelle arbitraire (car la différence de deux primitives a une dérivée nulle sur l'intervalle $I$, donc est constante par le TAF). On note l'ensemble des primitives (ou \textbf{intégrale indéfinie}) :
\[
    \boxed{\int f(x)\dd x = F(x) + C \quad (C \in \R)}.
\]
\end{definition}

\begin{astuce}[Le réflexe absolu : dériver sa primitive pour la vérifier !]
Contrairement au calcul de dérivée qui est purement mécanique, le calcul d'une primitive demande de l'intuition. Mais vous disposez d'un filet de sécurité infaillible : \textbf{une fois votre primitive $F(x)$ trouvée, dérivez-la au brouillon !} Si $F'(x)$ redonne bien $f(x)$, vous êtes certain à $100\%$ d'avoir tous les points de la question !
\end{astuce}

\subsection{Tableau des primitives immédiates et quasi-immédiates (formes composées)}

En lisant à l'envers le tableau des dérivées du Chapitre 3 (et la règle de dérivation des fonctions composées), on obtient le formulaire fondamental suivant :

\begin{theoreme}[Formulaire des primitives immédiates et composées ($C \in \R$)]
\begin{center}
\renewcommand{\arraystretch}{1.65}
\begin{tabular}{@{}ll@{}}
\toprule
\textbf{Forme élémentaire $\displaystyle\int f(x)\dd x$} & \textbf{Forme composée directe $\displaystyle\int f\bigl(u(x)\bigr)\,u'(x)\dd x$} \\
\midrule
$\displaystyle\int x^p\dd x = \frac{x^{p+1}}{p+1} + C \quad (p \neq -1)$ & $\displaystyle\int \bigl(u(x)\bigr)^p\,u'(x)\dd x = \frac{\bigl(u(x)\bigr)^{p+1}}{p+1} + C \quad (p \neq -1)$ \\
$\displaystyle\int \frac{1}{x}\dd x = \ln(|x|) + C \quad (x \neq 0)$ & $\displaystyle\int \frac{u'(x)}{u(x)}\dd x = \ln\bigl(|u(x)|\bigr) + C$ \\
$\displaystyle\int \frac{1}{\sqrt{x}}\dd x = 2\sqrt{x} + C \quad (x > 0)$ & $\displaystyle\int \frac{u'(x)}{\sqrt{u(x)}}\dd x = 2\sqrt{u(x)} + C \quad \left(\text{ou } \int \frac{u'}{2\sqrt{u}}\dd x = \sqrt{u} + C\right)$ \\
$\displaystyle\int \frac{1}{x^2}\dd x = -\frac{1}{x} + C \quad (x \neq 0)$ & $\displaystyle\int \frac{u'(x)}{\bigl(u(x)\bigr)^2}\dd x = -\frac{1}{u(x)} + C$ \\
$\displaystyle\int e^x\dd x = e^x + C$ & $\displaystyle\int u'(x)\,e^{u(x)}\dd x = e^{u(x)} + C$ \\
$\displaystyle\int \cos(x)\dd x = \sin(x) + C$ & $\displaystyle\int u'(x)\,\cos\bigl(u(x)\bigr)\dd x = \sin\bigl(u(x)\bigr) + C$ \\
$\displaystyle\int \sin(x)\dd x = -\cos(x) + C$ & $\displaystyle\int u'(x)\,\sin\bigl(u(x)\bigr)\dd x = -\cos\bigl(u(x)\bigr) + C$ \\
$\displaystyle\int \frac{1}{\cos^2(x)}\dd x = \int \bigl(1+\tan^2(x)\bigr)\dd x = \tan(x) + C$ & $\displaystyle\int \frac{u'(x)}{\cos^2(u(x))}\dd x = \tan\bigl(u(x)\bigr) + C$ \\
$\displaystyle\int \frac{1}{\sin^2(x)}\dd x = \int \bigl(1+\cot^2(x)\bigr)\dd x = -\cot(x) + C$ & $\displaystyle\int \frac{u'(x)}{\sin^2(u(x))}\dd x = -\cot\bigl(u(x)\bigr) + C$ \\
$\displaystyle\int \frac{1}{1 + x^2}\dd x = \arctan(x) + C$ & $\displaystyle\int \frac{u'(x)}{1 + \bigl(u(x)\bigr)^2}\dd x = \arctan\bigl(u(x)\bigr) + C$ \\
$\displaystyle\int \frac{1}{\sqrt{1 - x^2}}\dd x = \arcsin(x) + C \quad (|x| < 1)$ & $\displaystyle\int \frac{u'(x)}{\sqrt{1 - \bigl(u(x)\bigr)^2}}\dd x = \arcsin\bigl(u(x)\bigr) + C$ \\
\bottomrule
\end{tabular}
\end{center}
\end{theoreme}

% ----------------------------------------------------------------------------
\section{Intégration par substitution (changement de variable)}
% ----------------------------------------------------------------------------

La méthode d'\textbf{intégration par substitution} (ou changement de variable) est la traduction intégrale de la règle de dérivation des fonctions composées $\bigl(F(u(x))\bigr)' = F'(u(x))\cdot u'(x) = f(u(x))\,u'(x)$.

\begin{theoreme}[Règle d'intégration par substitution]
Pour calculer une primitive de la forme $\displaystyle\int f\bigl(u(x)\bigr)\,u'(x)\dd x$ :
\begin{enumerate}
    \item On pose la nouvelle variable $\boxed{u = u(x)}$ et on calcule sa différentielle $\boxed{\mathrm{d}u = u'(x)\dd x}$.
    \item On remplace toute l'expression en $x$ par une intégrale portant uniquement sur $u$ :
    \[
        \boxed{\int f\bigl(u(x)\bigr)\,u'(x)\dd x = \int f(u)\dd u = F(u) + C = F\bigl(u(x)\bigr) + C}.
    \]
    \item \emph{Variante inverse ($x = \varphi(u)$) :} Parfois (notamment pour éliminer une racine $\sqrt{a^2 - x^2}$ ou un dénominateur $(x^2+1)^5$), on pose $x = \varphi(u)$ (par ex. $x = a\sin u$ ou $x = \tan u$), on remplace $\mathrm{d}x = \varphi'(u)\dd u$, puis on revient à $x$ à la fin grâce à $u = \varphi^{-1}(x)$.
\end{enumerate}
\end{theoreme}

\begin{exemple}[Cinq primitives d'annales résolues par substitution (2021--2026)]
\textbf{1) UCLouvain -- Juillet 2021, Question 3b :} Calculer $\displaystyle\int \frac{\cos(x)}{1 + \sin^2(x)}\dd x$.\\
\emph{Résolution :} On repère au numérateur $\cos(x)$, qui est exactement la dérivée de $\sin(x)$ ! Posons $u = \sin(x) \implies \mathrm{d}u = \cos(x)\dd x$ :
\[
    \int \frac{\cos(x)}{1 + \sin^2(x)}\dd x = \int \frac{1}{1 + u^2}\dd u = \arctan(u) + C = \boxed{\arctan\bigl(\sin(x)\bigr) + C}.
\]

\medskip
\textbf{2) UCLouvain -- Septembre 2021, Question 3b (1\iere{} partie) :} Calculer $\displaystyle\int \frac{-2x}{\sqrt{1 - x^2}}\dd x$.\\
\emph{Résolution :} Le numérateur $-2x$ est la dérivée du radicande $1 - x^2$. Posons $u = 1 - x^2 \implies \mathrm{d}u = -2x\dd x$ :
\[
    \int \frac{-2x}{\sqrt{1 - x^2}}\dd x = \int \frac{1}{\sqrt{u}}\dd u = 2\sqrt{u} + C = \boxed{2\sqrt{1 - x^2} + C}.
\]

\medskip
\textbf{3) UCLouvain -- Septembre 2024, Question 3 (1\iere{} partie) et Septembre 2026, Question 3 (1\iere{} partie) :}
\begin{itemize}
    \item \textbf{Primitive de $\dfrac{1}{x\ln(x)}$ (Septembre 2024) :} On reconnaît $\frac{1/x}{\ln(x)} = \frac{u'(x)}{u(x)}$ avec $u = \ln(x) \implies \mathrm{d}u = \frac{1}{x}\dd x$ :
    \[
        \int \frac{1}{x\ln(x)}\dd x = \int \frac{1}{u}\dd u = \ln(|u|) + C = \boxed{\ln\bigl(|\ln(x)|\bigr) + C}.
    \]
    \item \textbf{Primitive de $\tan(x)$ (Septembre 2026 et Août 2023) :} Écrivons $\tan(x) = \dfrac{\sin(x)}{\cos(x)}$ et posons $u = \cos(x) \implies \mathrm{d}u = -\sin(x)\dd x$ :
    \[
        \int \tan(x)\dd x = \int \frac{\sin(x)}{\cos(x)}\dd x = -\int \frac{\mathrm{d}u}{u} = \boxed{-\ln\bigl(|\cos(x)|\bigr) + C}.
    \]
\end{itemize}

\medskip
\textbf{4) UCLouvain -- Juillet 2026, Question 3 (Substitution polynomiale $u = x^2 + 1$) :}\\
\emph{Énoncé :} Obtenir une primitive $g(x) = \displaystyle\int \frac{x^5}{(x^2 + 1)^5}\dd x$.\\
\emph{Résolution (Méthode 1 du solutionnaire -- par substitution $u = x^2 + 1$) :}
Séparons $x^5\dd x = x^4 \cdot (x\dd x)$ et posons $\boxed{u = x^2 + 1}$, ce qui donne $\mathrm{d}u = 2x\dd x \iff x\dd x = \frac{1}{2}\dd u$, et surtout $x^2 = u - 1 \implies x^4 = (u - 1)^2 = u^2 - 2u + 1$. L'intégrale se transforme en une simple somme de puissances de $u$ :
\begin{align*}
    g(x) &= \int \frac{(x^2)^2}{(x^2 + 1)^5}\,x\dd x = \frac{1}{2}\int \frac{(u - 1)^2}{u^5}\dd u = \frac{1}{2}\int \frac{u^2 - 2u + 1}{u^5}\dd u \\
    &= \frac{1}{2}\int \left(u^{-3} - 2u^{-4} + u^{-5}\right)\dd u = \frac{1}{2}\left[\frac{u^{-2}}{-2} - 2\frac{u^{-3}}{-3} + \frac{u^{-4}}{-4}\right] + C_1 \\
    &= -\frac{1}{4u^2} + \frac{1}{3u^3} - \frac{1}{8u^4} + C_1 = \boxed{-\frac{1}{4(x^2+1)^2} + \frac{1}{3(x^2+1)^3} - \frac{1}{8(x^2+1)^4} + C_1}.
\end{align*}
En mettant au même dénominateur $24u^4 = 24(x^2+1)^4$, on obtient la forme compacte :
\[
    g(x) = -\frac{6u^2 - 8u + 3}{24u^4} + C_1 = \boxed{-\frac{6x^4 + 4x^2 + 1}{24(x^2 + 1)^4} + C_1} \quad \left(\text{ou, avec } C_1 = \frac{1}{24} + C_2 : \;\; \boxed{\frac{x^6(x^2 + 4)}{24(x^2 + 1)^4} + C_2}\right).
\]
\emph{Note :} Le solutionnaire de Juillet 2026 présente également la substitution trigonométrique $x = \tan\theta \implies \mathrm{d}x = \frac{1}{\cos^2\theta}\dd\theta$ et $x^2+1 = \frac{1}{\cos^2\theta}$, qui transforme l'intégrale en $\int \sin^5\theta\cos^3\theta\dd\theta = \int \sin^5\theta(1-\sin^2\theta)\cos\theta\dd\theta = \frac{\sin^6\theta}{6} - \frac{\sin^8\theta}{8} + C_2$, où $\sin^2\theta = \frac{x^2}{x^2+1}$ !
\end{exemple}

% ----------------------------------------------------------------------------
\section{Intégration par parties (IPP)}
% ----------------------------------------------------------------------------

Lorsque l'intégrande est un \textbf{produit de deux fonctions de natures différentes} (par exemple un polynôme fois un logarithme $x\ln^2(x)$, un polynôme fois un sinus/cosinus $x^2\cos(x)$, un polynôme fois une exponentielle $x^2 e^{-x}$, ou une fonction réciproque $\arcsin(x)$), et qu'aucune des deux n'est la dérivée d'une partie de l'autre, la méthode reine est l'\textbf{intégration par parties}.

\begin{theoreme}[Formule d'intégration par parties (IPP)]
En intégrant la formule de dérivation d'un produit $\bigl(u(x)v(x)\bigr)' = u'(x)v(x) + u(x)v'(x)$, on obtient pour toutes fonctions dérivables $u(x)$ et $v(x)$ :
\[
    \boxed{\int u(x)\,v'(x)\dd x = u(x)\,v(x) - \int v(x)\,u'(x)\dd x} \qquad \text{notée en différentielles :} \qquad \boxed{\int u\dd v = u\,v - \int v\dd u}.
\]
Pour une intégrale définie entre $a$ et $b$ :
\[
    \boxed{\int_a^b u(x)\,v'(x)\dd x = \Bigl[u(x)\,v(x)\Bigr]_a^b - \int_a^b v(x)\,u'(x)\dd x}.
\]
\end{theoreme}

\begin{astuce}[Comment choisir qui est $u$ (à dériver) et qui est $\mathrm{d}v = v'\dd x$ (à intégrer) ? La règle ALPES !]
Le but de l'intégration par parties est que la nouvelle intégrale $\int v\dd u = \int v(x)u'(x)\dd x$ soit \textbf{plus simple} que l'intégrale de départ $\int u(x)v'(x)\dd x$. Pour choisir quelle fonction poser égale à $u(x)$ (celle qu'on va \textbf{dériver}), suivez l'ordre de priorité mnémotechnique \textbf{A -- L -- P -- E -- S} :
\begin{enumerate}
    \item \textbf{A}rc ($\arcsin, \arccos, \arctan$) : leur dérivée est algébrique ($\frac{1}{\sqrt{1-x^2}}, \frac{1}{1+x^2}$), il faut donc toujours les dériver ($u$) !
    \item \textbf{L}ogarithme ($\ln(x), \ln^2(x)$) : sa dérivée est $\frac{1}{x}$, on le choisit donc toujours comme $u(x)$ face à un polynôme ! (Exemple clé à connaître : $\displaystyle\int \ln(x)\dd x = x\ln(x) - \int x\cdot\frac{1}{x}\dd x = \boxed{x\ln(x) - x + C}$).
    \item \textbf{P}olynôme ($x, x^2, \dots$) : face à une exponentielle ou à un sinus/cosinus (mais pas face à $\ln$ ou $\arcsin$ !), on dérive le polynôme ($u = x^n$) pour faire baisser son degré jusqu'à $0$.
    \item \textbf{E}xponentielle ($e^{ax}$) et \textbf{S}inus/Cosinus ($\sin, \cos$) : faciles à intégrer ($\mathrm{d}v$), et si les deux sont présents ensemble (comme $\int e^{2x}\cos(2x)\dd x$ en Juillet 2021 Q3f), deux IPP successives font réapparaître l'intégrale de départ $\mathcal{I}$ (\textbf{IPP cyclique}).
\end{enumerate}
\end{astuce}

\begin{exemple}[Cinq primitives d'examen UCLouvain résolues par intégration par parties]
\textbf{1) IPP avec $\arcsin(x)$ (UCLouvain -- Septembre 2021, Question 3b) :}\\
\emph{Énoncé :} En utilisant la primitive $\displaystyle\int \frac{-2x}{\sqrt{1-x^2}}\dd x = 2\sqrt{1-x^2} + C$, calculer $\displaystyle\int \frac{x}{\sqrt{1-x^2}}\arcsin(x)\dd x$.\\
\emph{Résolution :} Posons $u(x) = \arcsin(x)$ (qu'on dérive : $u'(x) = \frac{1}{\sqrt{1-x^2}}$) et $v'(x) = \frac{x}{\sqrt{1-x^2}}$ (dont une primitive est $v(x) = -\sqrt{1-x^2}$ d'après la première partie). La formule d'IPP donne :
\begin{align*}
    \int \frac{x}{\sqrt{1-x^2}}\arcsin(x)\dd x &= -\sqrt{1-x^2}\,\arcsin(x) - \int \bigl(-\sqrt{1-x^2}\bigr)\frac{1}{\sqrt{1-x^2}}\dd x \\
    &= -\sqrt{1-x^2}\,\arcsin(x) + \int 1\dd x = \boxed{x - \sqrt{1-x^2}\,\arcsin(x) + C}.
\end{align*}

\medskip
\textbf{2) IPP avec fonction trigonométrique (UCLouvain -- Août 2023, Question 3) :}\\
\emph{Énoncé :} Obtenir la primitive $g(x) = \displaystyle\int \frac{x}{\cos^2(x)}\dd x$.\\
\emph{Résolution :} On reconnaît le produit du polynôme $u(x) = x$ (donc $\mathrm{d}u = \mathrm{d}x$) par $\mathrm{d}v = \frac{1}{\cos^2(x)}\dd x$, dont la primitive est immédiate : $v(x) = \tan(x)$. Par intégration par parties :
\[
    g(x) = x\tan(x) - \int \tan(x)\dd x = x\tan(x) - \int \frac{\sin(x)}{\cos(x)}\dd x = \boxed{x\tan(x) + \ln\bigl(|\cos(x)|\bigr) + C}.
\]

\medskip
\textbf{3) Double IPP vs substitution (UCLouvain -- Juillet 2023 -- Série 2, Question 4) :}\\
\emph{Énoncé :} Obtenir la primitive $g(x) = \displaystyle\int \Bigl(x^2\cos(x) - x\cos(x^2)\Bigr)\dd x$.\\
\emph{Résolution :} Séparons en deux intégrales :
\begin{itemize}
    \item Pour $\displaystyle\int x\cos(x^2)\dd x$, surtout pas d'IPP ! Comme $2x$ est la dérivée de $x^2$, on procède par \textbf{substitution} $u = x^2 \implies x\dd x = \frac{1}{2}\dd u$ :
    \[
        \int x\cos(x^2)\dd x = \frac{1}{2}\int \cos(u)\dd u = \frac{1}{2}\sin(x^2).
    \]
    \item Pour $\displaystyle\int x^2\cos(x)\dd x$, on effectue \textbf{deux intégrations par parties successives} pour abaisser le degré de $x^2$ :
    \begin{itemize}
        \item 1\iere{} IPP ($u = x^2 \implies \mathrm{d}u = 2x\dd x$ ; $\mathrm{d}v = \cos(x)\dd x \implies v = \sin(x)$) :
        \[
            \int x^2\cos(x)\dd x = x^2\sin(x) - 2\int x\sin(x)\dd x.
        \]
        \item 2\ieme{} IPP sur $\int x\sin(x)\dd x$ ($u = x \implies \mathrm{d}u = \mathrm{d}x$ ; $\mathrm{d}v = \sin(x)\dd x \implies v = -\cos(x)$) :
        \[
            \int x\sin(x)\dd x = -x\cos(x) - \int \bigl(-\cos(x)\bigr)\dd x = -x\cos(x) + \sin(x).
        \]
    \end{itemize}
    En rassemblant les deux parties :
    \[
        \boxed{g(x) = x^2\sin(x) + 2x\cos(x) - 2\sin(x) - \frac{1}{2}\sin(x^2) + C}.
    \]
\end{itemize}

\medskip
\textbf{4) Combinaison substitution + IPP avec logarithmes (UCLouvain -- Septembre 2024, Question 3) :}\\
\emph{Énoncé :} Obtenir la primitive $g(x) = \displaystyle\int \left(\frac{1}{x\ln(x)} - \frac{\ln(x)}{x^2}\right)\dd x$.\\
\emph{Résolution :}
\begin{itemize}
    \item La première intégrale vaut $\displaystyle\int \frac{1}{x\ln(x)}\dd x = \ln\bigl(|\ln(x)|\bigr)$ (par substitution $u = \ln(x)$).
    \item Pour la seconde intégrale $\displaystyle\int \frac{\ln(x)}{x^2}\dd x$, appliquons l'\textbf{IPP} en dérivant le logarithme ($u = \ln(x) \implies \mathrm{d}u = \frac{1}{x}\dd x$) et en intégrant la puissance ($\mathrm{d}v = \frac{1}{x^2}\dd x \implies v = -\frac{1}{x}$) :
    \[
        \int \frac{\ln(x)}{x^2}\dd x = -\frac{\ln(x)}{x} - \int \left(-\frac{1}{x}\right)\frac{1}{x}\dd x = -\frac{\ln(x)}{x} + \int \frac{1}{x^2}\dd x = -\frac{\ln(x)}{x} - \frac{1}{x} = -\frac{\ln(x) + 1}{x}.
    \]
    En soustrayant ce résultat à la première primitive, le signe $-$ devient $+$ :
    \[
        \boxed{g(x) = \ln\bigl(|\ln(x)|\bigr) + \frac{\ln(x) + 1}{x} + C}.
    \]
\end{itemize}

\medskip
\textbf{5) Double IPP sur un logarithme au carré (UCLouvain -- Septembre 2025, Question 3) :}\\
\emph{Énoncé :} Obtenir une primitive $g(x) = \displaystyle\int x\,\ln^2(x)\dd x$.\\
\emph{Résolution :} Suivons la règle ALPES en dérivant le logarithme et en intégrant $x$ :
\begin{itemize}
    \item \textbf{1\iere{} IPP :} Posons $u = \ln^2(x) \implies \mathrm{d}u = 2\ln(x)\frac{1}{x}\dd x$, et $\mathrm{d}v = x\dd x \implies v = \frac{x^2}{2}$ :
    \[
        g(x) = \frac{x^2}{2}\ln^2(x) - \int \frac{x^2}{2} \cdot \frac{2\ln(x)}{x}\dd x = \frac{x^2}{2}\ln^2(x) - \int x\ln(x)\dd x.
    \]
    \item \textbf{2\ieme{} IPP sur $\displaystyle\int x\ln(x)\dd x$ :} Posons $u = \ln(x) \implies \mathrm{d}u = \frac{1}{x}\dd x$, et $\mathrm{d}v = x\dd x \implies v = \frac{x^2}{2}$ :
    \[
        \int x\ln(x)\dd x = \frac{x^2}{2}\ln(x) - \int \frac{x^2}{2}\cdot\frac{1}{x}\dd x = \frac{x^2}{2}\ln(x) - \frac{1}{2}\int x\dd x = \frac{x^2}{2}\ln(x) - \frac{x^2}{4}.
    \]
\end{itemize}
En remplaçant et en factorisant par $\frac{x^2}{2}$, on retrouve la fonction étudiée en Septembre 2026 Q1 :
\[
    \boxed{g(x) = \frac{x^2}{2}\left(\ln^2(x) - \ln(x) + \frac{1}{2}\right) + C}.
\]
\end{exemple}

% ----------------------------------------------------------------------------
\section{Fractions rationnelles et primitives trigonométriques}
% ----------------------------------------------------------------------------

\subsection{Décomposition en éléments simples de $\dfrac{1}{1 - y^2}$ et fractions rationnelles}

Une identité algébrique très fréquemment utilisée aux examens de l'UCLouvain (Juillet 2021 Q3c, Septembre 2021 Q3c, Septembre 2026 Q3 -- où elle est même rappelée dans l'Aide !) est la \textbf{décomposition en fractions simples} de $\dfrac{2}{1 - t^2}$ (ou $\dfrac{2}{y^2 - 1}$) :

\begin{theoreme}[Primitives de $\dfrac{1}{1 - t^2}$ et de $\dfrac{1}{y^2 - 1}$ par décomposition en éléments simples]
Puisque $1 - t^2 = (1 - t)(1 + t)$, on a l'identité algébrique fondamentale :
\[
    \boxed{\frac{1}{1 - t^2} = \frac{1}{2}\left(\frac{1}{1 - t} + \frac{1}{1 + t}\right)} \qquad \text{et} \qquad \boxed{\frac{1}{y^2 - 1} = \frac{1}{2}\left(\frac{1}{y - 1} - \frac{1}{y + 1}\right)}.
\]
En intégrant terme à terme (attention au signe $-$ provenant de la dérivée de $1 - t$ !) :
\[
    \boxed{\int \frac{1}{1 - t^2}\dd t = \frac{1}{2}\ln(|1+t|) - \frac{1}{2}\ln(|1-t|) + C = \frac{1}{2}\ln\left(\left|\frac{1 + t}{1 - t}\right|\right) + C},
\]
\[
    \boxed{\int \frac{2}{y^2 - 1}\dd y = \int \left(\frac{1}{y-1} - \frac{1}{y+1}\right)\dd y = \ln\left(\left|\frac{y - 1}{y + 1}\right|\right) + C}.
\]
\end{theoreme}

\subsection{Les 4 techniques classiques pour les primitives trigonométriques}

\begin{theoreme}[Techniques d'intégration des expressions trigonométriques]
\begin{enumerate}
    \item \textbf{Présence d'une puissance impaire de $\sin(x)$ ou de $\cos(x)$ :}
    Si l'intégrande contient $\sin^{2k+1}(x)$ (comme $\frac{\sin^3(x)}{1+\cos(x)}$ en \textbf{Juillet 2022, Q3}), on isole un facteur $\sin(x)\dd x = -\mathrm{d}(\cos x)$ et on remplace le reste $\sin^{2k}(x) = \bigl(1 - \cos^2(x)\bigr)^k$ en posant $\boxed{u = \cos(x)}$.\\
    Symétriquement, pour $\cos^{2k+1}(x)$, on isole $\cos(x)\dd x = \mathrm{d}(\sin x)$, on remplace $\cos^{2k}(x) = \bigl(1 - \sin^2(x)\bigr)^k$ et on pose $\boxed{u = \sin(x)}$.
    \item \textbf{Puissances paires $\sin^2(x)$ ou $\cos^2(x)$ : linéarisation par les formules de Carnot !}
    (Juillet 2021 Q3f et Juillet 2024 Q3) :
    \[
        \boxed{\cos^2(x) = \frac{1 + \cos(2x)}{2} \implies \int \cos^2(x)\dd x = \frac{x}{2} + \frac{\sin(2x)}{4} + C},
    \]
    \[
        \boxed{\sin^2(x) = \frac{1 - \cos(2x)}{2} \implies \int \sin^2(x)\dd x = \frac{x}{2} - \frac{\sin(2x)}{4} + C}.
    \]
    \item \textbf{Dénominateur en $1 + \cos(x)$ ou $1 - \cos(x)$ : formule de l'angle demi ou binôme conjugué :}
    (Septembre 2022 Q3) : On utilise $1 + \cos(x) = 2\cos^2\left(\frac{x}{2}\right)$ (de primitive $\tan\left(\frac{x}{2}\right)$), ou l'on divise/factorise par $1+\cos(x)$ grâce à $\cos^2(x) - 1 = (\cos x + 1)(\cos x - 1)$.
    \item \textbf{Primitive de $\dfrac{1}{\cos(x)}$ ou de $\dfrac{1}{\sin(u)}$ (Juillet 2021 Q3c et Septembre 2026 Q3) :}
    On multiplie numérateur et dénominateur par $\cos(x)$ (resp. par $\sin(u)$) pour utiliser $\cos^2(x) = 1 - \sin^2(x)$ (resp. $\sin^2(u) = 1 - \cos^2(u)$) et se ramener à $\int \frac{\mathrm{d}t}{1 - t^2}$ :
    \[
        \int \frac{1}{\cos(x)}\dd x = \int \frac{\cos(x)}{1 - \sin^2(x)}\dd x \stackrel{t=\sin x}{=} \int \frac{\mathrm{d}t}{1 - t^2} = \boxed{\frac{1}{2}\ln\left(\frac{1 + \sin(x)}{1 - \sin(x)}\right) + C},
    \]
    \[
        \int \frac{1}{\sin(u)}\dd u = \int \frac{\sin(u)}{1 - \cos^2(u)}\dd u \stackrel{t=\cos u}{=} -\int \frac{\mathrm{d}t}{1 - t^2} = \boxed{\frac{1}{2}\ln\left(\frac{1 - \cos(u)}{1 + \cos(u)}\right) + C = \ln\left(\left|\tan\left(\frac{u}{2}\right)\right|\right) + C}.
    \]
\end{enumerate}
\end{theoreme}

\begin{exemple}[Quatre primitives trigonométriques d'annales UCLouvain (2022--2026)]
\textbf{1) UCLouvain -- Juillet 2022, Question 3 :} Calculer $I = \displaystyle\int \frac{\sin^3(x)}{1 + \cos(x)}\dd x$.\\
\emph{Résolution :} Isolons $\sin(x)\dd x$ et écrivons $\sin^2(x) = 1 - \cos^2(x) = \bigl(1 - \cos(x)\bigr)\bigl(1 + \cos(x)\bigr)$. Le dénominateur $1 + \cos(x)$ se simplifie immédiatement !
\begin{align*}
    I &= \int \frac{\bigl(1 - \cos^2(x)\bigr)\sin(x)}{1 + \cos(x)}\dd x = \int \bigl(1 - \cos(x)\bigr)\sin(x)\dd x \\
    &= \int \sin(x)\dd x - \int \cos(x)\sin(x)\dd x = \boxed{-\cos(x) + \frac{1}{2}\cos^2(x) + C_1} \quad \left(\text{ou } -\cos(x) - \frac{1}{2}\sin^2(x) + C_2\right).
\end{align*}

\medskip
\textbf{2) UCLouvain -- Septembre 2022, Question 3 :} Calculer $g(x) = \displaystyle\int \frac{\cos^2(x)}{1 + \cos(x)}\dd x$.\\
\emph{Résolution (Méthode 1 du solutionnaire -- Ajout et retrait de $1$ au numérateur) :}
Faisons apparaître $\cos^2(x) - 1 = (\cos x + 1)(\cos x - 1)$ au numérateur :
\begin{align*}
    g(x) &= \int \frac{\cos^2(x) - 1 + 1}{1 + \cos(x)}\dd x = \int \frac{(\cos x + 1)(\cos x - 1)}{1 + \cos(x)}\dd x + \int \frac{1}{1 + \cos(x)}\dd x \\
    &= \int \bigl(\cos(x) - 1\bigr)\dd x + \int \frac{1}{2\cos^2(x/2)}\dd x = \boxed{\sin(x) - x + \tan\left(\frac{x}{2}\right) + C},
\end{align*}
où nous avons utilisé la formule de Carnot $1 + \cos(x) = 2\cos^2\left(\frac{x}{2}\right)$ et posé $s = \frac{x}{2} \implies \mathrm{d}s = \frac{1}{2}\dd x$, d'où $\int \frac{\mathrm{d}x}{2\cos^2(x/2)} = \int \frac{\mathrm{d}s}{\cos^2(s)} = \tan(s) = \tan\left(\frac{x}{2}\right)$.

\medskip
\textbf{3) UCLouvain -- Juillet 2025, Question 3 :} Calculer $g(x) = \displaystyle\int \frac{1}{1 + \sin(2x)}\dd x$.\\
\emph{Rappel de l'énoncé :} $\sin(2x) = 2\sin(x)\cos(x) = \dfrac{2\tan(x)}{1 + \tan^2(x)}$.\\
\emph{Résolution (présentons 3 des 4 méthodes officielles !) :}
\begin{itemize}
    \item \emph{Méthode 1 (Avec l'aide en $\tan(x)$) :} En remplaçant $\sin(2x) = \frac{2\tan x}{1+\tan^2 x}$ et en utilisant $1 + \tan^2(x) = \frac{1}{\cos^2(x)}$ :
    \[
        g(x) = \int \frac{1}{1 + \dfrac{2\tan x}{1+\tan^2 x}}\dd x = \int \frac{1 + \tan^2(x)}{1 + 2\tan(x) + \tan^2(x)}\dd x = \int \frac{\dfrac{1}{\cos^2(x)}}{\bigl(1 + \tan(x)\bigr)^2}\dd x.
    \]
    En posant $u = 1 + \tan(x) \implies \mathrm{d}u = \frac{1}{\cos^2(x)}\dd x$, on obtient $\displaystyle\int \frac{\mathrm{d}u}{u^2} = -\frac{1}{u} + C = \boxed{-\frac{1}{1 + \tan(x)} + C}$.
    \item \emph{Méthode 2 (Multiplication par le conjugué $1 - \sin(2x)$) :} Comme $(1+\sin 2x)(1-\sin 2x) = 1 - \sin^2(2x) = \cos^2(2x)$ :
    \[
        g(x) = \int \frac{1 - \sin(2x)}{\cos^2(2x)}\dd x = \int \frac{1}{\cos^2(2x)}\dd x - \int \frac{\sin(2x)}{\cos^2(2x)}\dd x = \boxed{\frac{1}{2}\tan(2x) - \frac{1}{2\cos(2x)} + C_3 = \frac{\sin(2x) - 1}{2\cos(2x)} + C_3}.
    \]
    \item \emph{Méthode 3 (Factorisation de $1 + \sin(2x) = (\cos x + \sin x)^2$) :} Comme $(\cos x + \sin x)^2 = 2\cos^2\left(x - \frac{\pi}{4}\right)$, on a directement $g(x) = \frac{1}{2}\tan\left(x - \frac{\pi}{4}\right) + C_4 = -\frac{\cos(x)}{\cos(x)+\sin(x)} + C$.
\end{itemize}

\medskip
\textbf{4) UCLouvain -- Septembre 2026, Question 3 :} Calculer $g(x) = \displaystyle\int \left(\tan(x) - \frac{\sin(x)}{\sin(\cos(x))}\right)\dd x$.\\
\textbf{Aide :} $\dfrac{2}{1 - t^2} = \dfrac{1}{1 - t} + \dfrac{1}{1 + t}$.\\
\emph{Résolution :}
\begin{itemize}
    \item La primitive du premier terme est $\displaystyle\int \tan(x)\dd x = -\ln\bigl(|\cos(x)|\bigr)$.
    \item Pour le second terme $I_2 = \displaystyle\int \frac{\sin(x)}{\sin(\cos(x))}\dd x$, effectuons une première substitution $\boxed{u = \cos(x)} \implies \mathrm{d}u = -\sin(x)\dd x$ :
    \[
        I_2 = -\int \frac{\mathrm{d}u}{\sin(u)} = -\int \frac{\sin(u)}{\sin^2(u)}\dd u = -\int \frac{\sin(u)}{1 - \cos^2(u)}\dd u.
    \]
    Effectuons alors une seconde substitution $\boxed{t = \cos(u) = \cos(\cos(x))} \implies \mathrm{d}t = -\sin(u)\dd u$. Grâce à l'aide de l'énoncé :
    \[
        I_2 = \int \frac{\mathrm{d}t}{1 - t^2} = \frac{1}{2}\int \left(\frac{1}{1 - t} + \frac{1}{1 + t}\right)\dd t = \frac{1}{2}\ln\left(\left|\frac{1 + t}{1 - t}\right|\right) = \frac{1}{2}\ln\left(\left|\frac{1 + \cos(\cos(x))}{1 - \cos(\cos(x))}\right|\right).
    \]
    En soustrayant $I_2$ au premier terme (ou en utilisant la variante de l'angle demi $\int \frac{\mathrm{d}u}{\sin(u)} = \ln\bigl(|\tan(u/2)|\bigr)$) :
    \[
        \boxed{g(x) = -\ln\bigl(|\cos(x)|\bigr) - \frac{1}{2}\ln\left(\left|\frac{1 + \cos(\cos(x))}{1 - \cos(\cos(x))}\right|\right) + C = -\ln\bigl(|\cos(x)|\bigr) + \ln\left(\left|\tan\left(\frac{\cos(x)}{2}\right)\right|\right) + C'}.
    \]
\end{itemize}
\end{exemple}

% ----------------------------------------------------------------------------
\section{Exercices du Chapitre 6}
% ----------------------------------------------------------------------------

\emph{Les corrigés détaillés de tous ces exercices se trouvent au Chapitre 9, en fin d'ouvrage.}

\begin{enumerate}[label=\textbf{Exercice 6.\arabic* --}]
    \item \textbf{(Primitives par substitution directe)} Calculer les primitives suivantes :
    \begin{enumerate}
        \item $\displaystyle\int \frac{e^x}{1 + e^{2x}}\dd x$
        \item $\displaystyle\int \frac{\ln^3(x)}{x}\dd x$
        \item $\displaystyle\int \frac{x}{\sqrt{4 - x^2}}\dd x$ \quad et \quad $\displaystyle\int \frac{1}{\sqrt{4 - x^2}}\dd x$
        \item $\displaystyle\int \frac{\cos(\ln(x))}{x}\dd x$
    \end{enumerate}

    \item \textbf{(UCLouvain -- Septembre 2021, Q3c -- Substitution exponentielle et fractions simples)} Déterminer une primitive de la fonction :
    \[
        f(x) = \frac{2e^x}{(e^x)^2 - 1},
    \]
    en précisant sa condition d'existence.

    \item \textbf{(Primitives par intégration par parties)} Calculer les primitives suivantes :
    \begin{enumerate}
        \item $\displaystyle\int \arctan(x)\dd x$
        \item $\displaystyle\int \arcsin(x)\dd x$
        \item $\displaystyle\int (2x + 1)\,e^{-x}\dd x$
        \item $\displaystyle\int x^2\ln(x)\dd x$
    \end{enumerate}

    \item \textbf{(IPP cyclique -- Inspiré de Juillet 2021, Q3f)} En effectuant deux intégrations par parties successives, calculer les primitives suivantes :
    \begin{enumerate}
        \item $I_1 = \displaystyle\int e^x\cos(x)\dd x$
        \item $I_2 = \displaystyle\int e^{2x}\sin(x)\dd x$
    \end{enumerate}

    \item \textbf{(Primitives trigonométriques)} Calculer les primitives suivantes :
    \begin{enumerate}
        \item $\displaystyle\int \sin^3(x)\cos^2(x)\dd x$
        \item $\displaystyle\int \sin^2(x)\cos^2(x)\dd x$
        \item $\displaystyle\int \tan^2(x)\dd x$ \quad et \quad $\displaystyle\int \tan^3(x)\dd x$
    \end{enumerate}

    \item \textbf{(Défi multi-méthodes -- Inspiré de Juillet 2026, Q3)} Calculer la primitive :
    \[
        \int \frac{x^3}{(x^2 + 1)^3}\dd x
    \]
    par deux méthodes différentes : (a) par le changement de variable $u = x^2 + 1$, et (b) par intégration par parties en écrivant $\dfrac{x^3}{(x^2+1)^3} = x^2 \cdot \dfrac{x}{(x^2+1)^3}$. Vérifier que les deux résultats ne diffèrent que d'une constante.
\end{enumerate}

% ============================================================================
% CHAPITRE 7 : INTÉGRALES DÉFINIES, INÉGALITÉS, AIRES ET VOLUMES DE RÉVOLUTION
% ============================================================================
\chapter{Intégrales Définies, Inégalités, Aires et Volumes de Révolution}

Une fois maîtrisé le calcul des primitives (Chapitre 6), l'\textbf{intégrale définie} $\int_a^b f(x)\dd x$ permet de calculer des grandeurs géométriques et physiques exactes. À l'examen d'admission d'Analyse de l'UCLouvain, l'intégrale définie intervient régulièrement sous deux grandes formes dans les Questions 3 à 5 :
\begin{enumerate}
    \item \textbf{La démonstration d'inégalités fonctionnelles} (encadrements de $\ln(x)$, de $\int_1^x \ln(s)\dd s$ ou de $e^{-x}$) en intégrant une inégalité de départ entre $1$ et $x$ (ou entre $0$ et $x$) : \textbf{Juillet 2022 Q4, Septembre 2022 Q4, Juillet 2023 (Série 1) Q4, Juillet 2023 (Série 2) Q5, Août 2023 Q5}.
    \item \textbf{Le calcul du volume d'un solide de révolution} obtenu en faisant tourner une courbe ou une surface plane autour d'un axe : \textbf{Juillet 2021 Q3f, Septembre 2022 Q2, Juillet 2024 Q3 (volume d'un tore), Juillet 2026 Q4 (volume d'un œuf fossilisé)}.
\end{enumerate}

% ----------------------------------------------------------------------------
\section{L'intégrale définie et le Théorème Fondamental de l'Analyse}
% ----------------------------------------------------------------------------

\subsection{Sommes de Riemann, théorème fondamental et propriétés}

Géométriquement, pour une fonction continue $f : [a, b] \to \R$, l'\textbf{intégrale définie} $\int_a^b f(x)\dd x$ représente l'\textbf{aire algébrique} (c'est-à-dire comptée positivement là où $f(x) \ge 0$ et négativement là où $f(x) \le 0$) délimitée par la courbe $y = f(x)$, l'axe horizontal $Ox$ et les droites verticales $x = a$ et $x = b$.

\begin{theoreme}[Théorème Fondamental de l'Analyse (Newton-Leibniz)]
Soit $f : [a, b] \to \R$ une fonction continue sur l'intervalle $[a, b]$.
\begin{enumerate}
    \item \textbf{Dérivée d'une intégrale fonction de sa borne supérieure :} La fonction d'accumulation $\Phi : [a, b] \to \R$ définie par :
    \[
        \Phi(x) = \int_a^x f(s)\dd s
    \]
    est dérivable sur $[a, b]$ et est l'unique primitive de $f$ qui s'annule en $x = a$ :
    \[
        \boxed{\forall x \in [a, b], \quad \Phi'(x) = \frac{\mathrm{d}}{\mathrm{d}x}\left(\int_a^x f(s)\dd s\right) = f(x) \qquad \text{avec } \Phi(a) = 0}.
    \]
    \item \textbf{Formule de Newton-Leibniz :} Si $F$ est n'importe quelle primitive de $f$ sur $[a, b]$ ($F'(x) = f(x)$), alors :
    \[
        \boxed{\int_a^b f(x)\dd x = \Bigl[F(x)\Bigr]_a^b = F(b) - F(a)}.
    \]
\end{enumerate}
\end{theoreme}

\begin{theoreme}[Propriétés fondamentales de l'intégrale définie et valeur moyenne]
Pour toutes fonctions continues $f, g$ sur $[a, b]$ et constantes $\alpha, \beta \in \R$ :
\begin{enumerate}
    \item \textbf{Linéarité :} $\displaystyle\int_a^b \bigl(\alpha f(x) + \beta g(x)\bigr)\dd x = \alpha\int_a^b f(x)\dd x + \beta\int_a^b g(x)\dd x$.
    \item \textbf{Relation de Chasles et inversion des bornes :}
    \[
        \int_a^a f(x)\dd x = 0, \;\; \int_b^a f(x)\dd x = -\int_a^b f(x)\dd x, \;\; \int_a^c f(x)\dd x + \int_c^b f(x)\dd x = \int_a^b f(x)\dd x.
    \]
    \item \textbf{Simplifications par parité sur un intervalle symétrique $[-a, a]$ :}
    \begin{itemize}
        \item Si $f$ est \textbf{paire} ($f(-x) = f(x)$) : $\boxed{\displaystyle\int_{-a}^a f(x)\dd x = 2\int_0^a f(x)\dd x}$.
        \item Si $f$ est \textbf{impaire} ($f(-x) = -f(x)$) : $\boxed{\displaystyle\int_{-a}^a f(x)\dd x = 0}$.
    \end{itemize}
    \item \textbf{Positivité et croissance de l'intégrale (pour $a \le b$, l'ordre des bornes est crucial !) :}
    \[
        \boxed{\text{Si } a \le b \;\text{ et }\; \forall s \in [a, b], \; f(s) \le g(s), \quad \text{alors} \quad \int_a^b f(s)\dd s \le \int_a^b g(s)\dd s}.
    \]
    \item \textbf{Valeur moyenne d'une fonction $f$ sur $[a, b]$ (avec $a < b$) :} (Utilisée en Août 2023 Q2 pour l'inconfort moyen $\bar{c}$ dans le train de nuit) :
    \[
        \boxed{\bar{f} = \frac{1}{b - a}\int_a^b f(t)\dd t}.
    \]
\end{enumerate}
\end{theoreme}

% ----------------------------------------------------------------------------
\section{Démonstration d'inégalités par intégration ou par dérivation}
% ----------------------------------------------------------------------------

En 2022 et 2023, cinq examens d'admission consécutifs de l'UCLouvain ont comporté une question demandant de démontrer une inégalité (ou une double inégalité) fonctionnelle. Les solutionnaires officiels mettent en avant \textbf{deux méthodes complémentaires} qu'il faut absolument savoir rédiger.

\begin{theoreme}[Les deux méthodes reines pour démontrer une inégalité $f(x) \le g(x)$ pour $x \ge a$]
\begin{enumerate}
    \item \textbf{Méthode A -- Par intégration d'une inégalité connue sur $[a, x]$ (quand $x \ge a$) :}
    Si l'on connaît une inégalité simple entre deux fonctions dérivées $f'(s) \le g'(s)$ pour tout $s \in [a, x]$ (par exemple $\frac{1}{s^2} \le \frac{1}{s} \le 1$ pour $s \ge 1$, ou $1 - \ln(1+s) \le e^{-s}$ pour $s \ge 0$), on \textbf{intègre chaque membre par rapport à la variable muette $s$ entre $a$ et $x$} (avec $x \ge a$) :
    \[
        \int_a^x f'(s)\dd s \le \int_a^x g'(s)\dd s \iff f(x) - f(a) \le g(x) - g(a).
    \]
    Si de plus $f(a) = g(a)$, on obtient immédiatement $\boxed{f(x) \le g(x)}$ pour tout $x \ge a$ !\\
    \emph{Attention si $x \le a$ (par ex. $0 < x \le 1$ avec $a = 1$) :} Pour garder les bornes dans le sens croissant (borne inférieure $\le$ borne supérieure), on intègre entre $x$ et $a$ : $\int_x^1 f'(s)\dd s \le \int_x^1 g'(s)\dd s$ !
    \item \textbf{Méthode B -- Par étude du signe de la dérivée de la différence $\Delta(x) = g(x) - f(x)$ :}
    Pour prouver que $f(x) \le g(x)$ pour tout $x \ge a$, on pose $\Delta(x) = g(x) - f(x)$ :
    \begin{itemize}
        \item On vérifie que $\Delta(a) = g(a) - f(a) = 0$ (ou $\ge 0$).
        \item On calcule la dérivée $\Delta'(x) = g'(x) - f'(x)$ et on montre que $\Delta'(x) \ge 0$ pour tout $x \ge a$.
        \item Par conséquent, $\Delta$ est croissante sur $[a, +\infty[$, donc $\forall x \ge a,\; \Delta(x) \ge \Delta(a) \ge 0 \iff f(x) \le g(x)$.
    \end{itemize}
\end{enumerate}
\end{theoreme}

\begin{exemple}[Les quatre grandes questions d'inégalités des annales UCLouvain (2022--2023)]
\textbf{1) Encadrement fondamental de $\ln(x)$ et de $\int_1^x \ln(s)\dd s$ (UCLouvain -- Septembre 2022, Question 4 et Août 2023, Question 5) :}\\
\emph{Énoncé :}
\begin{enumerate}[label=(\alph*)]
    \item Soit $x \ge 1$ (on a donc $1 \le s \le s^2$ pour tout $s \in [1, x]$). Démontrer que :
    \[
        \boxed{1 - \frac{1}{x} \le \ln(x) \le x - 1}.
    \]
    \item (Août 2023, Q5) Démontrer que cette double inégalité reste vraie pour $0 < x \le 1$, et donc pour tout $x > 0$.
    \item (Septembre 2022, Q4) En utilisant ce résultat pour $x \ge 1$, démontrer que :
    \[
        \boxed{0 \le \int_1^x \ln(s)\dd s \le (x - 1)^2} \quad \left(\text{et même } \le \frac{1}{2}(x-1)^2\right).
    \]
\end{enumerate}
\emph{Résolution :}
\begin{enumerate}[label=(\alph*)]
    \item Pour tout $s \ge 1$, on a $1 \le s \le s^2$. En passant aux inverses (la fonction inverse $u \mapsto \frac{1}{u}$ étant strictement décroissante sur $\R_0^+$, le sens des inégalités s'inverse) :
    \[
        \forall s \ge 1, \quad \frac{1}{s^2} \le \frac{1}{s} \le 1.
    \]
    Intégrons cet encadrement par rapport à $s$ sur l'intervalle $[1, x]$ (avec $x \ge 1$) :
    \[
        \int_1^x \frac{1}{s^2}\dd s \le \int_1^x \frac{1}{s}\dd s \le \int_1^x 1\dd s \iff \left[-\frac{1}{s}\right]_1^x \le \Bigl[\ln(s)\Bigr]_1^x \le \Bigl[s\Bigr]_1^x,
    \]
    soit $-\frac{1}{x} - (-1) \le \ln(x) - 0 \le x - 1$, ce qui donne exactement $1 - \dfrac{1}{x} \le \ln(x) \le x - 1$.
    \item \textbf{Pour $0 < x \le 1$ (deux méthodes officielles d'Août 2023) :}
    \begin{itemize}
        \item \emph{Méthode 1 (par intégration sur $[x, 1]$) :} Pour $0 < s \le 1$, on a $s^2 \le s \le 1 \implies 1 \le \frac{1}{s} \le \frac{1}{s^2}$. En intégrant entre la plus petite borne $x$ et la plus grande borne $1$ :
        \[
            \int_x^1 1\dd s \le \int_x^1 \frac{1}{s}\dd s \le \int_x^1 \frac{1}{s^2}\dd s \iff 1 - x \le -\ln(x) \le \frac{1}{x} - 1.
        \]
        En multipliant toute la double inégalité par $-1$ (ce qui inverse les symboles $\le$), on retrouve exactement $1 - \frac{1}{x} \le \ln(x) \le x - 1$ !
        \item \emph{Méthode 2 (par changement de variable $y = \frac{1}{x} \ge 1$) :} Si $0 < x \le 1$, alors $y = \frac{1}{x} \ge 1$. Appliquons le résultat de (a) à $y \ge 1$ :
        \[
            1 - \frac{1}{y} \le \ln(y) \le y - 1 \iff 1 - x \le -\ln(x) \le \frac{1}{x} - 1 \iff 1 - \frac{1}{x} \le \ln(x) \le x - 1.
        \]
    \end{itemize}
    \item \textbf{Pour majorer $\int_1^x \ln(s)\dd s$ ($x \ge 1$) :}
    \begin{itemize}
        \item Pour $s \ge 1$, on a $0 \le \ln(s) \le s - 1$. En intégrant directement cette double inégalité entre $1$ et $x$ :
        \[
            0 \le \int_1^x \ln(s)\dd s \le \int_1^x (s - 1)\dd s = \left[\frac{(s-1)^2}{2}\right]_1^x = \frac{1}{2}(x - 1)^2 \le (x - 1)^2.
        \]
        \item \emph{Variante en calculant l'intégrale de $\ln(s)$ :} Comme $\int_1^x \ln(s)\dd s = \bigl[s\ln(s) - s\bigr]_1^x = x\ln(x) - (x - 1)$, il suffit de multiplier l'encadrement $\frac{x-1}{x} \le \ln(x) \le x - 1$ par $x > 0$ pour obtenir $x - 1 \le x\ln(x) \le x(x-1)$, et en soustrayant $x - 1$, on a $0 \le x\ln(x) - (x-1) \le (x-1)(x-1) = (x-1)^2$.
    \end{itemize}
\end{enumerate}

\medskip
\textbf{2) Encadrement fin de $\ln(x)$ par des puissances $x^a$ (UCLouvain -- Juillet 2022, Question 4) :}\\
\emph{Énoncé :} Soit $x > 1$ et $0 < a < 1$, de sorte que $s^{1-a} < s < s^{1+a}$ pour tout $s > 1$. Prouver que :
\[
    \frac{1}{a}\left(1 - \frac{1}{x^a}\right) < \ln(x) < \frac{1}{a}\bigl(x^a - 1\bigr).
\]
\emph{Résolution :} Pour tout $s \in \,]1, x]$, passons aux inverses dans $s^{1-a} < s < s^{1+a}$ :
\[
    \frac{1}{s^{1+a}} < \frac{1}{s} < \frac{1}{s^{1-a}} \iff s^{-1-a} < \frac{1}{s} < s^{a-1}.
\]
Intégrons chaque membre par rapport à $s$ entre $1$ et $x$ (avec $\int s^p\dd s = \frac{s^{p+1}}{p+1}$ pour $p = -1-a$ et $p = a-1$) :
\[
    \left[\frac{s^{-a}}{-a}\right]_1^x < \Bigl[\ln(s)\Bigr]_1^x < \left[\frac{s^a}{a}\right]_1^x \iff \frac{1}{a}\left(1 - \frac{1}{x^a}\right) < \ln(x) < \frac{1}{a}\bigl(x^a - 1\bigr).
\]

\medskip
\textbf{3) Inégalité exponentielle $e^y \ge \left(1 + \frac{y}{n}\right)^n$ (UCLouvain -- Juillet 2023 -- Série 2, Question 5) :}\\
\emph{Énoncé :} En exploitant le fait que $e^x \ge 1$ pour tout $x \ge 0$, prouver que $e^x \ge 1 + x$ pour tout $x \ge 0$. Grâce aux propriétés de l'exponentielle, en déduire ensuite que pour tout entier $n \ge 1$ et tout $y > 0$ :
\[
    e^y \ge \left(1 + \frac{y}{n}\right)^n.
\]
\emph{Résolution :}
\begin{itemize}
    \item Pour tout $w \ge 0$, on a $e^w \ge 1$. Intégrons cette inégalité entre $0$ et $x$ (pour $x \ge 0$) :
    \[
        \int_0^x e^w\dd w \ge \int_0^x 1\dd w \iff e^x - 1 \ge x \iff \boxed{e^x \ge 1 + x \quad (\forall x \ge 0)}.
    \]
    \item Appliquons cette inégalité au réel positif $x = \dfrac{y}{n} > 0$ :
    \[
        e^{y/n} \ge 1 + \frac{y}{n} > 0.
    \]
    Comme les deux membres sont positifs et que la fonction puissance $u \mapsto u^n$ ($n \ge 1$) est strictement croissante sur $\R^+$, élevons les deux membres à la puissance $n$ en utilisant $(e^{y/n})^n = e^{(y/n)\cdot n} = e^y$ :
    \[
        \bigl(e^{y/n}\bigr)^n \ge \left(1 + \frac{y}{n}\right)^n \iff \boxed{e^y \ge \left(1 + \frac{y}{n}\right)^n}.
    \]
\end{itemize}

\medskip
\textbf{4) Majoration de $e^{-x}$ par intégration de $\ln(1+s)$ (UCLouvain -- Juillet 2023 -- Série 1, Question 4) :}\\
\emph{Énoncé :} Sachant que pour tout $x \ge 0$, on a $1 - \ln(1+x) \le e^{-x}$, démontrer que pour tout $x \ge 0$ :
\[
    e^{-x} \le (1 + x)\ln(1 + x) + (1 - 2x).
\]
\emph{Résolution :} Intégrons l'inégalité $1 - \ln(1+s) \le e^{-s}$ par rapport à $s$ entre $0$ et $x$ (avec $x \ge 0$). Par le changement de variable $p = 1 + s \implies \mathrm{d}p = \mathrm{d}s$, une primitive de $\ln(1+s)$ est $p\ln(p) - p = (1+s)\ln(1+s) - (1+s)$. Dès lors :
\[
    \int_0^x \bigl(1 - \ln(1+s)\bigr)\dd s \le \int_0^x e^{-s}\dd s \iff \Bigl[s - (1+s)\ln(1+s) + (1+s)\Bigr]_0^x \le \Bigl[-e^{-s}\Bigr]_0^x,
\]
\[
    \Bigl[(1 + 2s) - (1+s)\ln(1+s)\Bigr]_0^x \le -e^{-x} - (-1) \iff (1 + 2x) - (1+x)\ln(1+x) - 1 \le 1 - e^{-x}.
\]
En isolant $e^{-x}$ à gauche, on obtient immédiatement $\boxed{e^{-x} \le (1+x)\ln(1+x) + (1 - 2x)}$.
\end{exemple}

% ----------------------------------------------------------------------------
\section{Calcul d'aires planes et volumes de solides de révolution}
% ----------------------------------------------------------------------------

\subsection{Aire comprise entre deux courbes et volume d'un solide de révolution}

\begin{theoreme}[Aire entre deux courbes et Volume d'un solide de révolution]
\begin{enumerate}
    \item \textbf{Aire d'une région plane comprise entre deux courbes $y = f(x)$ et $y = g(x)$ sur $[a, b]$ :}
    \[
        \boxed{\mathcal{A} = \int_a^b \bigl|f(x) - g(x)\bigr|\dd x = \int_a^b \bigl(y_{\text{haut}}(x) - y_{\text{bas}}(x)\bigr)\dd x}.
    \]
    \item \textbf{Volume d'un solide plein obtenu par révolution de $y = R(x)$ autour de l'axe $Ox$ sur $[a, b]$ :}
    En découpant le solide en tranches cylindriques élémentaires (disques) perpendiculaires à l'axe de rotation $Ox$, d'épaisseur $\mathrm{d}x$ et de rayon $R(x) = |f(x)|$, l'aire de chaque disque vaut $\pi\bigl(R(x)\bigr)^2$, d'où le volume total :
    \[
        \boxed{V = \pi \int_a^b \bigl(f(x)\bigr)^2\dd x = \pi \int_a^b \bigl(R(x)\bigr)^2\dd x}.
    \]
    \item \textbf{Volume d'un solide creux (anneau / tube) compris entre un rayon extérieur $y_+(x) \ge 0$ et un rayon intérieur $y_-(x) \ge 0$ autour de $Ox$ :}
    Chaque tranche est une couronne circulaire d'aire $\pi y_+^2(x) - \pi y_-^2(x)$, donc :
    \[
        \boxed{V = V_{\text{ext}} - V_{\text{int}} = \pi \int_a^b \Bigl(\bigl(y_+(x)\bigr)^2 - \bigl(y_-(x)\bigr)^2\Bigr)\dd x}.
    \]
\end{enumerate}
\end{theoreme}

\begin{center}
\begin{tikzpicture}[scale=1.0, >=Stealth]
    % Axe de révolution
    \draw[->, thick, darkgray] (-0.6,0) -- (6.2,0) node[right] {$x$ (axe de révolution)};
    \draw[->, thick, darkgray] (0,-2.0) -- (0,2.2) node[above] {$y$};

    % Profil supérieur et inférieur symétrique
    \draw[very thick, uclblue, domain=0.8:5.0, samples=80] plot (\x, {1.05 + 0.45*sin(deg(1.2*\x))}) node[above right] {$y = R(x)$};
    \draw[thick, uclblue!60, dashed, domain=0.8:5.0, samples=80] plot (\x, {-1.05 - 0.45*sin(deg(1.2*\x))});

    % Ellipses de révolution aux extrémités a et b
    \draw[thick, uclblue] (0.8,0) ellipse (0.22 and 1.42);
    \draw[thick, uclblue] (5.0,0) ellipse (0.22 and 0.92);

    % Tranche élémentaire (disque d'épaisseur dx) en x = 2.8
    \fill[uclaccent!25] (2.8,0) ellipse (0.20 and 0.95);
    \draw[thick, warnred] (2.8,0) ellipse (0.20 and 0.95);
    \draw[thick, warnred, <->] (2.8,0) -- (2.8,0.95) node[midway, right] {\small $R(x)$};
    \fill[warnred] (2.8,0) circle (1.5pt);

    \node[below] at (0.8,-1.45) {$x = a$};
    \node[below] at (5.0,-1.0) {$x = b$};
    \node[uclblue, font=\small\bfseries] at (3.0, -1.75) {Tranche élémentaire : $\mathrm{d}V = \pi \bigl(R(x)\bigr)^2\dd x$};
\end{tikzpicture}
\end{center}

\begin{exemple}[Trois volumes de révolution complets des annales UCLouvain (2021, 2024 et 2026)]
\textbf{1) Pièce mécanique exponentielle-trigonométrique (UCLouvain -- Juillet 2021, Question 3f) :}\\
\emph{Énoncé :} Une pièce mécanique de longueur $2\pi$ est définie par la révolution autour de l'axe des abscisses de la fonction $f : [0, 2\pi] \to \R, x \mapsto f(x) = \sin(x)\exp(x)$.
(i) Exprimer le volume $V$ de cette pièce sous forme d'une intégrale, et (ii) calculer sa valeur exacte. (\emph{Rappel :} $\cos(2\theta) = 1 - 2\sin^2(\theta)$).\\
\emph{Résolution :}
\begin{enumerate}[label=(\roman*)]
    \item Par la formule du volume de révolution autour de $Ox$ :
    \[
        \boxed{V = \pi \int_0^{2\pi} \bigl(\sin(x)\,e^x\bigr)^2\dd x = \pi \int_0^{2\pi} \sin^2(x)\,e^{2x}\dd x}.
    \]
    \item Linéarisons $\sin^2(x) = \dfrac{1 - \cos(2x)}{2}$ grâce à la formule de Carnot rappelée dans l'énoncé :
    \[
        V = \frac{\pi}{2}\int_0^{2\pi} e^{2x}\dd x - \frac{\pi}{2}\int_0^{2\pi} \cos(2x)\,e^{2x}\dd x = \frac{\pi}{2}\left[\frac{e^{2x}}{2}\right]_0^{2\pi} - \frac{\pi}{2}\int_0^{2\pi} \cos(2x)\,e^{2x}\dd x.
    \]
    Posons $y = 2x \implies \mathrm{d}y = 2\dd x$ (avec $y \in [0, 4\pi]$) dans la seconde intégrale : $\int_0^{2\pi} \cos(2x)e^{2x}\dd x = \frac{1}{2}\mathcal{I}$ où $\mathcal{I} = \int_0^{4\pi} \cos(y)e^y\dd y$. Calculons $\mathcal{I}$ par \textbf{double intégration par parties cyclique} :
    \begin{itemize}
        \item 1\iere{} IPP ($u = \cos y \implies \mathrm{d}u = -\sin y\dd y$ ; $\mathrm{d}v = e^y\dd y \implies v = e^y$) :
        \[
            \mathcal{I} = \Bigl[\cos(y)\,e^y\Bigr]_0^{4\pi} + \int_0^{4\pi} \sin(y)\,e^y\dd y = (e^{4\pi} - 1) + \int_0^{4\pi} \sin(y)\,e^y\dd y.
        \]
        \item 2\ieme{} IPP ($u = \sin y \implies \mathrm{d}u = \cos y\dd y$ ; $\mathrm{d}v = e^y\dd y \implies v = e^y$) :
        \[
            \int_0^{4\pi} \sin(y)\,e^y\dd y = \underbrace{\Bigl[\sin(y)\,e^y\Bigr]_0^{4\pi}}_{= 0 - 0 = 0} - \int_0^{4\pi} \cos(y)\,e^y\dd y = -\mathcal{I}.
        \]
    \end{itemize}
    On obtient l'équation $\mathcal{I} = e^{4\pi} - 1 - \mathcal{I} \iff 2\mathcal{I} = e^{4\pi} - 1 \iff \mathcal{I} = \dfrac{e^{4\pi} - 1}{2}$. Dès lors :
    \[
        V = \frac{\pi}{4}\bigl(e^{4\pi} - 1\bigr) - \frac{\pi}{4}\,\mathcal{I} = \frac{\pi}{4}\bigl(e^{4\pi} - 1\bigr) - \frac{\pi}{8}\bigl(e^{4\pi} - 1\bigr) = \boxed{\frac{\pi}{8}\bigl(e^{4\pi} - 1\bigr)}.
    \]
\end{enumerate}

\medskip
\textbf{2) Volume d'un anneau / tore de révolution -- Théorème de Guldin (UCLouvain -- Juillet 2024, Question 3) :}\\
\emph{Énoncé :} Un anneau (tore) est un solide de révolution dont le profil dans le plan $x$-$y$ est un cercle de rayon $a$ excentré d'une distance $R > a$ par rapport à l'axe de rotation $Ox$, d'équation $x^2 + (y - R)^2 = a^2$. Quel est le volume $V$ de cet anneau ? (\emph{Rappel :} $\cos^2(\theta) = \frac{1}{2}(1 + \cos(2\theta))$).\\
\emph{Résolution :}
\begin{itemize}
    \item \textbf{Équations des demi-cercles supérieur $y_+(x)$ et inférieur $y_-(x)$ pour $x \in [-a, a]$ :}
    \[
        (y - R)^2 = a^2 - x^2 \iff y - R = \pm\sqrt{a^2 - x^2} \iff y_\pm(x) = R \pm \sqrt{a^2 - x^2} \ge 0.
    \]
    \item \textbf{Mise en intégrale du volume creux $V = V_{\text{ext}} - V_{\text{int}}$ :}
    Puisque $y_+^2(x) - y_-^2(x) = \bigl(R + \sqrt{a^2-x^2}\bigr)^2 - \bigl(R - \sqrt{a^2-x^2}\bigr)^2 = 4R\sqrt{a^2 - x^2}$ :
    \[
        V = \pi \int_{-a}^a \Bigl(y_+^2(x) - y_-^2(x)\Bigr)\dd x = 4\pi R \int_{-a}^a \sqrt{a^2 - x^2}\dd x = 8\pi R \int_0^a \sqrt{a^2 - x^2}\dd x.
    \]
    \item \textbf{Calcul de $\displaystyle\int_0^a \sqrt{a^2 - x^2}\dd x$ par substitution trigonométrique $x = a\sin(\theta)$ :}
    Posons $x = a\sin(\theta)$ pour $\theta \in \left[0, \frac{\pi}{2}\right]$, de sorte que $\mathrm{d}x = a\cos(\theta)\dd\theta$ et $\sqrt{a^2 - x^2} = \sqrt{a^2(1-\sin^2\theta)} = a\cos(\theta)$ (car $\cos\theta \ge 0$ sur $[0, \pi/2]$). Grâce à la formule de Carnot :
    \begin{align*}
        \int_0^a \sqrt{a^2 - x^2}\dd x &= a^2 \int_0^{\pi/2} \cos^2(\theta)\dd\theta = \frac{a^2}{2}\int_0^{\pi/2} \bigl(1 + \cos(2\theta)\bigr)\dd\theta \\
        &= \frac{a^2}{2}\left[\theta + \frac{\sin(2\theta)}{2}\right]_0^{\pi/2} = \frac{\pi a^2}{4},
    \end{align*}
    (ce qui est bien le quart de l'aire $\pi a^2$ d'un disque de rayon $a$ !). En multipliant par $8\pi R$ :
    \[
        \boxed{V = 8\pi R \cdot \frac{\pi a^2}{4} = 2\pi^2 R a^2 = (2\pi R)\cdot(\pi a^2)},
    \]
    soit le produit du périmètre $2\pi R$ du cercle directeur par l'aire $\pi a^2$ de la section circulaire.
\end{itemize}

\medskip
\textbf{3) L'œuf gigantesque fossilisé (UCLouvain -- Juillet 2026, Question 4) :}\\
\emph{Énoncé résumé :} Un œuf fossilisé est un solide de révolution de hauteur $H$ dont le rayon horizontal à la hauteur $h \in [0, H]$ vaut :
\[
    R(h) = A\sqrt{\sin\left(\frac{\pi h}{H}\right) + \frac{1}{4}\sin\left(\frac{2\pi h}{H}\right)}.
\]
Calculer la formule littérale du volume $V_{\text{œuf}}$, puis sa valeur numérique pour $H = 1,00\text{ m}$ et $A = 0,50\text{ m}$, et la comparer au volume $V_{\text{sphère}}$ d'une sphère de rayon $0,50\text{ m}$.\\
\emph{Résolution :} La présence de la racine carrée dans $R(h)$ est un cadeau puisque la formule du volume de révolution élève $R(h)$ au carré !
\begin{align*}
    V_{\text{œuf}} &= \pi \int_0^H \bigl(R(h)\bigr)^2\dd h = \pi A^2 \int_0^H \left(\sin\left(\frac{\pi h}{H}\right) + \frac{1}{4}\sin\left(\frac{2\pi h}{H}\right)\right)\dd h \\
    &= \pi A^2 \left[-\frac{H}{\pi}\cos\left(\frac{\pi h}{H}\right) - \frac{1}{4}\cdot\frac{H}{2\pi}\cos\left(\frac{2\pi h}{H}\right)\right]_0^H = -H A^2 \left[\cos\left(\frac{\pi h}{H}\right) + \frac{1}{8}\cos\left(\frac{2\pi h}{H}\right)\right]_0^H \\
    &= -H A^2 \left[\left(-1 + \frac{1}{8}\right) - \left(1 + \frac{1}{8}\right)\right] = \boxed{2 H A^2}.
\end{align*}
Numériquement, pour $H = 1,00\text{ m}$ et $A = 0,50\text{ m}$ : $\boxed{V_{\text{œuf}} = 2 \times 1,00 \times (0,50)^2 = 0,50\text{ m}^3}$, tandis que le volume de la sphère de rayon $R = 0,50\text{ m}$ vaut $\boxed{V_{\text{sphère}} = \frac{4}{3}\pi R^3 = \frac{4}{3}\pi\left(\frac{1}{2}\right)^3 = \frac{\pi}{6} \simeq \frac{3,1416}{6} \simeq 0,52\text{ m}^3}$. L'œuf est donc légèrement \textbf{moins volumineux} que la sphère.
\end{exemple}

% ----------------------------------------------------------------------------
\section{Exercices du Chapitre 7}
% ----------------------------------------------------------------------------

\emph{Les corrigés détaillés de tous ces exercices se trouvent au Chapitre 9, en fin d'ouvrage.}

\begin{enumerate}[label=\textbf{Exercice 7.\arabic* --}]
    \item \textbf{(Calculs d'intégrales définies)} Calculer la valeur exacte des intégrales définies suivantes :
    \begin{enumerate}
        \item $\displaystyle\int_0^{\pi/4} \tan(x)\dd x$
        \item $\displaystyle\int_1^e x\ln(x)\dd x$
        \item $\displaystyle\int_{-1}^1 \frac{x^3\,\cos(x)}{1 + x^2}\dd x$
        \item $\displaystyle\int_0^{\pi/2} \cos^2(x)\sin(x)\dd x$
    \end{enumerate}

    \item \textbf{(UCLouvain -- Septembre 2021, Q3e -- Étude et majoration de $\dfrac{\sin(x)}{x}$)}
    \begin{enumerate}
        \item Soit $g : \left]-\frac{\pi}{2}, \frac{\pi}{2}\right[ \;\to \R, x \mapsto g(x) = \tan(x) - x$. Exploiter la dérivée de $g$ pour montrer que $g$ n'a qu'une seule racine sur $\left]-\frac{\pi}{2}, \frac{\pi}{2}\right[$ et déterminer les zones où $g$ est positive ou négative.
        \item Soit $f : \R \to \R$ définie par $f(x) = \dfrac{\sin(x)}{x}$ si $x \neq 0$ et $f(0) = 1$. Sur base de (a), donner les zones de croissance et de décroissance de $f$ sur $\left]-\frac{\pi}{2}, \frac{\pi}{2}\right[$.
        \item En traitant séparément les cas $|x| \le 1 < \frac{\pi}{2}$ et $|x| > 1$, prouver que $f(x) \le 1$ pour tout $x \in \R$ (et donc que $\sin(x) \le x$ pour tout $x \ge 0$).
    \end{enumerate}

    \item \textbf{(Inégalités polynomiales de Taylor pour $\cos(x)$ et $\sin(x)$ par intégrations successives)}
    En partant de l'inégalité $\cos(s) \le 1$ valable pour tout $s \ge 0$ et en intégrant successivement entre $0$ et $x$ (pour $x \ge 0$), démontrer la chaîne d'inégalités :
    \[
        \forall x \ge 0 : \qquad \sin(x) \le x, \qquad 1 - \frac{x^2}{2} \le \cos(x) \le 1, \qquad x - \frac{x^3}{6} \le \sin(x) \le x.
    \]

    \item \textbf{(Aire comprise entre deux courbes)} Calculer l'aire géométrique du domaine plan fermé délimité par la parabole d'équation $y = 4 - x^2$ et la droite d'équation $y = x + 2$.

    \item \textbf{(Démonstration des formules de volume des solides classiques par intégration)}
    En utilisant la formule du volume d'un solide de révolution $V = \pi \int_a^b y^2(x)\dd x$, retrouver :
    \begin{enumerate}
        \item Le volume d'un \textbf{cône circulaire droit} de rayon de base $R$ et de hauteur $H$ ($V = \frac{1}{3}\pi R^2 H$).
        \item Le volume d'une \textbf{sphère} de rayon $R$ ($V = \frac{4}{3}\pi R^3$).
        \item Le volume d'une \textbf{calotte sphérique} de hauteur $h \in [0, R]$ découpée au sommet d'une sphère de rayon $R$ ($V = \pi h^2\left(R - \frac{h}{3}\right)$).
    \end{enumerate}

    \item \textbf{(Volume d'un solide de révolution logarithmique)} On considère la région plane délimitée par la courbe $y = \dfrac{\ln(x)}{\sqrt{x}}$, l'axe $Ox$ et les droites $x = 1$ et $x = e$. Calculer le volume exact $V$ du solide engendré par la rotation de cette région autour de l'axe $Ox$.
\end{enumerate}

% ============================================================================
% CHAPITRE 8 : FORMULAIRE SYNTHÉTIQUE ET AIDE-MÉMOIRE (VERSION IMPRIMABLE 4 PAGES)
% ============================================================================
\cleardoublepage
\newgeometry{top=1.3cm,bottom=1.4cm,left=1.4cm,right=1.4cm}
\refstepcounter{chapter}
\addcontentsline{toc}{chapter}{\protect\numberline{\thechapter}Formulaire Synthétique et Aide-Mémoire (Imprimable -- 4 pages)}
\chaptermark{Formulaire Synthétique et Aide-Mémoire}

% ============================================================================
% PAGE 1/4 : FONCTIONS, DOMAINES, LIMITES ET ASYMPTOTES (1 PAGE A4 IMPRIMABLE)
% ============================================================================
\thispagestyle{plain}
\refstepcounter{section}
\addcontentsline{toc}{section}{\protect\numberline{\thesection}Fonctions, Domaines, Limites et Asymptotes}

\begin{center}
    {\LARGE\bfseries\color{uclblue} Fonctions, Domaines, Limites et Asymptotes}\\[2pt]
    {\small\color{darkgray} Examen d'entrée Ingénieur Civil -- UCLouvain (Fiche récapitulative d'Analyse 1/4)}
    \vspace{2pt}
    \hrule height 1pt
\end{center}
\vspace{-6pt}

{\fontsize{9.4pt}{11.4pt}\selectfont
\setlength{\columnsep}{18pt}
\begin{multicols}{2}

\noindent{\large\bfseries\color{uclblue} 1. Conditions d'Existence (Domaine $\D$)}
\begin{itemize}[leftmargin=1.2em, itemsep=1.5pt, topsep=2pt]
    \item \textbf{Quotient :} $\dfrac{N(x)}{D(x)} \implies \mathbf{D(x) \neq 0}$\\
    $\tan(u) \implies u \neq \frac{\pi}{2} + k\pi$, \quad $\cot(u) \implies u \neq k\pi$
    \item \textbf{Racine paire :} $\sqrt[2n]{u(x)} \implies \mathbf{u(x) \ge 0}$ \quad ($\frac{1}{\sqrt{u(x)}} \implies u(x) > 0$)
    \item \textbf{Logarithme :} $\ln\bigl(u(x)\bigr) \implies \mathbf{u(x) > 0}$
    \item \textbf{Cyclométriques :} $\arcsin(u), \arccos(u) \implies \mathbf{-1 \le u(x) \le 1}$
    \item \textbf{Équation irrationnelle :} $\sqrt{A(x)} = B(x) \iff B(x) \ge 0 \text{ et } A(x) = B^2(x)$
\end{itemize}

\vspace{3pt}
\noindent{\large\bfseries\color{uclblue} 2. Symétries et Périodicité}
\begin{itemize}[leftmargin=1.2em, itemsep=1.5pt, topsep=2pt]
    \item \textbf{Fonction paire ($\D$ sym.) :} $f(-x) = f(x)$ (axe $Oy$).\\
    Si $f$ est paire : $f'$ est impaire, $f''$ est paire, $\int_{-a}^a f(x)\dd x = 2\int_0^a f(x)\dd x$.
    \item \textbf{Fonction impaire ($\D$ sym.) :} $f(-x) = -f(x)$ (centre $O$).\\
    Si $f$ est impaire : $f'$ est paire, $f''$ est impaire, $\int_{-a}^a f(x)\dd x = 0$.
    \item \textbf{Périodique de période $T > 0$ :} $f(x+T) = f(x)$.\\
    $\sin(\omega x), \cos(\omega x) \to T = \frac{2\pi}{\omega}$ ; \; $\tan(\omega x), \sin^2(\omega x) \to T = \frac{\pi}{\omega}$.
\end{itemize}

\vspace{3pt}
\noindent{\large\bfseries\color{uclblue} 3. Logarithme $\ln(x)$ et Exponentielle $e^x$}
\begin{itemize}[leftmargin=1.2em, itemsep=1.5pt, topsep=2pt]
    \item \textbf{Réciprocité :} $y = e^x \iff x = \ln(y)$ \; ($x \in \R,\; y > 0$)\\
    $\ln(e^x) = x \; (\forall x \in \R)$ \quad et \quad $e^{\ln(y)} = y \; (\forall y > 0)$
    \item \textbf{Valeurs clés :} $\ln(1) = 0$, \; $\ln(e) = 1$, \; $e^0 = 1$, \; $e > 0$ ($\forall x, e^x > 0$)\\
    $e \approx 2,718$, \; $e^{-1} \approx 0,368$, \; $e^{-2} \approx 0,135$, \; $\ln(2) \approx 0,693$
    \item \textbf{Règles de $\ln$ ($a, b > 0$) :} $\ln(ab) = \ln a + \ln b$, \; $\ln\left(\frac{1}{a}\right) = -\ln a$,\\
    $\ln\left(\frac{a}{b}\right) = \ln a - \ln b$, \; $\ln(a^p) = p\ln a$, \; $\ln(x^2) = 2\ln(|x|)$ ($x \neq 0$)
    \item \textbf{Signe de $\ln(x)$ :} $<0$ sur $]0, 1[$, \; $=0$ en $x=1$, \; $>0$ sur $]1, +\infty[$
    \item \textbf{Règles de $\exp$ :} $e^{a+b} = e^a e^b$, \; $e^{-a} = \frac{1}{e^a}$, \; $e^{a-b} = \frac{e^a}{e^b}$, \; $(e^a)^p = e^{ap}$
\end{itemize}

\vspace{3pt}
\noindent{\large\bfseries\color{uclblue} 4. Fonction Réciproque $f^{-1}$ (ex: Lambert $W_0$)}
Si $f : I \to J$ est continue et strictement monotone sur un intervalle $I$ :
\begin{itemize}[leftmargin=1.2em, itemsep=1.5pt, topsep=2pt]
    \item $f^{-1} : J = f(I) \to I$ est continue et a le \textbf{même sens de variation} que $f$.
    \item \textbf{Dérivée :} $\bigl(f^{-1}\bigr)'(y) = \dfrac{1}{f'\bigl(f^{-1}(y)\bigr)}$ \quad (existe ssi $f'(x) \neq 0$).
\end{itemize}

\columnbreak

\noindent{\large\bfseries\color{uclblue} 5. Limites Fondamentales \& Croissances Comparées}
\begin{itemize}[leftmargin=1.2em, itemsep=2pt, topsep=2pt]
    \item \textbf{Définition ($\lim_{x \to +\infty} f(x) = c$) :}\\
    $\forall \varepsilon > 0, \;\exists M \;\text{ tel que }\; x > M \implies |f(x) - c| \le \varepsilon$.
    \item \textbf{Limites remarquables en $u \to 0$ :}
    \[
        \lim_{u \to 0} \frac{\sin(u)}{u} = 1, \quad \lim_{u \to 0} \frac{\tan(u)}{u} = 1, \quad \lim_{u \to 0} \frac{1 - \cos(u)}{u^2} = \frac{1}{2}
    \]
    \[
        \lim_{u \to 0} \frac{e^u - 1}{u} = 1, \qquad \lim_{u \to 0} \frac{\ln(1+u)}{u} = 1
    \]
    \item \textbf{Croissances comparées ($\forall p > 0$) :} ($\exp \gg x^p \gg \ln$)
    \[
        \lim_{x \to +\infty} \frac{e^x}{x^p} = +\infty, \qquad \lim_{x \to +\infty} x^p e^{-x} = 0
    \]
    \[
        \lim_{x \to +\infty} \frac{\ln(x)}{x^p} = 0, \qquad \lim_{y \to 0^+} y^p \ln(y) = 0
    \]
    \item \textbf{Fonction bornée ($|\sin u| \le 1$) :} Si $\lim v(x) = \pm\infty$, $\lim \dfrac{\sin(\dots)}{v(x)} = 0$.
\end{itemize}

\vspace{3pt}
\noindent{\large\bfseries\color{uclblue} 6. Continuité et Théorème des Valeurs Intermédiaires}
\begin{itemize}[leftmargin=1.2em, itemsep=1.5pt, topsep=2pt]
    \item \textbf{Continuité en $a \in \D$ :} $\lim_{x \to a^-} f(x) = \lim_{x \to a^+} f(x) = f(a)$.
    \item \textbf{TVI / Bolzano :} Si $f$ est continue sur $[a, b]$ et $f(a)\cdot f(b) < 0$ (changement de signe), alors $\exists c \in \,]a, b[$ tel que $\mathbf{f(c) = 0}$ (unique si $f$ est strictement monotone).
\end{itemize}

\vspace{3pt}
\noindent{\large\bfseries\color{uclblue} 7. Les Trois Types d'Asymptotes (AV, AH, AO)}
\begin{itemize}[leftmargin=1.2em, itemsep=1.5pt, topsep=2pt]
    \item \textbf{Asymptote Verticale ($\text{AV} \equiv x = a$) :} $\lim_{x \to a^\pm} f(x) = \pm\infty$.
    \item \textbf{Asymptote Horizontale ($\text{AH} \equiv y = b$) :} $\lim_{x \to \pm\infty} f(x) = b \in \R$.
    \item \textbf{Asymptote Oblique ($\text{AO} \equiv y = mx + p$, $m \neq 0$) :}
    \[
        m = \lim_{x \to \pm\infty} \frac{f(x)}{x} \in \R_0, \qquad p = \lim_{x \to \pm\infty} \bigl(f(x) - mx\bigr) \in \R
    \]
    \emph{Note :} Une fonction périodique non constante n'a ni AH ni AO.
\end{itemize}

\end{multicols}
}

\newpage
% ============================================================================
% PAGE 2/4 : DÉRIVÉES, L'HOSPITAL, ROLLE, TAF ET ÉTUDE DE FONCTION (1 PAGE A4)
% ============================================================================
\thispagestyle{plain}
\refstepcounter{section}
\addcontentsline{toc}{section}{\protect\numberline{\thesection}Dérivées, Règle de L'Hospital, Rolle, TAF et Étude de Fonction}

\begin{center}
    {\LARGE\bfseries\color{uclblue} Dérivées, L'Hospital, Rolle, TAF et Étude de Fonction}\\[2pt]
    {\small\color{darkgray} Examen d'entrée Ingénieur Civil -- UCLouvain (Fiche récapitulative d'Analyse 2/4)}
    \vspace{2pt}
    \hrule height 1pt
\end{center}
\vspace{-6pt}

{\fontsize{9.0pt}{10.8pt}\selectfont
\setlength{\columnsep}{18pt}
\begin{multicols}{2}

\noindent{\large\bfseries\color{uclblue} 1. Nombre Dérivé, Tangente et Règles Opératoires}
\begin{itemize}[leftmargin=1.2em, itemsep=1.2pt, topsep=1.5pt]
    \item \textbf{Définition :} $f'(a) = \lim_{x \to a} \dfrac{f(x) - f(a)}{x - a} = \lim_{h \to 0} \dfrac{f(a+h) - f(a)}{h}$
    \item \textbf{Tangente en $x_0$ :} $T_{x_0} \equiv \mathbf{y = f(x_0) + f'(x_0)\,(x - x_0)}$
    \item \textbf{Dérivable $\implies$ Continue} (mais réciproque fausse : $|x|$ ou $\sqrt{|x|}$ en $0$).
    \item \textbf{Produit :} $(u\cdot v)' = u'v + uv'$
    \item \textbf{Inverse \& Quotient :} $\left(\dfrac{1}{v}\right)' = -\dfrac{v'}{v^2}$, \qquad $\left(\dfrac{u}{v}\right)' = \dfrac{u'v - uv'}{v^2}$
    \item \textbf{Composée (Chaîne) :} $\Bigl(f\bigl(u(x)\bigr)\Bigr)' = f'\bigl(u(x)\bigr)\cdot u'(x)$
\end{itemize}

\vspace{2pt}
\noindent{\large\bfseries\color{uclblue} 2. Tableau des Dérivées Composées}
\begin{center}
\vspace{-4pt}
\renewcommand{\arraystretch}{1.35}
\setlength{\tabcolsep}{3.5pt}
\begin{tabular}{|l|l|}
\hline
\textbf{Fonction $f(u(x))$} & \textbf{Dérivée $\bigl(f(u(x))\bigr)'$} \\ \hline
$\bigl(u(x)\bigr)^n \;\; (n \in \R)$ & $n\,\bigl(u(x)\bigr)^{n-1} \cdot u'(x)$ \\ \hline
$\sqrt{u(x)}$ & $\dfrac{u'(x)}{2\sqrt{u(x)}}$ \\ \hline
$e^{u(x)}$ \; / \; $a^{u(x)}$ & $u'(x)\,e^{u(x)}$ \; / \; $(\ln a)\,u'(x)\,a^{u(x)}$ \\ \hline
$\ln\bigl(|u(x)|\bigr)$ & $\dfrac{u'(x)}{u(x)}$ \\ \hline
$\sin\bigl(u(x)\bigr)$ / $\cos\bigl(u(x)\bigr)$ & $u'(x)\cos\bigl(u(x)\bigr)$ / $-u'(x)\sin\bigl(u(x)\bigr)$ \\ \hline
$\tan\bigl(u(x)\bigr)$ & $\dfrac{u'(x)}{\cos^2(u(x))} = u'(x)\bigl(1 + \tan^2(u(x))\bigr)$ \\ \hline
$\cot\bigl(u(x)\bigr)$ & $-\dfrac{u'(x)}{\sin^2(u(x))}$ \\ \hline
$\arcsin\bigl(u(x)\bigr)$ & $\dfrac{u'(x)}{\sqrt{1 - u^2(x)}}$ \; ($\arccos : \text{signe } -$) \\ \hline
$\arctan\bigl(u(x)\bigr)$ & $\dfrac{u'(x)}{1 + u^2(x)}$ \\ \hline
\end{tabular}
\end{center}

\vspace{2pt}
\noindent{\large\bfseries\color{uclblue} 3. Règle de L'Hospital (Formes $\frac{0}{0}$ et $\frac{\infty}{\infty}$)}
Si $\lim_{x \to a} \dfrac{f(x)}{g(x)}$ est une F.I. $\left[\dfrac{0}{0}\right]$ ou $\left[\dfrac{\infty}{\infty}\right]$ et si $\lim_{x \to a} \dfrac{f'(x)}{g'(x)}$ existe :
\[
    \boxed{\lim_{x \to a} \frac{f(x)}{g(x)} \stackrel{\text{H}}{=} \lim_{x \to a} \frac{f'(x)}{g'(x)}}
\]
\emph{Astuces :} Simplifier l'expression avant de dériver (ex: $\ln(1+x)+\ln(1-x) = \ln(1-x^2)$) ; pour $[0\cdot\infty]$ comme $x\ln(x)$, écrire $\frac{\ln(x)}{1/x}$ ; pour $[u^v]$, écrire $e^{v\ln(u)}$.

\columnbreak

\noindent{\large\bfseries\color{uclblue} 4. Théorèmes de Rolle et des Accroissements Finis}
\begin{itemize}[leftmargin=1.2em, itemsep=1.2pt, topsep=1.5pt]
    \item \textbf{Théorème de Rolle :} Si $f : [a, b] \to \R$ est :
    \begin{enumerate}[leftmargin=1.5em, itemsep=0.5pt, topsep=0.5pt]
        \item \textbf{continue} sur l'intervalle fermé $[a, b]$,
        \item \textbf{dérivable} sur l'intervalle ouvert $]a, b[$,
        \item telle que $\mathbf{f(a) = f(b)}$,
    \end{enumerate}
    alors $\mathbf{\exists c \in \,]a, b[\; \text{ tel que } f'(c) = 0}$.
    \item \textbf{Théorème des Accroissements Finis (TAF) :} Si $f$ est \textbf{continue sur $[a, b]$} et \textbf{dérivable sur $]a, b[$}, alors $\mathbf{\exists c \in \,]a, b[}$ tel que :
    \[
        \boxed{f'(c) = \frac{f(b) - f(a)}{b - a}}
    \]
    \item \textbf{Astuce en cascade :} Si $f(x_0) = f(x_1) = f(x_2)$, Rolle sur $[x_0, x_1]$ et $[x_1, x_2]$ donne deux zéros distincts $c_1 < x_1 < c_2$ de $f'$, puis Rolle appliqué à $f'$ sur $[c_1, c_2]$ donne $d \in \,]c_1, c_2[$ tel que $\mathbf{f''(d) = 0}$ !
\end{itemize}

\vspace{2pt}
\noindent{\large\bfseries\color{uclblue} 5. Variations, Extrema, Concavité (Question 1)}
\begin{itemize}[leftmargin=1.2em, itemsep=1.2pt, topsep=1.5pt]
    \item \textbf{Croissance / Décroissance :} $f'(x) > 0 \implies \nearrow$, \quad $f'(x) < 0 \implies \searrow$.
    \item \textbf{Extrema locaux (en un point intérieur $x^*$ où $f'(x^*) = 0$) :}
    \begin{itemize}[leftmargin=1em, itemsep=0.5pt, topsep=0.5pt]
        \item $f'$ passe de $-$ à $+$ (ou $f''(x^*) > 0$) $\implies$ \textbf{Minimum local} ($\cup$)
        \item $f'$ passe de $+$ à $-$ (ou $f''(x^*) < 0$) $\implies$ \textbf{Maximum local} ($\cap$)
        \item $f'$ ne change pas de signe $\implies$ \textbf{Point d'inflexion à tangente horiz.}
    \end{itemize}
    \item \textbf{Concavité et Points d'Inflexion (PI) :}
    \begin{itemize}[leftmargin=1em, itemsep=0.5pt, topsep=0.5pt]
        \item $f''(x) > 0 \implies$ Concavité positive / vers le haut ($\cup$, convexe)
        \item $f''(x) < 0 \implies$ Concavité négative / vers le bas ($\cap$, concave)
        \item $f''(x_0) = 0$ \textbf{en changeant de signe} $\implies$ \textbf{Point d'inflexion (PI)}
    \end{itemize}
    \item \textbf{Ordre logique officiel du tableau de variations :}\\
    Ligne 1 : $x$ $\to$ Ligne 2 : $f'(x)$ $\to$ Ligne 3 : $f''(x)$ $\to$ Ligne 4 : $f(x)$ (avec valeurs exactes ET décimales approchées, flèches $\nearrow\searrow$, concavité $\cup\cap$).
\end{itemize}

\vspace{2pt}
\noindent{\large\bfseries\color{uclblue} 6. Méthode d'Optimisation Appliquée (Question 2)}
\begin{enumerate}[leftmargin=1.5em, itemsep=0.5pt, topsep=0.5pt]
    \item Identifier la fonction objectif à maximiser/minimiser et la contrainte.
    \item Éliminer les variables liées pour avoir une fonction d'\textbf{une seule variable} $\Phi(x)$ (en gardant les constantes sous forme de lettres !).
    \item Résoudre $\Phi'(x^*) = 0$ et vérifier le signe de $\Phi''(x^*)$ (ou le tableau de $\Phi'$).
    \item Remplacer les valeurs numériques à la toute fin et vérifier les unités.
\end{enumerate}

\end{multicols}
}

\newpage
% ============================================================================
% PAGE 3/4 : PRIMITIVES ET TECHNIQUES D'INTÉGRATION (1 PAGE A4 IMPRIMABLE)
% ============================================================================
\thispagestyle{plain}
\refstepcounter{section}
\addcontentsline{toc}{section}{\protect\numberline{\thesection}Primitives et Techniques d'Intégration}

\begin{center}
    {\LARGE\bfseries\color{uclblue} Primitives et Techniques d'Intégration (Question 3)}\\[2pt]
    {\small\color{darkgray} Examen d'entrée Ingénieur Civil -- UCLouvain (Fiche récapitulative d'Analyse 3/4)}
    \vspace{2pt}
    \hrule height 1pt
\end{center}
\vspace{-6pt}

{\fontsize{9.3pt}{11.3pt}\selectfont
\setlength{\columnsep}{18pt}
\begin{multicols}{2}

\noindent{\large\bfseries\color{uclblue} 1. Primitives Immédiates et Composées ($u = u(x)$)}
\begin{center}
\vspace{-4pt}
{\footnotesize
\renewcommand{\arraystretch}{1.42}
\setlength{\tabcolsep}{1.6pt}
\begin{tabular}{|l|l|}
\hline
\textbf{Intégrale $\displaystyle\int f(u)\,u'\dd x$} & \textbf{Primitive $F(u(x)) + C$} \\ \hline
$\displaystyle\int u^p\,u'\dd x \;\; (p \neq -1)$ & $\dfrac{u^{p+1}}{p+1} + C$ \\ \hline
$\displaystyle\int \frac{u'}{u}\dd x$ & $\ln(|u|) + C$ \\ \hline
$\displaystyle\int \frac{u'}{\sqrt{u}}\dd x$ \; / \; $\displaystyle\int \frac{u'}{u^2}\dd x$ & $2\sqrt{u} + C$ \; / \; $-\dfrac{1}{u} + C$ \\ \hline
$\displaystyle\int u'\,e^u\dd x$ & $e^u + C$ \\ \hline
$\displaystyle\int u'\cos(u)\dd x$ & $\sin(u) + C$ \\ \hline
$\displaystyle\int u'\sin(u)\dd x$ & $-\cos(u) + C$ \\ \hline
$\displaystyle\int \frac{u'}{\cos^2(u)}\dd x = \int u'(1+\tan^2 u)\dd x$ & $\tan(u) + C$ \\ \hline
$\displaystyle\int \frac{u'}{\sin^2(u)}\dd x = \int u'(1+\cot^2 u)\dd x$ & $-\cot(u) + C$ \\ \hline
$\displaystyle\int \frac{u'}{1 + u^2}\dd x$ & $\arctan(u) + C$ \\ \hline
$\displaystyle\int \frac{u'}{\sqrt{1 - u^2}}\dd x$ & $\arcsin(u) + C$ \\ \hline
$\displaystyle\int \tan(x)\dd x$ & $-\ln(|\cos x|) + C$ \\ \hline
$\displaystyle\int \cot(x)\dd x$ & $\ln(|\sin x|) + C$ \\ \hline
$\displaystyle\int \ln(x)\dd x$ & $x\ln(x) - x + C$ \\ \hline
\end{tabular}
}
\end{center}

\vspace{3pt}
\noindent{\large\bfseries\color{uclblue} 2. Intégration par Substitution (Changement de Variable)}
\begin{itemize}[leftmargin=1.2em, itemsep=1.5pt, topsep=2pt]
    \item \textbf{Principe :} Poser $\mathbf{u = u(x)}$ et remplacer $\mathbf{\mathrm{d}u = u'(x)\dd x}$ :
    \[
        \int f\bigl(u(x)\bigr)\,u'(x)\dd x = \int f(u)\dd u = F\bigl(u(x)\bigr) + C
    \]
    \item \textbf{Exemples types UCLouvain :}
    \begin{itemize}[leftmargin=1em, itemsep=1pt, topsep=1pt]
        \item $\int \frac{\cos x}{1+\sin^2 x}\dd x$ : $u = \sin x \implies \arctan(\sin x) + C$
        \item $\int \frac{1}{x\ln x}\dd x$ : $u = \ln x \implies \ln(|\ln x|) + C$
        \item $\int \frac{x^5}{(x^2+1)^5}\dd x$ : $u = x^2+1, x\dd x = \frac{1}{2}\dd u, x^4 = (u-1)^2$
        \item $\int \sqrt{a^2 - x^2}\dd x$ : $x = a\sin\theta, \mathrm{d}x = a\cos\theta\dd\theta \implies a^2\int \cos^2\theta\dd\theta$
    \end{itemize}
\end{itemize}

\columnbreak

\noindent{\large\bfseries\color{uclblue} 3. Intégration par Parties (IPP)}
\[
    \boxed{\int u(x)\,v'(x)\dd x = u(x)\,v(x) - \int v(x)\,u'(x)\dd x}
\]
\begin{itemize}[leftmargin=1.2em, itemsep=1.5pt, topsep=2pt]
    \item \textbf{Priorité pour le choix de $u(x)$ (à dériver) -- Règle ALPES :}
    \begin{enumerate}[leftmargin=1.4em, itemsep=1pt, topsep=1pt]
        \item \textbf{A}rc ($\arcsin, \arccos, \arctan$) $\to$ toujours dériver !
        \item \textbf{L}ogarithme ($\ln(x), \ln^2(x)$) $\to$ toujours dériver face à $x^n$ !
        \item \textbf{P}olynôme ($x, x^2$) $\to$ dériver face à $e^{ax}, \sin, \cos, \frac{1}{\cos^2 x}$.
        \item \textbf{E}xponentielle et \textbf{S}inus/Cosinus $\to$ intégrer ($\mathrm{d}v$) ; si produit $e^{ax}\cos(bx)$, faire \textbf{2 IPP cycliques} pour retrouver $\mathcal{I}$.
    \end{enumerate}
    \item \textbf{Exemples types UCLouvain :}
    \begin{itemize}[leftmargin=1em, itemsep=1pt, topsep=1pt]
        \item $\int \frac{x}{\cos^2 x}\dd x$ : $u=x, \mathrm{d}v=\frac{\mathrm{d}x}{\cos^2 x} \implies x\tan x + \ln(|\cos x|) + C$
        \item $\int \frac{\ln x}{x^2}\dd x$ : $u=\ln x, \mathrm{d}v=\frac{\mathrm{d}x}{x^2} \implies -\frac{\ln x + 1}{x} + C$
        \item $\int x\ln^2(x)\dd x$ : 2 IPP $\implies \frac{x^2}{2}\left(\ln^2(x) - \ln(x) + \frac{1}{2}\right) + C$
    \end{itemize}
\end{itemize}

\vspace{3pt}
\noindent{\large\bfseries\color{uclblue} 4. Fractions Simples et Primitives Trigonométriques}
\begin{itemize}[leftmargin=1.2em, itemsep=2pt, topsep=2pt]
    \item \textbf{Décomposition de $\dfrac{1}{1 - t^2}$ et $\dfrac{1}{y^2 - 1}$ :}
    \[
        \int \frac{\mathrm{d}t}{1 - t^2} = \frac{1}{2}\int \left(\frac{1}{1-t} + \frac{1}{1+t}\right)\dd t = \frac{1}{2}\ln\left(\left|\frac{1+t}{1-t}\right|\right) + C
    \]
    \[
        \int \frac{2\dd y}{y^2 - 1} = \int \left(\frac{1}{y-1} - \frac{1}{y+1}\right)\dd y = \ln\left(\left|\frac{y-1}{y+1}\right|\right) + C
    \]
    \item \textbf{Puissance impaire de $\sin$ ou $\cos$ :}\\
    Isoler $\sin(x)\dd x = -\mathrm{d}(\cos x)$ et remplacer $\sin^2(x) = 1 - \cos^2(x)$ (avec $u = \cos x$).
    \item \textbf{Puissance paire $\sin^2(x)$ ou $\cos^2(x)$ (Carnot) :}
    \[
        \cos^2(x) = \frac{1 + \cos(2x)}{2}, \qquad \sin^2(x) = \frac{1 - \cos(2x)}{2}
    \]
    \item \textbf{Dénominateur $1 + \cos(x)$ ou $1 + \sin(2x)$ :}
    \[
        1 + \cos(x) = 2\cos^2\left(\frac{x}{2}\right), \qquad 1 + \sin(2x) = \bigl(\cos x + \sin x\bigr)^2
    \]
    \item \textbf{Primitives de $\dfrac{1}{\cos(x)}$ et $\dfrac{1}{\sin(u)}$ :}
    \[
        \int \frac{\mathrm{d}x}{\cos(x)} = \int \frac{\cos(x)\dd x}{1 - \sin^2(x)} = \frac{1}{2}\ln\left(\frac{1+\sin x}{1-\sin x}\right) + C
    \]
    \[
        \int \frac{\mathrm{d}u}{\sin(u)} = \frac{1}{2}\ln\left(\frac{1-\cos u}{1+\cos u}\right) + C = \ln\left(\left|\tan\frac{u}{2}\right|\right) + C
    \]
\end{itemize}

\end{multicols}
}

\newpage
% ============================================================================
% PAGE 4/4 : INTÉGRALES DÉFINIES, INÉGALITÉS, AIRES ET VOLUMES (1 PAGE A4)
% ============================================================================
\thispagestyle{plain}
\refstepcounter{section}
\addcontentsline{toc}{section}{\protect\numberline{\thesection}Intégrales Définies, Inégalités, Aires et Volumes de Révolution}

\begin{center}
    {\LARGE\bfseries\color{uclblue} Intégrales Définies, Inégalités, Aires et Volumes}\\[2pt]
    {\small\color{darkgray} Examen d'entrée Ingénieur Civil -- UCLouvain (Fiche récapitulative d'Analyse 4/4)}
    \vspace{2pt}
    \hrule height 1pt
\end{center}
\vspace{-6pt}

{\fontsize{9.4pt}{11.4pt}\selectfont
\setlength{\columnsep}{18pt}
\begin{multicols}{2}

\noindent{\large\bfseries\color{uclblue} 1. Théorème Fondamental et Intégrale Définie}
\begin{itemize}[leftmargin=1.2em, itemsep=2pt, topsep=2pt]
    \item \textbf{Formule de Newton-Leibniz (pour $F'(x) = f(x)$) :}
    \[
        \boxed{\int_a^b f(x)\dd x = \Bigl[F(x)\Bigr]_a^b = F(b) - F(a)}
    \]
    \item \textbf{Dérivée par rapport à la borne supérieure $x$ :}
    \[
        \frac{\mathrm{d}}{\mathrm{d}x}\left(\int_a^x f(s)\dd s\right) = f(x)
    \]
    \item \textbf{Valeur moyenne de $f$ sur $[a, b]$ :}
    \[
        \bar{f} = \frac{1}{b - a}\int_a^b f(t)\dd t
    \]
    \item \textbf{Changement de variable dans une intégrale définie :}\\
    Si $u = u(x)$, penser à \textbf{changer aussi les bornes} ($x=a \to u=u(a)$ et $x=b \to u=u(b)$) :
    \[
        \int_a^b f\bigl(u(x)\bigr)\,u'(x)\dd x = \int_{u(a)}^{u(b)} f(u)\dd u
    \]
\end{itemize}

\vspace{4pt}
\noindent{\large\bfseries\color{uclblue} 2. Démonstration d'Inégalités Fonctionnelles}
\begin{itemize}[leftmargin=1.2em, itemsep=2pt, topsep=2pt]
    \item \textbf{Méthode 1 : Par intégration sur $[a, x]$ (avec $x \ge a$) :}
    \[
        \forall s \in [a, x], \; f'(s) \le g'(s) \implies \int_a^x f'(s)\dd s \le \int_a^x g'(s)\dd s
    \]
    \begin{itemize}[leftmargin=1em, itemsep=1.5pt, topsep=1pt]
        \item \textbf{Inégalité de $\ln(x)$ ($\forall x > 0$) :} Pour $s \ge 1$, $\frac{1}{s^2} \le \frac{1}{s} \le 1$. En intégrant sur $[1, x]$ ($x \ge 1$, ou sur $[x, 1]$ si $0 < x \le 1$) :
        \[
            \boxed{1 - \frac{1}{x} \le \ln(x) \le x - 1 \quad (\forall x > 0)}
        \]
        \item \textbf{Inégalité de $e^x$ ($\forall x \ge 0$) :} Pour $s \ge 0$, $e^s \ge 1 \implies \int_0^x e^s\dd s \ge \int_0^x 1\dd s$ :
        \[
            \boxed{e^x \ge 1 + x \;\; (\forall x \ge 0)} \implies e^y \ge \left(1 + \frac{y}{n}\right)^n \; (y > 0)
        \]
    \end{itemize}
    \item \textbf{Méthode 2 : Par étude du signe de la dérivée :}\\
    Pour montrer $f(x) \le g(x)$ sur $[a, +\infty[$, poser $\Delta(x) = g(x) - f(x)$, vérifier $\Delta(a) = 0$ et montrer $\Delta'(x) = g'(x) - f'(x) \ge 0$ pour tout $x \ge a$.
\end{itemize}

\columnbreak

\noindent{\large\bfseries\color{uclblue} 3. Aires Planes et Volumes de Révolution}
\begin{itemize}[leftmargin=1.2em, itemsep=2pt, topsep=2pt]
    \item \textbf{Aire entre deux courbes sur $[a, b]$ :}
    \[
        \mathcal{A} = \int_a^b \bigl(y_{\text{haut}}(x) - y_{\text{bas}}(x)\bigr)\dd x
    \]
    \item \textbf{Volume d'un solide plein de révolution autour de $Ox$ :}\\
    Rotation du profil de rayon $R(x) = |f(x)|$ pour $x \in [a, b]$ :
    \[
        \boxed{V = \pi \int_a^b \bigl(R(x)\bigr)^2\dd x = \pi \int_a^b \bigl(f(x)\bigr)^2\dd x}
    \]
    \item \textbf{Volume d'un solide creux (entre $y_+(x) \ge y_-(x) \ge 0$) :}
    \[
        \boxed{V = \pi \int_a^b \Bigl(\bigl(y_+(x)\bigr)^2 - \bigl(y_-(x)\bigr)^2\Bigr)\dd x}
    \]
    \item \textbf{Volumes classiques à connaître :}
    \begin{itemize}[leftmargin=1em, itemsep=1.5pt, topsep=1pt]
        \item Cylindre de rayon $R$ et hauteur $H$ : $V = \pi R^2 H$
        \item Cône de rayon $R$ et hauteur $H$ : $V = \dfrac{1}{3}\pi R^2 H$
        \item Sphère de rayon $R$ : $V = \dfrac{4}{3}\pi R^3$ \quad (Surface : $S = 4\pi R^2$)
        \item Tore (anneau de rayon de tube $a$ à distance $R > a$) :\\
        $V = (2\pi R)\cdot(\pi a^2) = 2\pi^2 R a^2$
    \end{itemize}
\end{itemize}

\vspace{4pt}
\noindent{\large\bfseries\color{uclblue} 4. Équations Différentielles \& Fonctions Liées}
Reconnaître une dérivée composée pour intégrer directement :
\begin{itemize}[leftmargin=1.2em, itemsep=2pt, topsep=2pt]
    \item $f'(x) = -\bigl(f(x)\bigr)^2 \iff -\dfrac{f'(x)}{f^2(x)} = 1 \iff \left(\dfrac{1}{f(x)}\right)' = 1$\\
    $\implies \dfrac{1}{f(x)} = x + C \implies f(x) = \dfrac{1}{x + C}$
    \item $\dfrac{g'(x)}{g(x)} = f(x)f'(x) \iff \bigl(\ln g(x)\bigr)' = \left(\dfrac{1}{2}f^2(x)\right)'$\\
    $\implies \ln\left(\dfrac{g(x)}{g(0)}\right) = \dfrac{1}{2}\bigl(f^2(x) - f^2(0)\bigr)$
    \item $f''(x) = -f(x) \implies \bigl(f^2(x) + f'^2(x)\bigr)' = 2f'(x)\bigl(f(x) + f''(x)\bigr) = 0$\\
    $\implies f^2(x) + f'^2(x) = k^2 = \text{constante}$
\end{itemize}

\end{multicols}
}

\restoregeometry

% ============================================================================
% ANNEXE A : ÉLÉMENTS DE LOGIQUE ET MÉTHODES DE DÉMONSTRATION EN ANALYSE
% ============================================================================
\chapter*{Annexe A : Logique et Méthodes de Démonstration en Analyse}
\addcontentsline{toc}{chapter}{Annexe A : Logique et Méthodes de Démonstration en Analyse}
\markboth{Annexe A : Logique et Démonstrations}{Annexe A : Logique et Démonstrations}

\emph{Bien que la logique formelle ne fasse pas l'objet d'un chapitre séparé au programme de l'examen d'admission, plusieurs questions d'Analyse de l'UCLouvain (telles que Juillet 2021 Q3d, Juillet 2023 Série 1 Q5, Juillet 2024 Q5, Septembre 2024 Q5 ou Juillet 2025 Q4) évaluent directement la capacité à manipuler rigoureusement une définition avec quantificateurs ($\forall, \exists$), à distinguer une implication ($\implies$) d'une équivalence ($\iff$), à utiliser une contraposée ou à construire un contre-exemple.}

% ----------------------------------------------------------------------------
\section*{A.1 Connecteurs logiques (« non », « et », « ou »), implication et équivalence}
% ----------------------------------------------------------------------------

\begin{theoreme}[Opérateurs logiques, implication ($\implies$) et équivalence ($\iff$)]
Soient $P$ et $Q$ deux assertions mathématiques :
\begin{enumerate}
    \item \textbf{Conjonction (« $P \text{ et } Q$ ») vs Disjonction (« $P \text{ ou } Q$ ») :}
    \begin{itemize}
        \item « $P \text{ et } Q$ » (noté par une accolade système) correspond à l'\textbf{intersection} $\cap$ des ensembles de validité : les deux conditions doivent être vraies simultanément.
        \item « $P \text{ ou } Q$ » correspond à l'\textbf{union} $\cup$ : il suffit qu'au moins l'une des deux conditions soit vraie (par exemple : $A(x)\cdot B(x) = 0 \iff A(x) = 0 \text{ \textbf{ou} } B(x) = 0$).
    \end{itemize}
    \item \textbf{Lois de De Morgan (négation d'un « et » et d'un « ou ») :}
    \[
        \boxed{\text{non}(P \text{ ou } Q) \iff (\text{non } P) \text{ et } (\text{non } Q)} \qquad \text{et} \qquad \boxed{\text{non}(P \text{ et } Q) \iff (\text{non } P) \text{ ou } (\text{non } Q)}.
    \]
    \item \textbf{Implication ($P \implies Q$) et sa contraposée ($\text{non } Q \implies \text{non } P$) :}
    L'implication $P \implies Q$ (« si $P$, alors $Q$ ») a toujours exactement la même valeur de vérité que sa \textbf{contraposée} :
    \[
        \boxed{(P \implies Q) \iff (\text{non } Q \implies \text{non } P)}.
    \]
    \emph{Exemple fondamental en analyse :} L'implication « $f \text{ dérivable en } a \implies f \text{ continue en } a$ » a pour contraposée « $f \text{ non continue en } a \implies f \text{ non dérivable en } a$ » (toutes deux vraies). En revanche, la \textbf{réciproque} « $f \text{ continue en } a \implies f \text{ dérivable en } a$ » est \textbf{fausse} (contre-exemple : $f(x) = |x|$ en $a = 0$).
    \item \textbf{Équivalence ($P \iff Q$, « si et seulement si ») :} Vraie lorsque $P \implies Q$ \textbf{et} $Q \implies P$ sont toutes deux vraies.
\end{enumerate}
\end{theoreme}

\begin{attention}[Trois fausses équivalences à éviter absolument en Analyse]
\begin{enumerate}
    \item \textbf{Condition d'extremum :} « $x^*$ est un extremum local intérieur de $f \implies f'(x^*) = 0$ » est une \textbf{implication simple}, pas une équivalence ! La réciproque est fausse : on peut avoir $f'(x^*) = 0$ sans que $x^*$ soit un extremum (si $f'$ ne change pas de signe en $x^*$, comme $x^3$ en $0$ ou la fonction de Septembre 2026 Q1 en $x = 1$, qui est un point d'inflexion à tangente horizontale).
    \item \textbf{Élévation au carré :} $\sqrt{1+x} = x \implies 1 + x = x^2$ n'est une équivalence que si l'on ajoute la condition de signe $x \ge 0$ (Juillet 2024 Q1a). De même, $(f(x+\pi))^2 = (f(x))^2 \iff f(x+\pi) = \pm f(x)$ n'implique pas $f(x+\pi) = f(x)$, contrairement au cube $(f(x+\pi))^3 = (f(x))^3 \iff f(x+\pi) = f(x)$ car $y \mapsto y^3$ est bijective sur $\R$ (Septembre 2024 Q5).
    \item \textbf{Théorème de Rolle :} Si les hypothèses du théorème de Rolle sont satisfaites sur $[a, b]$, alors $\exists c \in \,]a, b[$ tel que $f'(c) = 0$. Si l'une des hypothèses n'est pas satisfaite, le théorème de Rolle \textbf{ne s'applique pas} (ce qui ne prouve pas nécessairement que $f'$ ne s'annule pas, mais interdit d'invoquer Rolle pour le démontrer, comme en Juillet 2023 Série 1 Q5).
\end{enumerate}
\end{attention}

% ----------------------------------------------------------------------------
\section*{A.2 Quantificateurs ($\forall, \exists$), négation d'une définition et contre-exemples}
% ----------------------------------------------------------------------------

\begin{theoreme}[Règles de négation des quantificateurs $\forall$ et $\exists$]
Pour nier une proposition contenant un ou plusieurs quantificateurs, on remplace chaque $\forall$ par $\exists$, chaque $\exists$ par $\forall$, et on nie la propriété finale :
\[
    \boxed{\text{non}\bigl(\forall x \in E,\; P(x)\bigr) \iff \exists x \in E \text{ tel que } \text{non } P(x)},
\]
\[
    \boxed{\text{non}\bigl(\exists x \in E \text{ tel que } P(x)\bigr) \iff \forall x \in E,\; \text{non } P(x)}.
\]
\end{theoreme}

\begin{itemize}
    \item \textbf{Application 1 : Réfuter une propriété universelle par un contre-exemple (Septembre 2024 Q5b et Juillet 2025 Q4) :}
    \begin{itemize}
        \item La définition de « $f$ est croissante sur $I$ » est : $\forall x_1, x_2 \in I,\; (x_1 < x_2 \implies f(x_1) \le f(x_2))$.
        \item Sa négation exacte est : $\mathbf{\exists x_1, x_2 \in I \text{ tels que } x_1 < x_2 \text{ et } f(x_1) > f(x_2)}$. Pour prouver qu'une fonction n'est pas croissante sur $I$, il suffit donc de donner \textbf{un seul couple numérique précis} $(x_1, x_2)$ vérifiant $x_1 < x_2$ et $f(x_1) > f(x_2)$ (par ex. $x_1 = 0,09 < x_2 = 0,1$ en Juillet 2025 Q4).
    \end{itemize}
    \item \textbf{Application 2 : Nier une définition de limite $\varepsilon$-$M$ (Juillet 2021 Q3d-ii) :}
    \begin{itemize}
        \item La définition de « $\lim_{x \to +\infty} g(x) = 0$ » est :
        \[
            \forall \varepsilon > 0, \;\; \exists M \in \R \text{ tel que } \forall x > M, \;\; |g(x)| \le \varepsilon.
        \]
        \item En inversant chaque quantificateur de gauche à droite, la négation « $\text{non}\bigl(\lim_{x \to +\infty} g(x) = 0\bigr)$ » s'écrit :
        \[
            \boxed{\exists \varepsilon > 0 \text{ tel que } \forall M \in \R, \;\; \exists x > M \text{ vérifiant } |g(x)| > \varepsilon}.
        \]
        Pour démontrer que $g(x)$ ne tend pas vers $0$ en $+\infty$, on choisit donc une valeur fixe de $\varepsilon > 0$ (par ex. $\varepsilon = 1$) et on exhibe une suite de points $x_k \to +\infty$ (qui dépassent donc n'importe quel $M$) en lesquels $|g(x_k)| > 1$.
    \end{itemize}
\end{itemize}

% ============================================================================
% CHAPITRE 9 : CORRIGÉS DÉTAILLÉS DES EXERCICES DE FIN DE CHAPITRE
% ============================================================================
\chapter{Corrigés Détaillés des Exercices}

Ce chapitre rassemble les solutions complètes, rigoureuses et commentées de tous les exercices proposés à la fin des Chapitres 1 à 7.

% ----------------------------------------------------------------------------
\section{Corrigés du Chapitre 1 (Fonctions Réelles, Domaines et Fonctions Transcendantes)}
% ----------------------------------------------------------------------------

\paragraph{Corrigé de l'Exercice 1.1 (Domaines de définition d'annales UCLouvain) :}
\begin{enumerate}[label=(\alph*)]
    \item \textbf{Pour $f_1(x) = \dfrac{x}{\cos^2(x)}$ (Septembre 2023, Q1a) :} La seule condition d'existence est que le dénominateur ne s'annule pas : $\cos^2(x) \neq 0 \iff \cos(x) \neq 0 \iff x \neq \frac{\pi}{2} + k\pi$ ($k \in \Z$). D'où :
    \[
        \boxed{\D_1 = \R \setminus \left\{ \frac{\pi}{2} + k\pi \;\middle|\; k \in \Z \right\}}.
    \]
    \item \textbf{Pour $f_2(x) = \dfrac{1-x^2}{\cos^2(x)}$ (Septembre 2024, Q1a) :} Exactement la même condition d'existence $\cos^2(x) \neq 0 \iff \cos(x) \neq 0$, soit :
    \[
        \boxed{\D_2 = \R \setminus \left\{ \frac{\pi}{2} + k\pi \;\middle|\; k \in \Z \right\}}.
    \]
    \item \textbf{Pour $f_3(x) = x\ln^2(x)$ (Septembre 2025, Q1a) :} L'argument du logarithme népérien doit être strictement positif ($x > 0$), soit :
    \[
        \boxed{\D_3 = \,]0, +\infty[ = \R_0^+}.
    \]
    \emph{Remarque :} Comme $\lim_{x \to 0^+} x\ln^2(x) = 0$, on peut prolonger $f_3$ par continuité en $x = 0$ en posant $\tilde{f}_3(0) = 0$.
    \item \textbf{Pour $f_4(x) = (1-x^2)\ln(1-x) + x$ (Juillet 2026, Q1a) :} L'argument du logarithme doit vérifier $1 - x > 0 \iff x < 1$, soit :
    \[
        \boxed{\D_4 = \,]-\infty, 1[}.
    \]
\end{enumerate}

\paragraph{Corrigé de l'Exercice 1.2 (Racines et équations logarithmiques / irrationnelles) :}
\begin{enumerate}[label=(\alph*)]
    \item \textbf{Pour $f(x) = \ln(x^2 - 3x + 2)$ :}
    \begin{itemize}
        \item \textbf{CE :} $x^2 - 3x + 2 > 0 \iff (x-1)(x-2) > 0$. Le trinôme (à coefficient dominant $1 > 0$) est strictement positif à l'extérieur de ses racines $1$ et $2$, donc $\boxed{\D = \,]-\infty, 1[ \;\cup\; ]2, +\infty[}$.
        \item \textbf{Racines :} $f(x) = 0 \iff \ln(x^2 - 3x + 2) = 0 \iff x^2 - 3x + 2 = 1 \iff x^2 - 3x + 1 = 0$. Le discriminant vaut $\Delta = 9 - 4 = 5 > 0$, ce qui donne deux racines réelles qui appartiennent toutes deux à $\D$ (car en ces points $x^2 - 3x + 2 = 1 > 0$) :
        \[
            \boxed{x_{1,2} = \frac{3 \pm \sqrt{5}}{2}}.
        \]
    \end{itemize}
    \item \textbf{Pour $g(x) = \sqrt{2x+3} - x$ :}
    \begin{itemize}
        \item \textbf{CE :} $2x + 3 \ge 0 \iff x \ge -\frac{3}{2}$, d'où $\boxed{\D = \left[-\frac{3}{2}, +\infty\right[}$.
        \item \textbf{Racines :} $g(x) = 0 \iff \sqrt{2x+3} = x \iff \begin{cases} x \ge 0 \\ 2x + 3 = x^2 \end{cases} \iff \begin{cases} x \ge 0 \\ x^2 - 2x - 3 = (x-3)(x+1) = 0 \end{cases}$.
        La valeur $x = -1 < 0$ est une racine parasite (car $\sqrt{1} = 1 \neq -1$). L'unique racine est donc $\boxed{x = 3}$.
    \end{itemize}
    \item \textbf{Pour $h(x) = \dfrac{x^2}{2}\left(\ln^2(x) - \ln(x) + \dfrac{1}{2}\right)$ (Septembre 2026, Q1) :}
    \begin{itemize}
        \item \textbf{CE :} $x > 0 \implies \boxed{\D = \,]0, +\infty[}$.
        \item \textbf{Racines :} Sur $]0, +\infty[$, $\frac{x^2}{2} > 0$. Posons $u = \ln(x) \in \R$. L'équation $u^2 - u + \frac{1}{2} = 0$ a pour discriminant $\Delta = (-1)^2 - 4(1)\left(\frac{1}{2}\right) = -1 < 0$ (ou par complétion du carré : $\left(u - \frac{1}{2}\right)^2 + \frac{1}{4} \ge \frac{1}{4} > 0$). La fonction $h$ ne possède donc \textbf{aucune racine réelle} sur son domaine $\D = \,]0, +\infty[$.
    \end{itemize}
\end{enumerate}

\paragraph{Corrigé de l'Exercice 1.3 (UCLouvain -- Septembre 2024, Q4a) :}
Par hypothèse, $x_0$ est racine de $f(x) = \sqrt{1+x} - x$, ce qui signifie que $f(x_0) = 0 \iff \mathbf{\sqrt{1 + x_0} = x_0}$.
Évaluons alors la fonction $g(x) = \sqrt{1 + \sqrt{1+x}} - x$ en $x = x_0$ en remplaçant successivement $\sqrt{1+x_0}$ par $x_0$ :
\[
    g(x_0) = \sqrt{1 + \underbrace{\sqrt{1 + x_0}}_{= x_0}} - x_0 = \underbrace{\sqrt{1 + x_0}}_{= x_0} - x_0 = x_0 - x_0 = 0.
\]
Par conséquent, $x_0$ est bien racine de la fonction $g$. \hfill $\square$

\paragraph{Corrigé de l'Exercice 1.4 (Parité et symétries) :}
\begin{enumerate}[label=(\alph*)]
    \item $\D_1 = \R$ (symétrique). Pour tout $x \in \R$, $(-x)^2 = x^2$, donc $f_1(-x) = (-x)^2\sin\bigl((-x)^2\bigr) + \cos\bigl((-x)^2\bigr) = x^2\sin(x^2) + \cos(x^2) = f_1(x)$ : $f_1$ est \textbf{paire}.
    \item $\D_2 = \R$ (symétrique). Comme $\cos(-x) = \cos(x)$, on a $f_2(-x) = 3\cos(-x) - \bigl(\cos(-x)\bigr)^3 = 3\cos(x) - \cos^3(x) = f_2(x)$ : $f_2$ est \textbf{paire}.
    \item \textbf{CE :} $\frac{1+x}{1-x} > 0 \iff x \in \D_3 = \,]-1, 1[$ (intervalle symétrique par rapport à $0$). Pour tout $x \in \,]-1, 1[$ :
    \[
        f_3(-x) = \ln\left(\frac{1-x}{1+x}\right) = \ln\left(\left(\frac{1+x}{1-x}\right)^{-1}\right) = -\ln\left(\frac{1+x}{1-x}\right) = -f_3(x) \implies f_3 \text{ est \textbf{impaire}}.
    \]
    \item Comme $\sqrt{1+x^2} > \sqrt{x^2} = |x| \ge -x$, on a $x + \sqrt{1+x^2} > 0$ pour tout $x \in \R$, donc $\D_4 = \R$. En multipliant par le binôme conjugué :
    \[
        f_4(-x) = \ln\bigl(\sqrt{1+x^2} - x\bigr) = \ln\left(\frac{(\sqrt{1+x^2}-x)(\sqrt{1+x^2}+x)}{\sqrt{1+x^2}+x}\right) = \ln\left(\frac{1}{\sqrt{1+x^2}+x}\right) = -\ln\bigl(x + \sqrt{1+x^2}\bigr) = -f_4(x).
    \]
    La fonction $f_4$ (qui est l'argument sinus hyperbolique $\operatorname{argsinh}(x)$) est donc \textbf{impaire} sur $\R$.
\end{enumerate}

\paragraph{Corrigé de l'Exercice 1.5 (Simplifications algébriques avec $\ln$ et $\exp$) :}
\begin{enumerate}[label=(\alph*)]
    \item \textbf{CE :} $e^{2x} - 1 > 0 \iff e^{2x} > 1 \iff x > 0$ (ce qui assure aussi $1 - e^{-x} > 0$). Pour tout $x > 0$, factorisons $e^{2x} - 1 = (e^x - 1)(e^x + 1)$ et écrivons $1 - e^{-x} = 1 - \frac{1}{e^x} = \frac{e^x - 1}{e^x}$ :
    \[
        A(x) = \ln\left(\frac{(e^x-1)\cancel{(e^x+1)}}{\cancel{e^x+1}}\right) - \ln\left(\frac{e^x - 1}{e^x}\right) = \cancel{\ln(e^x - 1)} - \Bigl[\cancel{\ln(e^x - 1)} - \ln(e^x)\Bigr] = \ln(e^x) = \boxed{x}.
    \]
    \item En posant $u = e^x > 0$, l'équation devient $u^2 - 3u - 4 = (u - 4)(u + 1) = 0$. Comme $u = e^x > 0$, la racine négative $u = -1$ est impossible, et il reste $e^x = 4 \iff \boxed{x = \ln(4) = 2\ln(2)}$.
    \item \textbf{CE :} $x + 3 > 0$ et $x - 1 > 0 \iff \mathbf{x > 1}$. Sous cette condition, $\ln(x+3) + \ln(x-1) = \ln(5) \iff \ln\bigl((x+3)(x-1)\bigr) = \ln(5) \iff x^2 + 2x - 3 = 5 \iff x^2 + 2x - 8 = (x+4)(x-2) = 0$. La valeur $x = -4$ est rejetée par la CE ($x > 1$), d'où l'unique solution $\boxed{x = 2}$.
\end{enumerate}

\paragraph{Corrigé de l'Exercice 1.6 (UCLouvain -- Septembre 2026, Q4 -- Fonction de Lambert $W_0$) :}
\begin{enumerate}[label=(\alph*)]
    \item Calculons les images par $f(x) = x\,e^x$ des entiers simples $x \in \{0, 1, -1\} \subset I = [-1, +\infty[$ :
    \begin{itemize}
        \item $f(0) = 0 \cdot e^0 = 0 \implies \boxed{W_0(0) = 0}$.
        \item $f(1) = 1 \cdot e^1 = e \implies \boxed{W_0(e) = 1}$.
        \item $f(-1) = (-1)\cdot e^{-1} = -e^{-1} \implies \boxed{W_0(-e^{-1}) = -1}$.
    \end{itemize}
    \item La fonction $f(x) = x\,e^x$ est continue et strictement croissante sur $I = [-1, +\infty[$ (car $f'(x) = (1+x)e^x > 0$ sur $]-1, +\infty[$). Son ensemble image est $f(I) = \bigl[f(-1),\; \lim_{x \to +\infty} f(x)\bigr[ = [-e^{-1}, +\infty[$. Par les propriétés générales des fonctions réciproques :
    \begin{itemize}
        \item Le domaine de définition de $W_0$ est $\boxed{D = [-e^{-1}, +\infty[}$.
        \item $W_0$ est \textbf{strictement croissante} sur $D$ (comme $f$ sur $I$).
        \item $W_0$ est \textbf{continue} sur $D$ (comme réciproque d'une fonction continue strictement monotone sur un intervalle).
    \end{itemize}
\end{enumerate}

% ----------------------------------------------------------------------------
\section{Corrigés du Chapitre 2 (Limites, Continuité et Asymptotes)}
% ----------------------------------------------------------------------------

\paragraph{Corrigé de l'Exercice 2.1 (Calculs de limites sans L'Hospital) :}
\begin{enumerate}[label=(\alph*)]
    \item Forme $\left[\frac{0}{0}\right]$ en $x = 2$. Par les identités remarquables $x^3 - 8 = (x-2)(x^2 + 2x + 4)$ et $x^2 - 3x + 2 = (x-2)(x-1)$ :
    \[
        \lim_{x \to 2} \frac{\cancel{(x-2)}(x^2 + 2x + 4)}{\cancel{(x-2)}(x - 1)} = \frac{4 + 4 + 4}{2 - 1} = \boxed{12}.
    \]
    \item Forme $\left[\frac{0}{0}\right]$ en $x = 0$. Multiplions numérateur et dénominateur par le binôme conjugué $\sqrt{1+2x} + \sqrt{1-2x}$ :
    \[
        \lim_{x \to 0} \frac{(1 + 2x) - (1 - 2x)}{x\bigl(\sqrt{1+2x} + \sqrt{1-2x}\bigr)} = \lim_{x \to 0} \frac{4\cancel{x}}{\cancel{x}\bigl(\sqrt{1+2x} + \sqrt{1-2x}\bigr)} = \frac{4}{1 + 1} = \boxed{2}.
    \]
    \item Forme $[\infty - \infty]$ en $+\infty$. Par multiplication par le binôme conjugué $\sqrt{x^2+4x+1} + x$ et mise en évidence de $x = \sqrt{x^2}$ (car $x > 0$) :
    \[
        \lim_{x \to +\infty} \frac{(x^2 + 4x + 1) - x^2}{\sqrt{x^2+4x+1} + x} = \lim_{x \to +\infty} \frac{\cancel{x}\left(4 + \frac{1}{x}\right)}{\cancel{x}\left(\sqrt{1 + \frac{4}{x} + \frac{1}{x^2}} + 1\right)} = \frac{4 + 0}{\sqrt{1} + 1} = \boxed{2}.
    \]
    \item Comme $\ln(x^3) = 3\ln(x)$ pour $x > 0$, mettons le terme dominant $\ln(x)$ en évidence :
    \[
        \lim_{x \to +\infty} \frac{9\ln(x) - 2\cos(x^2)}{2\ln(x) + 5\sin(x)} = \lim_{x \to +\infty} \frac{9 - \dfrac{2\cos(x^2)}{\ln(x)}}{2 + \dfrac{5\sin(x)}{\ln(x)}} = \frac{9 - 0}{2 + 0} = \boxed{\frac{9}{2}},
    \]
    puisque $\cos(x^2)$ et $\sin(x)$ sont bornés entre $-1$ et $1$ tandis que $\ln(x) \to +\infty$.
\end{enumerate}

\paragraph{Corrigé de l'Exercice 2.2 (Limites trigonométriques et exponentielles -- Annales UCLouvain) :}
\begin{enumerate}[label=(\alph*)]
    \item \textbf{(Juillet 2021, Q3e)} En posant $y = \sin(x) \to 0$ lorsque $x \to 0$, décomposons le quotient en deux limites fondamentales :
    \[
        \lim_{x \to 0} \frac{x\bigl(e^{\sin(x)} - 1\bigr)}{\sin^2(x)} = \left(\lim_{x \to 0} \frac{x}{\sin(x)}\right) \cdot \left(\lim_{y \to 0} \frac{e^y - 1}{y}\right) = 1 \cdot 1 = \boxed{1}.
    \]
    \item \textbf{(Juillet 2025, Q5)} De même, en multipliant et divisant par $x$ :
    \[
        \lim_{x \to 0} \frac{\sin(x^2)}{x\tan(x)} = \left(\lim_{x \to 0} \frac{\sin(x^2)}{x^2}\right) \cdot \left(\lim_{x \to 0} \frac{x}{\tan(x)}\right) = 1 \cdot 1 = \boxed{1}.
    \]
    \item Forme $[0 \cdot (-\infty)]$. Écrivons $x\ln(\sin x) = x\ln\left(x\cdot\frac{\sin x}{x}\right) = x\ln(x) + x\ln\left(\frac{\sin x}{x}\right)$. Or $\lim_{x \to 0^+} x\ln(x) = 0$ (croissance comparée) et $\lim_{x \to 0^+} \frac{\sin x}{x} = 1 \implies \lim_{x \to 0^+} x\ln\left(\frac{\sin x}{x}\right) = 0 \cdot \ln(1) = 0$, d'où la limite $\boxed{0}$.
\end{enumerate}

\paragraph{Corrigé de l'Exercice 2.3 (Définition formelle de la limite $\varepsilon$-$M$) :}
\begin{enumerate}[label=(\alph*)]
    \item Soit $\varepsilon > 0$ quelconque. Pour tout $x > 0$, on a :
    \[
        \left|\frac{2x + \cos(x)}{x} - 2\right| = \left|\frac{\cos(x)}{x}\right| = \frac{|\cos(x)|}{x} \le \frac{1}{x}.
    \]
    En choisissant $\boxed{M = \dfrac{1}{\varepsilon}}$, on a pour tout $x > M$ : $|f(x) - 2| \le \frac{1}{x} < \varepsilon$, ce qui démontre que la limite vaut $2$.
    \item Soit $\varepsilon > 0$ quelconque. Pour tout $x > 1$ (de sorte que $\ln(x) > 0$), on a :
    \[
        \left|\frac{\ln(x) + 3}{\ln(x)} - 1\right| = \frac{3}{\ln(x)}.
    \]
    Cette quantité est inférieure ou égale à $\varepsilon$ dès que $\ln(x) \ge \frac{3}{\varepsilon} \iff x \ge e^{3/\varepsilon}$. En posant $\boxed{M = e^{3/\varepsilon}}$, l'implication $x > M \implies |f(x) - 1| \le \varepsilon$ est vérifiée.
\end{enumerate}

\paragraph{Corrigé de l'Exercice 2.4 (Continuité et prolongement par continuité) :}
Calculons la limite à gauche et la limite à droite de $f$ en $x = 0$ :
\[
    \lim_{x \to 0^-} f(x) = \lim_{x \to 0^-} \left(3\cdot\frac{\sin(3x)}{3x}\right) = 3 \cdot 1 = 3, \qquad \lim_{x \to 0^+} f(x) = \lim_{x \to 0^+} \bigl(\sqrt{1 + bx} + 2\bigr) = \sqrt{1 + 0} + 2 = 3.
\]
Les limites à gauche et à droite sont égales à $3$ quelle que soit la valeur de $b \in \R$. Pour que $f$ soit continue en $0$, il faut et il suffit que $f(0) = a$ soit égal à cette limite commune : $\boxed{a = 3 \text{ et } b \in \R \text{ quelconque}}$.

\paragraph{Corrigé de l'Exercice 2.5 (Recherche complète d'asymptotes) :}
\begin{enumerate}[label=(\alph*)]
    \item \textbf{Pour $f_1(x) = \dfrac{2x^2 - 3x + 5}{x - 2}$ :} $\D_1 = \R \setminus \{2\}$.
    \begin{itemize}
        \item En $x = 2$, le numérateur vaut $2(4) - 6 + 5 = 7 \neq 0$, donc $\lim_{x \to 2^\pm} f_1(x) = \pm\infty \implies \boxed{\text{AV} \equiv x = 2}$.
        \item En $x \to \pm\infty$ : par division euclidienne de $2x^2 - 3x + 5$ par $x - 2$, on a $2x^2 - 3x + 5 = (2x + 1)(x - 2) + 7$, soit $f_1(x) = 2x + 1 + \dfrac{7}{x - 2}$. Comme $\lim_{x \to \pm\infty} \frac{7}{x-2} = 0$, il n'y a pas d'AH et on obtient directement l'asymptote oblique en $\pm\infty$ : $\boxed{\text{AO} \equiv y = 2x + 1}$.
    \end{itemize}
    \item \textbf{Pour $f_2(x) = \sqrt{x^2 + 2x}$ :} $x(x+2) \ge 0 \iff \D_2 = \,]-\infty, -2] \cup [0, +\infty[$. Pas d'AV (la fonction est continue et nulle aux bornes fermées $-2$ et $0$). En $\pm\infty$, complétons le carré : $x^2 + 2x = (x+1)^2 - 1$.
    \begin{itemize}
        \item En $+\infty$ : $m = \lim_{x \to +\infty} \frac{\sqrt{x^2+2x}}{x} = 1$, et $p = \lim_{x \to +\infty} \bigl(\sqrt{x^2+2x} - x\bigr) = \lim_{x \to +\infty} \frac{2x}{\sqrt{x^2+2x}+x} = 1 \implies \boxed{\text{AO}_{+\infty} \equiv y = x + 1}$.
        \item En $-\infty$ (attention : pour $x < 0$, $\sqrt{x^2} = |x| = -x$ !) : $m = \lim_{x \to -\infty} \frac{|x|\sqrt{1+2/x}}{x} = -1$, et $p = \lim_{x \to -\infty} \bigl(\sqrt{x^2+2x} + x\bigr) = \lim_{x \to -\infty} \frac{2x}{\sqrt{x^2+2x}-x} = \frac{2}{-1 - 1} = -1 \implies \boxed{\text{AO}_{-\infty} \equiv y = -x - 1}$.
    \end{itemize}
    \item \textbf{Pour $f_3(x) = \dfrac{e^x - 1}{e^x + 1}$ :} $\D_3 = \R$ (car $e^x + 1 > 1 > 0$). Pas d'AV.
    \begin{itemize}
        \item En $-\infty$ : $\lim_{x \to -\infty} e^x = 0 \implies \lim_{x \to -\infty} f_3(x) = \frac{0 - 1}{0 + 1} = -1 \implies \boxed{\text{AH}_{-\infty} \equiv y = -1}$.
        \item En $+\infty$ : en factorisant $e^x$, $\lim_{x \to +\infty} \frac{1 - e^{-x}}{1 + e^{-x}} = \frac{1 - 0}{1 + 0} = 1 \implies \boxed{\text{AH}_{+\infty} \equiv y = 1}$.
    \end{itemize}
\end{enumerate}

\paragraph{Corrigé de l'Exercice 2.6 (Application du théorème des valeurs intermédiaires) :}
Posons $h(x) = \cos(x) - x$, définie et continue sur $\R$.
\begin{itemize}
    \item \textbf{Existence dans $\left]0, \frac{\pi}{2}\right[$ :} Sur $\left[0, \frac{\pi}{2}\right]$, $h$ est continue, $h(0) = \cos(0) - 0 = 1 > 0$ et $h\left(\frac{\pi}{2}\right) = \cos\left(\frac{\pi}{2}\right) - \frac{\pi}{2} = -\frac{\pi}{2} < 0$. Comme $h(0)\cdot h\left(\frac{\pi}{2}\right) < 0$, le théorème des valeurs intermédiaires (Bolzano) garantit l'existence d'au moins une racine $x^* \in \left]0, \frac{\pi}{2}\right[$ telle que $h(x^*) = 0 \iff \cos(x^*) = x^*$.
    \item \textbf{Unicité sur tout $\R$ :} La fonction $h$ est dérivable sur $\R$ et $h'(x) = -\sin(x) - 1 \le 0$ (car $\sin(x) \ge -1$), avec $h'(x) < 0$ sauf en des points isolés $x = -\frac{\pi}{2} + 2k\pi$. La fonction $h$ est donc \textbf{strictement décroissante sur tout $\R$} : elle ne peut s'annuler qu'en un unique point $x^* \in \left]0, \frac{\pi}{2}\right[$.
\end{itemize}

% ----------------------------------------------------------------------------
\section{Corrigés du Chapitre 3 (Dérivées et Grands Théorèmes d'Analyse)}
% ----------------------------------------------------------------------------

\paragraph{Corrigé de l'Exercice 3.1 (Calculs de dérivées d'annales UCLouvain) :}
\begin{enumerate}[label=(\alph*)]
    \item \textbf{Pour $f_1(x) = x^2\sin(x^2) + \cos(x^2)$ (Juillet 2023 -- Série 1, Q1a-b) :}
    \[
        f'_1(x) = \cancel{2x\sin(x^2)} + x^2\cdot 2x\cos(x^2) - \cancel{2x\sin(x^2)} = \boxed{2x^3\cos(x^2)},
    \]
    \[
        f''_1(x) = 6x^2\cos(x^2) + 2x^3\bigl(-2x\sin(x^2)\bigr) = \boxed{2x^2\Bigl(3\cos(x^2) - 2x^2\sin(x^2)\Bigr)}.
    \]
    \item \textbf{Pour $f_2(x) = 3\cos(x) - \cos^3(x)$ (Juillet 2023 -- Série 2, Q1a-b) :}
    \[
        f'_2(x) = -3\sin(x) - 3\cos^2(x)\bigl(-\sin(x)\bigr) = -3\sin(x)\bigl(1 - \cos^2(x)\bigr) = \boxed{-3\sin^3(x)},
    \]
    \[
        f''_2(x) = \boxed{-9\sin^2(x)\cos(x)}.
    \]
    \item \textbf{Pour $f_3(x) = \dfrac{1 - x^2}{\cos^2(x)}$ (Septembre 2024, Q1b) :} Par la règle du quotient :
    \[
        f'_3(x) = \frac{-2x\cos^2(x) - (1 - x^2)\cdot 2\cos(x)\bigl(-\sin(x)\bigr)}{\cos^4(x)} = \frac{-2x\cos(x) + 2(1-x^2)\sin(x)}{\cos^3(x)} = \frac{g(x)}{\cos^2(x)},
    \]
    avec $\boxed{g(x) = 2\Bigl(\tan(x)\bigl(1 - x^2\bigr) - x\Bigr)}$.
    \item \textbf{Pour $f_4(x) = \dfrac{x^2}{2}\left(\ln^2(x) - \ln(x) + \dfrac{1}{2}\right)$ (Septembre 2026, Q1b-c) :} Vu à l'Étude 5 du Chapitre 4 :
    \[
        \boxed{f'_4(x) = x\ln^2(x)} \qquad \text{et} \qquad \boxed{f''_4(x) = \ln^2(x) + 2\ln(x) = \ln(x)\bigl(\ln(x) + 2\bigr)}.
    \]
\end{enumerate}

\paragraph{Corrigé de l'Exercice 3.2 (Dérivées de fonctions réciproques et composées) :}
\begin{enumerate}[label=(\alph*)]
    \item Pour $g_1(x) = \arcsin(2x-1)$ (dérivable pour $-1 < 2x-1 < 1 \iff 0 < x < 1$) :
    \[
        g'_1(x) = \frac{(2x-1)'}{\sqrt{1 - (2x-1)^2}} = \frac{2}{\sqrt{1 - (4x^2 - 4x + 1)}} = \frac{2}{\sqrt{4x - 4x^2}} = \boxed{\frac{1}{\sqrt{x(1 - x)}}}.
    \]
    \item Pour $g_2(x) = \arctan\left(\dfrac{1+x}{1-x}\right)$ ($x \neq 1$), posons $u(x) = \frac{1+x}{1-x} \implies u'(x) = \frac{1(1-x) - (1+x)(-1)}{(1-x)^2} = \frac{2}{(1-x)^2}$ :
    \[
        g'_2(x) = \frac{u'(x)}{1 + u^2(x)} = \frac{\dfrac{2}{(1-x)^2}}{1 + \dfrac{(1+x)^2}{(1-x)^2}} = \frac{2}{(1-x)^2 + (1+x)^2} = \frac{2}{2 + 2x^2} = \boxed{\frac{1}{1 + x^2}}.
    \]
    \emph{Interprétation :} On retrouve exactement la dérivée de $\arctan(x)$ ! En effet, par la formule trigonométrique $\tan\left(\alpha + \frac{\pi}{4}\right) = \frac{1 + \tan\alpha}{1 - \tan\alpha}$, la fonction $g_2(x)$ ne diffère de $\arctan(x)$ que d'une constante ($\frac{\pi}{4}$ sur $]-\infty, 1[$ et $-\frac{3\pi}{4}$ sur $]1, +\infty[$).
    \item Pour $g_3(x) = x^x = e^{x\ln(x)}$ ($x > 0$), la règle de dérivation de $e^{u(x)}$ avec $u(x) = x\ln(x) \implies u'(x) = \ln(x) + 1$ donne :
    \[
        \boxed{g'_3(x) = \bigl(\ln(x) + 1\bigr)\,e^{x\ln(x)} = \bigl(\ln(x) + 1\bigr)\,x^x}.
    \]
\end{enumerate}

\paragraph{Corrigé de l'Exercice 3.3 (Limites par la règle de L'Hospital) :}
\begin{enumerate}[label=(\alph*)]
    \item Trois applications successives de la règle de L'Hospital (formes $\left[\frac{0}{0}\right]$ à chaque étape) donnent :
    \[
        \lim_{x \to 0} \frac{e^x - e^{-x} - 2x}{x - \sin(x)} \stackrel{\text{H}}{=} \lim_{x \to 0} \frac{e^x + e^{-x} - 2}{1 - \cos(x)} \quad \left[\frac{0}{0}\right] \stackrel{\text{H}}{=} \lim_{x \to 0} \frac{e^x - e^{-x}}{\sin(x)} \quad \left[\frac{0}{0}\right] \stackrel{\text{H}}{=} \lim_{x \to 0} \frac{e^x + e^{-x}}{\cos(x)} = \frac{1 + 1}{1} = \boxed{2}.
    \]
    \item Forme $\left[\frac{0}{0}\right]$ en $x = \frac{\pi}{2}$. Deux applications de la règle de L'Hospital donnent :
    \[
        \lim_{x \to \pi/2} \frac{\ln(\sin x)}{(\pi - 2x)^2} \stackrel{\text{H}}{=} \lim_{x \to \pi/2} \frac{\dfrac{\cos x}{\sin x}}{-4(\pi - 2x)} = \lim_{x \to \pi/2} \left(\frac{1}{\sin x} \cdot \frac{\cos x}{-4(\pi - 2x)}\right) \stackrel{\text{H}}{=} 1 \cdot \lim_{x \to \pi/2} \frac{-\sin x}{8} = \boxed{-\frac{1}{8}}.
    \]
    \item Forme $[0^0]$. Écrivons $x^{\sin(x)} = e^{\sin(x)\ln(x)}$ et calculons la limite de l'exposant par L'Hospital :
    \[
        \lim_{x \to 0^+} \sin(x)\ln(x) = \lim_{x \to 0^+} \left(\frac{\sin(x)}{x} \cdot \frac{\ln(x)}{1/x}\right) \stackrel{\text{H}}{=} 1 \cdot \lim_{x \to 0^+} \frac{1/x}{-1/x^2} = \lim_{x \to 0^+} (-x) = 0 \implies \lim_{x \to 0^+} x^{\sin(x)} = e^0 = \boxed{1}.
    \]
    \item Forme $[1^\infty]$. Écrivons $\left(1 + \frac{y}{x}\right)^x = e^{x\ln(1 + y/x)}$. En posant $u = \frac{1}{x} \to 0^+$ (ou par L'Hospital sur $\frac{\ln(1+y/x)}{1/x}$) :
    \[
        \lim_{x \to +\infty} x\ln\left(1 + \frac{y}{x}\right) = \lim_{u \to 0^+} \frac{\ln(1 + yu)}{u} \stackrel{\text{H}}{=} \lim_{u \to 0^+} \frac{y}{1 + yu} = y \implies \boxed{\lim_{x \to +\infty} \left(1 + \frac{y}{x}\right)^x = e^y}.
    \]
\end{enumerate}

\paragraph{Corrigé de l'Exercice 3.4 (UCLouvain -- Juillet 2023 -- Série 1, Q5 -- Applicabilité de Rolle) :}
\begin{enumerate}[label=(\alph*)]
    \item $f(x) = 1 - \sqrt{|x|}$ : \textbf{NON} -- Bien que $f$ soit continue sur $[-1, 1]$ et que $f(-1) = f(1) = 0$, $f$ \textbf{n'est pas dérivable en $x = 0 \in \,]-1, 1[$} (point de rebroussement à tangente verticale en $0$).
    \item $g(x) = x - \dfrac{1}{x}$ : \textbf{NON} -- $g$ \textbf{n'est pas définie (ni continue) en $x = 0 \in [-1, 1]$} (et de plus $g(-1) = 0 \neq g(1)$).
    \item $h(x) = \ln(1 + x^2)$ : \textbf{OUI} -- $h$ est définie et continue sur $[-1, 1]$, dérivable sur $]-1, 1[$, et $h(-1) = \ln(2) = h(1)$.
    \item $w(x) = x - x^5$ : \textbf{OUI} -- $w$ est un polynôme donc continue sur $[-1, 1]$ et dérivable sur $]-1, 1[$, et $w(-1) = -1 - (-1) = 0 = 1 - 1 = w(1)$.
\end{enumerate}

\paragraph{Corrigé de l'Exercice 3.5 (Application de Rolle et du TVI au polynôme $P(x) = x^4 - 4x + 1$) :}
\begin{enumerate}[label=(\alph*)]
    \item Supposons par l'absurde que $P(x)$ possède au moins $3$ racines réelles distinctes $r_1 < r_2 < r_3$. Comme $P$ est un polynôme (continu et dérivable partout), le théorème de Rolle appliqué sur $[r_1, r_2]$ et sur $[r_2, r_3]$ impliquerait l'existence de deux racines distinctes $c_1 \in \,]r_1, r_2[$ et $c_2 \in \,]r_2, r_3[$ de la dérivée $P'(x)$. Or $P'(x) = 4x^3 - 4 = 4(x^3 - 1)$ ne s'annule qu'en \textbf{un seul point réel} ($x = 1$) ! Cette contradiction prouve que $P$ a au plus $2$ racines réelles.
    \item Comme $P(0) = 1 > 0$, $P(1) = 1 - 4 + 1 = -2 < 0$ et $P(2) = 16 - 8 + 1 = 9 > 0$, le théorème des valeurs intermédiaires garantit une première racine dans $]0, 1[$ et une seconde racine dans $]1, 2[$. Il y en a donc exactement deux.
\end{enumerate}

\paragraph{Corrigé de l'Exercice 3.6 (Inégalités par le théorème des accroissements finis) :}
\begin{enumerate}[label=(\alph*)]
    \item Si $a = b$, l'inégalité $0 \le 0$ est évidente. Si $a \neq b$ (supposons $a < b$), appliquons le TAF à la fonction $f(t) = \sin(t)$ sur $[a, b]$ : il existe $c \in \,]a, b[$ tel que $\sin(b) - \sin(a) = \cos(c)\,(b - a)$. En passant aux valeurs absolues et en utilisant $|\cos(c)| \le 1$, on obtient $\bigl|\sin(b) - \sin(a)\bigr| = |\cos(c)|\cdot|b - a| \le |b - a|$.
    \item Soit $x > 0$. Appliquons le TAF à $f(t) = \arctan(t)$ sur l'intervalle $[0, x]$ : il existe $c \in \,]0, x[$ tel que :
    \[
        \frac{\arctan(x) - \arctan(0)}{x - 0} = f'(c) = \frac{1}{1 + c^2}.
    \]
    Or, puisque $0 < c < x$, on a $1 < 1 + c^2 < 1 + x^2$, d'où en inversant : $\dfrac{1}{1 + x^2} < \dfrac{1}{1 + c^2} < 1$. En remplaçant $\frac{1}{1+c^2} = \frac{\arctan(x)}{x}$ et en multipliant par $x > 0$, on obtient $\boxed{\dfrac{x}{1 + x^2} < \arctan(x) < x}$.
\end{enumerate}

% ----------------------------------------------------------------------------
\section{Corrigés du Chapitre 4 (Étude Complète de Fonctions)}
% ----------------------------------------------------------------------------

\paragraph{Corrigé de l'Exercice 4.1 (UCLouvain -- Juillet 2021, Q1 -- $f(x) = \ln(1/\cos x)$) :}
\begin{enumerate}[label=(\alph*)]
    \item \textbf{Domaine :} Vu au Chapitre 1 : $\cos(x) > 0 \iff \boxed{\D = \bigcup_{k\in\Z} \left]-\frac{\pi}{2} + 2k\pi,\; \frac{\pi}{2} + 2k\pi\right[}$. La fonction est paire et périodique de période $2\pi$.
    \item \textbf{Dérivées :} $f$ est dérivable sur tout $\D$. En écrivant $f(x) = -\ln(\cos x)$ :
    \[
        \boxed{f'(x) = -\frac{-\sin(x)}{\cos(x)} = \tan(x)} \qquad \text{et} \qquad \boxed{f''(x) = \frac{1}{\cos^2(x)} > 0 \quad (\forall x \in \D)}.
    \]
    \item \textbf{Extrema et asymptotes :} $f'(x) = \tan(x) = 0$ sur $\D \iff \boxed{x = 2k\pi}$ ($k \in \Z$). Comme $f''(x) = \frac{1}{\cos^2(x)} \ge 1 > 0$ sur tout $\D$, la concavité est toujours tournée vers le haut ($\cup$) et les points $x = 2k\pi$ sont des \textbf{minima locaux} (de valeur $f(2k\pi) = \ln(1) = 0$).
    Étant périodique non constante, $f$ n'a ni AH ni AO. Aux bords de chaque intervalle de $\D$, $\cos(x) \to 0^+ \implies \frac{1}{\cos x} \to +\infty \implies f(x) \to +\infty$, d'où une infinité d'\textbf{asymptotes verticales} $\boxed{\text{AV} \equiv x = \frac{\pi}{2} + k\pi}$ ($k \in \Z$).
    \item \textbf{Tableau sur $[-\pi, \pi[$ et lien avec $\bar{f}(x) = \ln\bigl(1/\cos^2(x)\bigr)$ :}
    Sur $[-\pi, \pi[$, $f$ n'est définie que sur $\left]-\frac{\pi}{2}, \frac{\pi}{2}\right[$, décroît de $+\infty$ (AV en $-\frac{\pi}{2}$) à $0$ (minimum en $x = 0$), puis croît de $0$ à $+\infty$ (AV en $\frac{\pi}{2}$) avec une concavité positive ($\cup$). Sur $[-4\pi, 4\pi]$, cette branche en « U » se répète tous les $2\pi$ (autour de $-4\pi, -2\pi, 0, 2\pi, 4\pi$).
    Pour $\bar{f}(x) = \ln\bigl(1/\cos^2(x)\bigr)$, la condition d'existence est simplement $\cos^2(x) > 0 \iff \cos(x) \neq 0 \iff x \neq \frac{\pi}{2} + k\pi$ : $\bar{f}$ est donc définie sur \textbf{tous} les intervalles $\left]-\frac{\pi}{2}+k\pi, \frac{\pi}{2}+k\pi\right[$, est \textbf{périodique de période $\pi$} (au lieu de $2\pi$), et vaut $\bar{f}(x) = 2\ln\left(\frac{1}{|\cos x|}\right) = 2f(x)$ sur $\dom f$ (il y a deux fois plus de branches en « U », toutes étirées verticalement d'un facteur $2$).
\end{enumerate}

\paragraph{Corrigé de l'Exercice 4.2 (UCLouvain -- Juillet 2022, Q1 -- $f(x) = \ln(1/(1+\cos x))$) :}
\begin{itemize}
    \item \textbf{Domaine, parité, périodicité et racines :} $1 + \cos(x) > 0 \iff \cos(x) \neq -1 \iff \boxed{\D = \R \setminus \{\pi + 2k\pi \mid k \in \Z\}}$. $f$ est paire et $2\pi$-périodique. Racines : $f(x) = -\ln(1+\cos x) = 0 \iff 1 + \cos(x) = 1 \iff \cos(x) = 0 \iff \boxed{x = \frac{\pi}{2} + k\pi}$ ($k \in \Z$). Aux points exclus $x = \pi + 2k\pi$, $1 + \cos(x) \to 0^+ \implies f(x) \to +\infty$ ($\text{AV} \equiv x = \pi + 2k\pi$).
    \item \textbf{Dérivées première et seconde :} En dérivant $f(x) = -\ln(1+\cos x)$ :
    \[
        \boxed{f'(x) = \frac{\sin(x)}{1 + \cos(x)} = \tan\left(\frac{x}{2}\right)},
    \]
    \[
        \boxed{f''(x) = \frac{\cos(x)(1+\cos x) - \sin(x)(-\sin x)}{(1+\cos x)^2} = \frac{\cos(x) + \cos^2(x) + \sin^2(x)}{(1+\cos x)^2} = \frac{1 + \cos(x)}{(1+\cos x)^2} = \frac{1}{1 + \cos(x)} \ge \frac{1}{2} > 0}.
    \]
    \item \textbf{Extrema et tableau sur $]-\pi, \pi[$ :} $f'(x) = 0 \text{ dans } \D \iff \sin(x) = 0 \text{ et } \cos(x) \neq -1 \iff \boxed{x = 2k\pi}$. Comme $f''(x) > 0$ partout sur $\D$, la concavité est toujours positive ($\cup$), il n'y a pas de maximum, et les points $x = 2k\pi$ sont des \textbf{minima locaux} de valeur $f(0) = -\ln(2) \simeq -0,693$. Sur $]-\pi, \pi[$, $f$ décroît de $+\infty$ (AV en $-\pi$) jusqu'à $-\ln(2)$ en $x=0$ (en s'annulant en $-\frac{\pi}{2}$), puis remonte à $+\infty$ (AV en $\pi$) en s'annulant en $\frac{\pi}{2}$.
\end{itemize}

\paragraph{Corrigé de l'Exercice 4.3 (UCLouvain -- Juillet 2023 -- Série 1, Q1 -- $f(x) = x^2\sin(x^2)+\cos(x^2)$) :}
\begin{enumerate}[label=(\alph*)]
    \item $f$ est définie et paire sur $\R$. D'après l'Exercice 3.1a, $\boxed{f'(x) = 2x^3\cos(x^2)}$. Les points critiques ($f'(x) = 0$) sont :
    \[
        x = 0 \quad \text{ou} \quad \cos(x^2) = 0 \iff x^2 = \frac{\pi}{2} + k\pi = (2n-1)\frac{\pi}{2} \;\;(n \in \N_0) \iff \boxed{x = 0 \;\text{ ou }\; x = \pm\sqrt{(2n-1)\frac{\pi}{2}} \;\;(n \in \N_0)}.
    \]
    \item Dérivée seconde : $\boxed{f''(x) = 2x^2\bigl(3\cos(x^2) - 2x^2\sin(x^2)\bigr)}$.
    \item \textbf{Tableau sur $[0, 2\sqrt{\pi}]$ :} Sur cet intervalle, $x^2$ varie de $0$ à $4\pi$, donc $\cos(x^2)$ s'annule en $x^2 \in \left\{\frac{\pi}{2}, \frac{3\pi}{2}, \frac{5\pi}{2}, \frac{7\pi}{2}\right\}$. En ces points où $\cos(x^2) = 0$ et $\sin(x^2) = (-1)^{n-1} = \pm 1$, on a $f(x) = \pm x^2$ et $f''(x) = \mp 4x^4$ :
    \begin{center}
    \renewcommand{\arraystretch}{1.4}
    \setlength{\tabcolsep}{4pt}
    \begin{tabular}{|c||c|c|c|c|c|c|c|c|c|c|}
    \hline
    $x^2$ & $0$ & & $\dfrac{\pi}{2}$ & & $\dfrac{3\pi}{2}$ & & $\dfrac{5\pi}{2}$ & & $\dfrac{7\pi}{2}$ & $4\pi$ \\ \hline
    $x$ & $0$ & & $\sqrt{\dfrac{\pi}{2}}$ & & $\sqrt{\dfrac{3\pi}{2}}$ & & $\sqrt{\dfrac{5\pi}{2}}$ & & $\sqrt{\dfrac{7\pi}{2}}$ & $2\sqrt{\pi}$ \\ \hline\hline
    $f'(x)$ & $0$ & $+$ & $0$ & $-$ & $0$ & $+$ & $0$ & $-$ & $0$ & $+\;(\text{en } 2\sqrt{\pi})$ \\ \hline
    $f''(x)$ & $0$ & & $-\pi^2 < 0$ & & $+(3\pi)^2 > 0$ & & $-(5\pi)^2 < 0$ & & $+(7\pi)^2 > 0$ & $+$ \\ \hline
    $f(x)$ & $\mathbf{1}$ (min plat) & $\nearrow$ & $\mathbf{\dfrac{\pi}{2}}$ (max) & $\searrow$ & $\mathbf{-\dfrac{3\pi}{2}}$ (min) & $\nearrow$ & $\mathbf{\dfrac{5\pi}{2}}$ (max) & $\searrow$ & $\mathbf{-\dfrac{7\pi}{2}}$ (min) & $\nearrow\; 1$ \\ \hline
    \end{tabular}
    \end{center}
    (En $x = 0$, par parité, $f'(x) = 2x^3\cos(x^2)$ est négative à gauche de $0$ et positive à droite de $0$, donc $x = 0$ est bien un minimum local très plat de valeur $f(0) = 1$).
\end{enumerate}

\paragraph{Corrigé de l'Exercice 4.4 (UCLouvain -- Juillet 2023 -- Série 2, Q1 -- $f(x) = 3\cos(x) - \cos^3(x)$) :}
\begin{enumerate}[label=(\alph*)]
    \item $f$ est définie, paire et $2\pi$-périodique sur $\R$. D'après l'Exercice 3.1b, $\boxed{f'(x) = -3\sin^3(x)}$.
    \begin{itemize}
        \item \textbf{Candidats extrema :} $f'(x) = 0 \iff \sin(x) = 0 \iff \boxed{x = k\pi \;\;(k \in \Z)}$.
        \item \textbf{Racines de $f$ :} $f(x) = \cos(x)\bigl(3 - \cos^2(x)\bigr) = 0 \iff \cos(x) = 0$ (car $3 - \cos^2(x) \ge 2 > 0$) $\iff \boxed{x = \frac{\pi}{2} + k\pi \;\;(k \in \Z)}$.
    \end{itemize}
    \item Dérivée seconde : $\boxed{f''(x) = -9\sin^2(x)\cos(x)}$. Elle s'annule en $x = k\pi$ (sans changer de signe car $\sin^2(x) \ge 0$, ce qui explique que les extrema en $k\pi$ sont « très plats » avec $f' = f'' = 0$ !) et en $x = \frac{\pi}{2} + k\pi$ (en changeant de signe : \textbf{points d'inflexion} qui coïncident avec les racines de $f$, de pente $f'(\frac{\pi}{2}+k\pi) = \mp 3$).
    \item \textbf{Tableau sur $[0, 3\pi]$ :}
    \begin{center}
    \renewcommand{\arraystretch}{1.35}
    \setlength{\tabcolsep}{4.2pt}
    \begin{tabular}{|c||c|c|c|c|c|c|c|c|c|c|c|c|c|}
    \hline
    $x$ & $0$ & & $\dfrac{\pi}{2}$ & & $\pi$ & & $\dfrac{3\pi}{2}$ & & $2\pi$ & & $\dfrac{5\pi}{2}$ & & $3\pi$ \\ \hline\hline
    $f'(x)$ & $0$ & $-$ & $-3$ & $-$ & $0$ & $+$ & $3$ & $+$ & $0$ & $-$ & $-3$ & $-$ & $0$ \\ \hline
    $f''(x)$ & $0$ & $-$ & $0$ & $+$ & $0$ & $+$ & $0$ & $-$ & $0$ & $-$ & $0$ & $+$ & $0$ \\ \hline
    $f(x)$ & $\mathbf{2}$ (max) & $\searrow\cap$ & $\mathbf{0}$ (PI) & $\searrow\cup$ & $\mathbf{-2}$ (min) & $\nearrow\cup$ & $\mathbf{0}$ (PI) & $\nearrow\cap$ & $\mathbf{2}$ (max) & $\searrow\cap$ & $\mathbf{0}$ (PI) & $\searrow\cup$ & $\mathbf{-2}$ (min) \\ \hline
    \end{tabular}
    \end{center}
\end{enumerate}

\paragraph{Corrigé de l'Exercice 4.5 (UCLouvain -- Septembre 2023, Q1 -- $f(x) = x/\cos^2(x)$) :}
\begin{itemize}
    \item \textbf{Domaine :} $\D = \R \setminus \left\{\frac{\pi}{2} + k\pi \mid k \in \Z\right\}$ (fonction impaire).
    \item \textbf{Dérivée première :}
    \[
        f'(x) = \frac{1\cdot\cos^2(x) - x\cdot 2\cos(x)(-\sin(x))}{\cos^4(x)} = \boxed{\frac{\cos(x) + 2x\sin(x)}{\cos^3(x)} = \frac{1 + 2x\tan(x)}{\cos^2(x)}}.
    \]
    \item \textbf{Tableau de valeurs sur $[0, \pi]$ :}
    \begin{center}
    \renewcommand{\arraystretch}{1.35}
    \begin{tabular}{|c||c|c|c|c|c|}
    \hline
    $x$ & $0$ & $\dfrac{\pi}{4}$ & $\dfrac{\pi}{2}$ & $\dfrac{3\pi}{4}$ & $\pi$ \\ \hline\hline
    $f(x)$ & $0$ (racine) & $\dfrac{\pi}{2}$ & $\text{AV}\;(+\infty)$ & $\dfrac{3\pi}{2}$ & $\pi$ \\ \hline
    $f'(x)$ & $1 > 0$ & $2 + \pi > 0$ & $+\infty \mid -\infty$ & $2 - 3\pi < 0$ & $1 > 0$ \\ \hline
    \end{tabular}
    \end{center}
    \item \textbf{Existence d'un minimum sur $\left]\frac{3\pi}{4}, \pi\right[$ :} Sur $\left[\frac{3\pi}{4}, \pi\right]$, la dérivée $f'$ est continue (car $\cos(x) \neq 0$), vaut $f'\left(\frac{3\pi}{4}\right) = 2 - 3\pi < 0$ et $f'(\pi) = 1 > 0$. Par le théorème des valeurs intermédiaires, $f'$ s'annule en passant du signe $-$ au signe $+$ en un point de $\left]\frac{3\pi}{4}, \pi\right[$, où $f$ admet donc un \textbf{minimum local}.
\end{itemize}

\paragraph{Corrigé de l'Exercice 4.6 (UCLouvain -- Juillet 2026, Q1 -- $f(x) = (1-x^2)\ln(1-x)+x$) :}
\begin{itemize}
    \item \textbf{Domaine et limite en $1^-$ :} $\D = \,]-\infty, 1[$ et $\lim_{x \to 1^-} f(x) = 2\lim_{y \to 0^+} y\ln(y) + 1 = \mathbf{1}$ (vu au Chapitre 2).
    \item \textbf{Dérivée première $f'(x)$ et ses racines :}
    \[
        f'(x) = -2x\ln(1-x) + (1-x^2)\frac{-1}{1-x} + 1 = -2x\ln(1-x) - (1+x) + 1 = \boxed{-x\bigl(2\ln(1-x) + 1\bigr)}.
    \]
    $f'(x) = 0 \iff x = 0 \;\text{ ou }\; \ln(1-x) = -\frac{1}{2} \iff \boxed{x_1 = 0 \;\text{ ou }\; x_2 = 1 - e^{-1/2} \simeq 1 - 0,6065 = 0,3935}$.
    (De plus, $\lim_{x \to 1^-} f'(x) = -1(2(-\infty)+1) = +\infty$ : demi-tangente verticale au point limite $(1, 1)$).
    \item \textbf{Dérivée seconde $f''(x)$ :} En dérivant $f'(x) = -2x\ln(1-x) - x$ :
    \[
        f''(x) = -2\ln(1-x) - 2x\frac{-1}{1-x} - 1 = \boxed{\frac{2x}{1-x} - \bigl(2\ln(1-x) + 1\bigr)}.
    \]
    \item \textbf{Valeurs remarquables et tableau de variations sur $[-1, 1[$ :}
    \begin{itemize}
        \item En $x = -1$ : $f(-1) = 0\cdot\ln(2) - 1 = -1$ ; $f'(-1) = 2\ln(2)+1 \simeq 2,3862 > 0$ ; $f''(-1) = -1 - (2\ln 2 + 1) \simeq -3,3862 < 0$.
        \item En $x_1 = 0$ : $f(0) = 0$ (racine et maximum local car $f'(0) = 0$ et $f''(0) = -1 < 0$).
        \item En $x_2 = 1 - e^{-1/2} \simeq 0,3935$ : comme $2\ln(1-x_2) + 1 = 0$, on a $f'(x_2) = 0$, $f''(x_2) = \frac{2x_2}{1-x_2} = 2(e^{1/2} - 1) \simeq 1,2976 > 0$ (minimum local), et $f(x_2) = -\frac{1}{2}(1-x_2^2) + x_2 \simeq -0,0291 < 0$.
    \end{itemize}
    \begin{center}
    \renewcommand{\arraystretch}{1.35}
    \setlength{\tabcolsep}{4pt}
    \begin{tabular}{|c||c|c|c|c|c|c|c|}
    \hline
    $x$ & $-1$ & & $x_1 = 0$ & & $x_2 \simeq 0,3935$ & & $1$ \\ \hline\hline
    $f'(x)$ & $2,386$ & $+$ & $0$ & $-$ & $0$ & $+$ & $\lim = +\infty$ \\ \hline
    $f''(x)$ & $-3,386$ & $-$ & $-1$ & $-\;0\;+$ & $+1,298$ & $+$ & $\lim = +\infty$ \\ \hline
    $f(x)$ & $-1$ & $\nearrow\;\cap$ & $\mathbf{0}$ (max \& racine) & $\searrow\;(\text{PI})$ & $\mathbf{-0,0291}$ (min) & $\nearrow\;\cup$ & $\lim = 1$ (2\ieme{} racine entre $x_2$ et $1$) \\ \hline
    \end{tabular}
    \end{center}
\end{itemize}

% ----------------------------------------------------------------------------
\section{Corrigés du Chapitre 5 (Problèmes d'Optimisation Appliquée)}
% ----------------------------------------------------------------------------

\paragraph{Corrigé de l'Exercice 5.1 (UCLouvain -- Juillet 2022, Q2 -- Le drone-lanceur) :}
\begin{enumerate}[label=(\alph*)]
    \item Comme $x(t) = vt$, l'altitude $H$ est atteinte à l'instant $\boxed{T = \dfrac{H}{v}}$.
    \item Intégrons la puissance instantanée $Q(t) = K + \dfrac{Cv^3}{vt + x_0}$ entre $t = 0$ et $t = T = \dfrac{H}{v}$ :
    \[
        E(v) = \int_0^{H/v} \left(K + \frac{Cv^3}{vt + x_0}\right)\dd t = \left[Kt + Cv^2\ln(vt + x_0)\right]_0^{H/v} = \boxed{K\frac{H}{v} + Cv^2\ln\left(1 + \frac{H}{x_0}\right)}.
    \]
    \item Dérivons $E(v)$ par rapport à la vitesse $v > 0$ (les constantes $K, C, H, x_0$ étant fixées) :
    \[
        E'(v) = -\frac{KH}{v^2} + 2Cv\ln\left(1 + \frac{H}{x_0}\right) = 0 \iff v^3 = \frac{KH}{2C\ln(1 + H/x_0)} \iff \boxed{v^* = \sqrt[3]{\frac{KH}{2C\ln(1 + H/x_0)}}}.
    \]
    Comme $E''(v) = \frac{2KH}{v^3} + 2C\ln(1+H/x_0) > 0$, il s'agit bien d'un minimum, et la consommation minimale vaut :
    \[
        \boxed{E(v^*) = \frac{3}{2^{2/3}}(KH)^{2/3}\left(C\ln\left(1 + \frac{H}{x_0}\right)\right)^{1/3}}.
    \]
\end{enumerate}

\paragraph{Corrigé de l'Exercice 5.2 (UCLouvain -- Septembre 2023, Q2 -- Le train de nuit) :}
\begin{enumerate}[label=(\alph*)]
    \item Sur $\left[-\frac{\pi}{2}, \frac{\pi}{2}\right]$, $\ell'(t) = b\cos(t) \ge 0$ (car $b > 0$ et $\cos t \ge 0$), donc $|\ell'(t)| = \ell'(t)$. Ainsi :
    \begin{align*}
        \bar{c}(a) &= \frac{1}{\pi}\int_{-\pi/2}^{\pi/2} \left(a + b + b\sin(t) + \frac{\ell'(t)}{\ell(t)}\right)\dd t = \frac{1}{\pi}\Bigl[(a+b)t - b\cos(t) + \ln\bigl(a + b(1+\sin t)\bigr)\Bigr]_{-\pi/2}^{\pi/2} \\
        &= \frac{1}{\pi}\Bigl[(a+b)\pi - 0 + \ln(a + 2b) - \ln(a)\Bigr] = \boxed{a + b + \frac{1}{\pi}\bigl(\ln(a + 2b) - \ln(a)\bigr)}.
    \end{align*}
    \item Dérivons $\bar{c}(a)$ par rapport à $a > 0$ :
    \[
        \bar{c}'(a) = 1 + \frac{1}{\pi}\left(\frac{1}{a + 2b} - \frac{1}{a}\right) = 1 - \frac{2b}{\pi a(a+2b)} = 0 \iff \pi a^2 + 2b\pi a - 2b = 0.
    \]
    Le discriminant réduit vaut $\Delta' = b^2\pi^2 + 2b\pi > 0$. Comme $a > 0$, on retient l'unique racine positive :
    \[
        \boxed{a^* = \frac{-b\pi + \sqrt{b^2\pi^2 + 2b\pi}}{\pi} = -b + \sqrt{b^2 + \frac{2b}{\pi}}}.
    \]
    \item Avec la fréquence $\lambda > 0$ sur $\left[-\frac{\pi}{2\lambda}, \frac{\pi}{2\lambda}\right]$ (de durée $\frac{\pi}{\lambda}$), le même calcul donne $\bar{c}(a) = a + b + \frac{\lambda}{\pi}\bigl(\ln(a+2b) - \ln(a)\bigr)$, ce qui revient à remplacer $\frac{1}{\pi}$ par $\frac{\lambda}{\pi}$, d'où $\boxed{a^* = -b + \sqrt{b^2 + \frac{2b\lambda}{\pi}}}$.
\end{enumerate}

\paragraph{Corrigé de l'Exercice 5.3 (UCLouvain -- Septembre 2024, Q2 -- Le thermos de café) :}
\begin{enumerate}[label=(\alph*)]
    \item La dérivée de $c(t) = \frac{D}{1-\beta/\alpha}\bigl(e^{-\beta t} - e^{-\alpha t}\bigr)$ par rapport à $t$ s'annule au temps $t_{\max}$ tel que :
    \[
        \frac{\mathrm{d}c}{\mathrm{d}t}(t_{\max}) = \frac{D}{1 - \beta/\alpha}\bigl(\alpha e^{-\alpha t_{\max}} - \beta e^{-\beta t_{\max}}\bigr) = 0 \iff \alpha e^{-\alpha t_{\max}} = \beta e^{-\beta t_{\max}}.
    \]
    Comme calculé à l'Exemple du Chapitre 1, cela donne :
    \[
        t_{\max} = \frac{1}{\alpha - \beta}\ln\left(\frac{\alpha}{\beta}\right) = \frac{\ln(5)}{0,8} \simeq \frac{1,61}{0,8} = 2,0125\text{ h} = 2,0125 \times 60 = 120,75\text{ min} \simeq 121\text{ minutes}.
    \]
    Pour que ce pic (qui a lieu $121\text{ min}$ après avoir bu le café) se produise $60\text{ min}$ après le début de l'examen, il faut boire le café $121 - 60 = \boxed{61\text{ minutes avant le début de l'examen}}$.
    \item La dérivée seconde en $t_{\max}$ vaut, en remplaçant $\alpha e^{-\alpha t_{\max}} = \beta e^{-\beta t_{\max}}$ :
    \[
        \left.\frac{\mathrm{d}^2 c}{\mathrm{d}t^2}\right|_{t_{\max}} = \frac{D}{1 - \beta/\alpha}\bigl(-\alpha^2 e^{-\alpha t_{\max}} + \beta^2 e^{-\beta t_{\max}}\bigr) = \frac{D}{1 - \beta/\alpha}\beta e^{-\beta t_{\max}}(\beta - \alpha) = -D\alpha\beta e^{-\beta t_{\max}}.
    \]
    Or la concentration de caféine dans le sang est une grandeur physique positive ($c(t) > 0$ pour $t > 0$) et $1 - \beta/\alpha = 0,8 > 0$ avec $e^{-\beta t} > e^{-\alpha t}$, ce qui impose $D > 0$. Comme $D, \alpha, \beta, e^{-\beta t_{\max}} > 0$, on a bien $\left.\frac{\mathrm{d}^2 c}{\mathrm{d}t^2}\right|_{t_{\max}} < 0$ : c'est un maximum !
\end{enumerate}

\paragraph{Corrigé de l'Exercice 5.4 (UCLouvain -- Juillet 2025, Q2 -- La ligue des « legendgénieurs ») :}
\begin{itemize}
    \item Comme $E = kA^2$ vaut $1$ pour $A = 560$, on a $k = \frac{1}{560^2}$, soit $E(A) = \left(\dfrac{A}{560}\right)^2$.
    \item La contrainte de budget est $A + V = 560 \implies V = 560 - A$. En remplaçant $V$ et $E$ dans la durée de vie $D = \left(7 + \frac{V}{10}\right)E$, on obtient une fonction de la seule variable $A \in [0, 560]$ :
    \[
        D(A) = \left(7 + \frac{560 - A}{10}\right)\left(\frac{A}{560}\right)^2 = \frac{1}{560^2}\left(63\,A^2 - \frac{1}{10}A^3\right).
    \]
    \item Dérivons par rapport à $A$ :
    \[
        D'(A) = \frac{1}{560^2}\left(126\,A - \frac{3}{10}A^2\right) = \frac{A}{560^2}\left(126 - \frac{3}{10}A\right) = 0 \implies A = 0 \text{ (min)} \;\text{ ou }\; \boxed{A^* = \frac{1260}{3} = 420\text{ bitcoins}},
    \]
    d'où $\boxed{V^* = 560 - 420 = 140\text{ bitcoins}}$ (avec $D''(420) = \frac{-126}{560^2} < 0$, donc maximum).
    \item La durée de vie maximale vaut $D(420) = \left(7 + \frac{140}{10}\right)\left(\frac{420}{560}\right)^2 = 21 \times \left(\frac{3}{4}\right)^2 = \boxed{\dfrac{189}{16}\text{ jours}}$.
    Effectuons la division euclidienne : $189 = 11 \times 16 + 13$, soit $\boxed{J = 11\text{ jours}}$ et un reste de $\frac{13}{16}\text{ jour} = \frac{13}{16} \times 24\text{ h} = \frac{39}{2}\text{ h} = 19,5\text{ h}$, soit $\boxed{H = 19\text{ heures}}$ et $\boxed{M = 30\text{ minutes}}$.
\end{itemize}

\paragraph{Corrigé de l'Exercice 5.5 (UCLouvain -- Septembre 2025, Q2 -- Le moteur propulseur spatial) :}
\begin{enumerate}[label=(\alph*)]
    \item L'accélération en fonction du flux $F > 0$ est $A(F) = k\,V_{\max}\,F\,\exp(-F/F_c)$. En dérivant par rapport à $F$ :
    \[
        \frac{\mathrm{d}A}{\mathrm{d}F} = k V_{\max}\left(1 - \frac{F}{F_c}\right)\exp(-F/F_c) = 0 \iff \boxed{F^* = F_c = 100\text{ g/s}}.
    \]
    Comme $\frac{\mathrm{d}A}{\mathrm{d}F} > 0$ pour $F < F_c$ et $< 0$ pour $F > F_c$, c'est un maximum global. Convertissons en $\text{kg/min}$ : $100\text{ g/s} = 100 \times 60\text{ g/min} = 6000\text{ g/min} = \boxed{6\text{ kg/min}}$.
    \item Pour $A(F) = k\,V_{\max}\,F\,\exp\bigl[-(F/F_c)^{1/2}\bigr]$, la règle du produit et de la chaîne donne :
    \[
        \frac{\mathrm{d}A}{\mathrm{d}F} = k V_{\max}\left(1 - \frac{1}{2}\left(\frac{F}{F_c}\right)^{1/2}\right)\exp\bigl[-(F/F_c)^{1/2}\bigr] = 0 \iff \left(\frac{F}{F_c}\right)^{1/2} = 2 \iff \boxed{F^* = 4F_c = 400\text{ g/s}}.
    \]
\end{enumerate}

\paragraph{Corrigé de l'Exercice 5.6 (Cylindre de volume maximal inscrit dans une sphère de rayon $R$) :}
Dans une coupe méridienne passant par l'axe du cylindre, le théorème de Pythagore relie le rayon $r$ du cylindre, sa demi-hauteur $\frac{h}{2}$ et le rayon $R$ de la sphère : $r^2 + \left(\frac{h}{2}\right)^2 = R^2 \implies r^2 = R^2 - \frac{h^2}{4}$.
Remplaçons directement $r^2$ dans le volume du cylindre $V = \pi r^2 h$ pour obtenir une fonction polynomiale de la seule variable $h \in \,]0, 2R[$ :
\[
    V(h) = \pi\left(R^2 h - \frac{h^3}{4}\right) \implies V'(h) = \pi\left(R^2 - \frac{3h^2}{4}\right) = 0 \iff \boxed{h^* = \frac{2R}{\sqrt{3}}},
\]
(avec $V''(h) = -\frac{3\pi}{2}h < 0$, donc maximum), d'où $r^* = \sqrt{R^2 - \frac{R^2}{3}} = \boxed{R\sqrt{\frac{2}{3}}}$ et $\boxed{V_{\max} = \frac{4\pi R^3}{3\sqrt{3}} = \frac{1}{\sqrt{3}}V_{\text{sphère}}}$.

% ----------------------------------------------------------------------------
\section{Corrigés du Chapitre 6 (Calcul Intégral et Primitives)}
% ----------------------------------------------------------------------------

\paragraph{Corrigé de l'Exercice 6.1 (Primitives par substitution directe) :}
\begin{enumerate}[label=(\alph*)]
    \item Posons $u = e^x \implies \mathrm{d}u = e^x\dd x$ : $\displaystyle\int \frac{e^x}{1 + (e^x)^2}\dd x = \int \frac{\mathrm{d}u}{1 + u^2} = \arctan(u) + C = \boxed{\arctan(e^x) + C}$.
    \item Posons $u = \ln(x) \implies \mathrm{d}u = \frac{1}{x}\dd x$ : $\displaystyle\int \ln^3(x)\frac{\mathrm{d}x}{x} = \int u^3\dd u = \frac{u^4}{4} + C = \boxed{\frac{\ln^4(x)}{4} + C}$.
    \item Pour la première, posons $u = 4 - x^2 \implies x\dd x = -\frac{1}{2}\dd u$ : $\displaystyle\int \frac{x\dd x}{\sqrt{4-x^2}} = -\frac{1}{2}\int \frac{\mathrm{d}u}{\sqrt{u}} = \boxed{-\sqrt{4 - x^2} + C}$.\\
    Pour la seconde, factorisons $4$ sous la racine et posons $u = \frac{x}{2} \implies \mathrm{d}x = 2\dd u$ :
    \[
        \int \frac{\mathrm{d}x}{2\sqrt{1 - (x/2)^2}} = \int \frac{\mathrm{d}u}{\sqrt{1 - u^2}} = \arcsin(u) + C = \boxed{\arcsin\left(\frac{x}{2}\right) + C}.
    \]
    \item Posons $u = \ln(x) \implies \mathrm{d}u = \frac{1}{x}\dd x$ : $\displaystyle\int \cos(\ln x)\frac{\mathrm{d}x}{x} = \int \cos(u)\dd u = \sin(u) + C = \boxed{\sin\bigl(\ln(x)\bigr) + C}$.
\end{enumerate}

\paragraph{Corrigé de l'Exercice 6.2 (UCLouvain -- Septembre 2021, Q3c) :}
\textbf{CE :} $(e^x)^2 \neq 1 \iff e^x \neq 1$ (car $e^x > 0$) $\iff \mathbf{x \neq 0}$.
Posons le changement de variable $y = e^x > 0 \implies \mathrm{d}y = e^x\dd x$. Par décomposition en fractions simples :
\[
    \int \frac{2e^x}{(e^x)^2 - 1}\dd x = \int \frac{2}{y^2 - 1}\dd y = \int \left(\frac{1}{y - 1} - \frac{1}{y + 1}\right)\dd y = \ln(|y - 1|) - \ln(y + 1) + C = \boxed{\ln\left(\frac{|e^x - 1|}{e^x + 1}\right) + C}.
\]

\paragraph{Corrigé de l'Exercice 6.3 (Primitives par intégration par parties) :}
\begin{enumerate}[label=(\alph*)]
    \item IPP avec $u = \arctan(x) \implies \mathrm{d}u = \frac{1}{1+x^2}\dd x$ et $\mathrm{d}v = 1\dd x \implies v = x$ :
    \[
        \int \arctan(x)\dd x = x\arctan(x) - \int \frac{x}{1 + x^2}\dd x = \boxed{x\arctan(x) - \frac{1}{2}\ln(1 + x^2) + C}.
    \]
    \item IPP avec $u = \arcsin(x) \implies \mathrm{d}u = \frac{1}{\sqrt{1-x^2}}\dd x$ et $\mathrm{d}v = 1\dd x \implies v = x$ :
    \[
        \int \arcsin(x)\dd x = x\arcsin(x) - \int \frac{x}{\sqrt{1 - x^2}}\dd x = \boxed{x\arcsin(x) + \sqrt{1 - x^2} + C}.
    \]
    \item IPP avec $u = 2x + 1 \implies \mathrm{d}u = 2\dd x$ et $\mathrm{d}v = e^{-x}\dd x \implies v = -e^{-x}$ :
    \[
        \int (2x + 1)e^{-x}\dd x = -(2x + 1)e^{-x} + 2\int e^{-x}\dd x = -(2x + 1)e^{-x} - 2e^{-x} + C = \boxed{-(2x + 3)e^{-x} + C}.
    \]
    \item IPP avec $u = \ln(x) \implies \mathrm{d}u = \frac{1}{x}\dd x$ et $\mathrm{d}v = x^2\dd x \implies v = \frac{x^3}{3}$ :
    \[
        \int x^2\ln(x)\dd x = \frac{x^3}{3}\ln(x) - \int \frac{x^3}{3}\cdot\frac{1}{x}\dd x = \frac{x^3}{3}\ln(x) - \frac{1}{3}\int x^2\dd x = \boxed{\frac{x^3}{3}\left(\ln(x) - \frac{1}{3}\right) + C}.
    \]
\end{enumerate}

\paragraph{Corrigé de l'Exercice 6.4 (IPP cyclique) :}
\begin{enumerate}[label=(\alph*)]
    \item Pour $I_1 = \int e^x\cos(x)\dd x$ :
    1\iere{} IPP ($u = \cos x, \mathrm{d}v = e^x\dd x$) : $I_1 = e^x\cos(x) + \int e^x\sin(x)\dd x$.\\
    2\ieme{} IPP ($u = \sin x, \mathrm{d}v = e^x\dd x$) : $\int e^x\sin(x)\dd x = e^x\sin(x) - \int e^x\cos(x)\dd x = e^x\sin(x) - I_1$.\\
    D'où $2I_1 = e^x(\cos x + \sin x) \implies \boxed{I_1 = \frac{1}{2}e^x\bigl(\cos(x) + \sin(x)\bigr) + C}$.
    \item De même pour $I_2 = \int e^{2x}\sin(x)\dd x$ :
    1\iere{} IPP ($u = \sin x, \mathrm{d}v = e^{2x}\dd x \implies v = \frac{1}{2}e^{2x}$) : $I_2 = \frac{1}{2}e^{2x}\sin(x) - \frac{1}{2}\int e^{2x}\cos(x)\dd x$.\\
    2\ieme{} IPP ($u = \cos x, \mathrm{d}v = e^{2x}\dd x$) : $\int e^{2x}\cos(x)\dd x = \frac{1}{2}e^{2x}\cos(x) + \frac{1}{2}I_2$.\\
    En remplaçant : $I_2 = \frac{1}{2}e^{2x}\sin(x) - \frac{1}{4}e^{2x}\cos(x) - \frac{1}{4}I_2 \iff \frac{5}{4}I_2 = \frac{1}{4}e^{2x}(2\sin x - \cos x) \implies \boxed{I_2 = \frac{1}{5}e^{2x}\bigl(2\sin(x) - \cos(x)\bigr) + C}$.
\end{enumerate}

\paragraph{Corrigé de l'Exercice 6.5 (Primitives trigonométriques) :}
\begin{enumerate}[label=(\alph*)]
    \item Puissance impaire de $\sin(x)$ : posons $u = \cos(x) \implies \mathrm{d}u = -\sin(x)\dd x$ avec $\sin^2(x) = 1 - u^2$ :
    \[
        \int (1 - \cos^2 x)\cos^2(x)\sin(x)\dd x = -\int (1 - u^2)u^2\dd u = \int (u^4 - u^2)\dd u = \boxed{\frac{\cos^5(x)}{5} - \frac{\cos^3(x)}{3} + C}.
    \]
    \item Par duplication $\sin(x)\cos(x) = \frac{1}{2}\sin(2x)$, puis Carnot $\sin^2(2x) = \frac{1 - \cos(4x)}{2}$ :
    \[
        \int \sin^2(x)\cos^2(x)\dd x = \frac{1}{4}\int \sin^2(2x)\dd x = \frac{1}{8}\int \bigl(1 - \cos(4x)\bigr)\dd x = \boxed{\frac{x}{8} - \frac{\sin(4x)}{32} + C}.
    \]
    \item En utilisant $1 + \tan^2(x) = \dfrac{1}{\cos^2(x)}$ :
    \[
        \int \tan^2(x)\dd x = \int \left(\frac{1}{\cos^2(x)} - 1\right)\dd x = \boxed{\tan(x) - x + C},
    \]
    \[
        \int \tan^3(x)\dd x = \int \tan(x)\left(\frac{1}{\cos^2(x)} - 1\right)\dd x = \boxed{\frac{1}{2}\tan^2(x) + \ln\bigl(|\cos(x)|\bigr) + C}.
    \]
\end{enumerate}

\paragraph{Corrigé de l'Exercice 6.6 (Défi multi-méthodes sur $\int \frac{x^3}{(x^2+1)^3}\dd x$) :}
\begin{enumerate}[label=(\alph*)]
    \item \textbf{Par substitution $u = x^2 + 1$ ($\mathrm{d}u = 2x\dd x$ et $x^2 = u - 1$) :}
    \[
        \int \frac{x^2}{(x^2+1)^3}\,x\dd x = \frac{1}{2}\int \frac{u - 1}{u^3}\dd u = \frac{1}{2}\int \bigl(u^{-2} - u^{-3}\bigr)\dd u = -\frac{1}{2u} + \frac{1}{4u^2} + C_1 = \boxed{-\frac{2x^2 + 1}{4(x^2 + 1)^2} + C_1}.
    \]
    \item \textbf{Par intégration par parties :} Posons $p = x^2 \implies \mathrm{d}p = 2x\dd x$ et $\mathrm{d}q = \frac{x}{(x^2+1)^3}\dd x \implies q = -\frac{1}{4(x^2+1)^2}$ :
    \[
        \int \frac{x^3}{(x^2+1)^3}\dd x = -\frac{x^2}{4(x^2+1)^2} + \frac{1}{2}\int \frac{x}{(x^2+1)^2}\dd x = -\frac{x^2}{4(x^2+1)^2} - \frac{1}{4(x^2+1)} + C_2 = \boxed{-\frac{2x^2 + 1}{4(x^2 + 1)^2} + C_2}.
    \]
    Les deux méthodes conduisent exactement à la même expression !
\end{enumerate}

% ----------------------------------------------------------------------------
\section{Corrigés du Chapitre 7 (Intégrales Définies, Inégalités, Aires et Volumes)}
% ----------------------------------------------------------------------------

\paragraph{Corrigé de l'Exercice 7.1 (Calculs d'intégrales définies) :}
\begin{enumerate}[label=(\alph*)]
    \item $\displaystyle\int_0^{\pi/4} \tan(x)\dd x = \Bigl[-\ln(\cos x)\Bigr]_0^{\pi/4} = -\ln\left(\frac{\sqrt{2}}{2}\right) + \ln(1) = -\ln\left(2^{-1/2}\right) = \boxed{\frac{1}{2}\ln(2)}$.
    \item Par l'IPP connue $\int x\ln(x)\dd x = \frac{x^2}{2}\ln(x) - \frac{x^2}{4}$ :
    \[
        \int_1^e x\ln(x)\dd x = \left[\frac{x^2}{2}\ln(x) - \frac{x^2}{4}\right]_1^e = \left(\frac{e^2}{2} - \frac{e^2}{4}\right) - \left(0 - \frac{1}{4}\right) = \boxed{\frac{e^2 + 1}{4}}.
    \]
    \item La fonction $f(x) = \dfrac{x^3\cos(x)}{1+x^2}$ est continue et \textbf{impaire} sur l'intervalle symétrique $[-1, 1]$ (car $(-x)^3 = -x^3$ et $\cos(-x) = \cos(x)$). Son intégrale sur $[-1, 1]$ est donc immédiatement nulle : $\boxed{0}$.
    \item En posant $u = \cos(x) \implies \mathrm{d}u = -\sin(x)\dd x$ :
    \[
        \int_0^{\pi/2} \cos^2(x)\sin(x)\dd x = \left[-\frac{\cos^3(x)}{3}\right]_0^{\pi/2} = 0 - \left(-\frac{1}{3}\right) = \boxed{\frac{1}{3}}.
    \]
\end{enumerate}

\paragraph{Corrigé de l'Exercice 7.2 (UCLouvain -- Septembre 2021, Q3e -- Majoration de $\sin(x)/x$) :}
\begin{enumerate}[label=(\alph*)]
    \item Sur $\left]-\frac{\pi}{2}, \frac{\pi}{2}\right[$, $g'(x) = \frac{1}{\cos^2(x)} - 1 = \frac{1-\cos^2(x)}{\cos^2(x)} = \tan^2(x) \ge 0$, et $g'(x)$ ne s'annule qu'en $x = 0$. Donc $g$ est \textbf{strictement croissante} sur $\left]-\frac{\pi}{2}, \frac{\pi}{2}\right[$. Comme $g(0) = 0$, $x = 0$ est son unique racine, avec $\boxed{g(x) < 0 \text{ sur } \left]-\frac{\pi}{2}, 0\right[}$ et $\boxed{g(x) > 0 \text{ sur } \left]0, \frac{\pi}{2}\right[}$.
    \item Pour $x \in \left]-\frac{\pi}{2}, \frac{\pi}{2}\right[ \setminus \{0\}$, on a $\cos(x) > 0$ et :
    \[
        f'(x) = \frac{x\cos(x) - \sin(x)}{x^2} = -\frac{\cos(x)}{x^2}\left(\frac{\sin(x)}{\cos(x)} - x\right) = -\frac{\cos(x)}{x^2}\,g(x).
    \]
    Puisque $-\frac{\cos(x)}{x^2} < 0$, $f'(x)$ est du signe opposé à $g(x)$ : $f'(x) > 0$ sur $\left]-\frac{\pi}{2}, 0\right[$ et $f'(x) < 0$ sur $\left]0, \frac{\pi}{2}\right[$. La fonction $f$ est donc \textbf{strictement croissante sur $\left]-\frac{\pi}{2}, 0\right]$ et strictement décroissante sur $\left[0, \frac{\pi}{2}\right[$}.
    \item \begin{itemize}
        \item \textbf{Si $|x| \le 1$ :} Comme $1 < \frac{\pi}{2} \approx 1,57$, on a $x \in [-1, 1] \subset \left]-\frac{\pi}{2}, \frac{\pi}{2}\right[$. D'après (b), $f$ atteint son maximum global sur $\left]-\frac{\pi}{2}, \frac{\pi}{2}\right[$ en $x = 0$, donc $f(x) \le f(0) = 1$.
        \item \textbf{Si $|x| > 1$ :}Comme $|\sin(x)| \le 1$, on a $f(x) \le |f(x)| = \dfrac{|\sin(x)|}{|x|} \le \dfrac{1}{|x|} < 1$.
    \end{itemize}
    Dans tous les cas, $\boxed{\forall x \in \R,\; f(x) \le 1}$.
\end{enumerate}

\paragraph{Corrigé de l'Exercice 7.3 (Inégalités polynomiales de Taylor pour $\cos(x)$ et $\sin(x)$) :}
\begin{itemize}
    \item \textbf{Étape 1 :} Pour tout $s \ge 0$, $\cos(s) \le 1$. En intégrant sur $[0, x]$ ($x \ge 0$) : $\int_0^x \cos(s)\dd s \le \int_0^x 1\dd s \iff \boxed{\sin(x) \le x}$.
    \item \textbf{Étape 2 :} Pour tout $s \ge 0$, $\sin(s) \le s$. En intégrant sur $[0, x]$ ($x \ge 0$) :
    \[
        \int_0^x \sin(s)\dd s \le \int_0^x s\dd s \iff \Bigl[-\cos(s)\Bigr]_0^x \le \frac{x^2}{2} \iff 1 - \cos(x) \le \frac{x^2}{2} \iff \boxed{1 - \frac{x^2}{2} \le \cos(x) \le 1}.
    \]
    \item \textbf{Étape 3 :} Intégrons à nouveau $1 - \frac{s^2}{2} \le \cos(s)$ sur $[0, x]$ ($x \ge 0$) :
    \[
        \int_0^x \left(1 - \frac{s^2}{2}\right)\dd s \le \int_0^x \cos(s)\dd s \iff \boxed{x - \frac{x^3}{6} \le \sin(x) \le x}.
    \]
\end{itemize}

\paragraph{Corrigé de l'Exercice 7.4 (Aire comprise entre une parabole et une droite) :}
Cherchons d'abord les points d'intersection des deux courbes : $4 - x^2 = x + 2 \iff x^2 + x - 2 = (x + 2)(x - 1) = 0 \iff x = -2$ ou $x = 1$. Sur $[-2, 1]$, la parabole $y_{\text{haut}} = 4 - x^2$ est au-dessus de la droite $y_{\text{bas}} = x + 2$. L'aire vaut donc :
\[
    \mathcal{A} = \int_{-2}^1 \bigl((4 - x^2) - (x + 2)\bigr)\dd x = \int_{-2}^1 \bigl(2 - x - x^2\bigr)\dd x = \left[2x - \frac{x^2}{2} - \frac{x^3}{3}\right]_{-2}^1 = \left(2 - \frac{1}{2} - \frac{1}{3}\right) - \left(-4 - 2 + \frac{8}{3}\right) = \boxed{\frac{9}{2}}.
\]

\paragraph{Corrigé de l'Exercice 7.5 (Volumes du cône, de la sphère et de la calotte sphérique) :}
\begin{enumerate}[label=(\alph*)]
    \item \textbf{Cône de rayon $R$ et hauteur $H$ :} Il est engendré par la rotation autour de $Ox$ du segment de droite $y(x) = \frac{R}{H}x$ pour $x \in [0, H]$ :
    \[
        V_{\text{cône}} = \pi \int_0^H \left(\frac{R}{H}x\right)^2\dd x = \frac{\pi R^2}{H^2}\left[\frac{x^3}{3}\right]_0^H = \boxed{\frac{1}{3}\pi R^2 H}.
    \]
    \item \textbf{Sphère de rayon $R$ :} Engendrée par la rotation du demi-cercle $y(x) = \sqrt{R^2 - x^2}$ sur $[-R, R]$ (fonction paire) :
    \[
        V_{\text{sphère}} = \pi \int_{-R}^R \bigl(R^2 - x^2\bigr)\dd x = 2\pi \int_0^R \bigl(R^2 - x^2\bigr)\dd x = 2\pi\left[R^2 x - \frac{x^3}{3}\right]_0^R = 2\pi\left(R^3 - \frac{R^3}{3}\right) = \boxed{\frac{4}{3}\pi R^3}.
    \]
    \item \textbf{Calotte sphérique de hauteur $h$ (entre $x = R - h$ et $x = R$) :}
    \[
        V_{\text{calotte}} = \pi \int_{R-h}^R \bigl(R^2 - x^2\bigr)\dd x = \pi\left[R^2 x - \frac{x^3}{3}\right]_{R-h}^R = \pi\left[\frac{2}{3}R^3 - R^2(R-h) + \frac{(R-h)^3}{3}\right] = \boxed{\pi h^2\left(R - \frac{h}{3}\right)}.
    \]
\end{enumerate}

\paragraph{Corrigé de l'Exercice 7.6 (Volume d'un solide de révolution logarithmique) :}
Par la formule du volume de révolution autour de $Ox$ sur $[1, e]$ :
\[
    V = \pi \int_1^e \left(\frac{\ln(x)}{\sqrt{x}}\right)^2\dd x = \pi \int_1^e \frac{\ln^2(x)}{x}\dd x.
\]
Posons le changement de variable $u = \ln(x) \implies \mathrm{d}u = \dfrac{1}{x}\dd x$. Lorsque $x$ varie de $1$ à $e$, $u = \ln(x)$ varie de $\ln(1) = 0$ à $\ln(e) = 1$ :
\[
    V = \pi \int_0^1 u^2\dd u = \pi\left[\frac{u^3}{3}\right]_0^1 = \boxed{\frac{\pi}{3}}.
\]


\end{document}
